%% ============================================================
%%  Série Avaliação na Prática — FGV CLEAR
%%  Introdução à Inferência Causal
%% ============================================================

% !TEX program = pdflatex
% !TEX encoding = UTF-8 Unicode

\documentclass[12pt, a4paper, openany]{book}

%% ── Codificação e Idioma ────────────────────────────────────
\usepackage[utf8]{inputenc}
\usepackage[T1]{fontenc}
\usepackage[brazil]{babel}

%% ── Fonte ───────────────────────────────────────────────────
\usepackage{tgpagella}

%% ── Matemática ──────────────────────────────────────────────
\usepackage{amsmath, amssymb, amsfonts, dsfont, amsthm}

%% ── Figuras e Tabelas ───────────────────────────────────────
\usepackage{graphicx, float}
\usepackage[labelfont=bf, font=small, justification=justified]{caption}
\usepackage{subcaption}
\usepackage{booktabs, multirow, multicol, array}
\usepackage{longtable}
\usepackage{tabularx}

%% ── Layout ──────────────────────────────────────────────────
\usepackage{geometry}
\geometry{a4paper, left=2.5cm, right=2.5cm, top=2.5cm, bottom=2.5cm}
\usepackage{setspace}
\onehalfspacing
\usepackage{fancyhdr}
\usepackage{emptypage}
\usepackage{enumitem}

%% ── TikZ e PGF ──────────────────────────────────────────────
\usepackage{csquotes}
\usepackage{tikz}
\usepackage{pgfplots}
\pgfplotsset{compat=1.18}
\usetikzlibrary{arrows.meta, positioning, decorations.pathreplacing}

%% ── Cores (paleta CLEAR) ────────────────────────────────────
\usepackage{xcolor}
\definecolor{clearblue}{HTML}{1B3A6E}
\definecolor{clearlightblue}{HTML}{5CA4A9}
\definecolor{cleardarkblue}{HTML}{003F72}
\definecolor{cleargray}{HTML}{6E6E6E}
\definecolor{salmonbg}{RGB}{252,235,230}
\definecolor{salmonframe}{RGB}{230,160,140}
\definecolor{bluebg}{RGB}{235,245,251}
\definecolor{blueframe}{RGB}{58,125,181}
\definecolor{orangebg}{RGB}{254,245,231}
\definecolor{orangeframe}{RGB}{230,126,34}
\definecolor{redbg}{RGB}{250,230,225}
\definecolor{redframe}{RGB}{180,40,40}
\definecolor{codebg}{RGB}{248,248,248}
\definecolor{codeframe}{RGB}{200,200,200}
\definecolor{codegray}{rgb}{0.45,0.45,0.45}
\definecolor{redbg}{RGB}{250,230,225}
\definecolor{redframe}{RGB}{180,40,40}
\definecolor{greenbg}{RGB}{225,242,225}
\definecolor{greenframe}{RGB}{60,140,80}

%% ── Hyperlinks e Bibliografia ───────────────────────────────
\usepackage{hyperref}
\hypersetup{
  colorlinks = true,
  linkcolor  = clearlightblue,
  urlcolor   = clearlightblue,
  citecolor  = clearlightblue,
  pdftitle   = {Introdução à Inferência Causal},
  pdfauthor  = {FGV CLEAR}
}
\usepackage{natbib}
\bibliographystyle{plainnat}

%% ── TColorbox ───────────────────────────────────────────────
\usepackage[many]{tcolorbox}
\tcbuselibrary{breakable, skins}

%% ── Código ──────────────────────────────────────────────────
\usepackage{listings}

%% ── Cabeçalhos e Rodapés ────────────────────────────────────
\pagestyle{fancy}
\fancyhf{}
\fancyhead[R]{\small\color{clearblue}\textit{Guia CLEAR $|$ Introdu\c{c}\~ao \`a Infer\^encia Causal}}
\fancyfoot[C]{\thepage}
\renewcommand{\headrulewidth}{0.4pt}
\renewcommand{\headrule}{\hbox to\headwidth{\color{clearblue}\leaders\hrule height \headrulewidth\hfill}}

%% ── Títulos de Seção ────────────────────────────────────────
\usepackage{titlesec}
\titleformat{\chapter}[display]
  {\normalfont\huge\bfseries\color{clearblue}}
  {\chaptertitlename\ \thechapter}{20pt}{\Huge}
\titleformat{\section}
  {\normalfont\Large\bfseries\color{clearblue}}{\thesection}{1em}{}
\titleformat{\subsection}
  {\normalfont\large\bfseries\color{cleardarkblue}}{\thesubsection}{1em}{}
\titleformat{\subsubsection}
  {\normalfont\normalsize\bfseries\color{cleardarkblue}}{\thesubsubsection}{1em}{}

%% ── Numeração ───────────────────────────────────────────────
\numberwithin{equation}{chapter}
\numberwithin{table}{chapter}
\numberwithin{figure}{chapter}

%% ── Caixas Temáticas ────────────────────────────────────────
\newtcolorbox{destaquebox}[1][Destaque]{
  enhanced, breakable,
  colback=salmonbg, colframe=salmonframe, coltitle=black, fonttitle=\bfseries,
  title=#1, boxrule=0.5pt, arc=2pt,
  left=6pt, right=6pt, top=4pt, bottom=4pt
}
\newtcolorbox{notabox}[1][Nota]{
  enhanced, breakable,
  colback=bluebg, colframe=blueframe, coltitle=white, fonttitle=\bfseries,
  title=#1, boxrule=0.5pt, arc=2pt,
  left=6pt, right=6pt, top=4pt, bottom=4pt
}
\newtcolorbox{conexaobox}[1][Conex\~ao com a s\'erie]{
  enhanced, breakable,
  colback=orangebg, colframe=orangeframe, coltitle=black, fonttitle=\bfseries,
  title=#1, boxrule=0.5pt, arc=2pt,
  left=6pt, right=6pt, top=4pt, bottom=4pt
}
\newtcolorbox{exemplobox}[1][Exemplo]{
  enhanced, breakable,
  colback=greenbg, colframe=greenframe, coltitle=white, fonttitle=\bfseries,
  title=#1, boxrule=0.5pt, arc=2pt,
  left=6pt, right=6pt, top=4pt, bottom=4pt
}
\newtcolorbox{hipotesebox}[1][Hip\'otese]{
  enhanced, breakable,
  colback=salmonbg, colframe=salmonframe, coltitle=black, fonttitle=\bfseries,
  title=#1, boxrule=0.5pt, arc=2pt,
  left=6pt, right=6pt, top=4pt, bottom=4pt
}
\newtcolorbox{atencaobox}[1][Aten\c{c}\~ao]{
  enhanced, breakable,
  colback=redbg, colframe=redframe, coltitle=white, fonttitle=\bfseries,
  title=#1, boxrule=0.5pt, arc=2pt,
  left=6pt, right=6pt, top=4pt, bottom=4pt
}

%% ── Estilo de Listings ──────────────────────────────────────
\lstdefinestyle{Rstyle}{
  language         = R,
  commentstyle     = \color{codegray}\itshape,
  keywordstyle     = \color{clearblue}\bfseries,
  numberstyle      = \tiny\color{codegray},
  stringstyle      = \color{clearlightblue},
  basicstyle       = \ttfamily\small,
  breakatwhitespace= false,
  breaklines       = true,
  captionpos       = b,
  keepspaces       = true,
  numbers          = left,
  numbersep        = 8pt,
  showspaces       = false,
  showstringspaces = false,
  showtabs         = false,
  tabsize          = 2,
  frame            = single,
  rulecolor        = \color{codeframe},
  backgroundcolor  = \color{codebg},
  inputencoding    = utf8,
  extendedchars    = true,
  literate =
    {á}{{\'{a}}}1 {à}{{\`{a}}}1 {â}{{\^{a}}}1 {ã}{{\~{a}}}1
    {é}{{\'{e}}}1 {ê}{{\^{e}}}1
    {í}{{\'{i}}}1
    {ó}{{\'{o}}}1 {ô}{{\^{o}}}1 {õ}{{\~{o}}}1
    {ú}{{\'{u}}}1
    {ç}{{\c{c}}}1
    {Á}{{\'{A}}}1 {Â}{{\^{A}}}1 {Ã}{{\~{A}}}1
    {É}{{\'{E}}}1 {Ê}{{\^{E}}}1
    {Í}{{\'{I}}}1
    {Ó}{{\'{O}}}1 {Ô}{{\^{O}}}1 {Õ}{{\~{O}}}1
    {Ú}{{\'{U}}}1
    {Ç}{{\c{C}}}1
}
\lstset{style=Rstyle}

%% ── Teoremas ────────────────────────────────────────────────
\newtheorem{definition}{Defini\c{c}\~ao}[chapter]
\newtheorem{proposition}{Proposi\c{c}\~ao}[chapter]
\newtheorem{remark}{Observa\c{c}\~ao}[chapter]

%% ── Atalhos Matemáticos ─────────────────────────────────────
\newcommand{\E}{\mathbb{E}}
\newcommand{\Var}{\mathrm{Var}}
\newcommand{\Cov}{\mathrm{Cov}}
\newcommand{\indep}{\perp\!\!\!\perp}
\newcommand{\iid}{\overset{\text{i.i.d.}}{\sim}}

%% ── Metadados de Figuras/Quadros ────────────────────────────
\newcommand{\figmeta}[3]{%
  \par\vspace{4pt}%
  \begin{minipage}{\linewidth}%
  \footnotesize%
  \textbf{Fonte:} #1%
  \ifx&#2&\else\\\textbf{Elabora\c{c}\~ao:} #2\fi%
  \ifx&#3&\else\\\textbf{Nota:} #3\fi%
  \end{minipage}%
  \vspace{6pt}%
}

%% ── Profundidade de numera\c{c}\~ao ─────────────────────────
\setcounter{secnumdepth}{2}
\setcounter{tocdepth}{2}

% ============================================================
\begin{document}

\frontmatter

% ============================================================
% CAPA
% ============================================================
\thispagestyle{empty}
\begin{center}
\vspace*{3cm}

{\Large\bfseries\color{clearlightblue} S\'ERIE: AVALIA\c{C}\~AO NA PR\'ATICA}

\vspace{1.8cm}

{\Huge\bfseries\color{clearblue} Como ir de correlação para causalidade }

\vspace{2.2cm}

{\large FGV CLEAR $|$ FGV EESP}

\vspace{0.5cm}

{\large Maio de 2026}

\vfill
\end{center}
\newpage

% ============================================================
% FICHA TÉCNICA
% ============================================================
\thispagestyle{empty}
\vspace*{1.5cm}

{\Large\bfseries\color{clearblue} Ficha T\'ecnica}

\bigskip

\noindent\textbf{Equipe t\'ecnica:} Caio de Souza Castro, Luigi Garzon.

\bigskip

\noindent\textbf{Apresenta\c{c}\~ao.} Este guia integra a s\'erie de publica\c{c}\~oes \textit{Avalia\c{c}\~ao na Pr\'atica}, desenvolvida pelo FGV CLEAR com o objetivo de ampliar o acesso a conhecimentos sobre monitoramento e avalia\c{c}\~ao com foco em pol\'iticas p\'ublicas.

\medskip

\noindent Fundado em 2015, o FGV CLEAR tem se dedicado ao fortalecimento da cultura de gest\~ao orientada por evid\^encias no Brasil e em pa\'ises lus\'ofonos. Com sede na Escola de Economia de S\~ao Paulo da Funda\c{c}\~ao Getulio Vargas (FGV EESP), o FGV CLEAR atua como centro regional da Iniciativa CLEAR (\textit{Centers for Learning on Evaluation and Results}).

\medskip

\noindent Saiba mais sobre o FGV CLEAR e acesse outras publica\c{c}\~oes em: \textbf{www.fgvclear.org}

\vfill

\noindent{\small\color{cleargray}%
\textbf{Cita\c{c}\~ao recomendada:} FGV CLEAR. (\the\year). \textit{Introdu\c{c}\~ao \`a Infer\^encia Causal}. S\'erie Avalia\c{c}\~ao na Pr\'atica. S\~ao Paulo: FGV CLEAR.

\medskip

\noindent Os textos desta publica\c{c}\~ao s\~ao de responsabilidade exclusiva dos autores. As opini\~oes expressas n\~ao refletem necessariamente a posi\c{c}\~ao das institui\c{c}\~oes parceiras.}

\newpage

% ============================================================
% SUMÁRIO
% ============================================================
% ─── SOBRE ESTE GUIA ─────────────────────────────────────────────────────────
\thispagestyle{empty}
\section*{Sobre este guia}
O cigarro causa câncer ou são as pessoas que fumam que já tinham mais
propensão a adoecer? Quando o Bolsa Família foi associado a queda na
pobreza, foi o programa que reduziu a pobreza ou o crescimento econômico
do mesmo período que ajudou? Crianças que frequentam creche têm
desempenho escolar melhor depois porque a creche desenvolve habilidades
cognitivas ou porque famílias que matriculam os filhos em creche são,
sistematicamente, famílias que investem mais nos filhos de outras formas
que não medimos? Cada uma dessas perguntas tem a mesma estrutura: uma
correlação observada nos dados, e uma dúvida razoável sobre se ela
reflete causa ou outra coisa. Distinguir essas duas possibilidades está no
centro de praticamente toda decisão de política pública. E quase nunca é
trivial.

A literatura contemporânea trata o problema com um arcabouço conceitual
chamado resultados potenciais. A ideia é a seguinte: para cada indivíduo
da população, imagine que existem dois mundos possíveis. No mundo em que
ele recebeu o tratamento, observa-se um certo resultado. No mundo em que
ele não recebeu, outro. O efeito causal do tratamento, para aquele
indivíduo, é a diferença entre os dois resultados. O problema é que cada
pessoa vive apenas um dos dois mundos. Quem fumou não pode ser observado
como se não tivesse fumado, e vice-versa. Os dois resultados potenciais
existem como construção teórica, mas só um deles é realizado e
observável. Por isso, o efeito causal individual é, em última instância,
inobservável.

A engenharia inteira da inferência causal moderna é a tentativa de
recuperar médias desses efeitos comparando pessoas diferentes que vivem
cada qual um dos dois mundos. A pergunta deixa de ser ``o que teria
acontecido com este indivíduo no outro mundo?'' e passa a ser ``no
agregado, qual é a diferença entre o que de fato aconteceu com quem foi
tratado e o que aconteceu com pessoas comparáveis que não foram?''. Essa
transição parece pequena mas tem implicação grande: a comparação só é
válida sob hipóteses específicas, e diferentes hipóteses sustentam
diferentes métodos.

O arcabouço separa três etapas que costumam ser confundidas. Identificação
é a pergunta de quando o parâmetro de interesse pode ser recuperado dos
dados observados, e sob quais suposições. Esse é o passo onde se decide
se a pergunta tem resposta, antes mesmo de pensar em estimação. Estimação
é como calcular o parâmetro numa amostra finita, com decisões sobre forma
funcional, ponderação, controles, especificação. Inferência é como
quantificar a incerteza em torno do que foi estimado, levando em conta
que com outra amostra teríamos outro número. Cada uma dessas etapas tem
armadilhas próprias, e tratá-las como uma só coisa (como em ``rodar uma
regressão'') esconde decisões importantes e propaga erros.

Este guia é a base conceitual da série. Os demais guias (RDD, controle
sintético, pareamento, diferenças-em-diferenças) aplicam essas etapas a
contextos específicos, com hipóteses adaptadas a cada situação. Quem
domina a estrutura aqui apresentada lê os outros guias com mais clareza
sobre o que cada método assume, o que cada estimativa significa, e por
que duas estimativas aparentemente da mesma coisa, vindas de métodos
diferentes, podem dar respostas diferentes sem que nenhuma das duas
esteja errada.
\newpage

% ─── COMO LER ESTE GUIA ──────────────────────────────────────────────────────
\thispagestyle{empty}
\section*{Como ler este guia}
Três tipos de caixa estruturam a leitura, identificados por cor:
\vspace{0.3cm}
\begin{notabox}[Caixa azul --- pontos-chave]
Aparece ao fim de cada capítulo, com sínteses e recomendações para quem
precisa decidir.
\end{notabox}
\begin{exemplobox}[Caixa verde --- exemplo aplicado]
Mostra como as ideias do capítulo se traduzem em situações de política
pública. Os casos são fictícios, mas calibrados por padrões recorrentes
na literatura e na prática de gestão.
\end{exemplobox}
\begin{atencaobox}[Caixa vermelha --- cuidados essenciais]
Sinaliza armadilhas comuns --- pontos em que uma escolha mal fundamentada
compromete toda a análise.
\end{atencaobox}
\vspace{0.5cm}
\noindent O guia faz referência recorrente ao \textit{Guia CLEAR de
Monitoramento e Avaliação de Políticas Públicas: do diagnóstico à decisão}
\citep{Lima2025} --- doravante \textbf{Guia CLEAR de M\&A} ---, que trata
em profundidade as etapas do ciclo da política pública adotado na série.
Quando o texto menciona Teoria da Mudança, plano de monitoramento ou
tipologia de avaliações, é nesse guia que o leitor encontra o
desenvolvimento completo.
\newpage

\tableofcontents
\newpage

% ============================================================
% CONTEÚDO
% ============================================================

\mainmatter

\chapter{Introdução}
\label{cap:introducao}

A avaliação quantitativa de políticas públicas enfrenta uma tensão permanente
entre duas perguntas aparentemente semelhantes, mas fundamentalmente distintas.
A primeira é descritiva: \textit{programas sociais estão associados a melhores
resultados?} A segunda é causal: \textit{esses programas causam melhores
resultados?} A diferença entre as duas não é apenas filosófica. Ela determina
se uma política deve ser expandida, reformulada ou descontinuada --- e,
portanto, como recursos públicos escassos são alocados.

Este guia é sobre a segunda pergunta. Seu objetivo é oferecer ao leitor as
ferramentas conceituais e práticas para estimar efeitos causais a partir de
dados observacionais e experimentais. Para isso, seguimos a estrutura que a
econometria moderna consagrou como o caminho rigoroso entre uma pergunta de
política e um número calculado em dados: \textbf{identificação},
\textbf{estimação} e \textbf{inferência}.


\section{O Problema e a Solução}

Toda análise causal começa com o \textbf{problema fundamental da inferência
causal}: para cada unidade de observação --- seja ela uma pessoa, um município
ou uma escola ---, só é possível observar o resultado sob \textit{um} dos
tratamentos possíveis. O salário de quem foi para a faculdade não pode ser
observado simultaneamente no mundo contrafactual em que a mesma pessoa não foi.
A nota de um aluno em turma pequena não pode ser comparada, para o mesmo aluno,
à nota que teria tirado em turma regular.

Essa limitação é estrutural, não uma falha dos dados. Ela implica que efeitos
causais individuais são intrinsecamente não-observáveis, e que a inferência
causal é, por natureza, um exercício de raciocínio sobre o que
\textit{não} foi observado.

A solução desenvolvida pela econometria moderna não elimina esse problema ---
ela o contorna por meio de hipóteses sobre o processo que gerou os dados. Se
indivíduos similares foram expostos a tratamentos distintos por razões que não
se relacionam com seus resultados potenciais, a comparação entre eles é
informativa sobre o efeito causal. Experimentos aleatórios garantem essa
condição por design. Métodos quasi-experimentais --- regressão com
descontinuidade, diferenças em diferenças, pareamento, variáveis instrumentais
--- procuram encontrá-la nos dados observacionais.


\section{A Revolução da Credibilidade}

A formação da econometria aplicada moderna foi marcada por uma virada
metodológica que ficou conhecida como \textbf{revolução da credibilidade}
\citep{Angrist2010}. Até os anos 1980, a estimação de efeitos causais dependia
amplamente de modelos estruturais com fortes restrições derivadas da teoria
econômica. Esses modelos produziam estimativas, mas sua credibilidade dependia
de hipóteses que raramente podiam ser verificadas nos dados.

A partir do trabalho seminal de pesquisadores como David Card, Joshua Angrist e
Guido Imbens, a prática econométrica mudou de direção: em vez de depender de
restrições teóricas implausíveis, passou-se a buscar \textit{variação exógena}
nos dados --- situações em que o tratamento é atribuído de forma suficientemente
aleatória para permitir identificação causal. A formalização do arcabouço de
resultados potenciais \citep{Rubin1974, Rubin1978, Neyman1923} forneceu o
vocabulário teórico unificado para essa abordagem.

Em 2021, o Prêmio Nobel de Ciências Econômicas reconheceu essa contribuição.
Card foi premiado pela aplicação de métodos quasi-experimentais a questões do
mercado de trabalho; Angrist e Imbens, pelas contribuições metodológicas à
análise de relações causais com dados observacionais. A Série Avaliação na
Prática, da qual este guia faz parte, é uma expressão direta desse legado.


\section{A Cadeia de Análise Causal}

O caminho que este guia percorre pode ser descrito como uma cadeia com
três elos:

\begin{enumerate}[leftmargin=*, label=\textbf{\arabic*.}]
  \item \textbf{Identificação} — Definir o parâmetro de interesse, derivar as
    hipóteses sobre o processo gerador dos dados que o tornam identificável, e
    especificar o estimando: a quantidade dos dados que, sob essas hipóteses,
    coincide com o parâmetro.

  \item \textbf{Estimação} — Construir um estimador que calcule o estimando a
    partir de uma amostra finita. Discutir suas propriedades (consistência,
    viés, eficiência) e as hipóteses funcionais adicionais que ele impõe.

  \item \textbf{Inferência} — Quantificar a incerteza amostral em torno da
    estimativa. Escolher estimadores de variância adequados à estrutura dos
    dados, realizar testes de hipótese, construir intervalos de confiança e
    verificar a plausibilidade das hipóteses de identificação por meio de
    testes de falsificação.
\end{enumerate}

Os três elos respondem a perguntas distintas e se articulam sequencialmente:
a inferência pressupõe a estimação, que pressupõe a identificação. Um erro
num dos elos não é corrigido pelos demais.

\begin{atencaobox}[Erro-padrão pequeno não é garantia de causalidade]
Um resultado com erro-padrão muito pequeno e p-valor minúsculo pode ser
completamente destituído de interpretação causal se a identificação for
inválida. A precisão estatística mede quão bem estimamos o estimando ---
não se o estimando corresponde ao parâmetro econômico de interesse.
Esta distinção percorre todo o guia.
\end{atencaobox}


\section{Exemplos Condutores}

Ao longo do guia, utilizamos dois exemplos para ilustrar os conceitos.

O primeiro é a questão dos \textbf{retornos à educação}: qual é o efeito
causal de um ano adicional de escolaridade sobre o salário? O viés de seleção
é imediato e intuitivo --- pessoas com mais educação tendem a ter maior
habilidade inata e condições socioeconômicas mais favorecidas, o que elevaria
seus salários mesmo sem o diploma. O exemplo tem longa tradição na literatura
\citep{Mincer1974, Card1994, Card1995} e dados abundantes no Brasil por meio
da PNAD Contínua.

O segundo exemplo, apresentado em detalhe no Capítulo~\ref{sec:star}, é o
\textbf{Experimento STAR} --- um experimento aleatório conduzido no Tennessee
entre 1985 e 1989, que avaliou o efeito do tamanho da turma sobre o desempenho
escolar \citep{Finn1990, Krueger1999}. O STAR serve como laboratório empírico
ideal para ilustrar identificação por aleatorização, estimação com efeitos
fixos, erros-padrão clusterizados, bootstrap, múltiplos testes, inferência
por permutação e testes de falsificação.


\section{Estrutura do Guia}

O Capítulo~\ref{sec:identificacao} apresenta o problema de identificação
causal: o arcabouço de resultados potenciais, os parâmetros de interesse
(ATE, ATT, ATU), o viés de seleção e as hipóteses de identificação.

O Capítulo~\ref{sec:estimacao} trata da passagem do estimando ao estimador:
o princípio plug-in, as propriedades dos estimadores, o papel da forma
funcional, o IPW, os estimadores de painel e o estimador de Lin.

O Capítulo~\ref{sec:inferencia} cobre inferência estatística: erros-padrão,
testes de hipótese, intervalos de confiança, clustering, bootstrap e
inferência por permutação, encerrando com os principais testes de
falsificação.

O Capítulo~\ref{sec:star} aplica todos esses elementos ao Experimento STAR.

O Capítulo~\ref{cap:conclusao} sintetiza as principais lições e conecta
este guia aos demais da Série Avaliação na Prática.

\begin{notabox}[Pré-requisitos]
Este guia pressupõe familiaridade com estatística básica (esperança,
variância, distribuição normal, testes de hipótese) e com o modelo de
regressão linear simples. Conhecimentos de cálculo são úteis mas não
indispensáveis. O código em R requer conhecimento básico da linguagem;
os pacotes utilizados são apresentados com exemplos autocontidos.
\end{notabox}
% =============================================================================
% Seção 2 — Identificação
% Guia Clear de Introdução à Inferência Causal — FGV EESP Clear
% =============================================================================

\chapter{Identificação}
\label{cap:identificacao}
\label{sec:identificacao}

Toda análise causal começa com uma pergunta da teoria econômica e termina com um número calculado em uma amostra finita. O caminho entre os dois é mais longo do que parece. Ele atravessa quatro espaços conceituais distintos: o \textbf{parâmetro econômico} (definido pela teoria, vivendo no espaço do modelo), o \textbf{estimando} (um funcional da distribuição populacional observável), o \textbf{estimador} (uma função dos dados amostrais) e a \textbf{inferência} (a quantificação da incerteza amostral). Esta seção trata exclusivamente do segundo elo dessa cadeia: a passagem do parâmetro para o estimando. Esse elo se chama \textit{identificação}.

Considere um modelo teórico genérico em que uma variável de interesse $Y$ é função de uma variável de tratamento $X$ e de fatores não observados $\varepsilon$:
\begin{equation}
Y = f(X, \varepsilon).
\label{eq:modelo-generico}
\end{equation}
A teoria econômica nos diz que o objeto relevante para responder à pergunta de política é, tipicamente, o efeito de mudar $X$ sobre $Y$. Em termos potenciais, definimos $Y_i(x)$ como o resultado que o indivíduo $i$ \textit{teria} se recebesse o valor $x$ do tratamento. O efeito causal individual de mover $X$ de $0$ para $1$ é:
\begin{equation}
\delta_i = Y_i(1) - Y_i(0).
\label{eq:efeito-individual}
\end{equation}
Note que $\delta_i$ é definido pela diferença entre dois mundos contrafactuais: o mundo em que $i$ recebe o tratamento e o mundo em que $i$ não o recebe. Para cada indivíduo, observamos apenas \textit{um} desses mundos. Esse é o \textbf{problema fundamental da inferência causal}: o contrafactual nunca é observado.

Como, então, podemos aprender alguma coisa sobre objetos como o ATE (\textit{Average Treatment Effect} --- Efeito Médio do Tratamento), definido por
\begin{equation}
\text{ATE} = E[Y_i(1) - Y_i(0)],
\label{eq:ate}
\end{equation}
ou o ATT (\textit{Average Treatment Effect on the Treated} --- Efeito Médio do Tratamento sobre os Tratados),
\begin{equation}
\text{ATT} = E[Y_i(1) - Y_i(0) \mid D_i = 1]?
\label{eq:att}
\end{equation}
Esses parâmetros vivem no espaço do modelo econômico. Eles são funções de objetos --- os resultados potenciais $Y_i(1)$ e $Y_i(0)$ --- que nunca aparecem juntos nos dados. O que efetivamente observamos, em uma amostra de $n$ indivíduos, é o trio $(Y_i, D_i, W_i)$, onde $D_i \in \{0,1\}$ indica se $i$ recebeu o tratamento e $Y_i = D_i \cdot Y_i(1) + (1 - D_i) \cdot Y_i(0)$ é o resultado realizado. A distribuição populacional desses três objetos --- chame-a de $F(Y, D, W)$ --- é, em princípio, aprendível. Funcionais dessa distribuição --- como $E[Y \mid D=1] - E[Y \mid D=0]$ --- são \textit{estimandos}. \textbf{Identificação é o conjunto de hipóteses sobre o processo gerador dos dados que estabelece uma equivalência entre o parâmetro econômico e algum estimando observável.}

Antes de prosseguir, precisamos isolar uma hipótese frequentemente tratada como invisível: a SUTVA (\textit{Stable Unit Treatment Value Assumption} --- Hipótese de Estabilidade do Valor do Tratamento). A SUTVA tem duas componentes. Primeiro, ela exige \textit{ausência de interferência}: o resultado potencial de $i$ depende apenas do tratamento de $i$, não do tratamento dos demais. Formalmente, escrever $Y_i(d_i)$ em vez de $Y_i(d_1, \dots, d_n)$ pressupõe que o vetor completo de tratamentos colapsa em uma única coordenada. Segundo, ela exige \textit{ausência de versões múltiplas do tratamento}: $D_i = 1$ tem o mesmo significado para todos os indivíduos. Sem SUTVA, os próprios resultados potenciais não estão bem definidos como funções de uma única variável $D_i$, e a discussão sobre identificação precisa ser refeita. Iremos voltar a esse ponto na Subseção~\ref{subsec:interferencia}.

Estabelecida essa cláusula, o restante desta seção desenvolve a identificação em um modelo linear simples (Subseção~\ref{subsec:linear-binario}), estende o argumento para painel (Subseção~\ref{subsec:painel}), discute o que muda quando o tratamento não é binário (Subseção~\ref{subsec:nao-binario}) e fecha com um panorama dos métodos quasi-experimentais que substituem a hipótese de exogeneidade (Subseção~\ref{subsec:alem-rct}). Em cada caso, seguimos a mesma ordem: o parâmetro econômico, a hipótese de identificação, o estimando, o argumento substantivo de plausibilidade.

% =============================================================================
\section{Modelo Linear com Tratamento Binário}
\label{subsec:linear-binario}

Para introduzir as ideias com o mínimo de fricção algébrica, começamos por um caso quase trivial: $X$ é binário e não há covariadas. Postulamos que a função $f$ em \eqref{eq:modelo-generico} é \textit{linear} em $X$:
\begin{equation}
Y_i = \alpha + \beta \cdot D_i + \varepsilon_i.
\label{eq:modelo-linear-binario}
\end{equation}
É importante notar que assumir linearidade já é uma hipótese sobre o mundo. Em geral, a relação entre $Y$ e $X$ pode ser não linear, e o efeito do tratamento pode variar entre indivíduos. No caso particular em que $D_i \in \{0, 1\}$ \textit{e} não há covariadas, contudo, a linearidade é praticamente uma identidade: dois pontos --- a média de $Y$ entre tratados e a média entre controles --- definem uma reta. Os parâmetros $\alpha$ e $\beta$ apenas reorganizam a informação que já existe nessas duas médias. O termo $\varepsilon_i$ recolhe tudo o que afeta $Y_i$ além de $D_i$.

Comparando \eqref{eq:modelo-linear-binario} com a definição de resultados potenciais, vemos que
\begin{equation}
Y_i(0) = \alpha + \varepsilon_i, \qquad Y_i(1) = \alpha + \beta + \varepsilon_i,
\label{eq:potenciais-linear}
\end{equation}
o que implica $\delta_i = Y_i(1) - Y_i(0) = \beta$. \textit{Neste modelo}, o efeito individual é constante e igual a $\beta$. Esta é uma segunda hipótese, ainda mais forte do que a linearidade --- ela exclui efeitos heterogêneos por construção. Iremos relaxá-la nas Subseções~\ref{subsubsec:covariadas} e~\ref{subsec:nao-binario}. Por ora, $\beta$ é, simultaneamente, o ATE, o ATT e o efeito individual de qualquer pessoa na população.

A pergunta de identificação se reduz, portanto, a: \textit{sob quais hipóteses o parâmetro $\beta$ é igual a algum funcional da distribuição observável de $(Y, D)$?} As próximas subseções respondem a essa pergunta sob diferentes premissas.

% -----------------------------------------------------------------------------
\subsection{Hipótese de Exogeneidade}
\label{subsubsec:exogeneidade}

A passagem do parâmetro econômico para o estimando depende, no modelo linear, de uma única hipótese central: a hipótese de \textbf{exogeneidade}. Ela diz que o tratamento e o erro são independentes em média:
\begin{equation}
E[\varepsilon_i \mid D_i] = 0.
\label{eq:exogeneidade}
\end{equation}
Em palavras, a expectativa condicional de tudo o que afeta $Y_i$ além de $D_i$ não muda quando $D_i$ muda. Quem é tratado e quem não é tratado têm, em média, os mesmos $\varepsilon_i$.

Antes da álgebra, vale a intuição. Suponha que estamos avaliando o efeito de um curso de qualificação profissional --- o PRONATEC (Programa Nacional de Acesso ao Ensino Técnico e Emprego), por exemplo --- sobre a renda futura. O termo $\varepsilon_i$ contém motivação, conexões familiares, habilidade cognitiva, condições do mercado local. Exogeneidade exige que tratados e controles tenham, na média, o mesmo nível dessas variáveis. Quando o tratamento é resultado de uma escolha do próprio indivíduo --- como inscrever-se voluntariamente no curso ---, essa hipótese é fortíssima: os mais motivados podem se autosselecionar para o tratamento, violando \eqref{eq:exogeneidade}.

\textbf{Exemplo numérico hipotético.} Considere uma população de seis indivíduos, três tratados (A, B, C) e três controles (1, 2, 3), com resultados observados conforme a Tabela~\ref{tab:exemplo-binario}.

\begin{table}[h]
\centering
\begin{tabular}{lcc}
\toprule
Indivíduo & $D_i$ & $Y_i$ \\
\midrule
A & 1 & 12 \\
B & 1 & 14 \\
C & 1 & 13 \\
1 & 0 & 8 \\
2 & 0 & 10 \\
3 & 0 & 9 \\
\bottomrule
\end{tabular}
\caption{Dados hipotéticos: três tratados e três controles.}
\label{tab:exemplo-binario}
\end{table}

A média entre tratados é $\bar{Y}_1 = 13$; entre controles, $\bar{Y}_0 = 9$. A diferença $\bar{Y}_1 - \bar{Y}_0 = 4$ é o estimador empírico (a versão amostral) do estimando $E[Y \mid D=1] - E[Y \mid D=0]$. \textit{Sob a hipótese de exogeneidade}, esse estimando é igual ao ATE.

A derivação algébrica torna esse argumento explícito. Aplicando a esperança condicional em \eqref{eq:modelo-linear-binario}:
\begin{align}
E[Y_i \mid D_i = 1] &= \alpha + \beta + E[\varepsilon_i \mid D_i = 1], \label{eq:cond-tratados}\\
E[Y_i \mid D_i = 0] &= \alpha + E[\varepsilon_i \mid D_i = 0]. \label{eq:cond-controles}
\end{align}
Subtraindo \eqref{eq:cond-controles} de \eqref{eq:cond-tratados}:
\begin{equation}
E[Y_i \mid D_i = 1] - E[Y_i \mid D_i = 0] = \beta + \big( E[\varepsilon_i \mid D_i = 1] - E[\varepsilon_i \mid D_i = 0] \big).
\label{eq:dec-bruta}
\end{equation}
Sob \eqref{eq:exogeneidade}, ambos os termos $E[\varepsilon_i \mid D_i = 1]$ e $E[\varepsilon_i \mid D_i = 0]$ são iguais a zero, e, portanto:
\begin{equation}
\beta = \underbrace{E[Y_i \mid D_i = 1] - E[Y_i \mid D_i = 0]}_{\text{estimando observável}}.
\label{eq:identificacao-rct}
\end{equation}
A equivalência entre o parâmetro econômico $\beta$ e o estimando $E[Y \mid D=1] - E[Y \mid D=0]$ está estabelecida. O lado esquerdo vive no modelo; o lado direito vive na distribuição observável de $(Y, D)$. A hipótese \eqref{eq:exogeneidade} é a \textit{ponte}.

A pergunta natural é: quando podemos defender que essa ponte é sólida? A resposta canônica é o RCT (\textit{Randomized Controlled Trial} --- Experimento Controlado Aleatorizado). Em um RCT, o pesquisador atribui $D_i$ por sorteio. Como o sorteio é independente de qualquer característica do indivíduo, ele é necessariamente independente de $\varepsilon_i$, e \eqref{eq:exogeneidade} vale por construção do desenho. Esse é o sentido em que se diz que o RCT é o \textit{padrão-ouro}: a hipótese de identificação não depende de argumentos substantivos sobre o comportamento dos agentes; ela é garantida pelo procedimento de atribuição.\footnote{Mesmo no RCT, exogeneidade pode falhar empiricamente por causa de \textit{attrition} diferencial, descumprimento do tratamento atribuído ou efeitos de equilíbrio geral. O ``padrão-ouro'' refere-se ao desenho ideal, não à execução perfeita.}

Vimos como, sob exogeneidade, o estimando da diferença de médias condicionais identifica o efeito causal. Mas o que acontece quando essa hipótese falha? É o que examinamos a seguir.

% -----------------------------------------------------------------------------
\subsection{Viés de Seleção}
\label{subsubsec:vies-selecao}

A maior parte dos dados de políticas públicas \textit{não} vem de RCTs. O tratamento é tipicamente alocado por algum mecanismo --- elegibilidade administrativa, autosseleção, focalização territorial --- que pode ser correlacionado com características que também afetam $Y$. Quando isso ocorre, $E[\varepsilon_i \mid D_i = 1] \neq E[\varepsilon_i \mid D_i = 0]$, e a equivalência \eqref{eq:identificacao-rct} se desfaz.

Para vermos o que sobra, vamos abrir o estimando observável da diferença de médias \textit{sem} impor exogeneidade. Reescrevendo $Y_i = D_i Y_i(1) + (1 - D_i) Y_i(0)$:
\begin{align}
E[Y_i \mid D_i = 1] - E[Y_i \mid D_i = 0]
  &= E[Y_i(1) \mid D_i = 1] - E[Y_i(0) \mid D_i = 0] \nonumber \\
  &= \big[ E[Y_i(1) \mid D_i = 1] - E[Y_i(0) \mid D_i = 1] \big] \nonumber \\
  &\quad + \big[ E[Y_i(0) \mid D_i = 1] - E[Y_i(0) \mid D_i = 0] \big]. \label{eq:dec-vies}
\end{align}
A primeira linha do lado direito é o \textbf{ATT}: o efeito médio do tratamento sobre os tratados. A segunda linha é o \textbf{viés de seleção}: a diferença entre o resultado potencial \textit{sob não-tratamento} de tratados e controles. Em notação compacta:
\begin{align}
E[Y_i \mid D_i = 1] - E[Y_i \mid D_i = 0]
&= \underbrace{E[Y_i(1) - Y_i(0) \mid D_i = 1]}_{\text{ATT}} \nonumber \\ 
&+ \underbrace{E[Y_i(0) \mid D_i = 1] - E[Y_i(0) \mid D_i = 0]}_{\text{viés de seleção}}.
\label{eq:ATT-mais-vies}
\end{align}
O viés de seleção tem uma interpretação direta. Ele é a diferença que existiria entre tratados e controles \textit{mesmo que o tratamento não tivesse sido aplicado}. Se os tratados são, por características pré-existentes, ``melhores'' em termos de $Y$ --- mais escolarizados, mais saudáveis, mais empregáveis ---, o viés é positivo: a diferença simples de médias superestima o ATT. Se os tratados são pior posicionados --- por exemplo, se o programa foca pessoas em situação de vulnerabilidade ---, o viés é negativo, e a diferença subestima o ATT. \textit{O sinal e a magnitude do viés são uma questão substantiva}, não estatística.

Tome como exemplo o Bolsa Família. Suponha que comparemos o desempenho escolar de crianças cujas famílias recebem o benefício com o de crianças cujas famílias não recebem, usando dados administrativos de larga escala. A diferença bruta tipicamente aparece negativa: crianças em famílias beneficiárias têm pior desempenho escolar. Será que isso significa que o programa \textit{prejudica} a aprendizagem? Quase certamente não. O programa é focalizado em famílias de baixa renda, com pais menos escolarizados, vivendo em municípios com piores escolas. Essas características também reduzem $Y_i(0)$ --- o resultado escolar que a criança teria \textit{na ausência} do programa. O viés de seleção é negativo, e domina o ATT. Vimos que precisão amostral é irrelevante neste caso: por mais grande que seja a amostra, a diferença bruta de médias \textit{não} identifica o efeito causal.

A álgebra de \eqref{eq:ATT-mais-vies} mostra a fonte da confusão entre correlação e causalidade. Sob exogeneidade, o segundo termo zera e a diferença bruta identifica o ATT (que, como $\beta$ é constante neste modelo, coincide com o ATE). Sem exogeneidade, restam duas componentes que somente em casos felizes têm o mesmo sinal. \textit{O trabalho do econometrista, quando o RCT não está disponível, é construir um argumento --- baseado em desenho de pesquisa, não em controles estatísticos arbitrários --- que torne o viés de seleção zero ou desprezível.} Esse é o tema da Subseção~\ref{subsec:alem-rct}.

Antes de chegar lá, ainda precisamos relaxar duas hipóteses que assumimos sem dizer: SUTVA e a ausência de covariadas.

% -----------------------------------------------------------------------------
\subsection{Interferência (Relaxando SUTVA)}
\label{subsec:interferencia}

Toda a argumentação anterior pressupôs que o resultado potencial de $i$ depende apenas do tratamento de $i$. Em muitos contextos de política pública, isso é uma simplificação séria. O resultado de uma criança vacinada depende de quantas \textit{outras} crianças foram vacinadas (efeito de imunidade de rebanho). O resultado de um trabalhador qualificado depende de quantos \textit{outros} trabalhadores foram qualificados na mesma cidade (efeito de equilíbrio no mercado de trabalho). O resultado de um aluno em uma turma depende do tratamento dos colegas (efeito de pares). Quando o tratamento de $j$ afeta o resultado de $i$, dizemos que há \textit{interferência} ou \textit{spillovers}.

Sob interferência, o resultado potencial não é mais $Y_i(d_i)$, mas
\begin{equation}
Y_i(d_1, d_2, \dots, d_n),
\label{eq:resultado-vetor}
\end{equation}
uma função do vetor completo de tratamentos atribuídos. Isso explode o espaço de resultados potenciais: para cada indivíduo, há $2^n$ contrafactuais possíveis. Sem alguma restrição sobre a forma da interferência, o problema é intratável.

A solução padrão é assumir um \textbf{mapa de exposição} (\textit{exposure mapping}). Postulamos que o vetor $(d_1, \dots, d_n)$ afeta $Y_i$ apenas por meio de uma função resumo $g_i(d_1, \dots, d_n)$ de baixa dimensão:
\begin{equation}
Y_i(d_1, \dots, d_n) = Y_i(d_i, g_i(d_1, \dots, d_n)).
\label{eq:exposure-mapping}
\end{equation}
Casos típicos: $g_i$ é a fração de tratados na vizinhança de $i$; $g_i$ é o número de pares tratados; $g_i$ é o tratamento da escola em que $i$ está matriculado.\footnote{A escolha do mapa de exposição é uma hipótese substantiva. Diferentes mapas implicam diferentes parâmetros identificáveis, e a literatura recente (Aronow e Samii 2017, Manski 2013, Sävje 2024) discute extensamente como escolher e validar essas restrições.}

Considere o exemplo da vacinação contra o sarampo no SUS. O resultado de interesse é se a criança $i$ contrai a doença. Mas a probabilidade de contrair sarampo não depende apenas de a criança $i$ estar ou não vacinada --- depende crucialmente da cobertura vacinal na sua comunidade. Se assumirmos um mapa de exposição em que $g_i$ é a taxa de cobertura vacinal no município de $i$, podemos escrever
\begin{equation}
Y_i = Y_i(D_i, G_i),
\label{eq:vacinacao}
\end{equation}
onde $G_i$ é a cobertura municipal. O parâmetro de interesse passa a ser o efeito direto $E[Y_i(1, g) - Y_i(0, g)]$ (efeito da vacinação individual mantendo a cobertura fixa) e/ou o efeito indireto $E[Y_i(0, g') - Y_i(0, g)]$ (efeito de aumentar a cobertura sobre quem não foi vacinado). Esses são objetos diferentes, com interpretações de política diferentes.

Sob \eqref{eq:exposure-mapping} e uma versão estendida da hipótese de exogeneidade --- $(Y_i(d, g))_{d, g} \perp (D_i, G_i)$ --- a identificação procede analogamente ao caso anterior:
\begin{equation}
E[Y_i(1, g) - Y_i(0, g)] = E[Y_i \mid D_i = 1, G_i = g] - E[Y_i \mid D_i = 0, G_i = g].
\label{eq:identificacao-spillover}
\end{equation}
Note-se que a randomização individual --- que garante exogeneidade no sentido de \eqref{eq:exogeneidade} --- \textit{não} garante exogeneidade no sentido estendido. Quando interferência está presente, o desenho experimental precisa ser repensado. Designs com randomização em dois níveis (saturação variável entre grupos) são uma das soluções padrão.\footnote{Em programas de transferência de renda como o Bolsa Família, há evidência de spillovers no consumo local: famílias não-beneficiárias em municípios com alta cobertura podem ter sua renda alterada por efeitos de equilíbrio no comércio local. Para uma discussão, ver Angelucci e De Giorgi (2009) sobre o caso do PROGRESA mexicano.}

A SUTVA é, portanto, uma hipótese substantiva como qualquer outra. Iremos manter SUTVA implícita no restante da seção, mas o leitor deve manter em mente que ela está sendo assumida e que a sua violação altera os parâmetros identificáveis.

% -----------------------------------------------------------------------------
\subsection{Covariadas}
\label{subsubsec:covariadas}

O modelo \eqref{eq:modelo-linear-binario} ignorou a existência de variáveis observáveis adicionais. Em aplicações reais, dispomos de um vetor $W_i$ de covariadas --- gênero, idade, escolaridade, características municipais. Há duas razões para incorporá-las. A primeira é \textit{eficiência}: condicionar em $W_i$ reduz a variância residual e estreita os intervalos de confiança. A segunda é \textit{identificação}: quando exogeneidade marginal falha, exogeneidade \textit{condicional} pode ainda ser plausível.

Postulamos um modelo com efeitos potencialmente heterogêneos por valores de $W$:
\begin{equation}
Y_i(x, w) = Y_i(0, w) + \beta(w) \cdot x.
\label{eq:modelo-cov}
\end{equation}
O efeito do tratamento depende de $w$: pode ser maior para mulheres do que para homens, maior em municípios pobres do que em ricos. A versão escalar do modelo é:
\begin{equation}
Y_i = \alpha(W_i) + \beta(W_i) \cdot D_i + \varepsilon_i,
\label{eq:modelo-cov-escalar}
\end{equation}
com a hipótese de exogeneidade adaptada para condicionalidade:
\begin{equation}
E[\varepsilon_i \mid D_i, W_i] = 0.
\label{eq:exog-condicional}
\end{equation}
A intuição é que, dentro de cada estrato definido por $W_i = w$, tratados e controles são comparáveis em $\varepsilon$. Se sortearmos quem recebe o curso PRONATEC \textit{dentro de} cada combinação de gênero, idade e escolaridade prévia, a comparação interna é válida mesmo que essas características predigam quem se inscreve.

Sob \eqref{eq:exog-condicional}, repete-se o argumento da Subseção~\ref{subsubsec:exogeneidade} \textit{dentro de cada $w$}:
\begin{equation}
\beta(w) = E[Y_i \mid D_i = 1, W_i = w] - E[Y_i \mid D_i = 0, W_i = w].
\label{eq:beta-w}
\end{equation}
O efeito médio do tratamento na população é, portanto, identificado como uma média ponderada dos efeitos condicionais:
\begin{equation}
\text{ATE} = E[\beta(W_i)] = \int \beta(w) \, dF_W(w).
\label{eq:ate-medio}
\end{equation}

Quando assumimos a forma mais forte $\beta(w) = \beta$ para todo $w$ --- isto é, efeitos \textit{constantes} no estrato ---, a identificação se simplifica drasticamente: basta uma única regressão de $Y$ em $D$ e $W$. Mas essa hipótese é restritiva. Considere o efeito do Bolsa Família sobre a frequência escolar. Estudos sugerem que o efeito é maior para meninas do que para meninos, e maior em municípios com pior infraestrutura escolar. Se rodarmos uma regressão linear simples assumindo $\beta(w) = \beta$, estamos forçando todas essas heterogeneidades a colapsarem em um único número --- e esse número é uma média ponderada que depende da composição da amostra, não do valor de política em populações específicas.

Para entender o que a regressão linear de $Y$ em $D$ controlando por $W$ identifica, sob \eqref{eq:exog-condicional} e a hipótese adicional $\beta(w) = \beta$, é útil invocar a versão populacional do teorema FWL (Frisch-Waugh-Lovell):
\begin{equation}
\beta = \frac{E\big[(D_i - E[D_i \mid W_i])(Y_i - E[Y_i \mid W_i])\big]}{E\big[(D_i - E[D_i \mid W_i])^2\big]}.
\label{eq:fwl-populacional}
\end{equation}
A leitura é a seguinte: $\beta$ é a covariância entre $D$ e $Y$, depois que ambos foram \textit{purgados} de sua componente previsível por $W$, normalizada pela variância de $D$ purgado de $W$. Em outras palavras, $\beta$ usa apenas a variação em $D$ que \textit{não} é explicada por $W$. Esse procedimento de ``parcializar'' as covariadas (\textit{partialling out}) é a base intuitiva da regressão múltipla.

Quando $\beta(w)$ varia com $w$, a fórmula \eqref{eq:fwl-populacional} ainda dá um número, mas esse número é uma \textit{média ponderada} dos $\beta(w)$ com pesos proporcionais à variância condicional de $D$ dado $W$:
\begin{equation}
\beta_{OLS} = \frac{E\big[\,\text{Var}(D_i \mid W_i) \cdot \beta(W_i)\,\big]}{E\big[\,\text{Var}(D_i \mid W_i)\,\big]}.
\label{eq:beta-ponderado}
\end{equation}
Os pesos são $\text{Var}(D_i \mid W_i) = p(W_i)(1 - p(W_i))$, onde $p(w) = P(D_i = 1 \mid W_i = w)$ é o \textit{propensity score}. Estratos com $p$ próximo de $0{,}5$ recebem mais peso; estratos quase puros (todos tratados ou todos controles) recebem peso quase nulo. Esses pesos não correspondem necessariamente à composição da população alvo da política. Esse é um ponto a que retornamos na Subseção~\ref{subsec:nao-binario}.

Vimos como a hipótese de exogeneidade pode ser adaptada para o caso com covariadas. A próxima fronteira é o tempo: o que muda quando observamos os mesmos indivíduos em mais de um período?

% =============================================================================
\section{Mais de um Período}
\label{subsec:painel}

Dados em painel --- observações repetidas dos mesmos indivíduos ao longo do tempo --- abrem possibilidades novas de identificação, mas também introduzem complicações. Para fixar ideias, considere um modelo com $T$ períodos, indivíduo $i$, sem efeitos dinâmicos do tratamento (isto é, o efeito do tratamento em $t$ depende apenas de $D_{it}$, não de $D_{i, t-1}$):
\begin{equation}
Y_{it} = \alpha_t + \beta_t \cdot D_{it} + \varepsilon_{it}.
\label{eq:modelo-painel}
\end{equation}
O parâmetro $\beta_t$ varia com $t$: o efeito do PRONATEC em 2015 pode ser diferente do efeito em 2020. Os interceptos $\alpha_t$ capturam tendências comuns na evolução de $Y$ não atribuíveis ao tratamento --- ciclo econômico, inflação, mudanças demográficas. Por simplicidade, mantemos por enquanto a ausência de covariadas individuais.

% -----------------------------------------------------------------------------
\subsection{Identificação em Painel}
\label{subsubsec:identificacao-painel}

A hipótese de exogeneidade adaptada para painel exige que o erro seja independente do tratamento dentro de cada período:
\begin{equation}
E[\varepsilon_{it} \mid D_{it}] = 0, \qquad t = 1, \dots, T.
\label{eq:exog-painel}
\end{equation}
Sob essa hipótese, $\beta_t$ é identificado período a período pela diferença de médias:
\begin{equation}
\beta_t = E[Y_{it} \mid D_{it} = 1] - E[Y_{it} \mid D_{it} = 0].
\label{eq:beta-t}
\end{equation}
A pergunta interessante é: o que acontece se, em vez de estimar $\beta_t$ período a período, agruparmos os dados e rodarmos uma única regressão?

\textbf{(A) Modelo \textit{pooled} com dummies de tempo.} Especificamos
\begin{equation}
Y_{it} = \alpha_t + \beta \cdot D_{it} + u_{it},
\label{eq:pooled}
\end{equation}
onde $\beta$ é forçado a ser constante. Aplicando FWL populacional, $\beta$ é o coeficiente da regressão de $Y_{it}$ em $D_{it}$ depois que ambos foram parcializados pelas dummies de tempo. Parcializar por dummies de tempo equivale a tirar a média de cada período: $\tilde{Y}_{it} = Y_{it} - \bar{Y}_{\cdot t}$ e $\tilde{D}_{it} = D_{it} - \bar{D}_{\cdot t}$. Logo:
\begin{equation}
\beta_{pooled} = \frac{\sum_t E\big[\,(D_{it} - \bar{D}_{\cdot t})(Y_{it} - \bar{Y}_{\cdot t})\,\big]}{\sum_t E\big[\,(D_{it} - \bar{D}_{\cdot t})^2\,\big]}
= \sum_t \omega_t^{P} \cdot \beta_t,
\label{eq:pooled-pesos}
\end{equation}
com pesos
\begin{equation}
\omega_t^{P} = \frac{p_t (1 - p_t)}{\sum_s p_s (1 - p_s)},
\label{eq:pesos-pooled}
\end{equation}
onde $p_t = P(D_{it} = 1)$ é a fração de tratados no período $t$. Os pesos somam $1$ e são maiores em períodos com $p_t$ próximo de $0{,}5$.

\textbf{(B) Modelo de efeitos fixos.} Adicionamos efeitos fixos individuais $\mu_i$:
\begin{equation}
Y_{it} = \mu_i + \alpha_t + \beta \cdot D_{it} + u_{it},
\label{eq:efeitos-fixos}
\end{equation}
e estimamos por \textit{within transformation}, isto é, subtraindo as médias \textit{por indivíduo} (e, simetricamente, por tempo). Após a dupla parcialização, obtemos $\ddot{Y}_{it}$ e $\ddot{D}_{it}$, e
\begin{equation}
\beta_{FE} = \frac{\sum_{i,t} E[\ddot{D}_{it} \ddot{Y}_{it}]}{\sum_{i,t} E[\ddot{D}_{it}^2]}
= \sum_t \omega_t^{FE} \cdot \beta_t.
\label{eq:fe-pesos}
\end{equation}
Os pesos $\omega_t^{FE}$ são proporcionais à variância \textit{within} de $D_{it}$ no período $t$: períodos em que muitos indivíduos mudam de status de tratamento contribuem mais. Se um indivíduo nunca muda de status (sempre tratado ou nunca tratado), ele contribui zero para a estimativa.

\textbf{Comparação.} Os pesos $\omega_t^P$ e $\omega_t^{FE}$ \textit{diferem}. O modelo \textit{pooled} pondera períodos pela variância \textit{cross-section} do tratamento; o modelo de efeitos fixos pondera pela variância \textit{within} do tratamento. Em desenhos com adoção escalonada (\textit{staggered adoption}) --- por exemplo, um programa que entra em municípios em ondas distintas ao longo de anos ---, essas duas ponderações dão pesos muito diferentes aos efeitos de cada coorte de adoção. A literatura recente sobre DiD com adoção escalonada (de Chaisemartin e D'Haultfœuille 2020, Goodman-Bacon 2021) mostra que $\beta_{FE}$ pode atribuir pesos \textit{negativos} a alguns $\beta_t$, o que dificulta a interpretação econômica do estimando.

A leitura é importante: \textit{nem o pooled nem o FE identificam o ATE médio simples $\frac{1}{T}\sum_t \beta_t$.} Eles identificam médias ponderadas, com pesos definidos pela mecânica do estimador. Quando os $\beta_t$ são iguais entre períodos (efeito constante), as três quantidades coincidem. Quando há heterogeneidade temporal, os números podem divergir substancialmente.

% -----------------------------------------------------------------------------
\subsection{Covariadas em Painel}
\label{subsubsec:covariadas-painel}

Voltamos agora à inclusão de covariadas $W_i$ (por simplicidade, invariantes no tempo --- o caso com $W_{it}$ variável é uma extensão direta). Há duas especificações naturais:

\textbf{(A) \textit{Pooled} sem interação $W \times D$:}
\begin{equation}
Y_{it} = \alpha_t + \beta \cdot D_{it} + \gamma \cdot W_i + u_{it}.
\label{eq:pooled-w}
\end{equation}
A pergunta é: o que essa regressão identifica? Aplicando FWL populacional, $\beta$ é o coeficiente da regressão de $Y_{it}$ em $D_{it}$ depois que ambos foram parcializados por $(\alpha_t, W_i)$. O ponto é que $W_i$ \textit{não varia com $t$}, mas o seu coeficiente $\gamma$ é forçado a ser constante --- e, portanto, parcializar por $W_i$ remove apenas uma única projeção em $W$, não uma para cada período.

Para ver o viés explicitamente, suponha que o modelo verdadeiro tem coeficiente de $W$ que varia com $t$:
\begin{equation}
Y_{it} = \alpha_t + \beta \cdot D_{it} + \gamma_t \cdot W_i + \varepsilon_{it}.
\label{eq:modelo-verdadeiro-w}
\end{equation}
A regressão \eqref{eq:pooled-w} impõe $\gamma_t = \gamma$ constante. Aplicando FWL e fazendo a álgebra (omitimos os passos detalhados, que envolvem álgebra padrão de regressão), o estimando da regressão restrita pode ser decomposto como:
\begin{equation}
\beta_{pooled}^{(A)} = \beta + \underbrace{\frac{\sum_t \text{Cov}(D_{it}, W_i) \cdot (\gamma_t - \bar{\gamma})}{\sum_t \text{Var}(\tilde{D}_{it} \mid W_i)}}_{\text{viés}},
\label{eq:vies-painel}
\end{equation}
onde $\bar{\gamma} = \frac{1}{T}\sum_t \gamma_t$ e $\tilde{D}_{it}$ denota $D_{it}$ parcializado pelas dummies de tempo. O viés tem origem clara: se $W_i$ está correlacionado com $D_{it}$ no período $t$ (por exemplo, mulheres tendem a entrar mais no programa em alguns períodos), \textit{e} o efeito de $W_i$ sobre $Y_{it}$ varia com $t$ (por exemplo, o gap salarial de gênero muda ao longo do tempo), então parcializar por uma única projeção média de $W_i$ não é suficiente. O viés captura essa diferença.

\textbf{(B) \textit{Pooled} com interação $W \times D$ e $W \times t$:}
\begin{equation}
Y_{it} = \alpha_t + \beta \cdot D_{it} + \sum_t \gamma_t \cdot (W_i \cdot \mathbb{1}\{T = t\}) + u_{it}.
\label{eq:pooled-w-int}
\end{equation}
Agora permitimos que o coeficiente de $W_i$ varie com $t$. Aplicando FWL com $D_{it}$ sendo parcializado por todas as $T$ projeções $W_i \cdot \mathbb{1}\{T=t\}$, recuperamos o caso em que cada período é tratado simetricamente. Algebricamente, a soma sobre $t$ no numerador de \eqref{eq:vies-painel} se torna $\sum_t \text{Cov}(D_{it}, W_i)(\gamma_t - \gamma_t) = 0$, e o viés desaparece:
\begin{equation}
\beta_{pooled}^{(B)} = \beta.
\label{eq:beta-pooled-b}
\end{equation}

A intuição empírica é direta. Suponha que avaliamos o efeito do PRONATEC sobre renda em painel. As covariadas $W_i$ incluem gênero. Se rodarmos a especificação (A), forçamos o efeito de gênero sobre renda a ser constante ao longo dos anos. Se, na realidade, o gap salarial de gênero diminuiu ao longo do período observado (uma característica empírica documentada do mercado de trabalho brasileiro), e se o PRONATEC atendeu mais mulheres em períodos posteriores, a regressão (A) atribui parte da redução do gap ao efeito do programa, viesando a estimativa de $\beta$. A especificação (B), ao permitir que o coeficiente de gênero varie por ano, isola corretamente o efeito do programa.

A lição geral é a mesma da Subseção~\ref{subsubsec:covariadas}: \textit{a forma funcional importa para o que o estimador calcula}. Identificação não é apenas uma questão de hipóteses sobre o erro; ela inclui hipóteses sobre como as covariadas entram no modelo. Restrições funcionais incorretas podem introduzir viés mesmo quando a hipótese de exogeneidade é válida.

Vimos como o tempo enriquece e complica a identificação. Resta examinar o que muda quando o tratamento não é binário.

% =============================================================================
\section{Modelo com Tratamento Não-Binário}
\label{subsec:nao-binario}

Até aqui, $D_i \in \{0, 1\}$. Em muitas aplicações, o tratamento é um número real --- horas de capacitação, valor de subsídio, dose de vacina, anos de escolaridade. Quando $X$ é binário, postular linearidade não é restritivo: dois pontos definem uma reta. Quando $X$ é contínuo ou tem múltiplos valores, linearidade é uma hipótese substantiva sobre a forma da relação $Y = f(X, \varepsilon)$.

Considere um modelo linear com efeitos heterogêneos:
\begin{equation}
Y_i = \alpha + \beta_i \cdot X_i + \varepsilon_i,
\label{eq:linear-heterogeneo}
\end{equation}
onde $\beta_i$ é o efeito marginal do tratamento para o indivíduo $i$ --- uma variável aleatória. Sob exogeneidade $E[\varepsilon_i \mid X_i] = 0$ \textit{e} a hipótese adicional $E[\beta_i \mid X_i] = E[\beta_i]$ (independência entre $\beta_i$ e $X_i$), a regressão linear identifica
\begin{equation}
\beta_{OLS} = E[\beta_i].
\label{eq:beta-medio}
\end{equation}
Ou seja, OLS recupera o efeito \textit{médio} dos efeitos individuais. Se a independência entre $\beta_i$ e $X_i$ falhar --- por exemplo, se quem escolhe níveis altos de tratamento é justamente quem tem efeito mais alto ---, o estimando de OLS é uma média ponderada dos $\beta_i$ com pesos que refletem a covariância entre $\beta_i$ e $X_i$. O número que sai da regressão pode não ter interpretação econômica simples.

% -----------------------------------------------------------------------------
\subsection{Modelo Estrutural Não-Linear}
\label{subsubsec:nao-linear}

Vamos um passo além. Postulamos um modelo estrutural geral:
\begin{equation}
Y_i = h(X_i, \varepsilon_i),
\label{eq:estrutural}
\end{equation}
sem assumir linearidade em $X$ nem aditividade em $\varepsilon$. O parâmetro causal de interesse depende da natureza de $X$:

\textbf{Caso $X$ contínuo.} O parâmetro relevante é a derivada parcial:
\begin{equation}
\frac{\partial h}{\partial x}(x, \varepsilon_i),
\label{eq:derivada-parcial}
\end{equation}
isto é, o efeito marginal de uma mudança em $X$ sobre $Y$, mantendo $\varepsilon_i$ fixo. Como $\varepsilon_i$ não é observado, esse objeto também varia entre indivíduos.

\textbf{Caso $X$ discreto.} O parâmetro é a diferença:
\begin{equation}
h(x+1, \varepsilon_i) - h(x, \varepsilon_i),
\label{eq:diferenca-discreta}
\end{equation}
o efeito de aumentar $X$ em uma unidade.

Em ambos os casos, ao impormos a especificação linear $Y_i = \alpha + \beta X_i + \varepsilon_i$ a um modelo estrutural não linear, OLS identifica uma média ponderada da derivada/diferença, com pesos que dependem da distribuição marginal de $X$ \textit{e} da curvatura de $h$. Esse estimando \textit{não} coincide, em geral, com a Average Derivative,
\begin{equation}
\theta_{AD} = E\left[\frac{\partial h}{\partial x}(X_i, \varepsilon_i)\right],
\label{eq:average-derivative}
\end{equation}
que é o objeto de política mais natural --- o efeito marginal médio na população.

Para enxergar a diferença, considere um exemplo simples: o efeito de horas de capacitação sobre renda, com retornos decrescentes (a função $h$ é côncava em $X$). A Average Derivative dá peso uniforme a todos os indivíduos. OLS dá mais peso àqueles cujo $X_i$ está próximo da média de $X$ \textit{e} em regiões da função onde a inclinação é maior. Em uma função côncava, isso significa peso desproporcional em níveis baixos de tratamento, onde a curva é mais íngreme. O coeficiente OLS pode, portanto, superestimar substancialmente o efeito marginal médio na população.

Imbens e Newey (2009) mostraram como, em modelos triangulares com variáveis instrumentais, é possível identificar a Average Derivative usando funções de controle. Garzon e Possebom (2025) generalizam essa identificação para o caso de tratamento exógeno com pesos amostrais explícitos, comparando os estimandos de OLS e da Average Derivative em uma classe ampla de aplicações empíricas.

Pode-se ir ainda mais longe: identificar efeitos heterogêneos condicionando em valores específicos de $X$. Sob hipóteses adicionais sobre a monotonicidade da função $h$ em $\varepsilon$ e a continuidade da distribuição condicional, é possível recuperar quantis dos efeitos individuais (Chesher 2003). A discussão técnica desses resultados foge ao escopo deste guia introdutório, mas o ponto a registrar é o seguinte: \textit{o que OLS identifica em modelos estruturais não lineares não é necessariamente o que a teoria econômica nos pediu para identificar}. A escolha entre estimadores deve refletir uma escolha consciente sobre qual média ponderada de efeitos individuais é relevante para a pergunta de política --- e essa escolha precede, logicamente, qualquer discussão sobre eficiência ou erros-padrão.

Todos os argumentos desta subseção pressupõem que $X$ é exógeno. Quando essa hipótese falha, voltamos ao problema do viés de seleção, agora com complicações adicionais devidas à não linearidade.

% =============================================================================
\section{Além do RCT}
\label{subsec:alem-rct}

A maior parte da literatura aplicada de avaliação de políticas \textit{não} dispõe de RCTs. Quando o tratamento não é atribuído por sorteio, a hipótese de exogeneidade marginal $E[\varepsilon_i \mid D_i] = 0$ não é defensável diretamente. O que se faz, então, é substituí-la por uma hipótese de identificação \textit{alternativa} --- tipicamente mais fraca em algum sentido, mas defensável a partir do desenho da pesquisa e do conhecimento institucional. Cada uma dessas hipóteses dá origem a uma \textit{estratégia de identificação} distinta.

Esta subseção apresenta, em nível introdutório, quatro das estratégias mais usadas em econometria aplicada. O objetivo aqui é mapear o terreno, não desenvolver cada método em detalhe --- cada estratégia é tratada em profundidade em um Guia Clear próprio.

% -----------------------------------------------------------------------------
\subsection{Matching}
\label{subsubsec:matching}

Quando a exogeneidade marginal falha, mas é plausível que ela valha \textit{condicionalmente} a um conjunto de observáveis, recorremos à hipótese de \textbf{independência condicional} (CIA, \textit{Conditional Independence Assumption}):
\begin{equation}
\big(Y_i(1), Y_i(0)\big) \perp D_i \mid W_i.
\label{eq:cia}
\end{equation}
A CIA diz que, \textit{dentro de} estratos definidos por $W_i$, o tratamento é como se aleatório. Quem é tratado e quem não é são comparáveis, controladas as diferenças observáveis.

A intuição do \textit{matching} é reconstruir o contrafactual de cada tratado a partir dos controles que mais se parecem com ele em $W_i$. Se um beneficiário do PRONATEC tem $35$ anos, ensino médio completo e mora em município de porte médio, busca-se nos dados um não-beneficiário com características similares e usa-se o seu $Y$ como aproximação para $Y_i(0)$. A diferença entre os $Y$ observados, agregada na população de tratados, identifica o ATT sob CIA.

O ponto delicado é que a CIA é uma hipótese substantiva. Ela exige que \textit{todas} as variáveis relevantes para a seleção e para o resultado estejam em $W_i$ --- não basta uma lista longa, é preciso que ela seja \textit{suficiente}. Diferentemente da ortogonalidade no RCT, a CIA não pode ser garantida por desenho; ela depende inteiramente de argumentos institucionais sobre quais fatores governam a entrada no programa. Para um tratamento aprofundado da estratégia, ver o Guia Clear de Matching.

% -----------------------------------------------------------------------------
\subsection{Regressão Descontínua (RDD)}
\label{subsubsec:rdd}

Em muitas políticas, a elegibilidade é determinada por uma regra administrativa baseada em uma variável contínua --- renda per capita, nota em uma prova, idade. O Bolsa Família, por exemplo, usa a renda per capita familiar para determinar quem é elegível. Quem está logo abaixo do limiar recebe; quem está logo acima, não.

A intuição da RDD (\textit{Regression Discontinuity Design} --- Desenho de Regressão Descontínua) é que indivíduos imediatamente abaixo e imediatamente acima do limiar tendem a ser \textit{essencialmente comparáveis} em todas as características relevantes --- exceto pelo tratamento. A hipótese de identificação é \textit{exogeneidade local}: na vizinhança do cutoff, a atribuição do tratamento é como se aleatória. Formalmente, exige-se continuidade das esperanças condicionais $E[Y_i(0) \mid R_i]$ e $E[Y_i(1) \mid R_i]$ no ponto de corte, onde $R_i$ é a variável de \textit{running}.

O parâmetro identificado é local: o efeito do tratamento \textit{para indivíduos no cutoff}. Esse é, simultaneamente, o ponto forte e a limitação do método. Forte, porque a hipótese de identificação é defensável por desenho institucional --- não exige argumentos sobre seleção comportamental, apenas sobre a impossibilidade de manipulação precisa do \textit{running}. Limitado, porque o parâmetro tem validade externa restrita: ele não diz nada sobre o efeito do programa em populações longe do cutoff. Para um tratamento aprofundado, ver o Guia Clear de Regressão Descontínua.

% -----------------------------------------------------------------------------
\subsection{Diferenças em Diferenças (DiD)}
\label{subsubsec:did}

Quando dispomos de dados em dois ou mais períodos e o tratamento muda de status para alguns indivíduos entre períodos, podemos substituir a hipótese de exogeneidade pela hipótese de \textbf{tendências paralelas}:
\begin{equation}
E[Y_{it}(0) - Y_{i, t-1}(0) \mid D_i = 1] = E[Y_{it}(0) - Y_{i, t-1}(0) \mid D_i = 0].
\label{eq:tendencias-paralelas}
\end{equation}
Em palavras: \textit{na ausência do tratamento}, os grupos tratado e controle teriam evoluído em paralelo. Note que isso é mais fraco do que exogeneidade marginal --- não exige que tratados e controles sejam comparáveis em nível, apenas que sua \textit{tendência} contrafactual seja a mesma.

A intuição da DiD (\textit{Difference-in-Differences} --- Diferenças em Diferenças) é que a diferença \textit{temporal} dentro do grupo de controle serve como proxy para a tendência contrafactual do grupo tratado. A diferença das diferenças --- (tratado pós - tratado pré) menos (controle pós - controle pré) --- isola o efeito do tratamento sob a hipótese \eqref{eq:tendencias-paralelas}.

Uma aplicação canônica no Brasil é o estudo do programa Mais Médicos. A entrada do programa em municípios em momentos distintos permite construir comparações DiD entre municípios que receberam o programa cedo, tarde ou nunca. A literatura recente sobre DiD com adoção escalonada (ver discussão na Subseção~\ref{subsubsec:identificacao-painel}) mostra que a aplicação ingênua de regressões com efeitos fixos pode produzir estimandos com pesos negativos sobre alguns efeitos de coorte --- um cuidado importante na aplicação. Para um tratamento aprofundado, ver o Guia Clear de Diferenças em Diferenças.

% -----------------------------------------------------------------------------
\subsection{Controle Sintético}
\label{subsubsec:controle-sintetico}

Em alguns contextos, a unidade tratada é única ou rara --- um estado que adotou uma política específica, uma cidade que recebeu um choque único. Não há, nesses casos, um grupo de controle natural comparável em nível e tendência. A alternativa proposta pela literatura de \textbf{controle sintético} (Abadie e Gardeazabal 2003, Abadie, Diamond e Hainmueller 2010) é construir um controle artificial como combinação ponderada de unidades não tratadas, de modo que a unidade sintética replique a unidade tratada nas variáveis pré-tratamento.

A intuição é que, se um estado tratado pode ser bem aproximado --- nas covariadas e nos resultados pré-tratamento --- por uma média ponderada de outros estados, então essa mesma combinação ponderada serve como contrafactual no período pós-tratamento. O parâmetro identificado é, novamente, um ATT ---- agora sobre a unidade tratada específica, não sobre uma população. Para um tratamento aprofundado, ver o Guia Clear de Controle Sintético.

\bigskip

As quatro estratégias descritas nesta subseção são respostas a diferentes formas de falha da exogeneidade marginal. Em todas elas, a estrutura lógica permanece a mesma: parte-se do parâmetro econômico de interesse, formula-se uma hipótese de identificação plausível à luz do contexto institucional e do desenho da pesquisa, deriva-se o estimando que essa hipótese permite identificar e, somente então, escolhe-se um estimador alinhado a esse estimando. Essa ordem é a espinha dorsal de toda análise causal séria. Ela será revisitada, em diferentes formas, ao longo de todos os guias da série Clear.

% =============================================================================
% Fim da Seção 2
% =============================================================================
% =============================================================================
% Seção 3 — Estimação
% Guia Clear de Introdução à Inferência Causal — FGV EESP Clear
% =============================================================================

\chapter{Estimação}
\label{cap:estimacao}
\label{sec:estimacao}

A Seção~\ref{sec:identificacao} desenvolveu a primeira metade da cadeia que conecta a teoria econômica aos dados. Partindo de um parâmetro econômico de interesse --- o ATE, o ATT, o efeito sobre uma população elegível ---, identificamos hipóteses sobre o processo gerador dos dados que estabelecem uma equivalência entre esse parâmetro e um \textit{estimando}: um funcional da distribuição populacional observável. A diferença de médias condicionais $\E[Y_i \mid D_i = 1] - \E[Y_i \mid D_i = 0]$, o coeficiente da projeção linear de $Y$ em $(D, W)$, o salto da esperança condicional no \textit{cutoff} de uma RDD, a diferença das diferenças sob tendências paralelas: todos esses são estimandos. Eles vivem no espaço populacional. Eles são números fixos, definidos pela distribuição $F(Y, D, W)$.

A questão prática, porém, é que $F$ não é conhecida. Em uma aplicação real, dispomos de uma amostra finita --- $n$ realizações $(Y_i, D_i, W_i)_{i=1}^n$. A pergunta da estimação é direta: \textit{dado que sabemos qual estimando queremos calcular, como o calculamos a partir dos dados que efetivamente temos?}

Essa passagem do estimando para o estimador acrescenta um terceiro espaço à cadeia da Seção~\ref{sec:identificacao}. Vale repetir a distinção, porque é ela que organiza tudo o que vem a seguir:

\begin{itemize}
  \item O \textbf{estimando} $\theta = \phi(F)$ é uma propriedade da distribuição populacional. É um \textit{número fixo}.
  \item O \textbf{estimador} $\hat{\theta} = \phi(\hat{F}_n)$ é uma função dos dados amostrais. É uma \textit{variável aleatória}.
  \item O objetivo da estimação é construir $\hat{\theta}$ tal que ele se aproxime de $\theta$ à medida que $n$ cresce.
\end{itemize}

A passagem entre os dois é probabilística, não determinística. Mesmo o melhor estimador pode produzir, em uma amostra particular, um número distante do estimando. A quantificação dessa distância --- a incerteza amostral --- é o tema da Seção~\ref{sec:inferencia}. Esta seção concentra-se em outra pergunta: dado o estimando, \textit{qual estimador escolher} e \textit{o que ele efetivamente calcula}?

A seção está organizada em quatro movimentos. Primeiro, formalizamos o que é uma amostra e introduzimos o princípio plug-in, que unifica a maioria dos estimadores práticos sob um único raciocínio (Subseções~\ref{subsec:amostra} e~\ref{subsec:plug-in}). Em seguida, discutimos os critérios pelos quais escolhemos entre estimadores --- consistência, viés, eficiência, robustez --- e o trade-off central entre eficiência e robustez (Subseção~\ref{subsec:propriedades}). O terceiro movimento examina um problema que conecta diretamente esta seção à anterior: identificação não garante que o estimador calcule o estimando, porque a forma funcional importa; daí seguem alternativas mais robustas como o IPW e o duplamente robusto (Subseções~\ref{subsec:forma-funcional} e~\ref{subsec:ipw}). O quarto movimento trata de duas famílias de estimadores específicas e amplamente usadas: estimadores de painel (Subseção~\ref{subsec:painel-estimacao}) e o estimador de Lin para experimentos com covariadas (Subseção~\ref{subsec:lin}).

% =============================================================================
\section{Amostra}
\label{subsec:amostra}

A amostra é o material bruto de toda a estimação. Formalmente, é uma coleção de $n$ realizações $(Y_i, D_i, W_i)_{i=1}^n$ geradas pela distribuição populacional verdadeira $F$. Em geral, $F$ é desconhecida --- caso contrário, não haveria nada a estimar. Toda a inferência empírica parte da premissa de que dispomos apenas de uma janela parcial sobre $F$.

Essa janela tem duas características que estruturam o problema da estimação.

A primeira é que ela é \textbf{finita}. A amostra contém informação limitada sobre $F$. Por mais cuidadosa que seja a coleta, há aspectos de $F$ que ficam invisíveis em qualquer amostra de tamanho fixo: caudas raras da distribuição, subgrupos pouco representados, configurações pouco frequentes de $(D, W)$. Quanto maior $n$, mais próxima a amostra está de $F$. Mas, para qualquer $n$ finito, a aproximação é imperfeita.

A segunda é que ela é \textbf{aleatória}. Cada observação $(Y_i, D_i, W_i)$ é um sorteio de $F$ --- e, portanto, qualquer função dos dados é também uma variável aleatória. Se sorteássemos uma nova amostra de tamanho $n$ da mesma população, obteríamos um valor diferente de $\hat{\theta}$. Essa variabilidade entre amostras possíveis, e não a "incerteza sobre o mundo", é a fonte da incerteza amostral. O estimando $\theta$ não muda; o estimador $\hat{\theta}$ flutua entre amostras.

Para tornar concreto o que significa "fazer com a amostra o análogo do que faríamos com a população", é útil introduzir um objeto: a distribuição empírica.

Dada uma amostra de tamanho $n$, a \textbf{distribuição empírica} $\hat{F}_n$ é a distribuição que coloca massa $1/n$ em cada observação. Ela é o análogo amostral de $F$. Pela lei dos grandes números, $\hat{F}_n \to F$ quando $n \to \infty$ --- a amostra "converge" para a distribuição verdadeira à medida que cresce. Essa convergência é a base de toda estimação consistente.

A distribuição empírica substitui o desconhecido ($F$) por algo que conhecemos completamente ($\hat{F}_n$). O estimando $\theta = \phi(F)$ é um funcional aplicado a algo que não temos; o estimador natural é o mesmo funcional aplicado àquilo que temos: $\hat{\theta} = \phi(\hat{F}_n)$. Esse princípio é a espinha dorsal da próxima subseção.

Antes de prosseguir, vale ancorar a discussão. Suponha que queiramos estimar o efeito médio do Bolsa Família sobre a frequência escolar de crianças em idade obrigatória. A distribuição populacional $F$ é a distribuição conjunta de $(Y_i, D_i, W_i)$ --- frequência escolar, recebimento do benefício, vetor de características --- em toda a população brasileira elegível. Essa distribuição existe; ela não é, em princípio, desconhecida em sentido absoluto. Mas, na prática, observamos apenas um subconjunto finito. Se a amostra vem da PNAD Contínua, são algumas dezenas de milhares de observações; se vem de dados administrativos do Cadastro Único cruzados com o Censo Escolar, podem ser milhões. Em qualquer caso, $\hat{F}_n$ é tudo o que temos para aprender sobre $F$. A boa notícia é que, para perguntas razoavelmente simples sobre $F$, $\hat{F}_n$ é uma aproximação excelente --- e melhora à medida que $n$ cresce.

% =============================================================================
\section{Estimadores Plug-in}
\label{subsec:plug-in}

Estabelecida a relação entre $F$ e $\hat{F}_n$, há um princípio de construção de estimadores que segue quase imediatamente. Ele é simples, geral e organiza a maior parte da estimação aplicada:

\begin{equation}
\text{Se } \theta = \phi(F),\text{ então } \hat{\theta} = \phi(\hat{F}_n).
\label{eq:plug-in}
\end{equation}

Em palavras: \textit{faça com a amostra exatamente o que você faria com a população}. Esse é o \textbf{princípio plug-in}. Ele é o mais natural dos princípios de estimação --- e, talvez surpreendentemente, é também o mais geral. Como veremos, OLS, diferença de médias, IV, IPW e matching são todos estimadores plug-in de funcionais distintos.

Antes da generalidade, três exemplos concretos. Em cada um, vamos identificar o funcional $\phi$, aplicá-lo a $F$ para obter o estimando, e em seguida aplicá-lo a $\hat{F}_n$ para obter o estimador.

\paragraph{Exemplo 1: a média populacional.} Considere o estimando $\theta = \E[Y]$. Em termos de $F$:
\[
\theta = \int y \, dF(y).
\]
O estimador plug-in substitui $F$ por $\hat{F}_n$. Como $\hat{F}_n$ coloca massa $1/n$ em cada $Y_i$, a integral se reduz a uma soma:
\begin{equation}
\hat{\theta} = \int y \, d\hat{F}_n(y) = \frac{1}{n} \sum_{i=1}^n Y_i = \bar{Y}.
\label{eq:plug-in-media}
\end{equation}
A média amostral é o estimador plug-in da média populacional. Esse resultado, que parece trivial, contém o que há de mais importante: o estimador é definido pela mesma operação aplicada à distribuição empírica.

\paragraph{Exemplo 2: a diferença de médias.} O estimando central da Subseção~\ref{subsubsec:exogeneidade} é $\theta = \E[Y \mid D=1] - \E[Y \mid D=0]$. As duas esperanças condicionais são funcionais de $F$. O estimador plug-in substitui cada esperança condicional pela média amostral correspondente:
\begin{equation}
\hat{\theta} = \bar{Y}_1 - \bar{Y}_0 = \frac{1}{n_1} \sum_{i: D_i = 1} Y_i - \frac{1}{n_0} \sum_{i: D_i = 0} Y_i.
\label{eq:plug-in-dif-medias}
\end{equation}
Sob exogeneidade, esse estimador calcula o ATE. A justificativa é a Equação~\eqref{eq:identificacao-rct} aplicada à amostra.

\paragraph{Exemplo 3: o coeficiente de Wald.} Em um IV (\textit{Instrumental Variables} --- Variáveis Instrumentais) com instrumento binário, o estimando é
\[
\theta_{\text{Wald}} = \frac{\E[Y \mid Z=1] - \E[Y \mid Z=0]}{\E[D \mid Z=1] - \E[D \mid Z=0]}.
\]
Ele é a razão de dois funcionais que já sabemos como estimar. O estimador plug-in é a razão das diferenças de médias amostrais:
\begin{equation}
\hat{\theta}_{\text{Wald}} = \frac{\bar{Y}_{Z=1} - \bar{Y}_{Z=0}}{\bar{D}_{Z=1} - \bar{D}_{Z=0}}.
\label{eq:plug-in-wald}
\end{equation}

\paragraph{Exemplo 4: OLS.} O estimando populacional do OLS é o coeficiente da projeção linear de $Y$ em um vetor de regressores $X_i$ (que pode incluir $D_i$, $W_i$, e uma constante):
\[
\beta = \big(\E[X_i X_i']\big)^{-1} \E[X_i Y_i].
\]
Os dois momentos --- a matriz de segundos momentos $\E[X_i X_i']$ e o vetor $\E[X_i Y_i]$ --- são funcionais de $F$. Substituindo cada um pelo análogo amostral:
\begin{equation}
\hat{\beta} = \left(\frac{1}{n} X'X\right)^{-1} \left(\frac{1}{n} X'Y\right) = (X'X)^{-1} X'Y.
\label{eq:plug-in-ols}
\end{equation}
O estimador OLS é \textit{exatamente} o plug-in do coeficiente de projeção linear populacional. Esse é um ponto importante: OLS não é um método especial criado para causalidade --- é o estimador plug-in de um funcional bem definido de $F$. O que muda quando interpretamos OLS causalmente são as hipóteses de identificação que ligam esse funcional a um parâmetro econômico (Seção~\ref{sec:identificacao}), não o estimador em si.

\paragraph{Por que o princípio é poderoso.} O plug-in unifica estimadores aparentemente distintos. Mais importante: ele \textit{herda propriedades} da convergência $\hat{F}_n \to F$. Como veremos na próxima subseção, isso dá consistência quase de graça para funcionais contínuos. Ele também torna transparente o que cada estimador está calculando: para entender $\hat{\theta}$, basta perguntar qual $\phi$ ele aplica e a qual estimando $\phi(F)$ esse $\phi$ corresponde.

\paragraph{Ancoragem empírica.} No contexto do PRONATEC, o estimando da diferença de médias salarial entre participantes e não-participantes é $\E[Y \mid D=1] - \E[Y \mid D=0]$. O estimador plug-in é simplesmente a diferença dos salários médios na nossa amostra. Sob exogeneidade --- por exemplo, num desenho experimental em que vagas escassas foram alocadas por sorteio ---, esse estimador calcula o ATE. Sem exogeneidade, o estimador continua a calcular um número bem definido, mas esse número deixa de ser o ATE: ele passa a misturar efeito causal e viés de seleção, como mostrado em~\eqref{eq:ATT-mais-vies}.

O bloco de código abaixo ilustra a construção dos quatro estimadores plug-in em uma amostra simulada.

\begin{lstlisting}[caption={Estimadores plug-in básicos.}, label={lst:plug-in}]
# -----------------------------
# 1) Limpeza e pacotes
# -----------------------------
rm(list = ls())
set.seed(2026)
library(tidyverse)

# -----------------------------
# 2) Simulação de dados (PGD com efeito = 4)
# -----------------------------
n <- 1000
W <- rnorm(n, mean = 0, sd = 1)         # covariada
D <- rbinom(n, 1, plogis(0.3 * W))      # tratamento
Y0 <- 2 + 0.8 * W + rnorm(n, sd = 1)    # potencial sob controle
Y1 <- Y0 + 4 + 0.5 * W                  # potencial sob tratamento
Y  <- D * Y1 + (1 - D) * Y0             # observado

dados <- tibble(Y = Y, D = D, W = W)

# -----------------------------
# 3) Plug-in (1): média populacional
# -----------------------------
mu_hat <- mean(dados$Y)
mu_hat

# -----------------------------
# 4) Plug-in (2): diferença de médias
# -----------------------------
delta_hat <- with(dados, mean(Y[D == 1]) - mean(Y[D == 0]))
delta_hat

# -----------------------------
# 5) Plug-in (3): OLS de Y em D
# -----------------------------
ols <- lm(Y ~ D, data = dados)
coef(ols)["D"]
\end{lstlisting}

Observe que, neste PGD (Processo Gerador de Dados), o efeito médio verdadeiro é $\E[Y_i(1) - Y_i(0)] = 4$. A diferença simples de médias amostrais aproxima esse valor porque a relação entre $D$ e $W$ é fraca (a probabilidade de tratamento varia pouco com $W$). Em contextos com seleção mais forte, a diferença de médias pode estar arbitrariamente longe de $4$ --- e nenhum aumento de $n$ resolve o problema, porque o estimando que ela calcula deixa de ser o ATE.

Vimos como o princípio plug-in unifica a construção de estimadores. Resta perguntar: que critérios usamos para julgar se um estimador é \textit{bom}?


% =============================================================================
\section{Propriedades dos Estimadores}
\label{subsec:propriedades}

Para escolher entre estimadores --- ou para entender o que estamos perdendo ao escolher um em vez de outro ---, precisamos de critérios. As três propriedades centrais são consistência, viés e eficiência. Cada uma responde a uma pergunta distinta sobre o comportamento de $\hat{\theta}$.

\subsection*{Consistência}

A primeira pergunta é a mais básica: \textit{com amostras grandes, o estimador acerta?} Formalmente, $\hat{\theta}$ é \textbf{consistente} para $\theta$ se
\begin{equation}
\hat{\theta} \xrightarrow{p} \theta \quad \text{quando} \quad n \to \infty.
\label{eq:consistencia}
\end{equation}
A notação $\xrightarrow{p}$ indica convergência em probabilidade: a probabilidade de $\hat{\theta}$ se afastar mais do que qualquer $\varepsilon > 0$ de $\theta$ tende a zero conforme $n$ cresce. Em palavras: dados suficientes, o estimador se aproxima do estimando.

A consistência é a propriedade \textit{mínima} exigível. Um estimador inconsistente não melhora com mais dados --- ele converge sistematicamente para um valor diferente do estimando. Aumentar $n$ pode reduzir a variância em torno desse valor errado, mas nunca corrige o desvio. Um intervalo de confiança estreito em torno de um estimando inconsistente é um intervalo de confiança preciso em torno do número errado.

A boa notícia é que, para estimadores plug-in de funcionais contínuos, a consistência é quase automática. Como $\hat{F}_n \to F$ pela lei dos grandes números, e como $\phi$ é contínuo, $\phi(\hat{F}_n) \to \phi(F)$. Esse argumento --- conhecido como teorema da aplicação contínua --- mostra por que tantos estimadores aplicados são consistentes sob hipóteses fracas: eles herdam a convergência da própria distribuição empírica.

\subsection*{Viés}

A consistência diz o que acontece com $n$ infinito. Em qualquer aplicação real, $n$ é finito. Para amostras finitas, importa saber se o estimador é, em média, igual ao estimando:
\begin{equation}
\text{Bias}(\hat{\theta}) = \E[\hat{\theta}] - \theta.
\label{eq:vies}
\end{equation}

Um estimador é \textbf{não viesado} se $\text{Bias}(\hat{\theta}) = 0$ para todo $n$. Em palavras: se replicássemos o experimento infinitas vezes, a média dos estimadores seria igual ao estimando.

É possível ter um estimador consistente mas viesado em amostras finitas. Exemplo: a variância amostral $\frac{1}{n} \sum (Y_i - \bar{Y})^2$ é consistente para $\Var(Y)$, mas viesada (subestima a variância populacional). A versão corrigida $\frac{1}{n-1} \sum (Y_i - \bar{Y})^2$ é não viesada \textit{e} consistente. Em estimação causal, a maioria dos estimadores semiparamétricos é consistente mas pode ter viés de amostras finitas não desprezível, especialmente com $n$ pequeno ou modelos auxiliares mal especificados.

Para fixar a distinção, considere o seguinte exemplo numérico hipotético. Suponha que $Y_i$ é a renda mensal (em milhares de reais) de um trabalhador, e que o estimando de interesse é $\theta = \E[Y]$. Geramos $200$ amostras de tamanho $n = 30$ de uma população em que $\E[Y] = 3$. Calculamos a média amostral em cada uma. A média das $200$ médias amostrais será aproximadamente $3$ --- a média amostral é não viesada. Calculamos também o estimador alternativo $\tilde{\theta} = \frac{1}{n+1} \sum Y_i$. Este é consistente (à medida que $n \to \infty$, o fator $\frac{1}{n+1}$ converge para $\frac{1}{n}$, e portanto $\tilde{\theta} \to \E[Y]$), mas viesado --- a média das $200$ replicações será aproximadamente $3 \cdot \frac{30}{31} \approx 2{,}90$. Para $n = 30$, o viés é cerca de $-0{,}10$; para $n = 1000$, é cerca de $-0{,}003$. O viés desaparece assintoticamente, mas é real para amostras finitas.

\subsection*{Eficiência}

Dado que dois estimadores são consistentes para o mesmo estimando, qual escolher? A resposta é: o que converge \textit{mais rápido}. Formalmente, o estimador \textbf{mais eficiente} é aquele com menor variância assintótica:
\[
\sqrt{n}(\hat{\theta} - \theta) \xrightarrow{d} \mathcal{N}(0, V),
\]
onde $V$ é a variância assintótica. Entre estimadores consistentes, preferimos os de $V$ menor, porque eles produzem intervalos de confiança mais estreitos com a mesma quantidade de dados --- ou, equivalentemente, atingem a mesma precisão com menos dados.

Eficiência importa quando a amostra é pequena, quando o efeito a estimar é sutil, ou quando a variância determina o poder estatístico. Em experimentos planejados, a escolha do estimador eficiente pode reduzir substancialmente o tamanho de amostra necessário para detectar um efeito de magnitude pré-definida.

\subsection*{O trade-off central: eficiência vs. robustez}

Eficiência tem um custo: ela é tipicamente obtida ao impor mais estrutura sobre o problema. Quanto mais hipóteses um estimador faz sobre $F$, mais informação ele extrai dos dados --- e mais eficiente é, \textit{quando essas hipóteses são verdadeiras}. Quando as hipóteses falham, o estimador eficiente pode ser inconsistente.

Esse é o trade-off central da estimação aplicada. Para fixar ideias, podemos organizar os estimadores em três famílias:

\begin{itemize}
  \item \textbf{Paramétricos.} Impõem uma forma funcional completa para $F$ (ex: MLE sob normalidade dos erros). Eficientes quando o modelo está correto; \textit{inconsistentes} quando está errado. Exemplo: máxima verossimilhança em um probit que assume normalidade do termo latente.

  \item \textbf{Semiparamétricos.} Impõem restrições apenas em partes da distribuição --- tipicamente, na esperança condicional ou no propensity score, mas não na distribuição completa dos erros. Robustos a algumas formas de má-especificação. Exemplo: OLS, que impõe linearidade em $\E[Y \mid X]$ mas não exige nada sobre a distribuição dos resíduos.

  \item \textbf{Não-paramétricos.} Deixam $F$ inteiramente livre e estimam funcionais por plug-in puro. Exemplo: a diferença de médias amostrais, ou o matching exato. Mais robustos, mas com convergência mais lenta --- a ``maldição da dimensionalidade'' implica que, com várias covariadas contínuas, a precisão cai rapidamente.
\end{itemize}

A regra geral é: \textit{quanto mais o estimador se apoia em hipóteses sobre $F$, mais eficiente é se essas hipóteses valem, e mais frágil é se elas falham.}

\paragraph{Ancoragem empírica.} Considere a estimação do efeito do PRONATEC sobre o salário de egressos. A distribuição de salários no Brasil é fortemente assimétrica à direita: a maioria dos trabalhadores ganha pouco, e uma minoria ganha muito. Um modelo paramétrico que assume normalidade dos erros (e, consequentemente, simetria em $Y$ condicional aos regressores) é mais eficiente \textit{se a hipótese for válida}. Mas, como a distribuição empírica de salários claramente não é normal, esse estimador pode ser inconsistente --- ele converge para o melhor ajuste linear gaussiano à distribuição verdadeira, que não é o ATE. Um estimador não-paramétrico --- por exemplo, uma diferença de médias dentro de estratos definidos por $W$ --- é mais robusto, mas exige amostras maiores para atingir a mesma precisão.

A escolha entre estimadores não é, portanto, uma decisão estatística pura. Ela é uma decisão sobre quais hipóteses estamos dispostos a defender. Em geral, é melhor um estimador robusto com erro-padrão maior do que um estimador eficiente que pode estar calculando o objeto errado.

\bigskip

Os critérios discutidos nesta subseção são gerais. A próxima é sobre um problema específico --- e particularmente importante --- que liga as propriedades do estimador às hipóteses que conduzimos da identificação para a estimação: o problema da forma funcional.

% =============================================================================
\section{Forma Funcional e Alinhamento Estimador-Estimando}
\label{subsec:forma-funcional}

A Seção~\ref{sec:identificacao} estabeleceu que, sob certas hipóteses sobre o processo gerador de dados, o parâmetro econômico de interesse é igual a um estimando bem definido. A subseção anterior introduziu critérios para escolher entre estimadores de um mesmo estimando. Falta uma pergunta essencial: \textit{o estimador escolhido efetivamente calcula o estimando que a identificação selecionou?}

A resposta nem sempre é ``sim''. Um estimador pode ser plug-in de um funcional bem definido --- e portanto calcular consistentemente \textit{aquele} funcional --- sem que esse funcional coincida com o estimando que a identificação selecionou. O caso mais frequente em econometria aplicada é o do OLS com forma funcional incorreta. Esta subseção examina esse problema em detalhe.

\subsection*{O que OLS efetivamente calcula}

OLS estima o coeficiente da \textit{projeção linear} de $Y$ em $(D, W)$. Esse é o estimando de OLS --- não o ATE, não o ATT, não nenhum efeito causal por si só. A interpretação causal de OLS depende de duas camadas de hipóteses:
\begin{enumerate}
  \item \textbf{Identificação:} alguma hipótese sobre $(Y_i(0), Y_i(1), D_i, W_i)$ que faz com que o estimando do OLS coincida com um efeito causal de interesse. Sob CIA, por exemplo, o ATT é identificado como uma média ponderada de diferenças condicionais.
  \item \textbf{Forma funcional:} a hipótese de que a esperança condicional $\E[Y \mid D, W]$ é \textit{linear} em $(D, W)$. Isso é o que faz OLS calcular o estimando que a identificação selecionou.
\end{enumerate}

A primeira camada é o tema da Seção~\ref{sec:identificacao}. A segunda --- e este é o ponto desta subseção --- é uma hipótese \textit{adicional}, que nem sempre é discutida com a mesma seriedade.

\subsection*{Quando a linearidade é restritiva e quando não é}

A força da hipótese de linearidade depende da natureza dos regressores. Três casos vale distinguir:

\paragraph{Caso 1: $D$ binário, sem covariadas.} A linearidade é, neste caso, uma identidade. Dois pontos --- a média de $Y$ em $D = 1$ e a média em $D = 0$ --- definem uma reta. OLS de $Y$ em $D$ é \textit{exatamente} a diferença de médias. Não há hipótese funcional implícita.

\paragraph{Caso 2: $D$ binário, com covariadas contínuas.} Aqui a linearidade impõe restrições reais. OLS de $Y$ em $D$ e $W$ assume que $\E[Y \mid D, W] = \alpha + \beta D + \gamma W$. Se a relação entre $W$ e $Y$ for não linear --- por exemplo, retornos decrescentes na escolaridade ---, o estimador é inconsistente para o ATE.

\paragraph{Caso 3: $D$ contínuo.} A linearidade é uma hipótese substantiva. Ela impõe que o efeito marginal de $D$ é constante: $\partial \E[Y \mid D, W] / \partial D = \beta$ para todo valor de $D$. Em modelos com retornos decrescentes ou efeitos de saturação, essa hipótese é fortemente restritiva.

\subsection*{O que acontece quando a forma funcional está errada}

Suponha que vale CIA, $(Y_i(0), Y_i(1)) \indep D_i \mid W_i$. A identificação seleciona como estimando o ATT (ou o ATE, dependendo das ponderações), construído a partir de diferenças condicionais $\E[Y \mid D=1, W=w] - \E[Y \mid D=0, W=w]$. OLS com controle linear em $W$ calcula, em vez disso, uma média ponderada dessas diferenças condicionais, com pesos que dependem da distribuição de $W$ \textit{e} do modelo linear imposto.

\citet{Angrist1998} mostrou que, com tratamento heterogêneo, o estimador OLS pode ser escrito como
\begin{equation}
\beta_{OLS} = \frac{\E\big[\,\Var(D_i \mid W_i) \cdot \tau(W_i)\,\big]}{\E\big[\,\Var(D_i \mid W_i)\,\big]},
\label{eq:angrist-pesos}
\end{equation}
onde $\tau(w) = \E[Y_i(1) - Y_i(0) \mid W_i = w]$ é o efeito condicional. Os pesos $\Var(D_i \mid W_i) = p(W_i)(1 - p(W_i))$ são proporcionais à variância do tratamento dentro de cada estrato. Estratos com $p(w)$ próximo de $0{,}5$ recebem mais peso; estratos quase puros (todos tratados ou todos controles) recebem peso quase nulo.

Esses pesos são \textit{mecânicos} --- eles vêm da álgebra do estimador, não de qualquer escolha consciente do pesquisador. Eles não correspondem necessariamente aos pesos da população de interesse para a política. Pior: em desenhos com adoção escalonada (Subseção~\ref{subsec:painel-estimacao}), pesos negativos podem aparecer, atribuindo peso negativo a alguns efeitos condicionais.

\subsection*{Um exemplo numérico do desalinhamento}

Para tornar concreto o que significa ``OLS calcular algo diferente do ATE'', considere uma população dividida em dois estratos definidos por $W \in \{0, 1\}$, conforme a Tabela~\ref{tab:pesos-ols}.

\begin{table}[h]
\centering
\caption{População hipotética: dois estratos com diferentes propensões e efeitos.}
\label{tab:pesos-ols}
\begin{tabular}{lccc}
\toprule
Estrato & Tamanho & $p(w)$ & $\tau(w)$ \\
\midrule
$W = 0$ & $50\%$ & $0{,}5$ & $2$ \\
$W = 1$ & $50\%$ & $0{,}1$ & $8$ \\
\bottomrule
\end{tabular}
\end{table}

O ATE verdadeiro é $\E[\tau(W)] = 0{,}5 \cdot 2 + 0{,}5 \cdot 8 = 5$. Mas os pesos do OLS são proporcionais a $p(w)(1-p(w))$:
\[
p(0)(1-p(0)) = 0{,}25, \qquad p(1)(1-p(1)) = 0{,}09.
\]
Normalizando, o estrato $W=0$ recebe peso $0{,}25/(0{,}25+0{,}09) \approx 0{,}735$, e o estrato $W=1$ recebe peso $\approx 0{,}265$. O estimando do OLS é, portanto:
\[
\beta_{OLS} \approx 0{,}735 \cdot 2 + 0{,}265 \cdot 8 \approx 3{,}59.
\]
A diferença é substantiva: o OLS estima $3{,}59$, enquanto o ATE é $5$. O viés vem do fato de que OLS subponderou o estrato com efeito grande --- e fez isso porque, naquele estrato, há pouca variação no tratamento (pouco se ``aprende'' sobre $\tau$ a partir dele, na lógica do estimador de variância mínima).

O bloco abaixo replica esse cálculo em uma simulação.

\begin{lstlisting}[caption={OLS vs. ATE quando os pesos do estimador não coincidem com os de política.}, label={lst:angrist-pesos}]
# -----------------------------
# 1) Configuração
# -----------------------------
rm(list = ls())
set.seed(2026)
library(tidyverse)

n <- 50000
W <- rbinom(n, 1, 0.5)                             # estrato (50/50)
p <- ifelse(W == 1, 0.1, 0.5)                      # propensity score
D <- rbinom(n, 1, p)
tau <- ifelse(W == 1, 8, 2)                        # efeito heterogêneo
Y0 <- rnorm(n)
Y1 <- Y0 + tau
Y  <- D * Y1 + (1 - D) * Y0
dados <- tibble(Y = Y, D = D, W = W)

# -----------------------------
# 2) ATE verdadeiro (na população)
# -----------------------------
ate_verdadeiro <- mean(tau)
ate_verdadeiro    # ~ 5

# -----------------------------
# 3) Estimando do OLS com controle linear
# -----------------------------
ols <- lm(Y ~ D + W, data = dados)
coef(ols)["D"]    # ~ 3.59 — não é o ATE

# -----------------------------
# 4) Estimando alternativo: média de diferenças por estrato
# -----------------------------
ate_estratos <- dados %>%
  group_by(W) %>%
  summarise(diff = mean(Y[D == 1]) - mean(Y[D == 0])) %>%
  summarise(ate = mean(diff)) %>%
  pull(ate)
ate_estratos     # ~ 5
\end{lstlisting}

A lição é geral: \textit{identificação não é suficiente para que o estimador acerte o estimando}. Mesmo com CIA válida, OLS pode estimar uma média ponderada que não coincide com o ATE. O leitor pode verificar que, no caso paradigmático em que $\tau(w) = \tau$ é constante ou em que $p(w)$ é constante (como em um RCT bem desenhado), o estimando do OLS coincide com o ATE --- a dificuldade emerge da \textit{combinação} entre heterogeneidade do tratamento e heterogeneidade da propensão.

\subsection*{Ancoragem empírica}

Considere o efeito de anos de escolaridade sobre salário. Postular linearidade --- $\E[\log W \mid \text{anos}] = \alpha + \beta \cdot \text{anos}$ --- impõe que cada ano adicional tem o mesmo efeito percentual, independentemente do nível. Isso pode ser razoável como descrição da média populacional, mas esconde heterogeneidades importantes. O retorno do ensino médio completo (12 anos) versus incompleto (11 anos) costuma ser muito maior do que o retorno de um ano adicional na pós-graduação. Especificações com termos quadráticos, splines, ou efeitos por nível de escolaridade --- em vez de uma única inclinação linear --- conseguem capturar essa heterogeneidade.

A questão não é abandonar OLS. É reconhecer que OLS calcula um funcional bem definido, e que a interpretação causal desse funcional depende não apenas das hipóteses de identificação, mas também da plausibilidade da forma funcional. Quando a linearidade é restritiva, há duas alternativas: (i) flexibilizar o modelo (interações, polinômios, splines) ou (ii) abandonar a abordagem baseada na esperança condicional e usar estimadores que evitam essa hipótese. A próxima subseção explora a segunda rota.


% =============================================================================
\section{IPW: Ponderação pelo Inverso da Probabilidade}
\label{subsec:ipw}

A subseção anterior identificou um problema: sob CIA, o ATE é identificado, mas OLS com controles lineares pode calcular uma média ponderada distinta do ATE. O problema vem de modelar a relação entre $W$ e $Y$ por uma forma funcional que pode estar errada. O IPW (\textit{Inverse Probability Weighting} --- Ponderação pelo Inverso da Probabilidade) propõe uma rota alternativa: em vez de modelar $\E[Y \mid D, W]$, modela-se a probabilidade de tratamento $P(D = 1 \mid W) = p(W)$ --- o \textit{propensity score} (escore de propensão).

\subsection*{Intuição}

Sob CIA, dentro de cada estrato definido por $W$, tratados e controles são comparáveis. Mas a distribuição de $W$ entre tratados \textit{difere} da distribuição entre controles --- justamente porque tratamento e $W$ são correlacionados. A diferença bruta de médias mistura, portanto, dois efeitos: o efeito do tratamento e a diferença composicional em $W$.

A ideia do IPW é corrigir essa diferença composicional pela reponderação da amostra. Indivíduos com baixa probabilidade de tratamento que \textit{foram} tratados são ``raros'' --- e, se queremos reconstruir a distribuição populacional de $W$ entre os tratados \textit{como se a alocação fosse aleatória}, esses indivíduos devem receber peso alto. Simetricamente, indivíduos com alta probabilidade de tratamento que \textit{não} foram tratados também são raros no grupo de controle --- e também devem receber peso alto.

O resultado da reponderação é uma ``pseudo-população'' em que tratados e controles têm a mesma distribuição de $W$. Nessa pseudo-população, a diferença bruta de médias identifica o ATE.

IPW é uma operação de reponderação. Em vez de modelar como $Y$ depende de $W$ (rota OLS), o IPW modela como $D$ depende de $W$ (rota propensity score) e usa esse modelo para ``consertar'' a composição do grupo tratado e do grupo controle. O ganho é robustez à forma funcional de $\E[Y \mid D, W]$. O custo é dependência da forma funcional do propensity score.

\subsection*{O estimando de Horvitz-Thompson e o estimador IPW}

A formalização passa pelo estimando populacional de Horvitz-Thompson:
\begin{equation}
\theta_{HT} = \E\!\left[\frac{D_i Y_i}{p(W_i)} - \frac{(1 - D_i) Y_i}{1 - p(W_i)}\right].
\label{eq:horvitz-thompson}
\end{equation}

A álgebra para mostrar que $\theta_{HT} = \E[Y_i(1) - Y_i(0)] = \text{ATE}$ sob CIA e \textit{overlap} ($0 < p(W) < 1$ com probabilidade um) usa apenas a definição de resultados potenciais e a lei das esperanças iteradas. Considere o primeiro termo:
\begin{align*}
\E\!\left[\frac{D_i Y_i}{p(W_i)}\right]
  &= \E\!\left[\frac{D_i Y_i(1)}{p(W_i)}\right] \\
  &= \E\!\left[\E\!\left[\frac{D_i Y_i(1)}{p(W_i)} \,\Big|\, W_i\right]\right] \\
  &= \E\!\left[\frac{\E[D_i \mid W_i] \cdot \E[Y_i(1) \mid W_i]}{p(W_i)}\right] \\
  &= \E[\E[Y_i(1) \mid W_i]] = \E[Y_i(1)],
\end{align*}
onde a terceira igualdade usa a CIA (que dá $D_i \indep Y_i(1) \mid W_i$) e a quarta usa $\E[D_i \mid W_i] = p(W_i)$. O argumento simétrico para o segundo termo dá $\E[Y_i(0)]$. A diferença é o ATE.

O estimador IPW é o plug-in de $\theta_{HT}$. Como $p(W_i)$ não é observado, ele é substituído por uma estimativa $\hat{p}(W_i)$ --- tipicamente, um logit de $D$ em $W$:
\begin{equation}
\hat{\tau}_{IPW}^{ATE} = \frac{1}{n} \sum_{i=1}^n \!\left[\frac{D_i Y_i}{\hat{p}(W_i)} - \frac{(1 - D_i) Y_i}{1 - \hat{p}(W_i)}\right].
\label{eq:ipw-ate}
\end{equation}

Para o ATT, a derivação análoga produz
\begin{equation}
\hat{\tau}_{IPW}^{ATT} = \frac{1}{N_1} \sum_{i: D_i = 1} Y_i - \frac{1}{N_1} \sum_{i: D_i = 0} \frac{\hat{p}(W_i)}{1 - \hat{p}(W_i)} \cdot Y_i,
\label{eq:ipw-att}
\end{equation}
onde $N_1$ é o número de tratados. A intuição é a mesma: o grupo controle é reponderado para se parecer com o grupo tratado em $W$.

\subsection*{Vantagens e limitações}

\paragraph{Vantagens em relação ao OLS.}
\begin{itemize}
  \item \textbf{Robusto à forma funcional de $\E[Y \mid D, W]$.} IPW não impõe linearidade no resultado --- ele opera diretamente sobre as observações via reponderação.
  \item \textbf{Transparência.} Os pesos $\hat{p}(W_i)^{-1}$ e $(1 - \hat{p}(W_i))^{-1}$ tornam explícito quais observações estão ``carregando'' o estimador. Isso facilita diagnósticos.
\end{itemize}

\paragraph{Limitações.}
\begin{itemize}
  \item \textbf{Sensibilidade à má-especificação do propensity score.} Se o modelo para $p(W)$ está errado, o estimador é inconsistente. Trocamos uma hipótese funcional (sobre $Y$) por outra (sobre $D$).
  \item \textbf{Pesos extremos.} Quando $\hat{p}(W_i) \approx 0$ ou $\hat{p}(W_i) \approx 1$, os pesos explodem e a variância do estimador se torna enorme. Em amostras finitas, uma única observação com peso muito alto pode dominar a estimativa.
\end{itemize}

\paragraph{Trimming.} A solução prática para o problema dos pesos extremos é o \textit{trimming}: descartar observações com $\hat{p}(W_i)$ fora de um intervalo (tipicamente $[0{,}05; 0{,}95]$). Isso restringe a estimação a uma sub-população em que há \textit{overlap} adequado --- isto é, em que tanto tratados quanto controles têm representatividade não-trivial. O custo é redefinir a população de inferência: o estimando deixa de ser o ATE em toda a população e passa a ser o ATE na região de \textit{overlap}.

\subsection*{Doubly Robust}

Há uma terceira rota que combina o melhor das duas: estimadores \textit{doubly robust} (duplamente robustos), como o AIPW (\textit{Augmented IPW}):
\begin{equation}
\hat{\tau}_{AIPW} = \frac{1}{n} \sum_i \!\left[ \hat{m}_1(W_i) - \hat{m}_0(W_i) + \frac{D_i (Y_i - \hat{m}_1(W_i))}{\hat{p}(W_i)} - \frac{(1 - D_i)(Y_i - \hat{m}_0(W_i))}{1 - \hat{p}(W_i)} \right],
\label{eq:aipw}
\end{equation}
onde $\hat{m}_d(W_i)$ é uma estimativa de $\E[Y \mid D = d, W]$. A propriedade central do AIPW é que ele é consistente para o ATE se \textit{pelo menos um} dos dois modelos --- o de resultado ($m_d$) ou o de propensão ($p$) --- estiver correto. Erros em um modelo são corrigidos pelo outro. Por isso o nome ``duplamente robusto''.

Estimadores duplamente robustos são uma direção natural quando há incerteza sobre qual modelo auxiliar é mais confiável. O custo é maior complexidade computacional e a necessidade de ajustar dois modelos em vez de um. A discussão técnica detalhada está fora do escopo deste guia introdutório, mas a referência é o artigo seminal de \citet{robins1994}.

\subsection*{Implementação}

\begin{lstlisting}[caption={Estimação por IPW e comparação com OLS.}, label={lst:ipw}]
# -----------------------------
# 1) Configuração
# -----------------------------
rm(list = ls())
set.seed(2026)
library(tidyverse)

n <- 5000
W <- rnorm(n)
p_real <- plogis(0.8 * W)                     # propensity score verdadeiro
D <- rbinom(n, 1, p_real)
Y0 <- exp(W) + rnorm(n)                       # E[Y|W] não-linear (exponencial)
Y1 <- Y0 + 4
Y  <- D * Y1 + (1 - D) * Y0
dados <- tibble(Y = Y, D = D, W = W)

# -----------------------------
# 2) OLS com controle linear (forma funcional errada)
# -----------------------------
ols <- lm(Y ~ D + W, data = dados)
coef(ols)["D"]    # provavelmente longe de 4

# -----------------------------
# 3) IPW: estimar propensity score via logit
# -----------------------------
modelo_ps <- glm(D ~ W, data = dados, family = binomial)
dados$ps <- predict(modelo_ps, type = "response")

# -----------------------------
# 4) Trimming: descartar pesos extremos
# -----------------------------
dados_trim <- dados %>% filter(ps > 0.05, ps < 0.95)

# -----------------------------
# 5) Estimador IPW para o ATE
# -----------------------------
ipw_ate <- with(dados_trim,
  mean(D * Y / ps - (1 - D) * Y / (1 - ps))
)
ipw_ate           # mais próximo de 4
\end{lstlisting}

\paragraph{Conexão com a Seção~\ref{sec:identificacao}.} IPW é o estimador mais diretamente motivado pela CIA. Se a identificação via CIA é válida, o IPW é uma alternativa natural ao OLS --- mais robusto à forma funcional do resultado, mais transparente sobre quais observações estão identificando o efeito. Sua escolha entre OLS, IPW e AIPW depende do que é mais fácil de defender no contexto: a forma funcional de $\E[Y \mid D, W]$, a forma funcional de $p(W)$, ou nenhuma das duas com certeza --- caso em que a robustez dupla do AIPW é particularmente atraente.

% =============================================================================
\section{Estimação em Painel}
\label{subsec:painel-estimacao}

Os estimadores discutidos até aqui assumem que as observações são independentes. Em muitas aplicações, dispomos de algo mais rico: dados em \textbf{painel}, isto é, observações repetidas das mesmas unidades $i$ em múltiplos períodos $t$. A estrutura adicional permite estimadores que controlam por características fixas não observadas das unidades --- em troca de hipóteses adicionais sobre como o tempo entra no modelo.

Considere o modelo geral
\begin{equation}
Y_{it} = \alpha_i + \lambda_t + \beta D_{it} + \varepsilon_{it},
\label{eq:painel-geral}
\end{equation}
onde $\alpha_i$ são \textit{efeitos fixos de unidade} (capturam características invariantes no tempo da unidade $i$ --- história, geografia, cultura) e $\lambda_t$ são \textit{efeitos fixos de tempo} (capturam choques agregados comuns a todas as unidades em $t$ --- ciclo econômico, mudanças nacionais de política, eventos macro).

A pergunta de estimação é: como recuperamos $\beta$ a partir dos dados, dado que $\alpha_i$ e $\lambda_t$ não são observados? Há duas estratégias principais: o estimador TWFE e o estimador FD.

\subsection{TWFE: Two-Way Fixed Effects}
\label{subsubsec:twfe}

O estimador TWFE (\textit{Two-Way Fixed Effects} --- Efeitos Fixos de Dois Sentidos) é o OLS aplicado ao modelo~\eqref{eq:painel-geral}, tratando $\alpha_i$ e $\lambda_t$ como parâmetros a estimar. Há duas implementações algebricamente equivalentes.

\paragraph{Implementação A: regressão com dummies.} Inclui-se uma dummy para cada unidade e cada período no modelo:
\[
Y_{it} = \sum_i \alpha_i \mathds{1}\{i\} + \sum_t \lambda_t \mathds{1}\{t\} + \beta D_{it} + \varepsilon_{it},
\]
e estima-se por OLS. Com $N$ unidades e $T$ períodos, isso significa $N + T - 1$ dummies (omite-se uma para evitar colinearidade). Em painéis grandes, a abordagem é computacionalmente custosa, embora pacotes modernos (\texttt{fixest} em R) implementem a fatoração de forma eficiente.

\paragraph{Implementação B: \textit{within transformation}.} Define-se a transformação \textit{within} (de duplo desvio):
\begin{equation}
\ddot{Y}_{it} = Y_{it} - \bar{Y}_{i\cdot} - \bar{Y}_{\cdot t} + \bar{Y}_{\cdot\cdot},
\label{eq:within}
\end{equation}
onde $\bar{Y}_{i\cdot} = \frac{1}{T} \sum_t Y_{it}$ é a média da unidade $i$ ao longo do tempo, $\bar{Y}_{\cdot t} = \frac{1}{N} \sum_i Y_{it}$ é a média do período $t$ entre unidades, e $\bar{Y}_{\cdot\cdot}$ é a grande média. Aplica-se a mesma transformação a $D_{it}$, obtendo $\ddot{D}_{it}$. O estimador TWFE é então o OLS de $\ddot{Y}_{it}$ em $\ddot{D}_{it}$:
\begin{equation}
\hat{\beta}_{TWFE} = \frac{\sum_{i,t} \ddot{D}_{it} \ddot{Y}_{it}}{\sum_{i,t} \ddot{D}_{it}^2}.
\label{eq:twfe}
\end{equation}

A equivalência entre as duas implementações é uma consequência do teorema FWL (\textit{Frisch-Waugh-Lovell}): o coeficiente de $D_{it}$ na regressão completa é idêntico ao coeficiente de $\ddot{D}_{it}$ na regressão simples sobre os resíduos das dummies. A transformação within ``limpa'' $D_{it}$ de toda a variação explicada por $\alpha_i$ e $\lambda_t$, deixando apenas a variação que sobra após controlar por unidade e tempo.

\subsection*{Os pesos do TWFE}

A Equação~\eqref{eq:twfe} torna explícito o que o TWFE calcula: ele é o OLS aplicado à variação \textit{within} de $D_{it}$. Em particular, unidades que nunca mudam de status de tratamento (sempre tratadas ou nunca tratadas) contribuem zero para a estimativa, porque $\ddot{D}_{it} = 0$ para elas. A identificação vem inteiramente das unidades que \textit{mudam} de status.

Quando $\beta$ é constante (efeito homogêneo), o TWFE estima esse $\beta$ comum. Quando $\beta$ varia entre unidades ou entre períodos --- por exemplo, $\beta_{it}$ ---, o TWFE estima uma média ponderada com pesos que dependem da variância \textit{within} de $D_{it}$ em cada combinação $(i, t)$. Essa ponderação é mecânica: ela vem da álgebra do OLS, não de uma escolha do pesquisador.

A literatura recente sobre DiD mostra que, em desenhos com adoção escalonada (unidades entrando no tratamento em períodos diferentes), o TWFE pode atribuir \textit{pesos negativos} a alguns efeitos $\beta_{it}$. Isso ocorre porque unidades já tratadas funcionam, em períodos posteriores, como ``controles'' para unidades recém-tratadas --- e a comparação resultante pode produzir contribuições negativas. \citet{dechaisemartin2020} e \citet{goodmanbacon2021} fornecem decomposições explícitas. O Guia Clear de Diferenças em Diferenças trata o problema em profundidade; aqui, o ponto é registrar que o estimando do TWFE pode não ser o que se espera intuitivamente.

\paragraph{Ancoragem empírica.} Considere a estimação do efeito do programa Mais Médicos sobre indicadores de saúde municipal usando dados de painel da Atenção Primária. O TWFE controla por características fixas dos municípios ($\alpha_i$: estrutura demográfica histórica, capacidade administrativa, densidade populacional pré-existente) e por choques anuais comuns ($\lambda_t$: mudanças no SUS, ciclo econômico nacional, pandemias). A identificação vem da comparação entre municípios que receberam o programa em momentos distintos --- e essa comparação carrega as armadilhas dos pesos do TWFE quando o efeito é heterogêneo entre coortes de adoção.

\subsection{FD: First Differences}
\label{subsubsec:fd}

Uma alternativa ao TWFE é o estimador FD (\textit{First Differences} --- Primeiras Diferenças). Em vez de remover $\alpha_i$ pela transformação within, removemos pela diferenciação temporal. Subtraindo $Y_{i, t-1}$ de $Y_{it}$ no modelo~\eqref{eq:painel-geral}:
\begin{equation}
\Delta Y_{it} = \Delta \lambda_t + \beta \Delta D_{it} + \Delta \varepsilon_{it},
\label{eq:fd}
\end{equation}
onde $\Delta Y_{it} = Y_{it} - Y_{i, t-1}$. O efeito fixo $\alpha_i$ desaparece por construção: $\alpha_i - \alpha_i = 0$. O estimador FD é o OLS aplicado à equação~\eqref{eq:fd}.

Tanto o TWFE quanto o FD eliminam $\alpha_i$, mas usam variação \textit{diferente} dos dados para identificar $\beta$. O TWFE usa toda a variação within de $D_{it}$ ao longo dos $T$ períodos. O FD usa apenas a variação \textit{entre períodos adjacentes}.

\subsection*{Comparação: TWFE vs. FD}

\paragraph{Caso $T = 2$: equivalência.} Com apenas dois períodos, TWFE e FD são algebricamente idênticos. A demonstração é direta: com $T = 2$, $\bar{Y}_{i\cdot} = (Y_{i1} + Y_{i2})/2$, e a variação within é simplesmente $\pm \Delta Y_{it}/2$. O OLS within é o mesmo que o OLS na primeira diferença. Esse caso recupera o estimador DiD padrão.

\paragraph{Caso $T > 2$: divergência.} Com mais períodos, os dois estimadores diferem porque usam ponderações diferentes da variação temporal. A escolha entre eles depende da estrutura de correlação serial dos erros:
\begin{itemize}
  \item Se $\varepsilon_{it}$ é i.i.d. (sem correlação serial), o TWFE é mais eficiente. A transformação within não introduz correlação espúria, e usa toda a variação within.
  \item Se $\varepsilon_{it}$ segue um \textit{random walk} (passeio aleatório), $\Delta \varepsilon_{it}$ é i.i.d. --- e, portanto, o FD é mais eficiente. A diferenciação ``estabiliza'' os erros.
\end{itemize}

Essa diferença pode ser explorada como diagnóstico. Se TWFE e FD produzem estimativas substantivamente diferentes em uma aplicação, vale investigar duas possibilidades: (i) o modelo está mal especificado (forma funcional errada, tendências paralelas violadas, efeitos heterogêneos no tempo); ou (ii) a estrutura de correlação serial dos erros é tal que um dos dois estimadores é mais apropriado. A divergência não é, por si, evidência de problema --- mas é um sinal que merece investigação.

\begin{lstlisting}[caption={Estimadores TWFE e FD em painel.}, label={lst:painel}]
# -----------------------------
# 1) Configuração
# -----------------------------
rm(list = ls())
set.seed(2026)
library(tidyverse)
library(fixest)

# -----------------------------
# 2) Simulação de painel (N = 200, T = 5)
# -----------------------------
N <- 200; T <- 5
painel <- expand_grid(i = 1:N, t = 1:T) %>%
  mutate(alpha_i = rep(rnorm(N, sd = 1), each = T),
         lambda_t = rep(0.2 * (1:T), times = N),
         D_it = ifelse(t >= 3 & i <= 100, 1, 0),
         eps  = rnorm(N * T, sd = 1),
         Y    = alpha_i + lambda_t + 3 * D_it + eps)  # efeito = 3

# -----------------------------
# 3) TWFE via fixest
# -----------------------------
twfe <- feols(Y ~ D_it | i + t, data = painel)
coef(twfe)         # ~ 3

# -----------------------------
# 4) FD: primeira diferença
# -----------------------------
painel_fd <- painel %>%
  arrange(i, t) %>%
  group_by(i) %>%
  mutate(dY = Y - lag(Y),
         dD = D_it - lag(D_it),
         dt = t) %>%
  ungroup() %>%
  filter(!is.na(dY))

fd <- feols(dY ~ dD | dt, data = painel_fd)
coef(fd)           # ~ 3
\end{lstlisting}

Em painéis com adoção homogênea e efeito constante --- como neste exemplo --- TWFE e FD entregam estimativas similares. Em desenhos escalonados com efeitos heterogêneos, a divergência pode ser substancial e indicativa de problemas no estimando. O Guia Clear de DiD é a referência para esse diagnóstico.


% =============================================================================
\section{Estimador de Lin}
\label{subsec:lin}

A discussão até aqui abordou três contextos: estimação por plug-in em geral, alternativas ao OLS quando a forma funcional importa (IPW), e estimação em painel. A última subseção desta seção trata de um problema mais específico, mas extremamente comum na avaliação experimental: em um RCT, como incorporar covariadas pré-tratamento para ganhar precisão sem introduzir viés?

\subsection*{O problema}

Em um experimento aleatorizado, a hipótese de exogeneidade $E[\varepsilon_i \mid D_i] = 0$ é garantida por desenho. Em particular, a regressão simples $Y_i = \alpha + \tau D_i + \varepsilon_i$ identifica o ATE sem qualquer hipótese funcional adicional --- a aleatorização cuida de tudo. O estimador OLS sem covariadas é não viesado para o ATE em amostras finitas (sob aleatorização) e consistente.

Mas há, tipicamente, covariadas $W_i$ medidas \textit{antes} do tratamento --- características demográficas, valores pré-experimentais do resultado, indicadores de qualidade institucional. Incluir $W_i$ na regressão pode reduzir a variância residual e estreitar os intervalos de confiança, aumentando o poder estatístico. A pergunta é: \textit{qual é a forma certa de incluir $W$?}

A resposta intuitiva é a regressão linear:
\begin{equation}
Y_i = \alpha + \tau D_i + \gamma' W_i + \varepsilon_i.
\label{eq:ols-com-w}
\end{equation}
Essa especificação tem um problema sutil. Ela impõe que o coeficiente de $W_i$ é o \textit{mesmo} para tratados e controles. Em amostras finitas, isso pode introduzir um pequeno viés se a relação verdadeira entre $W$ e $Y$ diferir entre os dois grupos (ainda que a aleatorização garanta que $W$ tenha a mesma distribuição entre tratados e controles em expectativa). \citet{freedman2008} demonstrou esse viés analiticamente e mostrou que, em alguns casos, a inclusão ``ingênua'' de covariadas pode \textit{piorar} a precisão em vez de melhorá-la.

\subsection*{A solução de Lin}

\citet{Lin2013} propôs uma correção simples: incluir as covariadas \textit{centradas} (subtraídas da média amostral) e também a interação dessas covariadas centradas com o tratamento:
\begin{equation}
Y_i = \alpha + \tau D_i + \gamma' (W_i - \bar{W}) + \delta' D_i (W_i - \bar{W}) + \varepsilon_i.
\label{eq:lin}
\end{equation}

O coeficiente $\tau$ desta regressão é o estimador de Lin do ATE. A inclusão da interação $D_i \cdot (W_i - \bar{W})$ permite que o coeficiente das covariadas seja \textit{diferente} para tratados e controles. A centragem de $W$ na média amostral garante que o coeficiente $\tau$ tenha interpretação direta como o efeito médio do tratamento.

\paragraph{Equivalência com OLS por grupo.} O estimador de Lin é algebricamente equivalente a estimar separadamente um modelo linear de $Y$ em $W$ no grupo tratado e no grupo controle, e em seguida tomar a diferença das predições na média da amostra. Essa equivalência é uma aplicação direta do FWL: a interação faz com que cada grupo ``decida'' sua própria relação entre $W$ e $Y$, sem ser forçado a uma inclinação comum.

\subsection*{Propriedades}

O estimador de Lin tem três propriedades que o tornam atraente em RCTs:

\begin{enumerate}
  \item \textbf{Não viesado em amostras finitas} (sob aleatorização). A centragem de $W$ na média amostral garante que o viés introduzido pela inclusão de covariadas em amostras finitas seja eliminado.

  \item \textbf{Assintoticamente ao menos tão eficiente quanto OLS sem covariadas.} Lin mostrou que a variância assintótica de $\hat{\tau}_{Lin}$ é menor ou igual à de $\hat{\tau}_{OLS\text{-sem-}W}$, com igualdade quando $W$ é não-correlacionado com $Y$.

  \item \textbf{Assintoticamente ao menos tão eficiente quanto OLS com covariadas sem interação.} Quando a relação entre $W$ e $Y$ difere entre tratados e controles, Lin é estritamente mais eficiente. Quando não difere, as duas variâncias coincidem.
\end{enumerate}

A combinação dessas propriedades é poderosa: \textit{em RCTs, o estimador de Lin nunca piora em relação ao OLS simples, e tipicamente melhora}. O custo computacional adicional é trivial.

\subsection*{Quando usar Lin vs. OLS simples}

A recomendação prática é direta:
\begin{itemize}
  \item \textbf{Sempre que houver covariadas pré-tratamento medidas em um RCT}, considere Lin como padrão.
  \item O ganho de precisão é maior quando as covariadas têm \textit{alta correlação} com $Y$. Se as covariadas explicam pouco da variação no resultado, o ganho é marginal --- mas o custo continua sendo zero.
  \item Em desenhos quase-experimentais (DiD, RDD, IV), a lógica de Lin não se aplica diretamente, porque a aleatorização não garante a propriedade de não-viesamento em amostras finitas. Lin é uma ferramenta \textit{para experimentos}.
\end{itemize}

\paragraph{Ancoragem empírica.} Suponha que avaliamos uma intervenção experimental na rede pública de educação --- um piloto de aulas suplementares de matemática alocado por sorteio entre $300$ escolas. O resultado $Y$ é a nota média da escola na Prova Brasil. As covariadas pré-tratamento $W$ incluem o NSE (Nível Socioeconômico) médio dos alunos, a infraestrutura da escola (índice INEP), e a nota da escola na Prova Brasil do ciclo anterior. Essa última, em particular, costuma ter altíssima correlação com a nota futura --- tipicamente $\rho > 0{,}8$. Incorporando-a via Lin (em vez de OLS sem covariadas), o erro-padrão do efeito do programa pode cair em ordens substanciais, sem qualquer ameaça à validade causal.

\begin{lstlisting}[caption={Estimador de Lin em um RCT.}, label={lst:lin}]
# -----------------------------
# 1) Configuração
# -----------------------------
rm(list = ls())
set.seed(2026)
library(tidyverse)
library(estimatr)

# -----------------------------
# 2) Simulação de RCT (n = 500, efeito = 5)
# -----------------------------
n <- 500
W <- rnorm(n, mean = 50, sd = 10)
D <- rbinom(n, 1, 0.5)                          # aleatorização
Y0 <- 0.7 * W + rnorm(n, sd = 5)                # alta correlação entre W e Y
Y1 <- Y0 + 5 + 0.3 * (W - mean(W))              # efeito heterogêneo em W
Y  <- D * Y1 + (1 - D) * Y0
dados <- tibble(Y = Y, D = D, W = W)

# -----------------------------
# 3) OLS sem covariadas
# -----------------------------
ols_sem_w <- lm_robust(Y ~ D, data = dados)
coef(ols_sem_w)["D"]
ols_sem_w$std.error[2]

# -----------------------------
# 4) OLS com covariada (sem interação)
# -----------------------------
ols_com_w <- lm_robust(Y ~ D + W, data = dados)
coef(ols_com_w)["D"]
ols_com_w$std.error[2]

# -----------------------------
# 5) Estimador de Lin: covariada centrada + interação
# -----------------------------
dados$Wc <- dados$W - mean(dados$W)
lin <- lm_robust(Y ~ D + Wc + D:Wc, data = dados)
coef(lin)["D"]
lin$std.error[2]
\end{lstlisting}

Em uma execução típica deste código, o erro-padrão de $\hat{\tau}_{Lin}$ é substancialmente menor do que o de $\hat{\tau}_{OLS\text{-sem-}W}$ --- frequentemente em $40\%$ a $60\%$, dado o alto $\rho$ entre $W$ e $Y$. O estimador é tão simples quanto o OLS, e a única diferença operacional é a centragem da covariada e a inclusão do termo de interação.

% =============================================================================
\section*{Encerramento da Seção}
\addcontentsline{toc}{section}{Encerramento da Seção}

Esta seção tratou da passagem do estimando para o estimador. Começamos pela definição da amostra e pelo princípio plug-in, que unifica OLS, diferença de médias, IV e IPW sob um único raciocínio: substituir a distribuição populacional pela distribuição empírica. Em seguida, organizamos os critérios de escolha entre estimadores --- consistência, viés, eficiência --- e identificamos o trade-off central entre eficiência e robustez.

A discussão sobre forma funcional mostrou que identificação não é suficiente para garantir que o estimador acerte o estimando. Mesmo com CIA válida, OLS pode calcular uma média ponderada com pesos que não correspondem aos da população alvo da política. O IPW oferece uma alternativa baseada em modelar a probabilidade de tratamento em vez do resultado, com vantagens em robustez funcional e custos em sensibilidade ao propensity score. O AIPW combina os dois e oferece dupla robustez.

A estimação em painel introduziu a possibilidade de controlar por características fixas não observadas das unidades. O TWFE e o FD são duas implementações que diferem na forma como exploram a variação temporal, com propriedades distintas em amostras finitas e em desenhos com adoção escalonada. O estimador de Lin, por fim, ofereceu uma resposta específica para o problema da incorporação de covariadas em RCTs: uma modificação simples ao OLS que preserva o não-viesamento da aleatorização e tipicamente melhora a precisão.

A próxima seção --- Inferência --- aborda o último elo da cadeia: dado que temos um estimador $\hat{\theta}$, como quantificar a incerteza amostral em torno dele? Erros-padrão, intervalos de confiança, testes de hipótese e cuidados com clustering e correlação serial são os temas dessa discussão. A linha que percorre todas as seções é a mesma: parâmetro econômico $\to$ estimando $\to$ estimador $\to$ incerteza. Cada etapa exige hipóteses distintas e cada hipótese precisa ser defendida.

% =============================================================================
% Fim da Seção 3
% =============================================================================
%% ─────────────────────────────────────────────────────────────────
%%  Seção 4 — Inferência
%%  Arquivo: secao4_inferencia.tex
%%  Inserir no projeto com: %% ─────────────────────────────────────────────────────────────────
%%  Seção 4 — Inferência
%%  Arquivo: secao4_inferencia.tex
%%  Inserir no projeto com: %% ─────────────────────────────────────────────────────────────────
%%  Seção 4 — Inferência
%%  Arquivo: secao4_inferencia.tex
%%  Inserir no projeto com: \input{secao4_inferencia}
%%
%%  Pacotes adicionais necessários no preâmbulo principal:
%%    \usepackage{pgfplots}
%%    \pgfplotsset{compat=1.18}
%% ─────────────────────────────────────────────────────────────────

\chapter{Inferência}
\label{cap:inferencia}
\label{sec:inferencia}

Nas seções anteriores construímos uma cadeia que vai da teoria ao dado: definimos o parâmetro de interesse, derivamos as hipóteses de identificação que o tornam observável e escolhemos um estimador $\hat{\theta}$ que converge para o estimando à medida que o tamanho da amostra cresce. Mas obter um número $\hat{\theta}$ não é o fim do trabalho. Se sorteássemos uma nova amostra da mesma população e aplicássemos exatamente o mesmo procedimento, obteríamos um valor diferente. A \textbf{inferência estatística} é o conjunto de ferramentas para raciocinar sobre essa variabilidade amostral e fazer afirmações sobre o parâmetro verdadeiro $\theta$ com base em uma única amostra observada.

É importante distinguir desde o início duas fontes de aleatoriedade que coexistem em econometria aplicada. A \textbf{aleatoriedade de amostragem} (\textit{sampling-based}) trata as observações como realizações independentes e identicamente distribuídas (i.i.d.) de uma população maior: a amostra é um sorteio, e a variabilidade de $\hat{\theta}$ reflete quantas amostras diferentes poderiam ter sido coletadas. Já a \textbf{aleatoriedade de atribuição} (\textit{design-based}) fixa a população de unidades observadas e trata como fonte de incerteza apenas o sorteio de quem recebeu o tratamento. Em experimentos aleatórios com um cadastro fixo de participantes, por exemplo, não há uma ``população maior'' de quem sortear: o que se repete hipoteticamente é a atribuição do tratamento. As duas perspectivas conduzem a fórmulas de erro-padrão distintas, e usar a fórmula errada — mesmo que matematicamente válida em outro contexto — gera inferências incorretas.

O ponto central desta seção: \textbf{um erro-padrão pequeno não implica identificação válida}. O erro-padrão mede com que precisão estimamos o estimando — a quantidade que definimos como objeto de interesse dado os dados. Ele não diz nada sobre se o estimando é igual ao parâmetro econômico que nos importa. Uma estimativa extremamente precisa de uma quantidade sem interpretação causal é, em todos os sentidos relevantes, pior do que uma estimativa imprecisa de algo bem identificado. Essa distinção, aparentemente óbvia, é sistematicamente ignorada na prática.


%% ─────────────────────────────────────────────────────────────────
\section{Erros-Padrão}
\label{sec:erros_padrao}

O \textbf{erro-padrão} de um estimador $\hat{\theta}$ é o desvio-padrão da sua distribuição amostral:
\begin{equation}
  SE(\hat{\theta}) = \sqrt{\text{Var}(\hat{\theta})}.
  \label{eq:se_definicao}
\end{equation}

Em termos práticos, $SE(\hat{\theta})$ quantifica o quanto $\hat{\theta}$ varia de amostra para amostra. Quanto menor o erro-padrão, mais concentrada é a distribuição do estimador em torno do seu valor esperado — e, se o estimador é não-viesado, em torno do parâmetro verdadeiro $\theta$.

\subsection*{Variância do estimador de MQO sob hipóteses clássicas}

No modelo de regressão linear $Y_i = \boldsymbol{x}_i'\boldsymbol{\beta} + \varepsilon_i$, o estimador MQO $\hat{\boldsymbol{\beta}} = (X'X)^{-1}X'Y$ tem variância:
\begin{equation}
  \text{Var}(\hat{\boldsymbol{\beta}} \mid X) = \sigma^2 (X'X)^{-1},
  \label{eq:var_ols_classica}
\end{equation}
onde $\sigma^2 = \text{Var}(\varepsilon_i \mid X)$. Esse resultado requer duas hipóteses fortes: (i) \textbf{homoscedasticidade} — a variância dos erros é constante para todos os valores de $X$ — e (ii) \textbf{independência} dos erros entre observações. Em dados de corte transversal representativos de populações heterogêneas, a homoscedasticidade é frequentemente violada: a variância do salário, por exemplo, tende a crescer com os anos de escolaridade. Quando \eqref{eq:var_ols_classica} é usada nesse contexto, os erros-padrão reportados são inconsistentes.

\subsection*{Erros-padrão robustos à heteroscedasticidade}

A solução padrão para dados de corte transversal são os \textbf{erros-padrão robustos} de Huber-White \citep{Huber1967, White1980}, também chamados de estimador \textit{sandwich}:
\begin{equation}
  \widehat{\text{Var}}_{\text{HC}}(\hat{\boldsymbol{\beta}}) =
    \underbrace{(X'X)^{-1}}_{\text{pão}}
    \underbrace{\left(\sum_{i=1}^{N} x_i x_i' \hat{\varepsilon}_i^2\right)}_{\text{recheio}}
    \underbrace{(X'X)^{-1}}_{\text{pão}},
  \label{eq:sandwich}
\end{equation}
onde $\hat{\varepsilon}_i = Y_i - \boldsymbol{x}_i'\hat{\boldsymbol{\beta}}$ são os resíduos da regressão. O estimador \eqref{eq:sandwich} é consistente independentemente da forma da heteroscedasticidade, sem exigir nenhuma hipótese sobre $\text{Var}(\varepsilon_i \mid X_i)$. Na prática, utiliza-se a versão com correção de amostras finitas HC2 ou HC3, que melhoram o desempenho em amostras pequenas.

\begin{notabox}[Erros robustos como padrão]
Em econometria aplicada moderna, erros-padrão robustos à heteroscedasticidade são o padrão mínimo para dados de corte transversal. A pergunta não é mais ``devo usar erros robustos?'', mas ``que tipo de robustez é necessária?'' (heteroscedasticidade apenas, ou também clustering?). Em R, a função \texttt{estimatr::lm\_robust()} calcula erros HC2 por padrão; o pacote \texttt{sandwich} oferece controle total sobre o tipo de correção.
\end{notabox}

\begin{lstlisting}[style=Rstyle, caption={Erros-padrão clássicos vs.\ robustos em R}]
# -----------------------------
# 1) Configuração inicial
# -----------------------------
rm(list = ls())

library(estimatr)   # lm_robust
library(sandwich)   # vcovHC
library(lmtest)     # coeftest

set.seed(42)
n <- 500

# Simular dados com heteroscedasticidade
edu   <- runif(n, 4, 18)
sigma <- 0.5 * edu          # variância cresce com escolaridade
salario <- 800 + 120 * edu + rnorm(n, sd = sigma * 300)

dados <- data.frame(salario, edu)

# -----------------------------
# 2) MQO com erros clássicos
# -----------------------------
mod_classico <- lm(salario ~ edu, data = dados)
cat("=== Erros clássicos ===\n")
print(coeftest(mod_classico)[, 1:2])

# -----------------------------
# 3) MQO com erros HC2 (robustos)
# -----------------------------
mod_robusto <- lm_robust(salario ~ edu, data = dados, se_type = "HC2")
cat("\n=== Erros HC2 robustos ===\n")
print(summary(mod_robusto)$coefficients[, 1:2])

# -----------------------------
# 4) Comparação direta
# -----------------------------
se_classico <- sqrt(diag(vcov(mod_classico)))["edu"]
se_robusto  <- sqrt(diag(vcovHC(mod_classico, type = "HC2")))["edu"]

cat("\nSE clássico (edu): ", round(se_classico, 3), "\n")
cat("SE HC2 robusto (edu):", round(se_robusto, 3), "\n")
cat("Razão HC2/clássico:  ", round(se_robusto / se_classico, 2), "\n")
\end{lstlisting}

\subsection*{O que o erro-padrão captura — e o que não captura}

Há uma confusão recorrente na literatura aplicada: tratar um erro-padrão pequeno como evidência de que ``a estimativa está correta''. A Tabela~\ref{tab:se_captura} sintetiza o que o erro-padrão realmente mede.

\begin{table}[H]
\centering
\caption{O que o erro-padrão captura e o que ele não captura}
\label{tab:se_captura}
\begin{tabular}{ll}
\toprule
\textbf{O SE captura} & \textbf{O SE não captura} \\
\midrule
Variabilidade amostral de $\hat{\theta}$ & Viés de identificação \\
Incerteza sobre a estimativa do estimando & Viés de estimação em amostras finitas \\
Reduz com aumento de $n$ & Erro de forma funcional \\
Melhora com amostragem mais eficiente & Viés por variável omitida \\
& Violação do SUTVA ou da CIA \\
\bottomrule
\end{tabular}

\vspace{4pt}
\figmeta{Elaboração dos autores}{}{SE: erro-padrão; CIA: hipótese de independência condicional; SUTVA: Stable Unit Treatment Value Assumption.}
\end{table}

Será que um erro-padrão pequeno garante que nossa estimativa está correta? Claramente não. Com uma amostra suficientemente grande, podemos estimar com extrema precisão a diferença de salários entre graduados e não-graduados — mas se essa diferença refletir diferenças de habilidade inata (viés de seleção), estimamos com precisão a coisa errada. Um SE pequeno com identificação fraca não apenas não impressiona: é potencialmente mais enganoso do que um SE grande com identificação sólida.


%% ─────────────────────────────────────────────────────────────────
\section{Teste de Hipótese}
\label{sec:teste_hipotese}

Dado um estimador $\hat{\theta}$ com erro-padrão $SE(\hat{\theta})$, a segunda pergunta natural é: os dados são compatíveis com a hipótese de que $\theta$ assume um valor particular $\theta_0$? O \textbf{teste de hipótese} oferece um procedimento formal para responder a essa pergunta. Apresentamos a lógica em cinco passos.

\medskip
\noindent\textbf{Passo 1 — Hipótese nula.} Definimos a hipótese nula $H_0: \theta = \theta_0$. Na maioria das aplicações em avaliação de impacto, $\theta_0 = 0$ (``o tratamento não tem efeito''), mas qualquer valor é válido.

\noindent\textbf{Passo 2 — Hipótese alternativa.} A hipótese alternativa $H_1: \theta \neq \theta_0$ especifica o que acreditamos ser verdadeiro caso $H_0$ seja falsa. Testes bicaudais são padrão; testes unicaudais ($H_1: \theta > \theta_0$) são admissíveis quando há razão teórica clara para a direção do efeito.

\noindent\textbf{Passo 3 — Estatística de teste.} Construímos a estatística:
\begin{equation}
  t = \frac{\hat{\theta} - \theta_0}{SE(\hat{\theta})},
  \label{eq:estatistica_t}
\end{equation}
que mede em quantos erros-padrão $\hat{\theta}$ se distancia de $\theta_0$. Sob $H_0$ e condições de regularidade, $t$ tem distribuição assintótica $\mathcal{N}(0,1)$ (ou $t_{N-K}$ exata se os erros forem normais).

\noindent\textbf{Passo 4 — P-valor.} O p-valor é a probabilidade de observar, sob $H_0$, uma estatística de teste tão ou mais extrema do que a observada:
\begin{equation}
  \text{p-valor} = P\!\left(|T| \geq |t_{\text{obs}}| \mid H_0\right),
  \label{eq:pvalor}
\end{equation}
onde $T$ é uma variável aleatória com a distribuição de $t$ sob $H_0$.

\noindent\textbf{Passo 5 — Regra de decisão.} Rejeitamos $H_0$ ao nível de significância $\alpha$ se p-valor $< \alpha$. O nível $\alpha = 0{,}05$ é convencional, mas arbitrário — corresponde a aceitar uma taxa de 5\% de falsos positivos quando $H_0$ é verdadeira.

\medskip
A Figura~\ref{fig:dist_h0} ilustra a lógica graficamente.

% FIGURA 4.1: Distribuição da estatística de teste sob H0
% Curva normal centrada em zero, regiões de rejeição sombreadas em cada cauda (|t| > 1,96),
% seta indicando t_obs = 2,35. Eixo x: estatística de teste. Eixo y: densidade.

\begin{figure}[H]
\centering
\begin{tikzpicture}
\begin{axis}[
  height  = 5.5cm,
  width   = 11cm,
  xmin    = -4.2, xmax = 4.2,
  ymin    = 0,    ymax = 0.46,
  xlabel  = {Estatística de teste $t$},
  ylabel  = {Densidade},
  axis lines = left,
  xtick  = {-1.96, 0, 1.96},
  xticklabels = {$-1{,}96$, $0$, $1{,}96$},
  ytick  = \empty,
  tick label style = {font=\small},
  label style      = {font=\small},
  clip   = false,
]
  %% Região de rejeição esquerda
  \addplot[fill=clearred!35, draw=none, domain=-4.2:-1.96, samples=80]
    {exp(-x^2/2)/sqrt(2*3.14159)} \closedcycle;
  %% Região de rejeição direita
  \addplot[fill=clearred!35, draw=none, domain=1.96:4.2, samples=80]
    {exp(-x^2/2)/sqrt(2*3.14159)} \closedcycle;
  %% Curva principal
  \addplot[thick, clearblue, domain=-4.2:4.2, samples=120]
    {exp(-x^2/2)/sqrt(2*3.14159)};
  %% Linhas verticais nos valores críticos
  \draw[dashed, gray!60] (axis cs:-1.96, 0) -- (axis cs:-1.96, 0.062);
  \draw[dashed, gray!60] (axis cs:1.96,  0) -- (axis cs:1.96,  0.062);
  %% Seta para t_obs
  \draw[-{Stealth}, thick, clearred]
    (axis cs:2.35, 0.10) -- (axis cs:2.35, 0.015);
  \node[above, font=\small, clearred] at (axis cs:2.35, 0.10)
    {$t_{\text{obs}} = 2{,}35$};
  %% Rótulos das caudas
  \node[font=\scriptsize] at (axis cs:-2.85, 0.018) {$\alpha/2$};
  \node[font=\scriptsize] at (axis cs: 2.85, 0.018) {$\alpha/2$};
  %% Rótulo da região central
  \node[font=\small, gray!60] at (axis cs:0, 0.20) {$H_0$ não rejeitada};
\end{axis}
\end{tikzpicture}
\caption{Distribuição da estatística de teste sob $H_0$. As regiões sombreadas em vermelho
  correspondem à região de rejeição ao nível $\alpha = 0{,}05$ (cada cauda com $\alpha/2 = 2{,}5\%$).
  O valor observado $t_{\text{obs}} = 2{,}35$ cai na região de rejeição ($|t| > 1{,}96$), levando à
  rejeição de $H_0$.}
\label{fig:dist_h0}
\end{figure}

A Figura~\ref{fig:poder} ilustra a relação entre a distribuição sob $H_0$, a distribuição sob uma hipótese alternativa específica com efeito $\delta$, e o conceito de poder do teste.

% FIGURA 4.2: Distribuição sob H0 (cinza) e H1 (azul).
% H0 ~ N(0,1); H1 ~ N(delta, 1) com delta = 2.
% Área azul à direita do valor crítico = poder do teste.
% Área azul à esquerda do valor crítico = erro Tipo II (beta).

\begin{figure}[H]
\centering
\begin{tikzpicture}
\begin{axis}[
  height  = 5.5cm,
  width   = 11cm,
  xmin    = -4.0, xmax = 5.5,
  ymin    = 0,    ymax = 0.46,
  xlabel  = {Estatística de teste $t$},
  ylabel  = {Densidade},
  axis lines = left,
  xtick  = {0, 1.96},
  xticklabels = {$0$, $t_c{=}1{,}96$},
  ytick  = \empty,
  tick label style = {font=\small},
  label style      = {font=\small},
  clip   = false,
  legend style = {at={(0.97,0.97)}, anchor=north east,
                  font=\small, draw=gray!50},
]
  %% Curva H0 (cinza preenchido)
  \addplot[fill=gray!20, draw=none, domain=-4:5.5, samples=100]
    {exp(-x^2/2)/sqrt(2*3.14159)} \closedcycle;
  \addplot[thick, gray!70, domain=-4:5.5, samples=100]
    {exp(-x^2/2)/sqrt(2*3.14159)};

  %% Poder = área sob H1 à direita de t_c
  \addplot[fill=clearblue!35, draw=none, domain=1.96:5.5, samples=80]
    {exp(-(x-2)^2/2)/sqrt(2*3.14159)} \closedcycle;

  %% Erro Tipo II = área sob H1 à esquerda de t_c
  \addplot[fill=clearteal!25, draw=none, domain=-1:1.96, samples=80]
    {exp(-(x-2)^2/2)/sqrt(2*3.14159)} \closedcycle;

  %% Curva H1
  \addplot[thick, clearblue, domain=-1:5.5, samples=100]
    {exp(-(x-2)^2/2)/sqrt(2*3.14159)};

  %% Linha vertical no valor crítico
  \draw[dashed, gray] (axis cs:1.96, 0) -- (axis cs:1.96, 0.40);

  %% Rótulos
  \node[font=\small, gray!80] at (axis cs:-0.7, 0.20) {$H_0$};
  \node[font=\small, clearblue] at (axis cs:3.4, 0.20) {$H_1$};
  \node[font=\small, clearblue] at (axis cs:2.0, -0.025) {$\delta$};
  \node[font=\scriptsize, clearblue!80] at (axis cs:2.9, 0.06)
    {\textbf{Poder}};
  \node[font=\scriptsize, clearteal!80] at (axis cs:1.0, 0.06)
    {Erro Tipo II ($\beta$)};

  \addlegendentry{$H_0: \theta = 0$}
  \addlegendentry{$H_1: \theta = \delta$}
\end{axis}
\end{tikzpicture}
\caption{Distribuição da estatística de teste sob $H_0$ (cinza) e sob $H_1$ com efeito $\delta = 2$
  (azul). A área azul à direita do valor crítico $t_c = 1{,}96$ representa o \textbf{poder do teste}
  — probabilidade de rejeitar $H_0$ quando $H_1$ é verdadeira. A área azul-esverdeada à esquerda
  é a probabilidade de erro Tipo~II ($\beta = 1 - \text{poder}$). Para referência sobre o cálculo
  do poder e do tamanho mínimo de amostra, ver o \textit{Guia de Cálculo de Poder Estatístico}
  da Série Avaliação na Prática.}
\label{fig:poder}
\end{figure}

\subsection*{O que o p-valor diz — e o que não diz}

O p-valor é um dos objetos mais mal interpretados em toda a estatística aplicada. Vale ser rigoroso. O p-valor \textbf{é} a probabilidade de observar, em um mundo em que $H_0$ é verdadeira, uma estatística tão ou mais extrema quanto a que observamos. Ele \textbf{não é}:

\begin{itemize}[leftmargin=*]
  \item A probabilidade de que $H_0$ seja verdadeira.
  \item A probabilidade de que o resultado seja ``devido ao acaso''.
  \item Uma medida do tamanho do efeito ou de sua relevância econômica.
\end{itemize}

Um ponto frequentemente ignorado em avaliações com amostras grandes: com $N$ suficientemente grande, \textit{qualquer} efeito não-nulo, por menor que seja, produz um p-valor arbitrariamente próximo de zero. Um programa que aumenta a renda em R\$\,2 por mês será ``altamente significativo'' ($p < 0{,}001$) com dados de um milhão de pessoas — e ainda assim economicamente irrelevante. Significância estatística não implica significância econômica.

\begin{atencaobox}[Reporte sempre o tamanho do efeito]
Nunca reporte apenas o p-valor. Reporte sempre: (i) a estimativa pontual $\hat{\theta}$, (ii) o erro-padrão ou o intervalo de confiança, e (iii) uma discussão da relevância econômica do efeito estimado, incluindo comparação com benchmarks relevantes (efeito de políticas similares, custo do programa, magnitude esperada pela teoria). Um p-valor sozinho não permite avaliar se o resultado importa.
\end{atencaobox}

\subsection*{Estatística $F$ para hipóteses conjuntas}

Quando $H_0$ envolve múltiplas restrições simultâneas — por exemplo, $H_0: \beta_1 = \beta_2 = \cdots = \beta_k = 0$ —, testar cada coeficiente individualmente com estatísticas-$t$ separadas infla a taxa de falsos positivos. A solução é a \textbf{estatística $F$}:
\begin{equation}
  F = \frac{(R\hat{\boldsymbol{\beta}} - r)'[R(X'X)^{-1}R']^{-1}(R\hat{\boldsymbol{\beta}} - r)/q}%
           {\hat{\sigma}^2},
  \label{eq:estatistica_F}
\end{equation}
onde $R\boldsymbol{\beta} = r$ codifica as $q$ restrições da hipótese nula. Sob $H_0$ e erros normais, $F \sim F(q, N-K)$.

A estatística $F$ aparece em três contextos recorrentes em avaliação de impacto:

\begin{enumerate}[leftmargin=*, label=(\roman*)]
  \item \textbf{Significância conjunta dos controles}: testar $H_0: \boldsymbol{\gamma} = \mathbf{0}$
    no modelo $Y_i = \alpha + \beta D_i + \boldsymbol{\gamma}' X_i + \varepsilon_i$.
  \item \textbf{Força do instrumento em IV}: a estatística $F$ do primeiro estágio mede a relevância do instrumento. A regra de bolso clássica é $F < 10$ como sinal de instrumento fraco \citep{StockYogo2005}; limiares mais rígidos são discutidos em trabalhos recentes.
  \item \textbf{Verificação de equilíbrio em experimentos}: $H_0$ de que as covariadas pré-tratamento não predizem o tratamento. Rejeição é evidência de falha na randomização.
\end{enumerate}


%% ─── 4.2.1 Múltiplos Testes ──────────────────────────────────────
\subsection{Múltiplos Testes}
\label{sec:multiplos_testes}

Um corolário muitas vezes subestimado do raciocínio da seção anterior: quando realizamos múltiplos testes simultaneamente, a probabilidade de ao menos um falso positivo cresce rapidamente, mesmo que cada teste individual seja conduzido corretamente ao nível $\alpha$. Se realizamos $m$ testes independentes, a probabilidade de ao menos uma rejeição espúria sob $H_0$ é:
\begin{equation}
  P(\text{ao menos um falso positivo}) = 1 - (1 - \alpha)^m.
  \label{eq:multiplos_testes}
\end{equation}

Com $m = 20$ testes ao nível $\alpha = 0{,}05$, essa probabilidade já atinge $1 - 0{,}95^{20} \approx 64\%$. A Tabela~\ref{tab:multiplos_testes} ilustra como ela cresce com $m$.

\begin{table}[H]
\centering
\caption{Probabilidade de ao menos um falso positivo em função do número de testes ($\alpha = 0{,}05$)}
\label{tab:multiplos_testes}
\begin{tabular}{rr}
\toprule
Número de testes ($m$) & $P(\text{ao menos um falso positivo})$ \\
\midrule
 1  & 5{,}0\% \\
 5  & 22{,}6\% \\
10  & 40{,}1\% \\
20  & 64{,}2\% \\
50  & 92{,}3\% \\
100 & 99{,}4\% \\
\bottomrule
\end{tabular}

\vspace{4pt}
\figmeta{Elaboração dos autores}{}{Calculado pela equação~\eqref{eq:multiplos_testes} com testes independentes.}
\end{table}

Dois procedimentos de correção são amplamente utilizados:

\medskip
\noindent\textbf{Correção de Bonferroni.} Para controlar a \textit{taxa de erro familiar} (FWER — probabilidade de ao menos um falso positivo), substitui-se o limiar $\alpha$ por $\alpha/m$. É conservadora — garante controle da FWER mesmo com dependência entre os testes —, mas perde poder quando $m$ é grande, pois cada teste exige evidência muito forte para ser rejeitado.

\medskip
\noindent\textbf{Taxa de Falsas Descobertas (FDR).} O procedimento de Benjamini-Hochberg \citep{BenjaminiHochberg1995} controla a proporção esperada de falsas descobertas entre as rejeições ($FDR = E[\text{falsos positivos}/\text{rejeições}]$), e não o número absoluto. É menos conservador que Bonferroni e preferível quando $m$ é grande e se aceita alguma proporção de erros. O procedimento ordena os $m$ p-valores em ordem crescente $p_{(1)} \leq p_{(2)} \leq \cdots \leq p_{(m)}$ e rejeita todas as hipóteses $H_{(j)}$ em que $p_{(j)} \leq j \cdot \alpha / m$.

\begin{lstlisting}[style=Rstyle, caption={Correção para múltiplos testes em R}]
# -----------------------------
# 1) Simulação: 20 testes sob H0 (todos nulos)
# -----------------------------
rm(list = ls())
set.seed(100)

m  <- 20
n  <- 200

# Gerar 20 p-valores sob H0 verdadeira
pvalores <- replicate(m, {
  Y <- rnorm(n)
  D <- rbinom(n, 1, 0.5)
  coef(summary(lm(Y ~ D)))[2, 4]  # p-valor do tratamento
})

cat("Sem correção:    rejeições =", sum(pvalores < 0.05), "de", m, "\n")

# -----------------------------
# 2) Correção de Bonferroni
# -----------------------------
pvalores_bonf <- p.adjust(pvalores, method = "bonferroni")
cat("Bonferroni:      rejeições =", sum(pvalores_bonf < 0.05), "\n")

# -----------------------------
# 3) Benjamini-Hochberg (FDR)
# -----------------------------
pvalores_bh <- p.adjust(pvalores, method = "BH")
cat("Benjamini-Hochberg: rejeições =", sum(pvalores_bh < 0.05), "\n")

# -----------------------------
# 4) Probabilidade teórica vs. simulada
# -----------------------------
set.seed(1)
sim_prob <- mean(replicate(10000, any(runif(m) < 0.05)))
cat("\nP(ao menos 1 falso positivo) — simulada: ", round(sim_prob, 3), "\n")
cat("P(ao menos 1 falso positivo) — teórica:  ",
    round(1 - 0.95^m, 3), "\n")
\end{lstlisting}

No entanto, correções estatísticas apenas não resolvem o problema estrutural do \textit{p-hacking}: quando o pesquisador realiza muitos testes e reporta apenas os ``significativos'', os p-valores não têm mais a cobertura nominal. A melhor prática é o \textbf{pré-registro}: declarar previamente, em repositório público (e.g., AEA RCT Registry, OSF), quais testes serão realizados, qual é o resultado primário da avaliação e quais são secundários. Independentemente do procedimento de correção utilizado, o relato deve sempre indicar quantos testes foram realizados.


%% ─────────────────────────────────────────────────────────────────
\section{Intervalos de Confiança}
\label{sec:intervalos}

O teste de hipótese responde a uma pergunta binária: os dados são compatíveis com $\theta = \theta_0$? O \textbf{intervalo de confiança} é mais informativo: ele identifica o conjunto de todos os valores $\theta_0$ que \textit{não seriam} rejeitados por um teste ao nível $\alpha$. Para estimadores assintoticamente normais, o intervalo de confiança de nível $1-\alpha$ é:

\begin{equation}
  IC_{1-\alpha} = \left[
    \hat{\theta} - z_{\alpha/2} \cdot SE(\hat{\theta}),\;\;
    \hat{\theta} + z_{\alpha/2} \cdot SE(\hat{\theta})
  \right],
  \label{eq:ic}
\end{equation}

onde $z_{\alpha/2}$ é o quantil $(1-\alpha/2)$ da distribuição normal padrão. Para $\alpha = 0{,}05$: $z_{0{,}025} = 1{,}96$.

\subsection*{Interpretação frequentista correta}

A interpretação do intervalo de confiança merece atenção especial, pois é uma das fontes mais comuns de confusão em estatística aplicada. O intervalo de confiança \textit{não} é uma afirmação sobre a probabilidade de $\theta$ estar em um intervalo específico: $\theta$ é um número fixo (desconhecido), e portanto não tem distribuição. A interpretação correta é frequentista: se o mesmo procedimento de construção do intervalo fosse repetido em muitas amostras independentes da mesma população, 95\% dos intervalos resultantes conteriam o valor verdadeiro $\theta$.

Em outras palavras, a aleatoriedade está no intervalo, não no parâmetro. A frase ``existe 95\% de probabilidade de que $\theta$ esteja entre $a$ e $b$'' está formalmente errada; o que é correto é ``o intervalo $[a, b]$ foi produzido por um procedimento que acerta em 95\% das amostras''.

\subsection*{Dualidade com o teste de hipótese}

Há uma correspondência exata entre intervalos de confiança e testes de hipótese: \textbf{rejeitar $H_0: \theta = \theta_0$ ao nível $\alpha$ é equivalente a $\theta_0$ não pertencer ao IC de nível $1-\alpha$}. Essa dualidade não é apenas uma curiosidade matemática — ela tem implicação prática direta.

O intervalo de confiança carrega estritamente mais informação do que o p-valor: ele diz não apenas se $\theta_0 = 0$ é rejeitado, mas quais outros valores são compatíveis com os dados. Um IC de $[{-0{,}5},\; 12{,}3]$ diz que efeitos nulos \textit{e} efeitos moderados são igualmente compatíveis com a evidência — informação que o p-valor não transmite. Reportar intervalos de confiança deve ser padrão; reportar apenas p-valores é uma prática a ser evitada.

\begin{lstlisting}[style=Rstyle, caption={Intervalos de confiança: construção e simulação da cobertura}]
# -----------------------------
# 1) IC padrão
# -----------------------------
rm(list = ls())
library(estimatr)

set.seed(42)
n <- 400
D <- rbinom(n, 1, 0.5)
Y <- 2 + 0.8 * D + rnorm(n)

mod <- lm_robust(Y ~ D, se_type = "HC2")
cat("Estimativa: ", round(coef(mod)["D"], 3), "\n")
cat("IC 95%:    [", round(confint(mod)["D", 1], 3),
               ";", round(confint(mod)["D", 2], 3), "]\n")

# -----------------------------
# 2) Verificar cobertura empírica do IC
# Deveria ser ~95% se o procedimento está correto
# -----------------------------
set.seed(99)
theta_verdadeiro <- 0.8
B <- 5000

cobertura <- replicate(B, {
  D_s <- rbinom(n, 1, 0.5)
  Y_s <- 2 + theta_verdadeiro * D_s + rnorm(n)
  m_s <- lm_robust(Y_s ~ D_s, se_type = "HC2")
  ic  <- confint(m_s)["D_s", ]
  ic[1] <= theta_verdadeiro & theta_verdadeiro <= ic[2]
})

cat("\nCobertura empírica do IC 95%:", round(mean(cobertura) * 100, 1), "%\n")
cat("(Esperado: 95%)\n")
\end{lstlisting}


%% ─────────────────────────────────────────────────────────────────
\section{Clustering}
\label{sec:clustering}

Na seção anterior, os erros-padrão robustos de Huber-White corrigem a heteroscedasticidade, mas ainda assumem independência entre as observações. Em muitas aplicações em políticas públicas essa hipótese é violada: estudantes de uma mesma escola enfrentam o mesmo professor e a mesma infraestrutura; municípios de um mesmo estado compartilham choques econômicos e políticas estaduais. Ignorar essa dependência dentro de grupos faz com que os erros-padrão clássicos \textit{subestimem} a variabilidade real de $\hat{\beta}$, inflando artificialmente a significância das estimativas.

O problema pode ser descrito com precisão. Suponha que as observações sejam agrupadas em $G$ clusters $g = 1, \ldots, G$, com $n_g$ observações no cluster $g$ ($N = \sum_g n_g$). O modelo é:
\begin{equation}
  Y_{ig} = \boldsymbol{x}_{ig}'\boldsymbol{\beta} + \varepsilon_{ig},
  \label{eq:modelo_cluster}
\end{equation}
onde $\varepsilon_{ig}$ e $\varepsilon_{jg}$ podem ser correlacionados (mesmo cluster $g$), mas $\varepsilon_{ig}$ e $\varepsilon_{jh}$ são independentes para $g \neq h$. O estimador MQO $\hat{\boldsymbol{\beta}}$ permanece não-viesado e consistente sob \eqref{eq:modelo_cluster}, mas a variância clássica $(X'X)^{-1}\hat{\sigma}^2$ é inconsistente. A alternativa consistente é o \textbf{estimador \textit{sandwich} clusterizado}:
\begin{equation}
  \widehat{\text{Var}}_{\text{CL}}(\hat{\boldsymbol{\beta}}) =
    (X'X)^{-1}
    \underbrace{\left(\sum_{g=1}^{G} B_g B_g'\right)}_{\text{``carne'' clusterizada}}
    (X'X)^{-1},
  \quad
  B_g = \sum_{i \in g} \boldsymbol{x}_{ig}\, \hat{\varepsilon}_{ig}.
  \label{eq:sandwich_cl}
\end{equation}

A correção de amostras finitas multiplica \eqref{eq:sandwich_cl} por $\frac{G}{G-1} \cdot \frac{N-1}{N-K}$, análoga à divisão por $N-K$ (e não $N$) nos erros clássicos.

\subsection*{A regra central: clusterizar no nível de atribuição do tratamento}

A regra mais importante do \textit{clustering} é independente dos dados e da forma do modelo: \textbf{o cluster deve ser definido no nível em que o tratamento foi atribuído}. Se uma política é implementada no nível do município e os dados são individuais (pessoas dentro de municípios), deve-se clusterizar por município — independentemente de quantas observações individuais existam. Usar erros robustos sem cluster nessa situação é equivalente a agir como se os 50.000 moradores de um município fossem 50.000 informações independentes sobre o efeito da política, quando na realidade todos compartilham exatamente a mesma exposição ao tratamento.

\begin{table}[H]
\centering
\caption{Quando clusterizar: critérios práticos}
\label{tab:quando_clusterizar}
\begin{tabular}{p{6.2cm} p{6.0cm}}
\toprule
\textbf{Situação} & \textbf{Nível de cluster} \\
\midrule
Política atribuída por escola, dados por aluno    & Escola \\
Política atribuída por município, dados individuais & Município \\
Dados de painel com tratamento variando no tempo  & Unidade (+ período, se bidirecional) \\
Dados de PNAD com estrato de amostragem           & PSU/UPA \\
Experimento com randomização por grupo            & Grupo (mesmo que $G$ seja pequeno) \\
Corte transversal puro sem estrutura de grupo     & Desnecessário (HC suficiente) \\
\bottomrule
\end{tabular}

\vspace{4pt}
\figmeta{Elaboração dos autores}{}{PSU: \textit{Primary Sampling Unit}; UPA: Unidade Primária de Amostragem.}
\end{table}

\begin{lstlisting}[style=Rstyle, caption={Erros clusterizados com sandwich e estimatr}]
# -----------------------------
# 1) Dados com estrutura de cluster
# Política municipal; observações individuais
# -----------------------------
rm(list = ls())

library(estimatr)
library(sandwich)
library(lmtest)

set.seed(77)
G  <- 50   # municípios
ng <- 100  # indivíduos por município

municipio  <- rep(1:G, each = ng)
tratamento <- rep(rbinom(G, 1, 0.5), each = ng)  # atribuição por município
erro_mun   <- rep(rnorm(G, sd = 2), each = ng)   # choque comum ao município
erro_ind   <- rnorm(G * ng, sd = 1)              # choque individual

Y <- 5 + 1.5 * tratamento + erro_mun + erro_ind

dados_cl <- data.frame(Y, tratamento, municipio)

# -----------------------------
# 2) Comparação de erros-padrão
# -----------------------------
mod_base <- lm(Y ~ tratamento, data = dados_cl)

se_classico <- sqrt(diag(vcov(mod_base)))["tratamento"]
se_hc2      <- sqrt(vcovHC(mod_base, type = "HC2"))["tratamento", "tratamento"]
se_cluster  <- sqrt(
  vcovCL(mod_base, cluster = ~municipio)
)["tratamento", "tratamento"]

cat("SE clássico:   ", round(se_classico, 4), "\n")
cat("SE HC2:        ", round(se_hc2, 4), "\n")
cat("SE clusterizado:", round(se_cluster, 4), "\n")

# Com estimatr (mais direto)
mod_cl <- lm_robust(Y ~ tratamento, clusters = municipio, data = dados_cl,
                    se_type = "CR2")
se_cr2 <- summary(mod_cl)$coefficients["tratamento", "Std. Error"]
cat("\nSE CR2 (estimatr):", round(se_cr2, 4), "\n")
\end{lstlisting}

\subsection*{Problema de poucos clusters}

O estimador sandwich clusterizado é válido assintoticamente em $G \to \infty$, mas tem desempenho pobre em amostras finitas quando $G$ é pequeno. Com $G < 30$--50, ele tende a \textit{subestimar} a variância real, produzindo p-valores espúrios. Algumas alternativas:

\begin{itemize}[leftmargin=*]
  \item \textbf{Distribuição $t(G-1)$}: usar a distribuição $t$ com $G-1$ graus de liberdade em vez da normal assintótica. Solução simples e conservadora, adequada quando $G \geq 10$--15.
  \item \textbf{Correção Bell-McCaffrey}: ajuste de amostras finitas mais sofisticado, implementado em \texttt{estimatr} com \texttt{se\_type = \textquotedbl CR2\textquotedbl}.
  \item \textbf{\textit{Wild cluster bootstrap}}: descrito na próxima seção. Solução preferida quando $G < 20$.
\end{itemize}

\subsection*{Clustering bidirecional}

Em dados de painel, pode haver correlação tanto dentro de unidades ao longo do tempo (dependência serial) quanto dentro de períodos entre unidades (dependência transversal, por exemplo, por choques setoriais). O \textbf{clustering bidirecional} \citep{Cameron2011} clusteriza simultaneamente por unidade $i$ e por período $t$:
\begin{equation}
  \widehat{\text{Var}}_{\text{2D}} = \widehat{\text{Var}}_{i} + \widehat{\text{Var}}_{t} - \widehat{\text{Var}}_{it},
  \label{eq:cluster_2d}
\end{equation}
onde cada componente é o estimador sandwich clusterizado no respectivo nível. Requer número suficiente de clusters em ambas as dimensões.


%% ─────────────────────────────────────────────────────────────────
\section{Bootstrap}
\label{sec:bootstrap}

O \textit{bootstrap} é um método computacional para estimar a distribuição amostral de um estimador sem depender de fórmulas analíticas. A ideia, devida a \citet{Efron1979}, é simples: como a distribuição empírica $\hat{F}_n$ é nossa melhor estimativa da distribuição verdadeira $F$, podemos simular a variabilidade amostral reamostrado os dados observados.

O algoritmo básico tem quatro passos:

\begin{enumerate}[leftmargin=*, label=\textbf{\arabic*.}]
  \item Reamostrar com reposição: sortear $n$ observações da amostra original para formar uma amostra \textit{bootstrap} $\{(Y_i^*, X_i^*)\}_{i=1}^n$.
  \item Calcular o estimador nessa amostra: $\hat{\theta}^* = \hat{\theta}(\{Y_i^*, X_i^*\})$.
  \item Repetir os passos 1--2 por $B$ vezes, obtendo $\{\hat{\theta}^*_b\}_{b=1}^B$.
  \item Usar a distribuição empírica de $\{\hat{\theta}^*_b\}$ como aproximação da distribuição amostral de $\hat{\theta}$.
\end{enumerate}

$B = 999$ ou $B = 1{.}999$ repetições são tipicamente suficientes para ICs e p-valores práticos. O \textbf{IC percentil} de nível $1-\alpha$ é:
\begin{equation}
  IC^{\text{perc}}_{1-\alpha} = \left[\hat{\theta}^*_{(\alpha/2)},\;\; \hat{\theta}^*_{(1-\alpha/2)}\right],
  \label{eq:ic_bootstrap}
\end{equation}
onde $\hat{\theta}^*_{(q)}$ denota o quantil $q$ da distribuição bootstrap. O \textbf{IC \textit{bias-corrected}} (BC) ajusta adicionalmente pelo viés $\hat{\theta}^*_{(0{,}5)} - \hat{\theta}$, melhorando a cobertura quando o estimador é enviesado em amostras finitas.

\subsection*{\textit{Cluster bootstrap} e \textit{wild cluster bootstrap}}

O bootstrap padrão (reamostrar observações individuais) não preserva a estrutura de cluster: ao reamostrar observações aleatoriamente, destroem-se as correlações dentro de grupos. A solução é o \textbf{\textit{cluster bootstrap}}: reamostrar clusters inteiros, não observações individuais. Com $G$ clusters, sorteia-se com reposição $G$ clusters da coleção original, mantendo todas as observações de cada cluster sorteado.

Quando $G$ é muito pequeno (digamos, $G < 20$), mesmo o cluster bootstrap pode ter cobertura insatisfatória. Nesses casos, o \textbf{\textit{wild cluster bootstrap}} \citep{Cameron2008} é a solução preferida. Em vez de reamostrar clusters, ele mantém as observações originais e perturba os resíduos de cada cluster por um fator aleatório $w_g \in \{-1, +1\}$ com igual probabilidade (distribuição de Rademacher):
\begin{equation}
  Y_{ig}^* = \boldsymbol{x}_{ig}'\hat{\boldsymbol{\beta}} + w_g\, \hat{\varepsilon}_{ig},
  \label{eq:wild_bootstrap}
\end{equation}
e então reestima o modelo na amostra assim construída. Com $G$ clusters, existem $2^G$ permutações possíveis dos fatores $w_g$, produzindo uma distribuição de permutação exata quando $G$ é muito pequeno.

\begin{lstlisting}[style=Rstyle, caption={Bootstrap e wild cluster bootstrap em R}]
# -----------------------------
# 1) Bootstrap padrão (sem cluster)
# -----------------------------
rm(list = ls())
set.seed(42)

n     <- 300
D     <- rbinom(n, 1, 0.5)
Y     <- 2 + 1.0 * D + rnorm(n)
dados <- data.frame(Y, D)

B         <- 1999
boot_est  <- numeric(B)

for (b in 1:B) {
  idx          <- sample(n, replace = TRUE)
  boot_est[b]  <- coef(lm(Y ~ D, data = dados[idx, ]))["D"]
}

ic_boot <- quantile(boot_est, c(0.025, 0.975))
cat("IC bootstrap 95%: [", round(ic_boot[1], 3),
                      ";", round(ic_boot[2], 3), "]\n")

# -----------------------------
# 2) Wild cluster bootstrap com fewclusterboot
# -----------------------------
library(fwildclusterboot)  # instalar se necessário

G  <- 20
ng <- 50
municipio  <- rep(1:G, each = ng)
trat_g     <- rep(rbinom(G, 1, 0.5), each = ng)
Y2 <- 3 + 0.8 * trat_g + rep(rnorm(G, sd = 1.5), each = ng) + rnorm(G*ng)
dados2 <- data.frame(Y = Y2, trat = trat_g, mun = factor(municipio))

mod2 <- lm(Y ~ trat, data = dados2)
boot_res <- boottest(
  mod2,
  clustid   = "mun",
  param     = "trat",
  B         = 1999,
  seed      = 7
)
summary(boot_res)
\end{lstlisting}

\begin{atencaobox}[{\textit{Bootstrap} não substitui identificação}]
O bootstrap reduz a dependência de aproximações assintóticas e pode melhorar substancialmente a cobertura dos intervalos de confiança, especialmente com clusters pequenos. No entanto, ele não corrige um estimador inconsistente nem valida uma estratégia de identificação fraca. Se $\hat{\theta}$ converge para um valor diferente de $\theta$ (por viés de seleção, variável omitida etc.), reamostrar os dados produzirá uma distribuição bootstrap concentrada em torno do valor errado — com erros-padrão menores e intervalos de confiança que não cobrem $\theta$.
\end{atencaobox}


%% ─────────────────────────────────────────────────────────────────
\section{Inferência Baseada em Design}
\label{sec:design_based}

Todas as abordagens discutidas até aqui assumem, implicitamente, que as observações são sorteios de uma superpopulação maior. A incerteza vem de podermos ter obtido uma amostra diferente. Há, no entanto, situações em que essa perspectiva não faz sentido: quando utilizamos dados de \textit{todos} os municípios brasileiros, de \textit{todas} as escolas de uma rede, ou dos \textit{todos} os participantes de um experimento em que a lista de participantes foi fixada antes da randomização, não há população maior de que amostramos.

Na \textbf{inferência baseada em \textit{design}} (\textit{design-based inference}), a população de unidades é tratada como \textbf{fixa}. A única fonte de aleatoriedade é o \textbf{processo de atribuição do tratamento} — o sorteio de quem entre as unidades observadas recebeu o tratamento. Em vez de perguntar ``o que aconteceria se sorteássemos outra amostra?'', pergunta-se: ``o que aconteceria se, com estas mesmas unidades, o tratamento tivesse sido atribuído de forma diferente?''

\subsection*{Por que isso importa na prática}

Três razões tornam a perspectiva \textit{design-based} relevante em avaliação de políticas públicas no Brasil:

\begin{enumerate}[leftmargin=*, label=(\roman*)]
  \item \textbf{Dados administrativos censitários}: quando os dados cobrem a população inteira (todos os beneficiários do Bolsa Família, todos os funcionários públicos federais), a noção de ``amostra de uma superpopulação'' é forçada. A aleatoriedade relevante é a da política, não a da amostragem.
  \item \textbf{Experimentos com cadastro fixo}: em RCTs em que a lista de participantes é definida antes da randomização, o processo que gera a incerteza é exatamente o sorteio do tratamento.
  \item \textbf{Validade em amostras finitas}: a inferência por permutação é exata em amostras finitas, sem depender de aproximações assintóticas.
\end{enumerate}

\subsection*{Inferência por permutação}

O procedimento central da inferência \textit{design-based} é a \textbf{inferência por permutação} (\textit{randomization inference}) \citep{Fisher1935}. Ela parte da \textbf{hipótese nula afiada} (\textit{sharp null}):

\begin{hipotesebox}[Sharp null (Fisher)]
Para todo indivíduo $i$: $Y_i(1) = Y_i(0)$, ou seja, o tratamento não afeta nenhum indivíduo.
\end{hipotesebox}

Sob a hipótese nula afiada, os resultados potenciais $Y_i(0)$ e $Y_i(1)$ são iguais para todo $i$ — e portanto os resultados observados seriam os mesmos independentemente de quem recebesse o tratamento. Isso significa que, sob $H_0^{\text{sharp}}$, qualquer permutação dos rótulos de tratamento deveria produzir uma estatística de teste semelhante à observada. O p-valor exato é a fração das permutações que geram uma estatística tão ou mais extrema:
\begin{equation}
  \text{p-valor}_{\text{perm}} = \frac{1}{|\mathcal{W}|}
    \sum_{\mathbf{d} \in \mathcal{W}}
    \mathds{1}\!\left[|T(\mathbf{d})| \geq |T(\mathbf{d}_{\text{obs}})|\right],
  \label{eq:pvalue_perm}
\end{equation}
onde $\mathcal{W}$ é o conjunto de todas as atribuições possíveis e $T(\mathbf{d})$ é a estatística de teste calculada com a atribuição $\mathbf{d}$.

\subsection*{Exemplo numérico}

A Tabela~\ref{tab:permutacao} ilustra o procedimento com seis municípios, três dos quais recebem um programa de transferência de renda ($D_i = 1$). A estatística de teste é a diferença de médias entre tratados e controles.

\begin{table}[H]
\centering
\caption{Dados observados para o exemplo de inferência por permutação}
\label{tab:permutacao}
\begin{tabular}{ccc}
\toprule
Município & Tratamento ($D_i$) & Taxa de pobreza (\%) \\
\midrule
A & 1 & 18 \\
B & 1 & 21 \\
C & 1 & 16 \\
D & 0 & 27 \\
E & 0 & 24 \\
F & 0 & 29 \\
\bottomrule
\end{tabular}

\vspace{4pt}
\figmeta{Elaboração dos autores}{}{Dados hipotéticos. Tratados: municípios A, B, C. Controles: D, E, F.}
\end{table}

A diferença de médias observada é $T_{\text{obs}} = \bar{Y}_1 - \bar{Y}_0 = (18+21+16)/3 - (27+24+29)/3 = 18{,}33 - 26{,}67 = -8{,}33$. Existem $\binom{6}{3} = 20$ possíveis atribuições de 3 tratados entre 6 municípios. Sob $H_0^{\text{sharp}}$, calculamos a diferença de médias para cada uma:

\begin{lstlisting}[style=Rstyle, caption={Inferência por permutação: exemplo com 6 municípios}]
# -----------------------------
# 1) Dados observados
# -----------------------------
rm(list = ls())
library(combinat)

Y <- c(18, 21, 16, 27, 24, 29)   # taxa de pobreza
D <- c( 1,  1,  1,  0,  0,  0)   # tratamento observado
n <- length(Y)
n1 <- sum(D)

T_obs <- mean(Y[D == 1]) - mean(Y[D == 0])
cat("Diferença de médias observada:", round(T_obs, 2), "\n")

# -----------------------------
# 2) Gerar todas as C(6,3) = 20 permutações
# -----------------------------
perms <- combn(n, n1)  # colunas = cada possível conjunto de tratados
T_perm <- apply(perms, 2, function(trat_idx) {
  D_perm <- rep(0, n)
  D_perm[trat_idx] <- 1
  mean(Y[D_perm == 1]) - mean(Y[D_perm == 0])
})

cat("Distribuição de permutação:\n")
print(sort(round(T_perm, 2)))

# -----------------------------
# 3) P-valor exato (bilateral)
# -----------------------------
p_perm <- mean(abs(T_perm) >= abs(T_obs))
cat("\np-valor de permutação (bilateral):", round(p_perm, 3), "\n")
cat("(", sum(abs(T_perm) >= abs(T_obs)), "de", ncol(perms), "permutações)\n")
\end{lstlisting}

No exemplo, apenas algumas das 20 permutações produzem uma diferença tão extrema quanto $-8{,}33$. O p-valor exato é a proporção dessas permutações. Note que esse p-valor é válido em amostras finitas, sem qualquer hipótese distribucional sobre os erros.

\begin{notabox}[Sharp null vs.\ weak null]
A hipótese nula afiada ($Y_i(1) = Y_i(0)$ para todo $i$) é mais forte do que a hipótese nula fraca ($\mathbb{E}[Y_i(1) - Y_i(0)] = 0$). A inferência por permutação testa a primeira. Com heterogeneidade de efeitos (alguns positivos, outros negativos com média zero), o teste de permutação pode não rejeitar $H_0^{\text{sharp}}$ mesmo quando a $H_0^{\text{weak}}$ também é falsa. Esse ponto é discutido em \citet{Imbens2015}, Capítulo 5.
\end{notabox}


%% ─────────────────────────────────────────────────────────────────
\section{Testes de Falsificação}
\label{sec:falsificacao}

Em econometria causal, muitos dos testes mais informativos não testam o efeito do tratamento diretamente, mas a \textbf{plausibilidade das hipóteses de identificação}. Se a estratégia empírica requer tendências paralelas, a ausência de tendências divergentes antes do tratamento é evidência (circunstancial) a seu favor. Se a estratégia requer que um instrumento seja exógeno, a ausência de efeito sobre resultados que não deveriam ser afetados é reconfortante. Esses \textbf{testes de falsificação} (\textit{placebo tests}, \textit{falsification tests}) são uma das ferramentas mais valiosas do arsenal do avaliador.

\subsection*{1 — Testes de pré-tendência (DiD)}

No contexto de Diferenças em Diferenças, a identificação requer que o grupo tratado e o grupo controle teriam evoluído em paralelo na ausência do tratamento. Esse contrafactual não é testável diretamente, mas podemos verificar se os grupos evoluíam em paralelo \textit{antes} do tratamento.

O procedimento padrão é o \textbf{estudo de evento} (\textit{event study}): regride-se o resultado sobre dummies de período relativo ao tratamento, separadas por grupo:
\begin{equation}
  Y_{it} = \alpha_i + \lambda_t + \sum_{k \neq -1} \delta_k \cdot D_i \cdot \mathds{1}[t = T_i + k] + \varepsilon_{it},
  \label{eq:event_study}
\end{equation}
onde $k = -1$ é o período de referência (imediatamente antes do tratamento). Os coeficientes $\delta_k$ para $k < 0$ são os \textbf{coeficientes de pré-tendência}: sob tendências paralelas, devem ser estatisticamente indistinguíveis de zero.

% FIGURA 4.3: Gráfico de event study.
% Eixo x: períodos relativos ao tratamento (k = -4, -3, -2, -1, 0, 1, 2, 3, 4)
% Eixo y: coeficiente delta_k com intervalo de confiança de 95%
% Linha vertical em k = -0.5 (cutoff) e linha horizontal em 0
% Coeficientes pré-tratamento (k < 0) próximos de zero; coeficientes pós-tratamento (k >= 0) positivos

\begin{lstlisting}[style=Rstyle, caption={Estudo de evento: coeficientes de pré-tendência}]
# -----------------------------
# 1) Dados de painel simulados com pré-tendências satisfeitas
# -----------------------------
rm(list = ls())
library(dplyr)
library(ggplot2)
library(lfe)  # para efeitos fixos (ou usar fixest)

set.seed(42)
N  <- 100   # unidades
T  <- 10    # períodos (tratamento no período 6)

dados_painel <- expand.grid(id = 1:N, t = 1:T)
dados_painel <- dados_painel %>%
  mutate(
    tratado = id <= 50,        # metade das unidades é tratada
    efeito  = ifelse(tratado & t >= 6, 2.0, 0),
    Y = 5 + 0.3 * t +          # tendência comum
        rnorm(N * T, sd = 1) +  # erro individual
        efeito
  )

# -----------------------------
# 2) Estudo de evento
# -----------------------------
dados_painel <- dados_painel %>%
  mutate(k = t - 6)  # período relativo ao tratamento

# Dummies de interação tratado × período (excluindo k = -1)
library(fixest)
mod_es <- feols(Y ~ i(k, tratado, ref = -1) | id + t,
                data = dados_painel)

# -----------------------------
# 3) Gráfico de event study
# -----------------------------
coefs <- as.data.frame(coef(mod_es))
iplot(mod_es,
      main   = "Estudo de evento: coeficientes por período",
      xlab   = "Período relativo ao tratamento",
      ylab   = expression(delta[k]),
      pt.pch = 16)
abline(h = 0, lty = 2, col = "gray50")
abline(v = -0.5, lty = 3, col = "gray30")
\end{lstlisting}

\begin{conexaobox}[Guia de Diferenças em Diferenças]
A discussão completa sobre tendências paralelas, eventos escalonados (\textit{staggered DiD}) e estimadores robustos à heterogeneidade de efeitos está no \textit{Guia de Diferenças em Diferenças} (2 volumes) da Série Avaliação na Prática. Em particular, o segundo volume trata dos estimadores de Callaway-Sant'Anna, Sun-Abraham e Borusyak-Jaravel-Spiess, que evitam os pesos negativos do estimador TWFE em designs escalonados.
\end{conexaobox}

\subsection*{2 — Placebo de resultado}

Um \textbf{placebo de resultado} aplica a mesma estratégia empírica a um resultado que, pela teoria, \textit{não} deveria ser afetado pelo tratamento. Um efeito estatisticamente significativo é evidência de que a estratégia pode estar captando algo além do tratamento de interesse.

\begin{exemplobox}[Placebo de resultado: Bolsa Família e consumo]
Suponha que se estime o efeito do Bolsa Família sobre o consumo de alimentos usando uma RDD no ponto de corte de elegibilidade (renda per capita = R\$\,218). Como placebo, aplica-se a mesma RDD ao consumo de bens de luxo (viagens internacionais, veículos de luxo) — itens que não deveriam ser afetados por uma transferência de baixo valor para famílias pobres. Um efeito significativo nesse placebo seria evidência de violação da hipótese de continuidade ao redor do corte.
\end{exemplobox}

\subsection*{3 — Placebo de tratamento}

Um \textbf{placebo de tratamento} aplica a estratégia empírica a um período anterior ao tratamento real, como se o tratamento tivesse ocorrido nesse período anterior. Se a estratégia é válida, não deveria haver efeito detectável nesse período contrafactual.

Em um painel com dados de 2010 a 2020 e tratamento em 2016, por exemplo, pode-se estimar o efeito como se o tratamento tivesse ocorrido em 2013, usando apenas dados de 2010 a 2015. Um efeito significativo em 2013 sugere que as diferenças pré-existentes entre grupos, e não o tratamento, explicam o resultado principal.

\subsection*{4 — Teste de balanceamento}

O \textbf{teste de balanceamento} verifica se covariadas pré-determinadas predizem o tratamento. Em um experimento aleatório, nenhuma característica pré-tratamento deveria prever quem recebe o tratamento (a não ser o próprio design de estratificação). Em um estudo observacional, o balanceamento não prova causalidade, mas desequilíbrio acentuado é evidência de seleção endógena.

\begin{lstlisting}[style=Rstyle, caption={Teste de balanceamento via regressão}]
# -----------------------------
# 1) Regredir tratamento em covariadas pré-tratamento
# Um F-teste conjunto testa se alguma covariada prediz D
# -----------------------------
rm(list = ls())
library(estimatr)

set.seed(55)
n <- 500

# Covariadas pré-tratamento
idade   <- rnorm(n, 35, 10)
renda   <- rnorm(n, 3000, 800)
sexo    <- rbinom(n, 1, 0.5)

# Tratamento: (a) aleatório ou (b) com seleção
D_aleat  <- rbinom(n, 1, 0.5)
D_selec  <- rbinom(n, 1, plogis(-1 + 0.03 * renda / 100))

# -----------------------------
# 2) F-teste de balanceamento
# -----------------------------
bal_aleat <- lm(D_aleat ~ idade + renda + sexo)
bal_selec <- lm(D_selec ~ idade + renda + sexo)

cat("=== Tratamento aleatório ===\n")
print(summary(bal_aleat)$fstatistic)

cat("\n=== Tratamento com seleção ===\n")
print(summary(bal_selec)$fstatistic)

# Tabela de diferenças por grupo
for (D_var in list(D_aleat, D_selec)) {
  cat("\nMédias por grupo:\n")
  cat("  Renda tratados:   ", round(mean(renda[D_var == 1])), "\n")
  cat("  Renda controles:  ", round(mean(renda[D_var == 0])), "\n")
}
\end{lstlisting}

\subsection*{Uma advertência obrigatória}

Os testes de falsificação são ferramentas poderosas, mas suas limitações precisam ser enunciadas com clareza. Passar em todos os placebos é evidência \textit{circunstancial} a favor da identificação — é um argumento para a credibilidade, não uma prova. Uma estratégia pode superar todos os testes disponíveis e ainda assim ter identificação inválida por razões não testáveis: um fator de confusão que não está nos dados, uma violação de SUTVA não observada, ou uma descontinuidade nas covariadas que coincide exatamente com o ponto de corte da política. Não confundir ausência de evidência contra a identificação com evidência a favor da identificação. Os testes de falsificação, como toda a inferência causal, devem ser lidos em conjunto com o argumento teórico e o conhecimento substantivo do fenômeno estudado.

\begin{conexaobox}[Testes de falsificação nos guias de métodos]
Cada guia da Série Avaliação na Prática discute os testes de falsificação específicos para o método em questão: o guia de \textit{RDD} trata da continuidade das covariadas e da densidade ao redor do corte (teste de McCrary); o guia de \textit{Matching} discute testes de qualidade do pareamento e equilíbrio; o guia de \textit{Variáveis Instrumentais} trata dos testes de superidentificação (Hansen-Sargan) e de instrumento fraco. Os procedimentos desta seção fornecem o arcabouço geral que fundamenta todos esses testes específicos.
\end{conexaobox}
%%
%%  Pacotes adicionais necessários no preâmbulo principal:
%%    \usepackage{pgfplots}
%%    \pgfplotsset{compat=1.18}
%% ─────────────────────────────────────────────────────────────────

\chapter{Inferência}
\label{cap:inferencia}
\label{sec:inferencia}

Nas seções anteriores construímos uma cadeia que vai da teoria ao dado: definimos o parâmetro de interesse, derivamos as hipóteses de identificação que o tornam observável e escolhemos um estimador $\hat{\theta}$ que converge para o estimando à medida que o tamanho da amostra cresce. Mas obter um número $\hat{\theta}$ não é o fim do trabalho. Se sorteássemos uma nova amostra da mesma população e aplicássemos exatamente o mesmo procedimento, obteríamos um valor diferente. A \textbf{inferência estatística} é o conjunto de ferramentas para raciocinar sobre essa variabilidade amostral e fazer afirmações sobre o parâmetro verdadeiro $\theta$ com base em uma única amostra observada.

É importante distinguir desde o início duas fontes de aleatoriedade que coexistem em econometria aplicada. A \textbf{aleatoriedade de amostragem} (\textit{sampling-based}) trata as observações como realizações independentes e identicamente distribuídas (i.i.d.) de uma população maior: a amostra é um sorteio, e a variabilidade de $\hat{\theta}$ reflete quantas amostras diferentes poderiam ter sido coletadas. Já a \textbf{aleatoriedade de atribuição} (\textit{design-based}) fixa a população de unidades observadas e trata como fonte de incerteza apenas o sorteio de quem recebeu o tratamento. Em experimentos aleatórios com um cadastro fixo de participantes, por exemplo, não há uma ``população maior'' de quem sortear: o que se repete hipoteticamente é a atribuição do tratamento. As duas perspectivas conduzem a fórmulas de erro-padrão distintas, e usar a fórmula errada — mesmo que matematicamente válida em outro contexto — gera inferências incorretas.

O ponto central desta seção: \textbf{um erro-padrão pequeno não implica identificação válida}. O erro-padrão mede com que precisão estimamos o estimando — a quantidade que definimos como objeto de interesse dado os dados. Ele não diz nada sobre se o estimando é igual ao parâmetro econômico que nos importa. Uma estimativa extremamente precisa de uma quantidade sem interpretação causal é, em todos os sentidos relevantes, pior do que uma estimativa imprecisa de algo bem identificado. Essa distinção, aparentemente óbvia, é sistematicamente ignorada na prática.


%% ─────────────────────────────────────────────────────────────────
\section{Erros-Padrão}
\label{sec:erros_padrao}

O \textbf{erro-padrão} de um estimador $\hat{\theta}$ é o desvio-padrão da sua distribuição amostral:
\begin{equation}
  SE(\hat{\theta}) = \sqrt{\text{Var}(\hat{\theta})}.
  \label{eq:se_definicao}
\end{equation}

Em termos práticos, $SE(\hat{\theta})$ quantifica o quanto $\hat{\theta}$ varia de amostra para amostra. Quanto menor o erro-padrão, mais concentrada é a distribuição do estimador em torno do seu valor esperado — e, se o estimador é não-viesado, em torno do parâmetro verdadeiro $\theta$.

\subsection*{Variância do estimador de MQO sob hipóteses clássicas}

No modelo de regressão linear $Y_i = \boldsymbol{x}_i'\boldsymbol{\beta} + \varepsilon_i$, o estimador MQO $\hat{\boldsymbol{\beta}} = (X'X)^{-1}X'Y$ tem variância:
\begin{equation}
  \text{Var}(\hat{\boldsymbol{\beta}} \mid X) = \sigma^2 (X'X)^{-1},
  \label{eq:var_ols_classica}
\end{equation}
onde $\sigma^2 = \text{Var}(\varepsilon_i \mid X)$. Esse resultado requer duas hipóteses fortes: (i) \textbf{homoscedasticidade} — a variância dos erros é constante para todos os valores de $X$ — e (ii) \textbf{independência} dos erros entre observações. Em dados de corte transversal representativos de populações heterogêneas, a homoscedasticidade é frequentemente violada: a variância do salário, por exemplo, tende a crescer com os anos de escolaridade. Quando \eqref{eq:var_ols_classica} é usada nesse contexto, os erros-padrão reportados são inconsistentes.

\subsection*{Erros-padrão robustos à heteroscedasticidade}

A solução padrão para dados de corte transversal são os \textbf{erros-padrão robustos} de Huber-White \citep{Huber1967, White1980}, também chamados de estimador \textit{sandwich}:
\begin{equation}
  \widehat{\text{Var}}_{\text{HC}}(\hat{\boldsymbol{\beta}}) =
    \underbrace{(X'X)^{-1}}_{\text{pão}}
    \underbrace{\left(\sum_{i=1}^{N} x_i x_i' \hat{\varepsilon}_i^2\right)}_{\text{recheio}}
    \underbrace{(X'X)^{-1}}_{\text{pão}},
  \label{eq:sandwich}
\end{equation}
onde $\hat{\varepsilon}_i = Y_i - \boldsymbol{x}_i'\hat{\boldsymbol{\beta}}$ são os resíduos da regressão. O estimador \eqref{eq:sandwich} é consistente independentemente da forma da heteroscedasticidade, sem exigir nenhuma hipótese sobre $\text{Var}(\varepsilon_i \mid X_i)$. Na prática, utiliza-se a versão com correção de amostras finitas HC2 ou HC3, que melhoram o desempenho em amostras pequenas.

\begin{notabox}[Erros robustos como padrão]
Em econometria aplicada moderna, erros-padrão robustos à heteroscedasticidade são o padrão mínimo para dados de corte transversal. A pergunta não é mais ``devo usar erros robustos?'', mas ``que tipo de robustez é necessária?'' (heteroscedasticidade apenas, ou também clustering?). Em R, a função \texttt{estimatr::lm\_robust()} calcula erros HC2 por padrão; o pacote \texttt{sandwich} oferece controle total sobre o tipo de correção.
\end{notabox}

\begin{lstlisting}[style=Rstyle, caption={Erros-padrão clássicos vs.\ robustos em R}]
# -----------------------------
# 1) Configuração inicial
# -----------------------------
rm(list = ls())

library(estimatr)   # lm_robust
library(sandwich)   # vcovHC
library(lmtest)     # coeftest

set.seed(42)
n <- 500

# Simular dados com heteroscedasticidade
edu   <- runif(n, 4, 18)
sigma <- 0.5 * edu          # variância cresce com escolaridade
salario <- 800 + 120 * edu + rnorm(n, sd = sigma * 300)

dados <- data.frame(salario, edu)

# -----------------------------
# 2) MQO com erros clássicos
# -----------------------------
mod_classico <- lm(salario ~ edu, data = dados)
cat("=== Erros clássicos ===\n")
print(coeftest(mod_classico)[, 1:2])

# -----------------------------
# 3) MQO com erros HC2 (robustos)
# -----------------------------
mod_robusto <- lm_robust(salario ~ edu, data = dados, se_type = "HC2")
cat("\n=== Erros HC2 robustos ===\n")
print(summary(mod_robusto)$coefficients[, 1:2])

# -----------------------------
# 4) Comparação direta
# -----------------------------
se_classico <- sqrt(diag(vcov(mod_classico)))["edu"]
se_robusto  <- sqrt(diag(vcovHC(mod_classico, type = "HC2")))["edu"]

cat("\nSE clássico (edu): ", round(se_classico, 3), "\n")
cat("SE HC2 robusto (edu):", round(se_robusto, 3), "\n")
cat("Razão HC2/clássico:  ", round(se_robusto / se_classico, 2), "\n")
\end{lstlisting}

\subsection*{O que o erro-padrão captura — e o que não captura}

Há uma confusão recorrente na literatura aplicada: tratar um erro-padrão pequeno como evidência de que ``a estimativa está correta''. A Tabela~\ref{tab:se_captura} sintetiza o que o erro-padrão realmente mede.

\begin{table}[H]
\centering
\caption{O que o erro-padrão captura e o que ele não captura}
\label{tab:se_captura}
\begin{tabular}{ll}
\toprule
\textbf{O SE captura} & \textbf{O SE não captura} \\
\midrule
Variabilidade amostral de $\hat{\theta}$ & Viés de identificação \\
Incerteza sobre a estimativa do estimando & Viés de estimação em amostras finitas \\
Reduz com aumento de $n$ & Erro de forma funcional \\
Melhora com amostragem mais eficiente & Viés por variável omitida \\
& Violação do SUTVA ou da CIA \\
\bottomrule
\end{tabular}

\vspace{4pt}
\figmeta{Elaboração dos autores}{}{SE: erro-padrão; CIA: hipótese de independência condicional; SUTVA: Stable Unit Treatment Value Assumption.}
\end{table}

Será que um erro-padrão pequeno garante que nossa estimativa está correta? Claramente não. Com uma amostra suficientemente grande, podemos estimar com extrema precisão a diferença de salários entre graduados e não-graduados — mas se essa diferença refletir diferenças de habilidade inata (viés de seleção), estimamos com precisão a coisa errada. Um SE pequeno com identificação fraca não apenas não impressiona: é potencialmente mais enganoso do que um SE grande com identificação sólida.


%% ─────────────────────────────────────────────────────────────────
\section{Teste de Hipótese}
\label{sec:teste_hipotese}

Dado um estimador $\hat{\theta}$ com erro-padrão $SE(\hat{\theta})$, a segunda pergunta natural é: os dados são compatíveis com a hipótese de que $\theta$ assume um valor particular $\theta_0$? O \textbf{teste de hipótese} oferece um procedimento formal para responder a essa pergunta. Apresentamos a lógica em cinco passos.

\medskip
\noindent\textbf{Passo 1 — Hipótese nula.} Definimos a hipótese nula $H_0: \theta = \theta_0$. Na maioria das aplicações em avaliação de impacto, $\theta_0 = 0$ (``o tratamento não tem efeito''), mas qualquer valor é válido.

\noindent\textbf{Passo 2 — Hipótese alternativa.} A hipótese alternativa $H_1: \theta \neq \theta_0$ especifica o que acreditamos ser verdadeiro caso $H_0$ seja falsa. Testes bicaudais são padrão; testes unicaudais ($H_1: \theta > \theta_0$) são admissíveis quando há razão teórica clara para a direção do efeito.

\noindent\textbf{Passo 3 — Estatística de teste.} Construímos a estatística:
\begin{equation}
  t = \frac{\hat{\theta} - \theta_0}{SE(\hat{\theta})},
  \label{eq:estatistica_t}
\end{equation}
que mede em quantos erros-padrão $\hat{\theta}$ se distancia de $\theta_0$. Sob $H_0$ e condições de regularidade, $t$ tem distribuição assintótica $\mathcal{N}(0,1)$ (ou $t_{N-K}$ exata se os erros forem normais).

\noindent\textbf{Passo 4 — P-valor.} O p-valor é a probabilidade de observar, sob $H_0$, uma estatística de teste tão ou mais extrema do que a observada:
\begin{equation}
  \text{p-valor} = P\!\left(|T| \geq |t_{\text{obs}}| \mid H_0\right),
  \label{eq:pvalor}
\end{equation}
onde $T$ é uma variável aleatória com a distribuição de $t$ sob $H_0$.

\noindent\textbf{Passo 5 — Regra de decisão.} Rejeitamos $H_0$ ao nível de significância $\alpha$ se p-valor $< \alpha$. O nível $\alpha = 0{,}05$ é convencional, mas arbitrário — corresponde a aceitar uma taxa de 5\% de falsos positivos quando $H_0$ é verdadeira.

\medskip
A Figura~\ref{fig:dist_h0} ilustra a lógica graficamente.

% FIGURA 4.1: Distribuição da estatística de teste sob H0
% Curva normal centrada em zero, regiões de rejeição sombreadas em cada cauda (|t| > 1,96),
% seta indicando t_obs = 2,35. Eixo x: estatística de teste. Eixo y: densidade.

\begin{figure}[H]
\centering
\begin{tikzpicture}
\begin{axis}[
  height  = 5.5cm,
  width   = 11cm,
  xmin    = -4.2, xmax = 4.2,
  ymin    = 0,    ymax = 0.46,
  xlabel  = {Estatística de teste $t$},
  ylabel  = {Densidade},
  axis lines = left,
  xtick  = {-1.96, 0, 1.96},
  xticklabels = {$-1{,}96$, $0$, $1{,}96$},
  ytick  = \empty,
  tick label style = {font=\small},
  label style      = {font=\small},
  clip   = false,
]
  %% Região de rejeição esquerda
  \addplot[fill=clearred!35, draw=none, domain=-4.2:-1.96, samples=80]
    {exp(-x^2/2)/sqrt(2*3.14159)} \closedcycle;
  %% Região de rejeição direita
  \addplot[fill=clearred!35, draw=none, domain=1.96:4.2, samples=80]
    {exp(-x^2/2)/sqrt(2*3.14159)} \closedcycle;
  %% Curva principal
  \addplot[thick, clearblue, domain=-4.2:4.2, samples=120]
    {exp(-x^2/2)/sqrt(2*3.14159)};
  %% Linhas verticais nos valores críticos
  \draw[dashed, gray!60] (axis cs:-1.96, 0) -- (axis cs:-1.96, 0.062);
  \draw[dashed, gray!60] (axis cs:1.96,  0) -- (axis cs:1.96,  0.062);
  %% Seta para t_obs
  \draw[-{Stealth}, thick, clearred]
    (axis cs:2.35, 0.10) -- (axis cs:2.35, 0.015);
  \node[above, font=\small, clearred] at (axis cs:2.35, 0.10)
    {$t_{\text{obs}} = 2{,}35$};
  %% Rótulos das caudas
  \node[font=\scriptsize] at (axis cs:-2.85, 0.018) {$\alpha/2$};
  \node[font=\scriptsize] at (axis cs: 2.85, 0.018) {$\alpha/2$};
  %% Rótulo da região central
  \node[font=\small, gray!60] at (axis cs:0, 0.20) {$H_0$ não rejeitada};
\end{axis}
\end{tikzpicture}
\caption{Distribuição da estatística de teste sob $H_0$. As regiões sombreadas em vermelho
  correspondem à região de rejeição ao nível $\alpha = 0{,}05$ (cada cauda com $\alpha/2 = 2{,}5\%$).
  O valor observado $t_{\text{obs}} = 2{,}35$ cai na região de rejeição ($|t| > 1{,}96$), levando à
  rejeição de $H_0$.}
\label{fig:dist_h0}
\end{figure}

A Figura~\ref{fig:poder} ilustra a relação entre a distribuição sob $H_0$, a distribuição sob uma hipótese alternativa específica com efeito $\delta$, e o conceito de poder do teste.

% FIGURA 4.2: Distribuição sob H0 (cinza) e H1 (azul).
% H0 ~ N(0,1); H1 ~ N(delta, 1) com delta = 2.
% Área azul à direita do valor crítico = poder do teste.
% Área azul à esquerda do valor crítico = erro Tipo II (beta).

\begin{figure}[H]
\centering
\begin{tikzpicture}
\begin{axis}[
  height  = 5.5cm,
  width   = 11cm,
  xmin    = -4.0, xmax = 5.5,
  ymin    = 0,    ymax = 0.46,
  xlabel  = {Estatística de teste $t$},
  ylabel  = {Densidade},
  axis lines = left,
  xtick  = {0, 1.96},
  xticklabels = {$0$, $t_c{=}1{,}96$},
  ytick  = \empty,
  tick label style = {font=\small},
  label style      = {font=\small},
  clip   = false,
  legend style = {at={(0.97,0.97)}, anchor=north east,
                  font=\small, draw=gray!50},
]
  %% Curva H0 (cinza preenchido)
  \addplot[fill=gray!20, draw=none, domain=-4:5.5, samples=100]
    {exp(-x^2/2)/sqrt(2*3.14159)} \closedcycle;
  \addplot[thick, gray!70, domain=-4:5.5, samples=100]
    {exp(-x^2/2)/sqrt(2*3.14159)};

  %% Poder = área sob H1 à direita de t_c
  \addplot[fill=clearblue!35, draw=none, domain=1.96:5.5, samples=80]
    {exp(-(x-2)^2/2)/sqrt(2*3.14159)} \closedcycle;

  %% Erro Tipo II = área sob H1 à esquerda de t_c
  \addplot[fill=clearteal!25, draw=none, domain=-1:1.96, samples=80]
    {exp(-(x-2)^2/2)/sqrt(2*3.14159)} \closedcycle;

  %% Curva H1
  \addplot[thick, clearblue, domain=-1:5.5, samples=100]
    {exp(-(x-2)^2/2)/sqrt(2*3.14159)};

  %% Linha vertical no valor crítico
  \draw[dashed, gray] (axis cs:1.96, 0) -- (axis cs:1.96, 0.40);

  %% Rótulos
  \node[font=\small, gray!80] at (axis cs:-0.7, 0.20) {$H_0$};
  \node[font=\small, clearblue] at (axis cs:3.4, 0.20) {$H_1$};
  \node[font=\small, clearblue] at (axis cs:2.0, -0.025) {$\delta$};
  \node[font=\scriptsize, clearblue!80] at (axis cs:2.9, 0.06)
    {\textbf{Poder}};
  \node[font=\scriptsize, clearteal!80] at (axis cs:1.0, 0.06)
    {Erro Tipo II ($\beta$)};

  \addlegendentry{$H_0: \theta = 0$}
  \addlegendentry{$H_1: \theta = \delta$}
\end{axis}
\end{tikzpicture}
\caption{Distribuição da estatística de teste sob $H_0$ (cinza) e sob $H_1$ com efeito $\delta = 2$
  (azul). A área azul à direita do valor crítico $t_c = 1{,}96$ representa o \textbf{poder do teste}
  — probabilidade de rejeitar $H_0$ quando $H_1$ é verdadeira. A área azul-esverdeada à esquerda
  é a probabilidade de erro Tipo~II ($\beta = 1 - \text{poder}$). Para referência sobre o cálculo
  do poder e do tamanho mínimo de amostra, ver o \textit{Guia de Cálculo de Poder Estatístico}
  da Série Avaliação na Prática.}
\label{fig:poder}
\end{figure}

\subsection*{O que o p-valor diz — e o que não diz}

O p-valor é um dos objetos mais mal interpretados em toda a estatística aplicada. Vale ser rigoroso. O p-valor \textbf{é} a probabilidade de observar, em um mundo em que $H_0$ é verdadeira, uma estatística tão ou mais extrema quanto a que observamos. Ele \textbf{não é}:

\begin{itemize}[leftmargin=*]
  \item A probabilidade de que $H_0$ seja verdadeira.
  \item A probabilidade de que o resultado seja ``devido ao acaso''.
  \item Uma medida do tamanho do efeito ou de sua relevância econômica.
\end{itemize}

Um ponto frequentemente ignorado em avaliações com amostras grandes: com $N$ suficientemente grande, \textit{qualquer} efeito não-nulo, por menor que seja, produz um p-valor arbitrariamente próximo de zero. Um programa que aumenta a renda em R\$\,2 por mês será ``altamente significativo'' ($p < 0{,}001$) com dados de um milhão de pessoas — e ainda assim economicamente irrelevante. Significância estatística não implica significância econômica.

\begin{atencaobox}[Reporte sempre o tamanho do efeito]
Nunca reporte apenas o p-valor. Reporte sempre: (i) a estimativa pontual $\hat{\theta}$, (ii) o erro-padrão ou o intervalo de confiança, e (iii) uma discussão da relevância econômica do efeito estimado, incluindo comparação com benchmarks relevantes (efeito de políticas similares, custo do programa, magnitude esperada pela teoria). Um p-valor sozinho não permite avaliar se o resultado importa.
\end{atencaobox}

\subsection*{Estatística $F$ para hipóteses conjuntas}

Quando $H_0$ envolve múltiplas restrições simultâneas — por exemplo, $H_0: \beta_1 = \beta_2 = \cdots = \beta_k = 0$ —, testar cada coeficiente individualmente com estatísticas-$t$ separadas infla a taxa de falsos positivos. A solução é a \textbf{estatística $F$}:
\begin{equation}
  F = \frac{(R\hat{\boldsymbol{\beta}} - r)'[R(X'X)^{-1}R']^{-1}(R\hat{\boldsymbol{\beta}} - r)/q}%
           {\hat{\sigma}^2},
  \label{eq:estatistica_F}
\end{equation}
onde $R\boldsymbol{\beta} = r$ codifica as $q$ restrições da hipótese nula. Sob $H_0$ e erros normais, $F \sim F(q, N-K)$.

A estatística $F$ aparece em três contextos recorrentes em avaliação de impacto:

\begin{enumerate}[leftmargin=*, label=(\roman*)]
  \item \textbf{Significância conjunta dos controles}: testar $H_0: \boldsymbol{\gamma} = \mathbf{0}$
    no modelo $Y_i = \alpha + \beta D_i + \boldsymbol{\gamma}' X_i + \varepsilon_i$.
  \item \textbf{Força do instrumento em IV}: a estatística $F$ do primeiro estágio mede a relevância do instrumento. A regra de bolso clássica é $F < 10$ como sinal de instrumento fraco \citep{StockYogo2005}; limiares mais rígidos são discutidos em trabalhos recentes.
  \item \textbf{Verificação de equilíbrio em experimentos}: $H_0$ de que as covariadas pré-tratamento não predizem o tratamento. Rejeição é evidência de falha na randomização.
\end{enumerate}


%% ─── 4.2.1 Múltiplos Testes ──────────────────────────────────────
\subsection{Múltiplos Testes}
\label{sec:multiplos_testes}

Um corolário muitas vezes subestimado do raciocínio da seção anterior: quando realizamos múltiplos testes simultaneamente, a probabilidade de ao menos um falso positivo cresce rapidamente, mesmo que cada teste individual seja conduzido corretamente ao nível $\alpha$. Se realizamos $m$ testes independentes, a probabilidade de ao menos uma rejeição espúria sob $H_0$ é:
\begin{equation}
  P(\text{ao menos um falso positivo}) = 1 - (1 - \alpha)^m.
  \label{eq:multiplos_testes}
\end{equation}

Com $m = 20$ testes ao nível $\alpha = 0{,}05$, essa probabilidade já atinge $1 - 0{,}95^{20} \approx 64\%$. A Tabela~\ref{tab:multiplos_testes} ilustra como ela cresce com $m$.

\begin{table}[H]
\centering
\caption{Probabilidade de ao menos um falso positivo em função do número de testes ($\alpha = 0{,}05$)}
\label{tab:multiplos_testes}
\begin{tabular}{rr}
\toprule
Número de testes ($m$) & $P(\text{ao menos um falso positivo})$ \\
\midrule
 1  & 5{,}0\% \\
 5  & 22{,}6\% \\
10  & 40{,}1\% \\
20  & 64{,}2\% \\
50  & 92{,}3\% \\
100 & 99{,}4\% \\
\bottomrule
\end{tabular}

\vspace{4pt}
\figmeta{Elaboração dos autores}{}{Calculado pela equação~\eqref{eq:multiplos_testes} com testes independentes.}
\end{table}

Dois procedimentos de correção são amplamente utilizados:

\medskip
\noindent\textbf{Correção de Bonferroni.} Para controlar a \textit{taxa de erro familiar} (FWER — probabilidade de ao menos um falso positivo), substitui-se o limiar $\alpha$ por $\alpha/m$. É conservadora — garante controle da FWER mesmo com dependência entre os testes —, mas perde poder quando $m$ é grande, pois cada teste exige evidência muito forte para ser rejeitado.

\medskip
\noindent\textbf{Taxa de Falsas Descobertas (FDR).} O procedimento de Benjamini-Hochberg \citep{BenjaminiHochberg1995} controla a proporção esperada de falsas descobertas entre as rejeições ($FDR = E[\text{falsos positivos}/\text{rejeições}]$), e não o número absoluto. É menos conservador que Bonferroni e preferível quando $m$ é grande e se aceita alguma proporção de erros. O procedimento ordena os $m$ p-valores em ordem crescente $p_{(1)} \leq p_{(2)} \leq \cdots \leq p_{(m)}$ e rejeita todas as hipóteses $H_{(j)}$ em que $p_{(j)} \leq j \cdot \alpha / m$.

\begin{lstlisting}[style=Rstyle, caption={Correção para múltiplos testes em R}]
# -----------------------------
# 1) Simulação: 20 testes sob H0 (todos nulos)
# -----------------------------
rm(list = ls())
set.seed(100)

m  <- 20
n  <- 200

# Gerar 20 p-valores sob H0 verdadeira
pvalores <- replicate(m, {
  Y <- rnorm(n)
  D <- rbinom(n, 1, 0.5)
  coef(summary(lm(Y ~ D)))[2, 4]  # p-valor do tratamento
})

cat("Sem correção:    rejeições =", sum(pvalores < 0.05), "de", m, "\n")

# -----------------------------
# 2) Correção de Bonferroni
# -----------------------------
pvalores_bonf <- p.adjust(pvalores, method = "bonferroni")
cat("Bonferroni:      rejeições =", sum(pvalores_bonf < 0.05), "\n")

# -----------------------------
# 3) Benjamini-Hochberg (FDR)
# -----------------------------
pvalores_bh <- p.adjust(pvalores, method = "BH")
cat("Benjamini-Hochberg: rejeições =", sum(pvalores_bh < 0.05), "\n")

# -----------------------------
# 4) Probabilidade teórica vs. simulada
# -----------------------------
set.seed(1)
sim_prob <- mean(replicate(10000, any(runif(m) < 0.05)))
cat("\nP(ao menos 1 falso positivo) — simulada: ", round(sim_prob, 3), "\n")
cat("P(ao menos 1 falso positivo) — teórica:  ",
    round(1 - 0.95^m, 3), "\n")
\end{lstlisting}

No entanto, correções estatísticas apenas não resolvem o problema estrutural do \textit{p-hacking}: quando o pesquisador realiza muitos testes e reporta apenas os ``significativos'', os p-valores não têm mais a cobertura nominal. A melhor prática é o \textbf{pré-registro}: declarar previamente, em repositório público (e.g., AEA RCT Registry, OSF), quais testes serão realizados, qual é o resultado primário da avaliação e quais são secundários. Independentemente do procedimento de correção utilizado, o relato deve sempre indicar quantos testes foram realizados.


%% ─────────────────────────────────────────────────────────────────
\section{Intervalos de Confiança}
\label{sec:intervalos}

O teste de hipótese responde a uma pergunta binária: os dados são compatíveis com $\theta = \theta_0$? O \textbf{intervalo de confiança} é mais informativo: ele identifica o conjunto de todos os valores $\theta_0$ que \textit{não seriam} rejeitados por um teste ao nível $\alpha$. Para estimadores assintoticamente normais, o intervalo de confiança de nível $1-\alpha$ é:

\begin{equation}
  IC_{1-\alpha} = \left[
    \hat{\theta} - z_{\alpha/2} \cdot SE(\hat{\theta}),\;\;
    \hat{\theta} + z_{\alpha/2} \cdot SE(\hat{\theta})
  \right],
  \label{eq:ic}
\end{equation}

onde $z_{\alpha/2}$ é o quantil $(1-\alpha/2)$ da distribuição normal padrão. Para $\alpha = 0{,}05$: $z_{0{,}025} = 1{,}96$.

\subsection*{Interpretação frequentista correta}

A interpretação do intervalo de confiança merece atenção especial, pois é uma das fontes mais comuns de confusão em estatística aplicada. O intervalo de confiança \textit{não} é uma afirmação sobre a probabilidade de $\theta$ estar em um intervalo específico: $\theta$ é um número fixo (desconhecido), e portanto não tem distribuição. A interpretação correta é frequentista: se o mesmo procedimento de construção do intervalo fosse repetido em muitas amostras independentes da mesma população, 95\% dos intervalos resultantes conteriam o valor verdadeiro $\theta$.

Em outras palavras, a aleatoriedade está no intervalo, não no parâmetro. A frase ``existe 95\% de probabilidade de que $\theta$ esteja entre $a$ e $b$'' está formalmente errada; o que é correto é ``o intervalo $[a, b]$ foi produzido por um procedimento que acerta em 95\% das amostras''.

\subsection*{Dualidade com o teste de hipótese}

Há uma correspondência exata entre intervalos de confiança e testes de hipótese: \textbf{rejeitar $H_0: \theta = \theta_0$ ao nível $\alpha$ é equivalente a $\theta_0$ não pertencer ao IC de nível $1-\alpha$}. Essa dualidade não é apenas uma curiosidade matemática — ela tem implicação prática direta.

O intervalo de confiança carrega estritamente mais informação do que o p-valor: ele diz não apenas se $\theta_0 = 0$ é rejeitado, mas quais outros valores são compatíveis com os dados. Um IC de $[{-0{,}5},\; 12{,}3]$ diz que efeitos nulos \textit{e} efeitos moderados são igualmente compatíveis com a evidência — informação que o p-valor não transmite. Reportar intervalos de confiança deve ser padrão; reportar apenas p-valores é uma prática a ser evitada.

\begin{lstlisting}[style=Rstyle, caption={Intervalos de confiança: construção e simulação da cobertura}]
# -----------------------------
# 1) IC padrão
# -----------------------------
rm(list = ls())
library(estimatr)

set.seed(42)
n <- 400
D <- rbinom(n, 1, 0.5)
Y <- 2 + 0.8 * D + rnorm(n)

mod <- lm_robust(Y ~ D, se_type = "HC2")
cat("Estimativa: ", round(coef(mod)["D"], 3), "\n")
cat("IC 95%:    [", round(confint(mod)["D", 1], 3),
               ";", round(confint(mod)["D", 2], 3), "]\n")

# -----------------------------
# 2) Verificar cobertura empírica do IC
# Deveria ser ~95% se o procedimento está correto
# -----------------------------
set.seed(99)
theta_verdadeiro <- 0.8
B <- 5000

cobertura <- replicate(B, {
  D_s <- rbinom(n, 1, 0.5)
  Y_s <- 2 + theta_verdadeiro * D_s + rnorm(n)
  m_s <- lm_robust(Y_s ~ D_s, se_type = "HC2")
  ic  <- confint(m_s)["D_s", ]
  ic[1] <= theta_verdadeiro & theta_verdadeiro <= ic[2]
})

cat("\nCobertura empírica do IC 95%:", round(mean(cobertura) * 100, 1), "%\n")
cat("(Esperado: 95%)\n")
\end{lstlisting}


%% ─────────────────────────────────────────────────────────────────
\section{Clustering}
\label{sec:clustering}

Na seção anterior, os erros-padrão robustos de Huber-White corrigem a heteroscedasticidade, mas ainda assumem independência entre as observações. Em muitas aplicações em políticas públicas essa hipótese é violada: estudantes de uma mesma escola enfrentam o mesmo professor e a mesma infraestrutura; municípios de um mesmo estado compartilham choques econômicos e políticas estaduais. Ignorar essa dependência dentro de grupos faz com que os erros-padrão clássicos \textit{subestimem} a variabilidade real de $\hat{\beta}$, inflando artificialmente a significância das estimativas.

O problema pode ser descrito com precisão. Suponha que as observações sejam agrupadas em $G$ clusters $g = 1, \ldots, G$, com $n_g$ observações no cluster $g$ ($N = \sum_g n_g$). O modelo é:
\begin{equation}
  Y_{ig} = \boldsymbol{x}_{ig}'\boldsymbol{\beta} + \varepsilon_{ig},
  \label{eq:modelo_cluster}
\end{equation}
onde $\varepsilon_{ig}$ e $\varepsilon_{jg}$ podem ser correlacionados (mesmo cluster $g$), mas $\varepsilon_{ig}$ e $\varepsilon_{jh}$ são independentes para $g \neq h$. O estimador MQO $\hat{\boldsymbol{\beta}}$ permanece não-viesado e consistente sob \eqref{eq:modelo_cluster}, mas a variância clássica $(X'X)^{-1}\hat{\sigma}^2$ é inconsistente. A alternativa consistente é o \textbf{estimador \textit{sandwich} clusterizado}:
\begin{equation}
  \widehat{\text{Var}}_{\text{CL}}(\hat{\boldsymbol{\beta}}) =
    (X'X)^{-1}
    \underbrace{\left(\sum_{g=1}^{G} B_g B_g'\right)}_{\text{``carne'' clusterizada}}
    (X'X)^{-1},
  \quad
  B_g = \sum_{i \in g} \boldsymbol{x}_{ig}\, \hat{\varepsilon}_{ig}.
  \label{eq:sandwich_cl}
\end{equation}

A correção de amostras finitas multiplica \eqref{eq:sandwich_cl} por $\frac{G}{G-1} \cdot \frac{N-1}{N-K}$, análoga à divisão por $N-K$ (e não $N$) nos erros clássicos.

\subsection*{A regra central: clusterizar no nível de atribuição do tratamento}

A regra mais importante do \textit{clustering} é independente dos dados e da forma do modelo: \textbf{o cluster deve ser definido no nível em que o tratamento foi atribuído}. Se uma política é implementada no nível do município e os dados são individuais (pessoas dentro de municípios), deve-se clusterizar por município — independentemente de quantas observações individuais existam. Usar erros robustos sem cluster nessa situação é equivalente a agir como se os 50.000 moradores de um município fossem 50.000 informações independentes sobre o efeito da política, quando na realidade todos compartilham exatamente a mesma exposição ao tratamento.

\begin{table}[H]
\centering
\caption{Quando clusterizar: critérios práticos}
\label{tab:quando_clusterizar}
\begin{tabular}{p{6.2cm} p{6.0cm}}
\toprule
\textbf{Situação} & \textbf{Nível de cluster} \\
\midrule
Política atribuída por escola, dados por aluno    & Escola \\
Política atribuída por município, dados individuais & Município \\
Dados de painel com tratamento variando no tempo  & Unidade (+ período, se bidirecional) \\
Dados de PNAD com estrato de amostragem           & PSU/UPA \\
Experimento com randomização por grupo            & Grupo (mesmo que $G$ seja pequeno) \\
Corte transversal puro sem estrutura de grupo     & Desnecessário (HC suficiente) \\
\bottomrule
\end{tabular}

\vspace{4pt}
\figmeta{Elaboração dos autores}{}{PSU: \textit{Primary Sampling Unit}; UPA: Unidade Primária de Amostragem.}
\end{table}

\begin{lstlisting}[style=Rstyle, caption={Erros clusterizados com sandwich e estimatr}]
# -----------------------------
# 1) Dados com estrutura de cluster
# Política municipal; observações individuais
# -----------------------------
rm(list = ls())

library(estimatr)
library(sandwich)
library(lmtest)

set.seed(77)
G  <- 50   # municípios
ng <- 100  # indivíduos por município

municipio  <- rep(1:G, each = ng)
tratamento <- rep(rbinom(G, 1, 0.5), each = ng)  # atribuição por município
erro_mun   <- rep(rnorm(G, sd = 2), each = ng)   # choque comum ao município
erro_ind   <- rnorm(G * ng, sd = 1)              # choque individual

Y <- 5 + 1.5 * tratamento + erro_mun + erro_ind

dados_cl <- data.frame(Y, tratamento, municipio)

# -----------------------------
# 2) Comparação de erros-padrão
# -----------------------------
mod_base <- lm(Y ~ tratamento, data = dados_cl)

se_classico <- sqrt(diag(vcov(mod_base)))["tratamento"]
se_hc2      <- sqrt(vcovHC(mod_base, type = "HC2"))["tratamento", "tratamento"]
se_cluster  <- sqrt(
  vcovCL(mod_base, cluster = ~municipio)
)["tratamento", "tratamento"]

cat("SE clássico:   ", round(se_classico, 4), "\n")
cat("SE HC2:        ", round(se_hc2, 4), "\n")
cat("SE clusterizado:", round(se_cluster, 4), "\n")

# Com estimatr (mais direto)
mod_cl <- lm_robust(Y ~ tratamento, clusters = municipio, data = dados_cl,
                    se_type = "CR2")
se_cr2 <- summary(mod_cl)$coefficients["tratamento", "Std. Error"]
cat("\nSE CR2 (estimatr):", round(se_cr2, 4), "\n")
\end{lstlisting}

\subsection*{Problema de poucos clusters}

O estimador sandwich clusterizado é válido assintoticamente em $G \to \infty$, mas tem desempenho pobre em amostras finitas quando $G$ é pequeno. Com $G < 30$--50, ele tende a \textit{subestimar} a variância real, produzindo p-valores espúrios. Algumas alternativas:

\begin{itemize}[leftmargin=*]
  \item \textbf{Distribuição $t(G-1)$}: usar a distribuição $t$ com $G-1$ graus de liberdade em vez da normal assintótica. Solução simples e conservadora, adequada quando $G \geq 10$--15.
  \item \textbf{Correção Bell-McCaffrey}: ajuste de amostras finitas mais sofisticado, implementado em \texttt{estimatr} com \texttt{se\_type = \textquotedbl CR2\textquotedbl}.
  \item \textbf{\textit{Wild cluster bootstrap}}: descrito na próxima seção. Solução preferida quando $G < 20$.
\end{itemize}

\subsection*{Clustering bidirecional}

Em dados de painel, pode haver correlação tanto dentro de unidades ao longo do tempo (dependência serial) quanto dentro de períodos entre unidades (dependência transversal, por exemplo, por choques setoriais). O \textbf{clustering bidirecional} \citep{Cameron2011} clusteriza simultaneamente por unidade $i$ e por período $t$:
\begin{equation}
  \widehat{\text{Var}}_{\text{2D}} = \widehat{\text{Var}}_{i} + \widehat{\text{Var}}_{t} - \widehat{\text{Var}}_{it},
  \label{eq:cluster_2d}
\end{equation}
onde cada componente é o estimador sandwich clusterizado no respectivo nível. Requer número suficiente de clusters em ambas as dimensões.


%% ─────────────────────────────────────────────────────────────────
\section{Bootstrap}
\label{sec:bootstrap}

O \textit{bootstrap} é um método computacional para estimar a distribuição amostral de um estimador sem depender de fórmulas analíticas. A ideia, devida a \citet{Efron1979}, é simples: como a distribuição empírica $\hat{F}_n$ é nossa melhor estimativa da distribuição verdadeira $F$, podemos simular a variabilidade amostral reamostrado os dados observados.

O algoritmo básico tem quatro passos:

\begin{enumerate}[leftmargin=*, label=\textbf{\arabic*.}]
  \item Reamostrar com reposição: sortear $n$ observações da amostra original para formar uma amostra \textit{bootstrap} $\{(Y_i^*, X_i^*)\}_{i=1}^n$.
  \item Calcular o estimador nessa amostra: $\hat{\theta}^* = \hat{\theta}(\{Y_i^*, X_i^*\})$.
  \item Repetir os passos 1--2 por $B$ vezes, obtendo $\{\hat{\theta}^*_b\}_{b=1}^B$.
  \item Usar a distribuição empírica de $\{\hat{\theta}^*_b\}$ como aproximação da distribuição amostral de $\hat{\theta}$.
\end{enumerate}

$B = 999$ ou $B = 1{.}999$ repetições são tipicamente suficientes para ICs e p-valores práticos. O \textbf{IC percentil} de nível $1-\alpha$ é:
\begin{equation}
  IC^{\text{perc}}_{1-\alpha} = \left[\hat{\theta}^*_{(\alpha/2)},\;\; \hat{\theta}^*_{(1-\alpha/2)}\right],
  \label{eq:ic_bootstrap}
\end{equation}
onde $\hat{\theta}^*_{(q)}$ denota o quantil $q$ da distribuição bootstrap. O \textbf{IC \textit{bias-corrected}} (BC) ajusta adicionalmente pelo viés $\hat{\theta}^*_{(0{,}5)} - \hat{\theta}$, melhorando a cobertura quando o estimador é enviesado em amostras finitas.

\subsection*{\textit{Cluster bootstrap} e \textit{wild cluster bootstrap}}

O bootstrap padrão (reamostrar observações individuais) não preserva a estrutura de cluster: ao reamostrar observações aleatoriamente, destroem-se as correlações dentro de grupos. A solução é o \textbf{\textit{cluster bootstrap}}: reamostrar clusters inteiros, não observações individuais. Com $G$ clusters, sorteia-se com reposição $G$ clusters da coleção original, mantendo todas as observações de cada cluster sorteado.

Quando $G$ é muito pequeno (digamos, $G < 20$), mesmo o cluster bootstrap pode ter cobertura insatisfatória. Nesses casos, o \textbf{\textit{wild cluster bootstrap}} \citep{Cameron2008} é a solução preferida. Em vez de reamostrar clusters, ele mantém as observações originais e perturba os resíduos de cada cluster por um fator aleatório $w_g \in \{-1, +1\}$ com igual probabilidade (distribuição de Rademacher):
\begin{equation}
  Y_{ig}^* = \boldsymbol{x}_{ig}'\hat{\boldsymbol{\beta}} + w_g\, \hat{\varepsilon}_{ig},
  \label{eq:wild_bootstrap}
\end{equation}
e então reestima o modelo na amostra assim construída. Com $G$ clusters, existem $2^G$ permutações possíveis dos fatores $w_g$, produzindo uma distribuição de permutação exata quando $G$ é muito pequeno.

\begin{lstlisting}[style=Rstyle, caption={Bootstrap e wild cluster bootstrap em R}]
# -----------------------------
# 1) Bootstrap padrão (sem cluster)
# -----------------------------
rm(list = ls())
set.seed(42)

n     <- 300
D     <- rbinom(n, 1, 0.5)
Y     <- 2 + 1.0 * D + rnorm(n)
dados <- data.frame(Y, D)

B         <- 1999
boot_est  <- numeric(B)

for (b in 1:B) {
  idx          <- sample(n, replace = TRUE)
  boot_est[b]  <- coef(lm(Y ~ D, data = dados[idx, ]))["D"]
}

ic_boot <- quantile(boot_est, c(0.025, 0.975))
cat("IC bootstrap 95%: [", round(ic_boot[1], 3),
                      ";", round(ic_boot[2], 3), "]\n")

# -----------------------------
# 2) Wild cluster bootstrap com fewclusterboot
# -----------------------------
library(fwildclusterboot)  # instalar se necessário

G  <- 20
ng <- 50
municipio  <- rep(1:G, each = ng)
trat_g     <- rep(rbinom(G, 1, 0.5), each = ng)
Y2 <- 3 + 0.8 * trat_g + rep(rnorm(G, sd = 1.5), each = ng) + rnorm(G*ng)
dados2 <- data.frame(Y = Y2, trat = trat_g, mun = factor(municipio))

mod2 <- lm(Y ~ trat, data = dados2)
boot_res <- boottest(
  mod2,
  clustid   = "mun",
  param     = "trat",
  B         = 1999,
  seed      = 7
)
summary(boot_res)
\end{lstlisting}

\begin{atencaobox}[{\textit{Bootstrap} não substitui identificação}]
O bootstrap reduz a dependência de aproximações assintóticas e pode melhorar substancialmente a cobertura dos intervalos de confiança, especialmente com clusters pequenos. No entanto, ele não corrige um estimador inconsistente nem valida uma estratégia de identificação fraca. Se $\hat{\theta}$ converge para um valor diferente de $\theta$ (por viés de seleção, variável omitida etc.), reamostrar os dados produzirá uma distribuição bootstrap concentrada em torno do valor errado — com erros-padrão menores e intervalos de confiança que não cobrem $\theta$.
\end{atencaobox}


%% ─────────────────────────────────────────────────────────────────
\section{Inferência Baseada em Design}
\label{sec:design_based}

Todas as abordagens discutidas até aqui assumem, implicitamente, que as observações são sorteios de uma superpopulação maior. A incerteza vem de podermos ter obtido uma amostra diferente. Há, no entanto, situações em que essa perspectiva não faz sentido: quando utilizamos dados de \textit{todos} os municípios brasileiros, de \textit{todas} as escolas de uma rede, ou dos \textit{todos} os participantes de um experimento em que a lista de participantes foi fixada antes da randomização, não há população maior de que amostramos.

Na \textbf{inferência baseada em \textit{design}} (\textit{design-based inference}), a população de unidades é tratada como \textbf{fixa}. A única fonte de aleatoriedade é o \textbf{processo de atribuição do tratamento} — o sorteio de quem entre as unidades observadas recebeu o tratamento. Em vez de perguntar ``o que aconteceria se sorteássemos outra amostra?'', pergunta-se: ``o que aconteceria se, com estas mesmas unidades, o tratamento tivesse sido atribuído de forma diferente?''

\subsection*{Por que isso importa na prática}

Três razões tornam a perspectiva \textit{design-based} relevante em avaliação de políticas públicas no Brasil:

\begin{enumerate}[leftmargin=*, label=(\roman*)]
  \item \textbf{Dados administrativos censitários}: quando os dados cobrem a população inteira (todos os beneficiários do Bolsa Família, todos os funcionários públicos federais), a noção de ``amostra de uma superpopulação'' é forçada. A aleatoriedade relevante é a da política, não a da amostragem.
  \item \textbf{Experimentos com cadastro fixo}: em RCTs em que a lista de participantes é definida antes da randomização, o processo que gera a incerteza é exatamente o sorteio do tratamento.
  \item \textbf{Validade em amostras finitas}: a inferência por permutação é exata em amostras finitas, sem depender de aproximações assintóticas.
\end{enumerate}

\subsection*{Inferência por permutação}

O procedimento central da inferência \textit{design-based} é a \textbf{inferência por permutação} (\textit{randomization inference}) \citep{Fisher1935}. Ela parte da \textbf{hipótese nula afiada} (\textit{sharp null}):

\begin{hipotesebox}[Sharp null (Fisher)]
Para todo indivíduo $i$: $Y_i(1) = Y_i(0)$, ou seja, o tratamento não afeta nenhum indivíduo.
\end{hipotesebox}

Sob a hipótese nula afiada, os resultados potenciais $Y_i(0)$ e $Y_i(1)$ são iguais para todo $i$ — e portanto os resultados observados seriam os mesmos independentemente de quem recebesse o tratamento. Isso significa que, sob $H_0^{\text{sharp}}$, qualquer permutação dos rótulos de tratamento deveria produzir uma estatística de teste semelhante à observada. O p-valor exato é a fração das permutações que geram uma estatística tão ou mais extrema:
\begin{equation}
  \text{p-valor}_{\text{perm}} = \frac{1}{|\mathcal{W}|}
    \sum_{\mathbf{d} \in \mathcal{W}}
    \mathds{1}\!\left[|T(\mathbf{d})| \geq |T(\mathbf{d}_{\text{obs}})|\right],
  \label{eq:pvalue_perm}
\end{equation}
onde $\mathcal{W}$ é o conjunto de todas as atribuições possíveis e $T(\mathbf{d})$ é a estatística de teste calculada com a atribuição $\mathbf{d}$.

\subsection*{Exemplo numérico}

A Tabela~\ref{tab:permutacao} ilustra o procedimento com seis municípios, três dos quais recebem um programa de transferência de renda ($D_i = 1$). A estatística de teste é a diferença de médias entre tratados e controles.

\begin{table}[H]
\centering
\caption{Dados observados para o exemplo de inferência por permutação}
\label{tab:permutacao}
\begin{tabular}{ccc}
\toprule
Município & Tratamento ($D_i$) & Taxa de pobreza (\%) \\
\midrule
A & 1 & 18 \\
B & 1 & 21 \\
C & 1 & 16 \\
D & 0 & 27 \\
E & 0 & 24 \\
F & 0 & 29 \\
\bottomrule
\end{tabular}

\vspace{4pt}
\figmeta{Elaboração dos autores}{}{Dados hipotéticos. Tratados: municípios A, B, C. Controles: D, E, F.}
\end{table}

A diferença de médias observada é $T_{\text{obs}} = \bar{Y}_1 - \bar{Y}_0 = (18+21+16)/3 - (27+24+29)/3 = 18{,}33 - 26{,}67 = -8{,}33$. Existem $\binom{6}{3} = 20$ possíveis atribuições de 3 tratados entre 6 municípios. Sob $H_0^{\text{sharp}}$, calculamos a diferença de médias para cada uma:

\begin{lstlisting}[style=Rstyle, caption={Inferência por permutação: exemplo com 6 municípios}]
# -----------------------------
# 1) Dados observados
# -----------------------------
rm(list = ls())
library(combinat)

Y <- c(18, 21, 16, 27, 24, 29)   # taxa de pobreza
D <- c( 1,  1,  1,  0,  0,  0)   # tratamento observado
n <- length(Y)
n1 <- sum(D)

T_obs <- mean(Y[D == 1]) - mean(Y[D == 0])
cat("Diferença de médias observada:", round(T_obs, 2), "\n")

# -----------------------------
# 2) Gerar todas as C(6,3) = 20 permutações
# -----------------------------
perms <- combn(n, n1)  # colunas = cada possível conjunto de tratados
T_perm <- apply(perms, 2, function(trat_idx) {
  D_perm <- rep(0, n)
  D_perm[trat_idx] <- 1
  mean(Y[D_perm == 1]) - mean(Y[D_perm == 0])
})

cat("Distribuição de permutação:\n")
print(sort(round(T_perm, 2)))

# -----------------------------
# 3) P-valor exato (bilateral)
# -----------------------------
p_perm <- mean(abs(T_perm) >= abs(T_obs))
cat("\np-valor de permutação (bilateral):", round(p_perm, 3), "\n")
cat("(", sum(abs(T_perm) >= abs(T_obs)), "de", ncol(perms), "permutações)\n")
\end{lstlisting}

No exemplo, apenas algumas das 20 permutações produzem uma diferença tão extrema quanto $-8{,}33$. O p-valor exato é a proporção dessas permutações. Note que esse p-valor é válido em amostras finitas, sem qualquer hipótese distribucional sobre os erros.

\begin{notabox}[Sharp null vs.\ weak null]
A hipótese nula afiada ($Y_i(1) = Y_i(0)$ para todo $i$) é mais forte do que a hipótese nula fraca ($\mathbb{E}[Y_i(1) - Y_i(0)] = 0$). A inferência por permutação testa a primeira. Com heterogeneidade de efeitos (alguns positivos, outros negativos com média zero), o teste de permutação pode não rejeitar $H_0^{\text{sharp}}$ mesmo quando a $H_0^{\text{weak}}$ também é falsa. Esse ponto é discutido em \citet{Imbens2015}, Capítulo 5.
\end{notabox}


%% ─────────────────────────────────────────────────────────────────
\section{Testes de Falsificação}
\label{sec:falsificacao}

Em econometria causal, muitos dos testes mais informativos não testam o efeito do tratamento diretamente, mas a \textbf{plausibilidade das hipóteses de identificação}. Se a estratégia empírica requer tendências paralelas, a ausência de tendências divergentes antes do tratamento é evidência (circunstancial) a seu favor. Se a estratégia requer que um instrumento seja exógeno, a ausência de efeito sobre resultados que não deveriam ser afetados é reconfortante. Esses \textbf{testes de falsificação} (\textit{placebo tests}, \textit{falsification tests}) são uma das ferramentas mais valiosas do arsenal do avaliador.

\subsection*{1 — Testes de pré-tendência (DiD)}

No contexto de Diferenças em Diferenças, a identificação requer que o grupo tratado e o grupo controle teriam evoluído em paralelo na ausência do tratamento. Esse contrafactual não é testável diretamente, mas podemos verificar se os grupos evoluíam em paralelo \textit{antes} do tratamento.

O procedimento padrão é o \textbf{estudo de evento} (\textit{event study}): regride-se o resultado sobre dummies de período relativo ao tratamento, separadas por grupo:
\begin{equation}
  Y_{it} = \alpha_i + \lambda_t + \sum_{k \neq -1} \delta_k \cdot D_i \cdot \mathds{1}[t = T_i + k] + \varepsilon_{it},
  \label{eq:event_study}
\end{equation}
onde $k = -1$ é o período de referência (imediatamente antes do tratamento). Os coeficientes $\delta_k$ para $k < 0$ são os \textbf{coeficientes de pré-tendência}: sob tendências paralelas, devem ser estatisticamente indistinguíveis de zero.

% FIGURA 4.3: Gráfico de event study.
% Eixo x: períodos relativos ao tratamento (k = -4, -3, -2, -1, 0, 1, 2, 3, 4)
% Eixo y: coeficiente delta_k com intervalo de confiança de 95%
% Linha vertical em k = -0.5 (cutoff) e linha horizontal em 0
% Coeficientes pré-tratamento (k < 0) próximos de zero; coeficientes pós-tratamento (k >= 0) positivos

\begin{lstlisting}[style=Rstyle, caption={Estudo de evento: coeficientes de pré-tendência}]
# -----------------------------
# 1) Dados de painel simulados com pré-tendências satisfeitas
# -----------------------------
rm(list = ls())
library(dplyr)
library(ggplot2)
library(lfe)  # para efeitos fixos (ou usar fixest)

set.seed(42)
N  <- 100   # unidades
T  <- 10    # períodos (tratamento no período 6)

dados_painel <- expand.grid(id = 1:N, t = 1:T)
dados_painel <- dados_painel %>%
  mutate(
    tratado = id <= 50,        # metade das unidades é tratada
    efeito  = ifelse(tratado & t >= 6, 2.0, 0),
    Y = 5 + 0.3 * t +          # tendência comum
        rnorm(N * T, sd = 1) +  # erro individual
        efeito
  )

# -----------------------------
# 2) Estudo de evento
# -----------------------------
dados_painel <- dados_painel %>%
  mutate(k = t - 6)  # período relativo ao tratamento

# Dummies de interação tratado × período (excluindo k = -1)
library(fixest)
mod_es <- feols(Y ~ i(k, tratado, ref = -1) | id + t,
                data = dados_painel)

# -----------------------------
# 3) Gráfico de event study
# -----------------------------
coefs <- as.data.frame(coef(mod_es))
iplot(mod_es,
      main   = "Estudo de evento: coeficientes por período",
      xlab   = "Período relativo ao tratamento",
      ylab   = expression(delta[k]),
      pt.pch = 16)
abline(h = 0, lty = 2, col = "gray50")
abline(v = -0.5, lty = 3, col = "gray30")
\end{lstlisting}

\begin{conexaobox}[Guia de Diferenças em Diferenças]
A discussão completa sobre tendências paralelas, eventos escalonados (\textit{staggered DiD}) e estimadores robustos à heterogeneidade de efeitos está no \textit{Guia de Diferenças em Diferenças} (2 volumes) da Série Avaliação na Prática. Em particular, o segundo volume trata dos estimadores de Callaway-Sant'Anna, Sun-Abraham e Borusyak-Jaravel-Spiess, que evitam os pesos negativos do estimador TWFE em designs escalonados.
\end{conexaobox}

\subsection*{2 — Placebo de resultado}

Um \textbf{placebo de resultado} aplica a mesma estratégia empírica a um resultado que, pela teoria, \textit{não} deveria ser afetado pelo tratamento. Um efeito estatisticamente significativo é evidência de que a estratégia pode estar captando algo além do tratamento de interesse.

\begin{exemplobox}[Placebo de resultado: Bolsa Família e consumo]
Suponha que se estime o efeito do Bolsa Família sobre o consumo de alimentos usando uma RDD no ponto de corte de elegibilidade (renda per capita = R\$\,218). Como placebo, aplica-se a mesma RDD ao consumo de bens de luxo (viagens internacionais, veículos de luxo) — itens que não deveriam ser afetados por uma transferência de baixo valor para famílias pobres. Um efeito significativo nesse placebo seria evidência de violação da hipótese de continuidade ao redor do corte.
\end{exemplobox}

\subsection*{3 — Placebo de tratamento}

Um \textbf{placebo de tratamento} aplica a estratégia empírica a um período anterior ao tratamento real, como se o tratamento tivesse ocorrido nesse período anterior. Se a estratégia é válida, não deveria haver efeito detectável nesse período contrafactual.

Em um painel com dados de 2010 a 2020 e tratamento em 2016, por exemplo, pode-se estimar o efeito como se o tratamento tivesse ocorrido em 2013, usando apenas dados de 2010 a 2015. Um efeito significativo em 2013 sugere que as diferenças pré-existentes entre grupos, e não o tratamento, explicam o resultado principal.

\subsection*{4 — Teste de balanceamento}

O \textbf{teste de balanceamento} verifica se covariadas pré-determinadas predizem o tratamento. Em um experimento aleatório, nenhuma característica pré-tratamento deveria prever quem recebe o tratamento (a não ser o próprio design de estratificação). Em um estudo observacional, o balanceamento não prova causalidade, mas desequilíbrio acentuado é evidência de seleção endógena.

\begin{lstlisting}[style=Rstyle, caption={Teste de balanceamento via regressão}]
# -----------------------------
# 1) Regredir tratamento em covariadas pré-tratamento
# Um F-teste conjunto testa se alguma covariada prediz D
# -----------------------------
rm(list = ls())
library(estimatr)

set.seed(55)
n <- 500

# Covariadas pré-tratamento
idade   <- rnorm(n, 35, 10)
renda   <- rnorm(n, 3000, 800)
sexo    <- rbinom(n, 1, 0.5)

# Tratamento: (a) aleatório ou (b) com seleção
D_aleat  <- rbinom(n, 1, 0.5)
D_selec  <- rbinom(n, 1, plogis(-1 + 0.03 * renda / 100))

# -----------------------------
# 2) F-teste de balanceamento
# -----------------------------
bal_aleat <- lm(D_aleat ~ idade + renda + sexo)
bal_selec <- lm(D_selec ~ idade + renda + sexo)

cat("=== Tratamento aleatório ===\n")
print(summary(bal_aleat)$fstatistic)

cat("\n=== Tratamento com seleção ===\n")
print(summary(bal_selec)$fstatistic)

# Tabela de diferenças por grupo
for (D_var in list(D_aleat, D_selec)) {
  cat("\nMédias por grupo:\n")
  cat("  Renda tratados:   ", round(mean(renda[D_var == 1])), "\n")
  cat("  Renda controles:  ", round(mean(renda[D_var == 0])), "\n")
}
\end{lstlisting}

\subsection*{Uma advertência obrigatória}

Os testes de falsificação são ferramentas poderosas, mas suas limitações precisam ser enunciadas com clareza. Passar em todos os placebos é evidência \textit{circunstancial} a favor da identificação — é um argumento para a credibilidade, não uma prova. Uma estratégia pode superar todos os testes disponíveis e ainda assim ter identificação inválida por razões não testáveis: um fator de confusão que não está nos dados, uma violação de SUTVA não observada, ou uma descontinuidade nas covariadas que coincide exatamente com o ponto de corte da política. Não confundir ausência de evidência contra a identificação com evidência a favor da identificação. Os testes de falsificação, como toda a inferência causal, devem ser lidos em conjunto com o argumento teórico e o conhecimento substantivo do fenômeno estudado.

\begin{conexaobox}[Testes de falsificação nos guias de métodos]
Cada guia da Série Avaliação na Prática discute os testes de falsificação específicos para o método em questão: o guia de \textit{RDD} trata da continuidade das covariadas e da densidade ao redor do corte (teste de McCrary); o guia de \textit{Matching} discute testes de qualidade do pareamento e equilíbrio; o guia de \textit{Variáveis Instrumentais} trata dos testes de superidentificação (Hansen-Sargan) e de instrumento fraco. Os procedimentos desta seção fornecem o arcabouço geral que fundamenta todos esses testes específicos.
\end{conexaobox}
%%
%%  Pacotes adicionais necessários no preâmbulo principal:
%%    \usepackage{pgfplots}
%%    \pgfplotsset{compat=1.18}
%% ─────────────────────────────────────────────────────────────────

\chapter{Inferência}
\label{cap:inferencia}
\label{sec:inferencia}

Nas seções anteriores construímos uma cadeia que vai da teoria ao dado: definimos o parâmetro de interesse, derivamos as hipóteses de identificação que o tornam observável e escolhemos um estimador $\hat{\theta}$ que converge para o estimando à medida que o tamanho da amostra cresce. Mas obter um número $\hat{\theta}$ não é o fim do trabalho. Se sorteássemos uma nova amostra da mesma população e aplicássemos exatamente o mesmo procedimento, obteríamos um valor diferente. A \textbf{inferência estatística} é o conjunto de ferramentas para raciocinar sobre essa variabilidade amostral e fazer afirmações sobre o parâmetro verdadeiro $\theta$ com base em uma única amostra observada.

É importante distinguir desde o início duas fontes de aleatoriedade que coexistem em econometria aplicada. A \textbf{aleatoriedade de amostragem} (\textit{sampling-based}) trata as observações como realizações independentes e identicamente distribuídas (i.i.d.) de uma população maior: a amostra é um sorteio, e a variabilidade de $\hat{\theta}$ reflete quantas amostras diferentes poderiam ter sido coletadas. Já a \textbf{aleatoriedade de atribuição} (\textit{design-based}) fixa a população de unidades observadas e trata como fonte de incerteza apenas o sorteio de quem recebeu o tratamento. Em experimentos aleatórios com um cadastro fixo de participantes, por exemplo, não há uma ``população maior'' de quem sortear: o que se repete hipoteticamente é a atribuição do tratamento. As duas perspectivas conduzem a fórmulas de erro-padrão distintas, e usar a fórmula errada — mesmo que matematicamente válida em outro contexto — gera inferências incorretas.

O ponto central desta seção: \textbf{um erro-padrão pequeno não implica identificação válida}. O erro-padrão mede com que precisão estimamos o estimando — a quantidade que definimos como objeto de interesse dado os dados. Ele não diz nada sobre se o estimando é igual ao parâmetro econômico que nos importa. Uma estimativa extremamente precisa de uma quantidade sem interpretação causal é, em todos os sentidos relevantes, pior do que uma estimativa imprecisa de algo bem identificado. Essa distinção, aparentemente óbvia, é sistematicamente ignorada na prática.


%% ─────────────────────────────────────────────────────────────────
\section{Erros-Padrão}
\label{sec:erros_padrao}

O \textbf{erro-padrão} de um estimador $\hat{\theta}$ é o desvio-padrão da sua distribuição amostral:
\begin{equation}
  SE(\hat{\theta}) = \sqrt{\text{Var}(\hat{\theta})}.
  \label{eq:se_definicao}
\end{equation}

Em termos práticos, $SE(\hat{\theta})$ quantifica o quanto $\hat{\theta}$ varia de amostra para amostra. Quanto menor o erro-padrão, mais concentrada é a distribuição do estimador em torno do seu valor esperado — e, se o estimador é não-viesado, em torno do parâmetro verdadeiro $\theta$.

\subsection*{Variância do estimador de MQO sob hipóteses clássicas}

No modelo de regressão linear $Y_i = \boldsymbol{x}_i'\boldsymbol{\beta} + \varepsilon_i$, o estimador MQO $\hat{\boldsymbol{\beta}} = (X'X)^{-1}X'Y$ tem variância:
\begin{equation}
  \text{Var}(\hat{\boldsymbol{\beta}} \mid X) = \sigma^2 (X'X)^{-1},
  \label{eq:var_ols_classica}
\end{equation}
onde $\sigma^2 = \text{Var}(\varepsilon_i \mid X)$. Esse resultado requer duas hipóteses fortes: (i) \textbf{homoscedasticidade} — a variância dos erros é constante para todos os valores de $X$ — e (ii) \textbf{independência} dos erros entre observações. Em dados de corte transversal representativos de populações heterogêneas, a homoscedasticidade é frequentemente violada: a variância do salário, por exemplo, tende a crescer com os anos de escolaridade. Quando \eqref{eq:var_ols_classica} é usada nesse contexto, os erros-padrão reportados são inconsistentes.

\subsection*{Erros-padrão robustos à heteroscedasticidade}

A solução padrão para dados de corte transversal são os \textbf{erros-padrão robustos} de Huber-White \citep{Huber1967, White1980}, também chamados de estimador \textit{sandwich}:
\begin{equation}
  \widehat{\text{Var}}_{\text{HC}}(\hat{\boldsymbol{\beta}}) =
    \underbrace{(X'X)^{-1}}_{\text{pão}}
    \underbrace{\left(\sum_{i=1}^{N} x_i x_i' \hat{\varepsilon}_i^2\right)}_{\text{recheio}}
    \underbrace{(X'X)^{-1}}_{\text{pão}},
  \label{eq:sandwich}
\end{equation}
onde $\hat{\varepsilon}_i = Y_i - \boldsymbol{x}_i'\hat{\boldsymbol{\beta}}$ são os resíduos da regressão. O estimador \eqref{eq:sandwich} é consistente independentemente da forma da heteroscedasticidade, sem exigir nenhuma hipótese sobre $\text{Var}(\varepsilon_i \mid X_i)$. Na prática, utiliza-se a versão com correção de amostras finitas HC2 ou HC3, que melhoram o desempenho em amostras pequenas.

\begin{notabox}[Erros robustos como padrão]
Em econometria aplicada moderna, erros-padrão robustos à heteroscedasticidade são o padrão mínimo para dados de corte transversal. A pergunta não é mais ``devo usar erros robustos?'', mas ``que tipo de robustez é necessária?'' (heteroscedasticidade apenas, ou também clustering?). Em R, a função \texttt{estimatr::lm\_robust()} calcula erros HC2 por padrão; o pacote \texttt{sandwich} oferece controle total sobre o tipo de correção.
\end{notabox}

\begin{lstlisting}[style=Rstyle, caption={Erros-padrão clássicos vs.\ robustos em R}]
# -----------------------------
# 1) Configuração inicial
# -----------------------------
rm(list = ls())

library(estimatr)   # lm_robust
library(sandwich)   # vcovHC
library(lmtest)     # coeftest

set.seed(42)
n <- 500

# Simular dados com heteroscedasticidade
edu   <- runif(n, 4, 18)
sigma <- 0.5 * edu          # variância cresce com escolaridade
salario <- 800 + 120 * edu + rnorm(n, sd = sigma * 300)

dados <- data.frame(salario, edu)

# -----------------------------
# 2) MQO com erros clássicos
# -----------------------------
mod_classico <- lm(salario ~ edu, data = dados)
cat("=== Erros clássicos ===\n")
print(coeftest(mod_classico)[, 1:2])

# -----------------------------
# 3) MQO com erros HC2 (robustos)
# -----------------------------
mod_robusto <- lm_robust(salario ~ edu, data = dados, se_type = "HC2")
cat("\n=== Erros HC2 robustos ===\n")
print(summary(mod_robusto)$coefficients[, 1:2])

# -----------------------------
# 4) Comparação direta
# -----------------------------
se_classico <- sqrt(diag(vcov(mod_classico)))["edu"]
se_robusto  <- sqrt(diag(vcovHC(mod_classico, type = "HC2")))["edu"]

cat("\nSE clássico (edu): ", round(se_classico, 3), "\n")
cat("SE HC2 robusto (edu):", round(se_robusto, 3), "\n")
cat("Razão HC2/clássico:  ", round(se_robusto / se_classico, 2), "\n")
\end{lstlisting}

\subsection*{O que o erro-padrão captura — e o que não captura}

Há uma confusão recorrente na literatura aplicada: tratar um erro-padrão pequeno como evidência de que ``a estimativa está correta''. A Tabela~\ref{tab:se_captura} sintetiza o que o erro-padrão realmente mede.

\begin{table}[H]
\centering
\caption{O que o erro-padrão captura e o que ele não captura}
\label{tab:se_captura}
\begin{tabular}{ll}
\toprule
\textbf{O SE captura} & \textbf{O SE não captura} \\
\midrule
Variabilidade amostral de $\hat{\theta}$ & Viés de identificação \\
Incerteza sobre a estimativa do estimando & Viés de estimação em amostras finitas \\
Reduz com aumento de $n$ & Erro de forma funcional \\
Melhora com amostragem mais eficiente & Viés por variável omitida \\
& Violação do SUTVA ou da CIA \\
\bottomrule
\end{tabular}

\vspace{4pt}
\figmeta{Elaboração dos autores}{}{SE: erro-padrão; CIA: hipótese de independência condicional; SUTVA: Stable Unit Treatment Value Assumption.}
\end{table}

Será que um erro-padrão pequeno garante que nossa estimativa está correta? Claramente não. Com uma amostra suficientemente grande, podemos estimar com extrema precisão a diferença de salários entre graduados e não-graduados — mas se essa diferença refletir diferenças de habilidade inata (viés de seleção), estimamos com precisão a coisa errada. Um SE pequeno com identificação fraca não apenas não impressiona: é potencialmente mais enganoso do que um SE grande com identificação sólida.


%% ─────────────────────────────────────────────────────────────────
\section{Teste de Hipótese}
\label{sec:teste_hipotese}

Dado um estimador $\hat{\theta}$ com erro-padrão $SE(\hat{\theta})$, a segunda pergunta natural é: os dados são compatíveis com a hipótese de que $\theta$ assume um valor particular $\theta_0$? O \textbf{teste de hipótese} oferece um procedimento formal para responder a essa pergunta. Apresentamos a lógica em cinco passos.

\medskip
\noindent\textbf{Passo 1 — Hipótese nula.} Definimos a hipótese nula $H_0: \theta = \theta_0$. Na maioria das aplicações em avaliação de impacto, $\theta_0 = 0$ (``o tratamento não tem efeito''), mas qualquer valor é válido.

\noindent\textbf{Passo 2 — Hipótese alternativa.} A hipótese alternativa $H_1: \theta \neq \theta_0$ especifica o que acreditamos ser verdadeiro caso $H_0$ seja falsa. Testes bicaudais são padrão; testes unicaudais ($H_1: \theta > \theta_0$) são admissíveis quando há razão teórica clara para a direção do efeito.

\noindent\textbf{Passo 3 — Estatística de teste.} Construímos a estatística:
\begin{equation}
  t = \frac{\hat{\theta} - \theta_0}{SE(\hat{\theta})},
  \label{eq:estatistica_t}
\end{equation}
que mede em quantos erros-padrão $\hat{\theta}$ se distancia de $\theta_0$. Sob $H_0$ e condições de regularidade, $t$ tem distribuição assintótica $\mathcal{N}(0,1)$ (ou $t_{N-K}$ exata se os erros forem normais).

\noindent\textbf{Passo 4 — P-valor.} O p-valor é a probabilidade de observar, sob $H_0$, uma estatística de teste tão ou mais extrema do que a observada:
\begin{equation}
  \text{p-valor} = P\!\left(|T| \geq |t_{\text{obs}}| \mid H_0\right),
  \label{eq:pvalor}
\end{equation}
onde $T$ é uma variável aleatória com a distribuição de $t$ sob $H_0$.

\noindent\textbf{Passo 5 — Regra de decisão.} Rejeitamos $H_0$ ao nível de significância $\alpha$ se p-valor $< \alpha$. O nível $\alpha = 0{,}05$ é convencional, mas arbitrário — corresponde a aceitar uma taxa de 5\% de falsos positivos quando $H_0$ é verdadeira.

\medskip
A Figura~\ref{fig:dist_h0} ilustra a lógica graficamente.

% FIGURA 4.1: Distribuição da estatística de teste sob H0
% Curva normal centrada em zero, regiões de rejeição sombreadas em cada cauda (|t| > 1,96),
% seta indicando t_obs = 2,35. Eixo x: estatística de teste. Eixo y: densidade.

\begin{figure}[H]
\centering
\begin{tikzpicture}
\begin{axis}[
  height  = 5.5cm,
  width   = 11cm,
  xmin    = -4.2, xmax = 4.2,
  ymin    = 0,    ymax = 0.46,
  xlabel  = {Estatística de teste $t$},
  ylabel  = {Densidade},
  axis lines = left,
  xtick  = {-1.96, 0, 1.96},
  xticklabels = {$-1{,}96$, $0$, $1{,}96$},
  ytick  = \empty,
  tick label style = {font=\small},
  label style      = {font=\small},
  clip   = false,
]
  %% Região de rejeição esquerda
  \addplot[fill=clearred!35, draw=none, domain=-4.2:-1.96, samples=80]
    {exp(-x^2/2)/sqrt(2*3.14159)} \closedcycle;
  %% Região de rejeição direita
  \addplot[fill=clearred!35, draw=none, domain=1.96:4.2, samples=80]
    {exp(-x^2/2)/sqrt(2*3.14159)} \closedcycle;
  %% Curva principal
  \addplot[thick, clearblue, domain=-4.2:4.2, samples=120]
    {exp(-x^2/2)/sqrt(2*3.14159)};
  %% Linhas verticais nos valores críticos
  \draw[dashed, gray!60] (axis cs:-1.96, 0) -- (axis cs:-1.96, 0.062);
  \draw[dashed, gray!60] (axis cs:1.96,  0) -- (axis cs:1.96,  0.062);
  %% Seta para t_obs
  \draw[-{Stealth}, thick, clearred]
    (axis cs:2.35, 0.10) -- (axis cs:2.35, 0.015);
  \node[above, font=\small, clearred] at (axis cs:2.35, 0.10)
    {$t_{\text{obs}} = 2{,}35$};
  %% Rótulos das caudas
  \node[font=\scriptsize] at (axis cs:-2.85, 0.018) {$\alpha/2$};
  \node[font=\scriptsize] at (axis cs: 2.85, 0.018) {$\alpha/2$};
  %% Rótulo da região central
  \node[font=\small, gray!60] at (axis cs:0, 0.20) {$H_0$ não rejeitada};
\end{axis}
\end{tikzpicture}
\caption{Distribuição da estatística de teste sob $H_0$. As regiões sombreadas em vermelho
  correspondem à região de rejeição ao nível $\alpha = 0{,}05$ (cada cauda com $\alpha/2 = 2{,}5\%$).
  O valor observado $t_{\text{obs}} = 2{,}35$ cai na região de rejeição ($|t| > 1{,}96$), levando à
  rejeição de $H_0$.}
\label{fig:dist_h0}
\end{figure}

A Figura~\ref{fig:poder} ilustra a relação entre a distribuição sob $H_0$, a distribuição sob uma hipótese alternativa específica com efeito $\delta$, e o conceito de poder do teste.

% FIGURA 4.2: Distribuição sob H0 (cinza) e H1 (azul).
% H0 ~ N(0,1); H1 ~ N(delta, 1) com delta = 2.
% Área azul à direita do valor crítico = poder do teste.
% Área azul à esquerda do valor crítico = erro Tipo II (beta).

\begin{figure}[H]
\centering
\begin{tikzpicture}
\begin{axis}[
  height  = 5.5cm,
  width   = 11cm,
  xmin    = -4.0, xmax = 5.5,
  ymin    = 0,    ymax = 0.46,
  xlabel  = {Estatística de teste $t$},
  ylabel  = {Densidade},
  axis lines = left,
  xtick  = {0, 1.96},
  xticklabels = {$0$, $t_c{=}1{,}96$},
  ytick  = \empty,
  tick label style = {font=\small},
  label style      = {font=\small},
  clip   = false,
  legend style = {at={(0.97,0.97)}, anchor=north east,
                  font=\small, draw=gray!50},
]
  %% Curva H0 (cinza preenchido)
  \addplot[fill=gray!20, draw=none, domain=-4:5.5, samples=100]
    {exp(-x^2/2)/sqrt(2*3.14159)} \closedcycle;
  \addplot[thick, gray!70, domain=-4:5.5, samples=100]
    {exp(-x^2/2)/sqrt(2*3.14159)};

  %% Poder = área sob H1 à direita de t_c
  \addplot[fill=clearblue!35, draw=none, domain=1.96:5.5, samples=80]
    {exp(-(x-2)^2/2)/sqrt(2*3.14159)} \closedcycle;

  %% Erro Tipo II = área sob H1 à esquerda de t_c
  \addplot[fill=clearteal!25, draw=none, domain=-1:1.96, samples=80]
    {exp(-(x-2)^2/2)/sqrt(2*3.14159)} \closedcycle;

  %% Curva H1
  \addplot[thick, clearblue, domain=-1:5.5, samples=100]
    {exp(-(x-2)^2/2)/sqrt(2*3.14159)};

  %% Linha vertical no valor crítico
  \draw[dashed, gray] (axis cs:1.96, 0) -- (axis cs:1.96, 0.40);

  %% Rótulos
  \node[font=\small, gray!80] at (axis cs:-0.7, 0.20) {$H_0$};
  \node[font=\small, clearblue] at (axis cs:3.4, 0.20) {$H_1$};
  \node[font=\small, clearblue] at (axis cs:2.0, -0.025) {$\delta$};
  \node[font=\scriptsize, clearblue!80] at (axis cs:2.9, 0.06)
    {\textbf{Poder}};
  \node[font=\scriptsize, clearteal!80] at (axis cs:1.0, 0.06)
    {Erro Tipo II ($\beta$)};

  \addlegendentry{$H_0: \theta = 0$}
  \addlegendentry{$H_1: \theta = \delta$}
\end{axis}
\end{tikzpicture}
\caption{Distribuição da estatística de teste sob $H_0$ (cinza) e sob $H_1$ com efeito $\delta = 2$
  (azul). A área azul à direita do valor crítico $t_c = 1{,}96$ representa o \textbf{poder do teste}
  — probabilidade de rejeitar $H_0$ quando $H_1$ é verdadeira. A área azul-esverdeada à esquerda
  é a probabilidade de erro Tipo~II ($\beta = 1 - \text{poder}$). Para referência sobre o cálculo
  do poder e do tamanho mínimo de amostra, ver o \textit{Guia de Cálculo de Poder Estatístico}
  da Série Avaliação na Prática.}
\label{fig:poder}
\end{figure}

\subsection*{O que o p-valor diz — e o que não diz}

O p-valor é um dos objetos mais mal interpretados em toda a estatística aplicada. Vale ser rigoroso. O p-valor \textbf{é} a probabilidade de observar, em um mundo em que $H_0$ é verdadeira, uma estatística tão ou mais extrema quanto a que observamos. Ele \textbf{não é}:

\begin{itemize}[leftmargin=*]
  \item A probabilidade de que $H_0$ seja verdadeira.
  \item A probabilidade de que o resultado seja ``devido ao acaso''.
  \item Uma medida do tamanho do efeito ou de sua relevância econômica.
\end{itemize}

Um ponto frequentemente ignorado em avaliações com amostras grandes: com $N$ suficientemente grande, \textit{qualquer} efeito não-nulo, por menor que seja, produz um p-valor arbitrariamente próximo de zero. Um programa que aumenta a renda em R\$\,2 por mês será ``altamente significativo'' ($p < 0{,}001$) com dados de um milhão de pessoas — e ainda assim economicamente irrelevante. Significância estatística não implica significância econômica.

\begin{atencaobox}[Reporte sempre o tamanho do efeito]
Nunca reporte apenas o p-valor. Reporte sempre: (i) a estimativa pontual $\hat{\theta}$, (ii) o erro-padrão ou o intervalo de confiança, e (iii) uma discussão da relevância econômica do efeito estimado, incluindo comparação com benchmarks relevantes (efeito de políticas similares, custo do programa, magnitude esperada pela teoria). Um p-valor sozinho não permite avaliar se o resultado importa.
\end{atencaobox}

\subsection*{Estatística $F$ para hipóteses conjuntas}

Quando $H_0$ envolve múltiplas restrições simultâneas — por exemplo, $H_0: \beta_1 = \beta_2 = \cdots = \beta_k = 0$ —, testar cada coeficiente individualmente com estatísticas-$t$ separadas infla a taxa de falsos positivos. A solução é a \textbf{estatística $F$}:
\begin{equation}
  F = \frac{(R\hat{\boldsymbol{\beta}} - r)'[R(X'X)^{-1}R']^{-1}(R\hat{\boldsymbol{\beta}} - r)/q}%
           {\hat{\sigma}^2},
  \label{eq:estatistica_F}
\end{equation}
onde $R\boldsymbol{\beta} = r$ codifica as $q$ restrições da hipótese nula. Sob $H_0$ e erros normais, $F \sim F(q, N-K)$.

A estatística $F$ aparece em três contextos recorrentes em avaliação de impacto:

\begin{enumerate}[leftmargin=*, label=(\roman*)]
  \item \textbf{Significância conjunta dos controles}: testar $H_0: \boldsymbol{\gamma} = \mathbf{0}$
    no modelo $Y_i = \alpha + \beta D_i + \boldsymbol{\gamma}' X_i + \varepsilon_i$.
  \item \textbf{Força do instrumento em IV}: a estatística $F$ do primeiro estágio mede a relevância do instrumento. A regra de bolso clássica é $F < 10$ como sinal de instrumento fraco \citep{StockYogo2005}; limiares mais rígidos são discutidos em trabalhos recentes.
  \item \textbf{Verificação de equilíbrio em experimentos}: $H_0$ de que as covariadas pré-tratamento não predizem o tratamento. Rejeição é evidência de falha na randomização.
\end{enumerate}


%% ─── 4.2.1 Múltiplos Testes ──────────────────────────────────────
\subsection{Múltiplos Testes}
\label{sec:multiplos_testes}

Um corolário muitas vezes subestimado do raciocínio da seção anterior: quando realizamos múltiplos testes simultaneamente, a probabilidade de ao menos um falso positivo cresce rapidamente, mesmo que cada teste individual seja conduzido corretamente ao nível $\alpha$. Se realizamos $m$ testes independentes, a probabilidade de ao menos uma rejeição espúria sob $H_0$ é:
\begin{equation}
  P(\text{ao menos um falso positivo}) = 1 - (1 - \alpha)^m.
  \label{eq:multiplos_testes}
\end{equation}

Com $m = 20$ testes ao nível $\alpha = 0{,}05$, essa probabilidade já atinge $1 - 0{,}95^{20} \approx 64\%$. A Tabela~\ref{tab:multiplos_testes} ilustra como ela cresce com $m$.

\begin{table}[H]
\centering
\caption{Probabilidade de ao menos um falso positivo em função do número de testes ($\alpha = 0{,}05$)}
\label{tab:multiplos_testes}
\begin{tabular}{rr}
\toprule
Número de testes ($m$) & $P(\text{ao menos um falso positivo})$ \\
\midrule
 1  & 5{,}0\% \\
 5  & 22{,}6\% \\
10  & 40{,}1\% \\
20  & 64{,}2\% \\
50  & 92{,}3\% \\
100 & 99{,}4\% \\
\bottomrule
\end{tabular}

\vspace{4pt}
\figmeta{Elaboração dos autores}{}{Calculado pela equação~\eqref{eq:multiplos_testes} com testes independentes.}
\end{table}

Dois procedimentos de correção são amplamente utilizados:

\medskip
\noindent\textbf{Correção de Bonferroni.} Para controlar a \textit{taxa de erro familiar} (FWER — probabilidade de ao menos um falso positivo), substitui-se o limiar $\alpha$ por $\alpha/m$. É conservadora — garante controle da FWER mesmo com dependência entre os testes —, mas perde poder quando $m$ é grande, pois cada teste exige evidência muito forte para ser rejeitado.

\medskip
\noindent\textbf{Taxa de Falsas Descobertas (FDR).} O procedimento de Benjamini-Hochberg \citep{BenjaminiHochberg1995} controla a proporção esperada de falsas descobertas entre as rejeições ($FDR = E[\text{falsos positivos}/\text{rejeições}]$), e não o número absoluto. É menos conservador que Bonferroni e preferível quando $m$ é grande e se aceita alguma proporção de erros. O procedimento ordena os $m$ p-valores em ordem crescente $p_{(1)} \leq p_{(2)} \leq \cdots \leq p_{(m)}$ e rejeita todas as hipóteses $H_{(j)}$ em que $p_{(j)} \leq j \cdot \alpha / m$.

\begin{lstlisting}[style=Rstyle, caption={Correção para múltiplos testes em R}]
# -----------------------------
# 1) Simulação: 20 testes sob H0 (todos nulos)
# -----------------------------
rm(list = ls())
set.seed(100)

m  <- 20
n  <- 200

# Gerar 20 p-valores sob H0 verdadeira
pvalores <- replicate(m, {
  Y <- rnorm(n)
  D <- rbinom(n, 1, 0.5)
  coef(summary(lm(Y ~ D)))[2, 4]  # p-valor do tratamento
})

cat("Sem correção:    rejeições =", sum(pvalores < 0.05), "de", m, "\n")

# -----------------------------
# 2) Correção de Bonferroni
# -----------------------------
pvalores_bonf <- p.adjust(pvalores, method = "bonferroni")
cat("Bonferroni:      rejeições =", sum(pvalores_bonf < 0.05), "\n")

# -----------------------------
# 3) Benjamini-Hochberg (FDR)
# -----------------------------
pvalores_bh <- p.adjust(pvalores, method = "BH")
cat("Benjamini-Hochberg: rejeições =", sum(pvalores_bh < 0.05), "\n")

# -----------------------------
# 4) Probabilidade teórica vs. simulada
# -----------------------------
set.seed(1)
sim_prob <- mean(replicate(10000, any(runif(m) < 0.05)))
cat("\nP(ao menos 1 falso positivo) — simulada: ", round(sim_prob, 3), "\n")
cat("P(ao menos 1 falso positivo) — teórica:  ",
    round(1 - 0.95^m, 3), "\n")
\end{lstlisting}

No entanto, correções estatísticas apenas não resolvem o problema estrutural do \textit{p-hacking}: quando o pesquisador realiza muitos testes e reporta apenas os ``significativos'', os p-valores não têm mais a cobertura nominal. A melhor prática é o \textbf{pré-registro}: declarar previamente, em repositório público (e.g., AEA RCT Registry, OSF), quais testes serão realizados, qual é o resultado primário da avaliação e quais são secundários. Independentemente do procedimento de correção utilizado, o relato deve sempre indicar quantos testes foram realizados.


%% ─────────────────────────────────────────────────────────────────
\section{Intervalos de Confiança}
\label{sec:intervalos}

O teste de hipótese responde a uma pergunta binária: os dados são compatíveis com $\theta = \theta_0$? O \textbf{intervalo de confiança} é mais informativo: ele identifica o conjunto de todos os valores $\theta_0$ que \textit{não seriam} rejeitados por um teste ao nível $\alpha$. Para estimadores assintoticamente normais, o intervalo de confiança de nível $1-\alpha$ é:

\begin{equation}
  IC_{1-\alpha} = \left[
    \hat{\theta} - z_{\alpha/2} \cdot SE(\hat{\theta}),\;\;
    \hat{\theta} + z_{\alpha/2} \cdot SE(\hat{\theta})
  \right],
  \label{eq:ic}
\end{equation}

onde $z_{\alpha/2}$ é o quantil $(1-\alpha/2)$ da distribuição normal padrão. Para $\alpha = 0{,}05$: $z_{0{,}025} = 1{,}96$.

\subsection*{Interpretação frequentista correta}

A interpretação do intervalo de confiança merece atenção especial, pois é uma das fontes mais comuns de confusão em estatística aplicada. O intervalo de confiança \textit{não} é uma afirmação sobre a probabilidade de $\theta$ estar em um intervalo específico: $\theta$ é um número fixo (desconhecido), e portanto não tem distribuição. A interpretação correta é frequentista: se o mesmo procedimento de construção do intervalo fosse repetido em muitas amostras independentes da mesma população, 95\% dos intervalos resultantes conteriam o valor verdadeiro $\theta$.

Em outras palavras, a aleatoriedade está no intervalo, não no parâmetro. A frase ``existe 95\% de probabilidade de que $\theta$ esteja entre $a$ e $b$'' está formalmente errada; o que é correto é ``o intervalo $[a, b]$ foi produzido por um procedimento que acerta em 95\% das amostras''.

\subsection*{Dualidade com o teste de hipótese}

Há uma correspondência exata entre intervalos de confiança e testes de hipótese: \textbf{rejeitar $H_0: \theta = \theta_0$ ao nível $\alpha$ é equivalente a $\theta_0$ não pertencer ao IC de nível $1-\alpha$}. Essa dualidade não é apenas uma curiosidade matemática — ela tem implicação prática direta.

O intervalo de confiança carrega estritamente mais informação do que o p-valor: ele diz não apenas se $\theta_0 = 0$ é rejeitado, mas quais outros valores são compatíveis com os dados. Um IC de $[{-0{,}5},\; 12{,}3]$ diz que efeitos nulos \textit{e} efeitos moderados são igualmente compatíveis com a evidência — informação que o p-valor não transmite. Reportar intervalos de confiança deve ser padrão; reportar apenas p-valores é uma prática a ser evitada.

\begin{lstlisting}[style=Rstyle, caption={Intervalos de confiança: construção e simulação da cobertura}]
# -----------------------------
# 1) IC padrão
# -----------------------------
rm(list = ls())
library(estimatr)

set.seed(42)
n <- 400
D <- rbinom(n, 1, 0.5)
Y <- 2 + 0.8 * D + rnorm(n)

mod <- lm_robust(Y ~ D, se_type = "HC2")
cat("Estimativa: ", round(coef(mod)["D"], 3), "\n")
cat("IC 95%:    [", round(confint(mod)["D", 1], 3),
               ";", round(confint(mod)["D", 2], 3), "]\n")

# -----------------------------
# 2) Verificar cobertura empírica do IC
# Deveria ser ~95% se o procedimento está correto
# -----------------------------
set.seed(99)
theta_verdadeiro <- 0.8
B <- 5000

cobertura <- replicate(B, {
  D_s <- rbinom(n, 1, 0.5)
  Y_s <- 2 + theta_verdadeiro * D_s + rnorm(n)
  m_s <- lm_robust(Y_s ~ D_s, se_type = "HC2")
  ic  <- confint(m_s)["D_s", ]
  ic[1] <= theta_verdadeiro & theta_verdadeiro <= ic[2]
})

cat("\nCobertura empírica do IC 95%:", round(mean(cobertura) * 100, 1), "%\n")
cat("(Esperado: 95%)\n")
\end{lstlisting}


%% ─────────────────────────────────────────────────────────────────
\section{Clustering}
\label{sec:clustering}

Na seção anterior, os erros-padrão robustos de Huber-White corrigem a heteroscedasticidade, mas ainda assumem independência entre as observações. Em muitas aplicações em políticas públicas essa hipótese é violada: estudantes de uma mesma escola enfrentam o mesmo professor e a mesma infraestrutura; municípios de um mesmo estado compartilham choques econômicos e políticas estaduais. Ignorar essa dependência dentro de grupos faz com que os erros-padrão clássicos \textit{subestimem} a variabilidade real de $\hat{\beta}$, inflando artificialmente a significância das estimativas.

O problema pode ser descrito com precisão. Suponha que as observações sejam agrupadas em $G$ clusters $g = 1, \ldots, G$, com $n_g$ observações no cluster $g$ ($N = \sum_g n_g$). O modelo é:
\begin{equation}
  Y_{ig} = \boldsymbol{x}_{ig}'\boldsymbol{\beta} + \varepsilon_{ig},
  \label{eq:modelo_cluster}
\end{equation}
onde $\varepsilon_{ig}$ e $\varepsilon_{jg}$ podem ser correlacionados (mesmo cluster $g$), mas $\varepsilon_{ig}$ e $\varepsilon_{jh}$ são independentes para $g \neq h$. O estimador MQO $\hat{\boldsymbol{\beta}}$ permanece não-viesado e consistente sob \eqref{eq:modelo_cluster}, mas a variância clássica $(X'X)^{-1}\hat{\sigma}^2$ é inconsistente. A alternativa consistente é o \textbf{estimador \textit{sandwich} clusterizado}:
\begin{equation}
  \widehat{\text{Var}}_{\text{CL}}(\hat{\boldsymbol{\beta}}) =
    (X'X)^{-1}
    \underbrace{\left(\sum_{g=1}^{G} B_g B_g'\right)}_{\text{``carne'' clusterizada}}
    (X'X)^{-1},
  \quad
  B_g = \sum_{i \in g} \boldsymbol{x}_{ig}\, \hat{\varepsilon}_{ig}.
  \label{eq:sandwich_cl}
\end{equation}

A correção de amostras finitas multiplica \eqref{eq:sandwich_cl} por $\frac{G}{G-1} \cdot \frac{N-1}{N-K}$, análoga à divisão por $N-K$ (e não $N$) nos erros clássicos.

\subsection*{A regra central: clusterizar no nível de atribuição do tratamento}

A regra mais importante do \textit{clustering} é independente dos dados e da forma do modelo: \textbf{o cluster deve ser definido no nível em que o tratamento foi atribuído}. Se uma política é implementada no nível do município e os dados são individuais (pessoas dentro de municípios), deve-se clusterizar por município — independentemente de quantas observações individuais existam. Usar erros robustos sem cluster nessa situação é equivalente a agir como se os 50.000 moradores de um município fossem 50.000 informações independentes sobre o efeito da política, quando na realidade todos compartilham exatamente a mesma exposição ao tratamento.

\begin{table}[H]
\centering
\caption{Quando clusterizar: critérios práticos}
\label{tab:quando_clusterizar}
\begin{tabular}{p{6.2cm} p{6.0cm}}
\toprule
\textbf{Situação} & \textbf{Nível de cluster} \\
\midrule
Política atribuída por escola, dados por aluno    & Escola \\
Política atribuída por município, dados individuais & Município \\
Dados de painel com tratamento variando no tempo  & Unidade (+ período, se bidirecional) \\
Dados de PNAD com estrato de amostragem           & PSU/UPA \\
Experimento com randomização por grupo            & Grupo (mesmo que $G$ seja pequeno) \\
Corte transversal puro sem estrutura de grupo     & Desnecessário (HC suficiente) \\
\bottomrule
\end{tabular}

\vspace{4pt}
\figmeta{Elaboração dos autores}{}{PSU: \textit{Primary Sampling Unit}; UPA: Unidade Primária de Amostragem.}
\end{table}

\begin{lstlisting}[style=Rstyle, caption={Erros clusterizados com sandwich e estimatr}]
# -----------------------------
# 1) Dados com estrutura de cluster
# Política municipal; observações individuais
# -----------------------------
rm(list = ls())

library(estimatr)
library(sandwich)
library(lmtest)

set.seed(77)
G  <- 50   # municípios
ng <- 100  # indivíduos por município

municipio  <- rep(1:G, each = ng)
tratamento <- rep(rbinom(G, 1, 0.5), each = ng)  # atribuição por município
erro_mun   <- rep(rnorm(G, sd = 2), each = ng)   # choque comum ao município
erro_ind   <- rnorm(G * ng, sd = 1)              # choque individual

Y <- 5 + 1.5 * tratamento + erro_mun + erro_ind

dados_cl <- data.frame(Y, tratamento, municipio)

# -----------------------------
# 2) Comparação de erros-padrão
# -----------------------------
mod_base <- lm(Y ~ tratamento, data = dados_cl)

se_classico <- sqrt(diag(vcov(mod_base)))["tratamento"]
se_hc2      <- sqrt(vcovHC(mod_base, type = "HC2"))["tratamento", "tratamento"]
se_cluster  <- sqrt(
  vcovCL(mod_base, cluster = ~municipio)
)["tratamento", "tratamento"]

cat("SE clássico:   ", round(se_classico, 4), "\n")
cat("SE HC2:        ", round(se_hc2, 4), "\n")
cat("SE clusterizado:", round(se_cluster, 4), "\n")

# Com estimatr (mais direto)
mod_cl <- lm_robust(Y ~ tratamento, clusters = municipio, data = dados_cl,
                    se_type = "CR2")
se_cr2 <- summary(mod_cl)$coefficients["tratamento", "Std. Error"]
cat("\nSE CR2 (estimatr):", round(se_cr2, 4), "\n")
\end{lstlisting}

\subsection*{Problema de poucos clusters}

O estimador sandwich clusterizado é válido assintoticamente em $G \to \infty$, mas tem desempenho pobre em amostras finitas quando $G$ é pequeno. Com $G < 30$--50, ele tende a \textit{subestimar} a variância real, produzindo p-valores espúrios. Algumas alternativas:

\begin{itemize}[leftmargin=*]
  \item \textbf{Distribuição $t(G-1)$}: usar a distribuição $t$ com $G-1$ graus de liberdade em vez da normal assintótica. Solução simples e conservadora, adequada quando $G \geq 10$--15.
  \item \textbf{Correção Bell-McCaffrey}: ajuste de amostras finitas mais sofisticado, implementado em \texttt{estimatr} com \texttt{se\_type = \textquotedbl CR2\textquotedbl}.
  \item \textbf{\textit{Wild cluster bootstrap}}: descrito na próxima seção. Solução preferida quando $G < 20$.
\end{itemize}

\subsection*{Clustering bidirecional}

Em dados de painel, pode haver correlação tanto dentro de unidades ao longo do tempo (dependência serial) quanto dentro de períodos entre unidades (dependência transversal, por exemplo, por choques setoriais). O \textbf{clustering bidirecional} \citep{Cameron2011} clusteriza simultaneamente por unidade $i$ e por período $t$:
\begin{equation}
  \widehat{\text{Var}}_{\text{2D}} = \widehat{\text{Var}}_{i} + \widehat{\text{Var}}_{t} - \widehat{\text{Var}}_{it},
  \label{eq:cluster_2d}
\end{equation}
onde cada componente é o estimador sandwich clusterizado no respectivo nível. Requer número suficiente de clusters em ambas as dimensões.


%% ─────────────────────────────────────────────────────────────────
\section{Bootstrap}
\label{sec:bootstrap}

O \textit{bootstrap} é um método computacional para estimar a distribuição amostral de um estimador sem depender de fórmulas analíticas. A ideia, devida a \citet{Efron1979}, é simples: como a distribuição empírica $\hat{F}_n$ é nossa melhor estimativa da distribuição verdadeira $F$, podemos simular a variabilidade amostral reamostrado os dados observados.

O algoritmo básico tem quatro passos:

\begin{enumerate}[leftmargin=*, label=\textbf{\arabic*.}]
  \item Reamostrar com reposição: sortear $n$ observações da amostra original para formar uma amostra \textit{bootstrap} $\{(Y_i^*, X_i^*)\}_{i=1}^n$.
  \item Calcular o estimador nessa amostra: $\hat{\theta}^* = \hat{\theta}(\{Y_i^*, X_i^*\})$.
  \item Repetir os passos 1--2 por $B$ vezes, obtendo $\{\hat{\theta}^*_b\}_{b=1}^B$.
  \item Usar a distribuição empírica de $\{\hat{\theta}^*_b\}$ como aproximação da distribuição amostral de $\hat{\theta}$.
\end{enumerate}

$B = 999$ ou $B = 1{.}999$ repetições são tipicamente suficientes para ICs e p-valores práticos. O \textbf{IC percentil} de nível $1-\alpha$ é:
\begin{equation}
  IC^{\text{perc}}_{1-\alpha} = \left[\hat{\theta}^*_{(\alpha/2)},\;\; \hat{\theta}^*_{(1-\alpha/2)}\right],
  \label{eq:ic_bootstrap}
\end{equation}
onde $\hat{\theta}^*_{(q)}$ denota o quantil $q$ da distribuição bootstrap. O \textbf{IC \textit{bias-corrected}} (BC) ajusta adicionalmente pelo viés $\hat{\theta}^*_{(0{,}5)} - \hat{\theta}$, melhorando a cobertura quando o estimador é enviesado em amostras finitas.

\subsection*{\textit{Cluster bootstrap} e \textit{wild cluster bootstrap}}

O bootstrap padrão (reamostrar observações individuais) não preserva a estrutura de cluster: ao reamostrar observações aleatoriamente, destroem-se as correlações dentro de grupos. A solução é o \textbf{\textit{cluster bootstrap}}: reamostrar clusters inteiros, não observações individuais. Com $G$ clusters, sorteia-se com reposição $G$ clusters da coleção original, mantendo todas as observações de cada cluster sorteado.

Quando $G$ é muito pequeno (digamos, $G < 20$), mesmo o cluster bootstrap pode ter cobertura insatisfatória. Nesses casos, o \textbf{\textit{wild cluster bootstrap}} \citep{Cameron2008} é a solução preferida. Em vez de reamostrar clusters, ele mantém as observações originais e perturba os resíduos de cada cluster por um fator aleatório $w_g \in \{-1, +1\}$ com igual probabilidade (distribuição de Rademacher):
\begin{equation}
  Y_{ig}^* = \boldsymbol{x}_{ig}'\hat{\boldsymbol{\beta}} + w_g\, \hat{\varepsilon}_{ig},
  \label{eq:wild_bootstrap}
\end{equation}
e então reestima o modelo na amostra assim construída. Com $G$ clusters, existem $2^G$ permutações possíveis dos fatores $w_g$, produzindo uma distribuição de permutação exata quando $G$ é muito pequeno.

\begin{lstlisting}[style=Rstyle, caption={Bootstrap e wild cluster bootstrap em R}]
# -----------------------------
# 1) Bootstrap padrão (sem cluster)
# -----------------------------
rm(list = ls())
set.seed(42)

n     <- 300
D     <- rbinom(n, 1, 0.5)
Y     <- 2 + 1.0 * D + rnorm(n)
dados <- data.frame(Y, D)

B         <- 1999
boot_est  <- numeric(B)

for (b in 1:B) {
  idx          <- sample(n, replace = TRUE)
  boot_est[b]  <- coef(lm(Y ~ D, data = dados[idx, ]))["D"]
}

ic_boot <- quantile(boot_est, c(0.025, 0.975))
cat("IC bootstrap 95%: [", round(ic_boot[1], 3),
                      ";", round(ic_boot[2], 3), "]\n")

# -----------------------------
# 2) Wild cluster bootstrap com fewclusterboot
# -----------------------------
library(fwildclusterboot)  # instalar se necessário

G  <- 20
ng <- 50
municipio  <- rep(1:G, each = ng)
trat_g     <- rep(rbinom(G, 1, 0.5), each = ng)
Y2 <- 3 + 0.8 * trat_g + rep(rnorm(G, sd = 1.5), each = ng) + rnorm(G*ng)
dados2 <- data.frame(Y = Y2, trat = trat_g, mun = factor(municipio))

mod2 <- lm(Y ~ trat, data = dados2)
boot_res <- boottest(
  mod2,
  clustid   = "mun",
  param     = "trat",
  B         = 1999,
  seed      = 7
)
summary(boot_res)
\end{lstlisting}

\begin{atencaobox}[{\textit{Bootstrap} não substitui identificação}]
O bootstrap reduz a dependência de aproximações assintóticas e pode melhorar substancialmente a cobertura dos intervalos de confiança, especialmente com clusters pequenos. No entanto, ele não corrige um estimador inconsistente nem valida uma estratégia de identificação fraca. Se $\hat{\theta}$ converge para um valor diferente de $\theta$ (por viés de seleção, variável omitida etc.), reamostrar os dados produzirá uma distribuição bootstrap concentrada em torno do valor errado — com erros-padrão menores e intervalos de confiança que não cobrem $\theta$.
\end{atencaobox}


%% ─────────────────────────────────────────────────────────────────
\section{Inferência Baseada em Design}
\label{sec:design_based}

Todas as abordagens discutidas até aqui assumem, implicitamente, que as observações são sorteios de uma superpopulação maior. A incerteza vem de podermos ter obtido uma amostra diferente. Há, no entanto, situações em que essa perspectiva não faz sentido: quando utilizamos dados de \textit{todos} os municípios brasileiros, de \textit{todas} as escolas de uma rede, ou dos \textit{todos} os participantes de um experimento em que a lista de participantes foi fixada antes da randomização, não há população maior de que amostramos.

Na \textbf{inferência baseada em \textit{design}} (\textit{design-based inference}), a população de unidades é tratada como \textbf{fixa}. A única fonte de aleatoriedade é o \textbf{processo de atribuição do tratamento} — o sorteio de quem entre as unidades observadas recebeu o tratamento. Em vez de perguntar ``o que aconteceria se sorteássemos outra amostra?'', pergunta-se: ``o que aconteceria se, com estas mesmas unidades, o tratamento tivesse sido atribuído de forma diferente?''

\subsection*{Por que isso importa na prática}

Três razões tornam a perspectiva \textit{design-based} relevante em avaliação de políticas públicas no Brasil:

\begin{enumerate}[leftmargin=*, label=(\roman*)]
  \item \textbf{Dados administrativos censitários}: quando os dados cobrem a população inteira (todos os beneficiários do Bolsa Família, todos os funcionários públicos federais), a noção de ``amostra de uma superpopulação'' é forçada. A aleatoriedade relevante é a da política, não a da amostragem.
  \item \textbf{Experimentos com cadastro fixo}: em RCTs em que a lista de participantes é definida antes da randomização, o processo que gera a incerteza é exatamente o sorteio do tratamento.
  \item \textbf{Validade em amostras finitas}: a inferência por permutação é exata em amostras finitas, sem depender de aproximações assintóticas.
\end{enumerate}

\subsection*{Inferência por permutação}

O procedimento central da inferência \textit{design-based} é a \textbf{inferência por permutação} (\textit{randomization inference}) \citep{Fisher1935}. Ela parte da \textbf{hipótese nula afiada} (\textit{sharp null}):

\begin{hipotesebox}[Sharp null (Fisher)]
Para todo indivíduo $i$: $Y_i(1) = Y_i(0)$, ou seja, o tratamento não afeta nenhum indivíduo.
\end{hipotesebox}

Sob a hipótese nula afiada, os resultados potenciais $Y_i(0)$ e $Y_i(1)$ são iguais para todo $i$ — e portanto os resultados observados seriam os mesmos independentemente de quem recebesse o tratamento. Isso significa que, sob $H_0^{\text{sharp}}$, qualquer permutação dos rótulos de tratamento deveria produzir uma estatística de teste semelhante à observada. O p-valor exato é a fração das permutações que geram uma estatística tão ou mais extrema:
\begin{equation}
  \text{p-valor}_{\text{perm}} = \frac{1}{|\mathcal{W}|}
    \sum_{\mathbf{d} \in \mathcal{W}}
    \mathds{1}\!\left[|T(\mathbf{d})| \geq |T(\mathbf{d}_{\text{obs}})|\right],
  \label{eq:pvalue_perm}
\end{equation}
onde $\mathcal{W}$ é o conjunto de todas as atribuições possíveis e $T(\mathbf{d})$ é a estatística de teste calculada com a atribuição $\mathbf{d}$.

\subsection*{Exemplo numérico}

A Tabela~\ref{tab:permutacao} ilustra o procedimento com seis municípios, três dos quais recebem um programa de transferência de renda ($D_i = 1$). A estatística de teste é a diferença de médias entre tratados e controles.

\begin{table}[H]
\centering
\caption{Dados observados para o exemplo de inferência por permutação}
\label{tab:permutacao}
\begin{tabular}{ccc}
\toprule
Município & Tratamento ($D_i$) & Taxa de pobreza (\%) \\
\midrule
A & 1 & 18 \\
B & 1 & 21 \\
C & 1 & 16 \\
D & 0 & 27 \\
E & 0 & 24 \\
F & 0 & 29 \\
\bottomrule
\end{tabular}

\vspace{4pt}
\figmeta{Elaboração dos autores}{}{Dados hipotéticos. Tratados: municípios A, B, C. Controles: D, E, F.}
\end{table}

A diferença de médias observada é $T_{\text{obs}} = \bar{Y}_1 - \bar{Y}_0 = (18+21+16)/3 - (27+24+29)/3 = 18{,}33 - 26{,}67 = -8{,}33$. Existem $\binom{6}{3} = 20$ possíveis atribuições de 3 tratados entre 6 municípios. Sob $H_0^{\text{sharp}}$, calculamos a diferença de médias para cada uma:

\begin{lstlisting}[style=Rstyle, caption={Inferência por permutação: exemplo com 6 municípios}]
# -----------------------------
# 1) Dados observados
# -----------------------------
rm(list = ls())
library(combinat)

Y <- c(18, 21, 16, 27, 24, 29)   # taxa de pobreza
D <- c( 1,  1,  1,  0,  0,  0)   # tratamento observado
n <- length(Y)
n1 <- sum(D)

T_obs <- mean(Y[D == 1]) - mean(Y[D == 0])
cat("Diferença de médias observada:", round(T_obs, 2), "\n")

# -----------------------------
# 2) Gerar todas as C(6,3) = 20 permutações
# -----------------------------
perms <- combn(n, n1)  # colunas = cada possível conjunto de tratados
T_perm <- apply(perms, 2, function(trat_idx) {
  D_perm <- rep(0, n)
  D_perm[trat_idx] <- 1
  mean(Y[D_perm == 1]) - mean(Y[D_perm == 0])
})

cat("Distribuição de permutação:\n")
print(sort(round(T_perm, 2)))

# -----------------------------
# 3) P-valor exato (bilateral)
# -----------------------------
p_perm <- mean(abs(T_perm) >= abs(T_obs))
cat("\np-valor de permutação (bilateral):", round(p_perm, 3), "\n")
cat("(", sum(abs(T_perm) >= abs(T_obs)), "de", ncol(perms), "permutações)\n")
\end{lstlisting}

No exemplo, apenas algumas das 20 permutações produzem uma diferença tão extrema quanto $-8{,}33$. O p-valor exato é a proporção dessas permutações. Note que esse p-valor é válido em amostras finitas, sem qualquer hipótese distribucional sobre os erros.

\begin{notabox}[Sharp null vs.\ weak null]
A hipótese nula afiada ($Y_i(1) = Y_i(0)$ para todo $i$) é mais forte do que a hipótese nula fraca ($\mathbb{E}[Y_i(1) - Y_i(0)] = 0$). A inferência por permutação testa a primeira. Com heterogeneidade de efeitos (alguns positivos, outros negativos com média zero), o teste de permutação pode não rejeitar $H_0^{\text{sharp}}$ mesmo quando a $H_0^{\text{weak}}$ também é falsa. Esse ponto é discutido em \citet{Imbens2015}, Capítulo 5.
\end{notabox}


%% ─────────────────────────────────────────────────────────────────
\section{Testes de Falsificação}
\label{sec:falsificacao}

Em econometria causal, muitos dos testes mais informativos não testam o efeito do tratamento diretamente, mas a \textbf{plausibilidade das hipóteses de identificação}. Se a estratégia empírica requer tendências paralelas, a ausência de tendências divergentes antes do tratamento é evidência (circunstancial) a seu favor. Se a estratégia requer que um instrumento seja exógeno, a ausência de efeito sobre resultados que não deveriam ser afetados é reconfortante. Esses \textbf{testes de falsificação} (\textit{placebo tests}, \textit{falsification tests}) são uma das ferramentas mais valiosas do arsenal do avaliador.

\subsection*{1 — Testes de pré-tendência (DiD)}

No contexto de Diferenças em Diferenças, a identificação requer que o grupo tratado e o grupo controle teriam evoluído em paralelo na ausência do tratamento. Esse contrafactual não é testável diretamente, mas podemos verificar se os grupos evoluíam em paralelo \textit{antes} do tratamento.

O procedimento padrão é o \textbf{estudo de evento} (\textit{event study}): regride-se o resultado sobre dummies de período relativo ao tratamento, separadas por grupo:
\begin{equation}
  Y_{it} = \alpha_i + \lambda_t + \sum_{k \neq -1} \delta_k \cdot D_i \cdot \mathds{1}[t = T_i + k] + \varepsilon_{it},
  \label{eq:event_study}
\end{equation}
onde $k = -1$ é o período de referência (imediatamente antes do tratamento). Os coeficientes $\delta_k$ para $k < 0$ são os \textbf{coeficientes de pré-tendência}: sob tendências paralelas, devem ser estatisticamente indistinguíveis de zero.

% FIGURA 4.3: Gráfico de event study.
% Eixo x: períodos relativos ao tratamento (k = -4, -3, -2, -1, 0, 1, 2, 3, 4)
% Eixo y: coeficiente delta_k com intervalo de confiança de 95%
% Linha vertical em k = -0.5 (cutoff) e linha horizontal em 0
% Coeficientes pré-tratamento (k < 0) próximos de zero; coeficientes pós-tratamento (k >= 0) positivos

\begin{lstlisting}[style=Rstyle, caption={Estudo de evento: coeficientes de pré-tendência}]
# -----------------------------
# 1) Dados de painel simulados com pré-tendências satisfeitas
# -----------------------------
rm(list = ls())
library(dplyr)
library(ggplot2)
library(lfe)  # para efeitos fixos (ou usar fixest)

set.seed(42)
N  <- 100   # unidades
T  <- 10    # períodos (tratamento no período 6)

dados_painel <- expand.grid(id = 1:N, t = 1:T)
dados_painel <- dados_painel %>%
  mutate(
    tratado = id <= 50,        # metade das unidades é tratada
    efeito  = ifelse(tratado & t >= 6, 2.0, 0),
    Y = 5 + 0.3 * t +          # tendência comum
        rnorm(N * T, sd = 1) +  # erro individual
        efeito
  )

# -----------------------------
# 2) Estudo de evento
# -----------------------------
dados_painel <- dados_painel %>%
  mutate(k = t - 6)  # período relativo ao tratamento

# Dummies de interação tratado × período (excluindo k = -1)
library(fixest)
mod_es <- feols(Y ~ i(k, tratado, ref = -1) | id + t,
                data = dados_painel)

# -----------------------------
# 3) Gráfico de event study
# -----------------------------
coefs <- as.data.frame(coef(mod_es))
iplot(mod_es,
      main   = "Estudo de evento: coeficientes por período",
      xlab   = "Período relativo ao tratamento",
      ylab   = expression(delta[k]),
      pt.pch = 16)
abline(h = 0, lty = 2, col = "gray50")
abline(v = -0.5, lty = 3, col = "gray30")
\end{lstlisting}

\begin{conexaobox}[Guia de Diferenças em Diferenças]
A discussão completa sobre tendências paralelas, eventos escalonados (\textit{staggered DiD}) e estimadores robustos à heterogeneidade de efeitos está no \textit{Guia de Diferenças em Diferenças} (2 volumes) da Série Avaliação na Prática. Em particular, o segundo volume trata dos estimadores de Callaway-Sant'Anna, Sun-Abraham e Borusyak-Jaravel-Spiess, que evitam os pesos negativos do estimador TWFE em designs escalonados.
\end{conexaobox}

\subsection*{2 — Placebo de resultado}

Um \textbf{placebo de resultado} aplica a mesma estratégia empírica a um resultado que, pela teoria, \textit{não} deveria ser afetado pelo tratamento. Um efeito estatisticamente significativo é evidência de que a estratégia pode estar captando algo além do tratamento de interesse.

\begin{exemplobox}[Placebo de resultado: Bolsa Família e consumo]
Suponha que se estime o efeito do Bolsa Família sobre o consumo de alimentos usando uma RDD no ponto de corte de elegibilidade (renda per capita = R\$\,218). Como placebo, aplica-se a mesma RDD ao consumo de bens de luxo (viagens internacionais, veículos de luxo) — itens que não deveriam ser afetados por uma transferência de baixo valor para famílias pobres. Um efeito significativo nesse placebo seria evidência de violação da hipótese de continuidade ao redor do corte.
\end{exemplobox}

\subsection*{3 — Placebo de tratamento}

Um \textbf{placebo de tratamento} aplica a estratégia empírica a um período anterior ao tratamento real, como se o tratamento tivesse ocorrido nesse período anterior. Se a estratégia é válida, não deveria haver efeito detectável nesse período contrafactual.

Em um painel com dados de 2010 a 2020 e tratamento em 2016, por exemplo, pode-se estimar o efeito como se o tratamento tivesse ocorrido em 2013, usando apenas dados de 2010 a 2015. Um efeito significativo em 2013 sugere que as diferenças pré-existentes entre grupos, e não o tratamento, explicam o resultado principal.

\subsection*{4 — Teste de balanceamento}

O \textbf{teste de balanceamento} verifica se covariadas pré-determinadas predizem o tratamento. Em um experimento aleatório, nenhuma característica pré-tratamento deveria prever quem recebe o tratamento (a não ser o próprio design de estratificação). Em um estudo observacional, o balanceamento não prova causalidade, mas desequilíbrio acentuado é evidência de seleção endógena.

\begin{lstlisting}[style=Rstyle, caption={Teste de balanceamento via regressão}]
# -----------------------------
# 1) Regredir tratamento em covariadas pré-tratamento
# Um F-teste conjunto testa se alguma covariada prediz D
# -----------------------------
rm(list = ls())
library(estimatr)

set.seed(55)
n <- 500

# Covariadas pré-tratamento
idade   <- rnorm(n, 35, 10)
renda   <- rnorm(n, 3000, 800)
sexo    <- rbinom(n, 1, 0.5)

# Tratamento: (a) aleatório ou (b) com seleção
D_aleat  <- rbinom(n, 1, 0.5)
D_selec  <- rbinom(n, 1, plogis(-1 + 0.03 * renda / 100))

# -----------------------------
# 2) F-teste de balanceamento
# -----------------------------
bal_aleat <- lm(D_aleat ~ idade + renda + sexo)
bal_selec <- lm(D_selec ~ idade + renda + sexo)

cat("=== Tratamento aleatório ===\n")
print(summary(bal_aleat)$fstatistic)

cat("\n=== Tratamento com seleção ===\n")
print(summary(bal_selec)$fstatistic)

# Tabela de diferenças por grupo
for (D_var in list(D_aleat, D_selec)) {
  cat("\nMédias por grupo:\n")
  cat("  Renda tratados:   ", round(mean(renda[D_var == 1])), "\n")
  cat("  Renda controles:  ", round(mean(renda[D_var == 0])), "\n")
}
\end{lstlisting}

\subsection*{Uma advertência obrigatória}

Os testes de falsificação são ferramentas poderosas, mas suas limitações precisam ser enunciadas com clareza. Passar em todos os placebos é evidência \textit{circunstancial} a favor da identificação — é um argumento para a credibilidade, não uma prova. Uma estratégia pode superar todos os testes disponíveis e ainda assim ter identificação inválida por razões não testáveis: um fator de confusão que não está nos dados, uma violação de SUTVA não observada, ou uma descontinuidade nas covariadas que coincide exatamente com o ponto de corte da política. Não confundir ausência de evidência contra a identificação com evidência a favor da identificação. Os testes de falsificação, como toda a inferência causal, devem ser lidos em conjunto com o argumento teórico e o conhecimento substantivo do fenômeno estudado.

\begin{conexaobox}[Testes de falsificação nos guias de métodos]
Cada guia da Série Avaliação na Prática discute os testes de falsificação específicos para o método em questão: o guia de \textit{RDD} trata da continuidade das covariadas e da densidade ao redor do corte (teste de McCrary); o guia de \textit{Matching} discute testes de qualidade do pareamento e equilíbrio; o guia de \textit{Variáveis Instrumentais} trata dos testes de superidentificação (Hansen-Sargan) e de instrumento fraco. Os procedimentos desta seção fornecem o arcabouço geral que fundamenta todos esses testes específicos.
\end{conexaobox}
%% ─────────────────────────────────────────────────────────────────
%%  Seção 5 — Aplicação Empírica: O Experimento STAR
%%  Arquivo: secao5_aplicacao.tex
%%  Inserir no projeto com: %% ─────────────────────────────────────────────────────────────────
%%  Seção 5 — Aplicação Empírica: O Experimento STAR
%%  Arquivo: secao5_aplicacao.tex
%%  Inserir no projeto com: %% ─────────────────────────────────────────────────────────────────
%%  Seção 5 — Aplicação Empírica: O Experimento STAR
%%  Arquivo: secao5_aplicacao.tex
%%  Inserir no projeto com: \input{secao5_aplicacao}
%%
%%  Dados: AER::STAR (pacote AER do R)
%%  Referência principal: Krueger (1999)
%% ─────────────────────────────────────────────────────────────────

\chapter{Aplicação Empírica: O Experimento STAR}
\label{sec:star}
\label{cap:star}

As seções anteriores construíram o arcabouço conceitual da inferência causal: resultados potenciais, identificação, estimação, erros-padrão, testes de hipótese, bootstrap e falsificação. O objetivo desta seção é usar um experimento clássico como laboratório empírico para ver esses conceitos funcionando juntos. O Experimento STAR --- \textit{Student/Teacher Achievement Ratio} --- oferece um dos melhores exemplos disponíveis de avaliação experimental em políticas educacionais. Ele é suficientemente limpo para ilustrar a lógica da identificação por aleatorização, mas suficientemente rico para revelar as complexidades práticas de inferência, múltiplos testes e falsificação.

O experimento foi conduzido no estado americano do Tennessee entre 1985 e 1989 \citep{Finn1990, Krueger1999}. A pergunta central é direta: turmas menores melhoram o desempenho escolar? Para respondê-la com rigor causal, o estado financiou um experimento em larga escala em que estudantes e professoras foram aleatoriamente designados a três tipos de turma: \textbf{turma pequena} (13--17 alunos), \textbf{turma regular} (22--26 alunos) e \textbf{turma regular com auxiliar} (22--26 alunos com uma professora auxiliar). O acompanhamento ocorreu ao longo de quatro anos, do jardim de infância (pré-escola, \textit{kindergarten}) ao terceiro ano do ensino fundamental (K--3).

O STAR é um caso de referência na literatura de avaliação de impacto por três razões. Primeira, é um dos poucos experimentos aleatórios de grande escala realizados em educação --- campo em que, por razões éticas e logísticas, experimentos são raros. Segunda, sua metodologia está bem documentada, o que permite verificar e replicar cada escolha empírica. Terceira, os dados tornaram-se públicos e foram explorados por dezenas de pesquisadores, criando uma literatura acumulada que serve de referência para discussão metodológica \citep{Chetty2011}. Neste guia, não temos o objetivo de contribuir com nova pesquisa sobre educação: usamos o STAR como plataforma para ilustrar, com um exemplo concreto, os métodos discutidos nas seções anteriores.


%% ─────────────────────────────────────────────────────────────────
\section{Dados e Desenho Experimental}
\label{sec:star_dados}

\subsection*{Estrutura da amostra}

O experimento acompanhou aproximadamente 11.600 alunos distribuídos em cerca de 300 turmas em 79 escolas do Tennessee. As 79 escolas foram selecionadas de forma a representar áreas rurais, suburbanas e urbanas do estado. Dentro de cada escola, alunos e professoras foram aleatoriamente designados a um dos três tipos de turma.

É importante distinguir três unidades de análise que coexistem nos dados:

\begin{itemize}[leftmargin=*]
  \item \textbf{Aluno ($i$)}: unidade de observação. Os dados contêm, para cada aluno, notas em diferentes disciplinas, características demográficas e indicadores de situação socioeconômica.
  \item \textbf{Turma ($c$)}: unidade de tratamento. O tratamento --- tipo de turma --- é atribuído no nível da turma, não do aluno individualmente. Todos os alunos de uma mesma turma recebem exatamente o mesmo tratamento.
  \item \textbf{Escola ($s$)}: unidade de randomização. A aleatorização foi feita \textit{dentro} de cada escola: dentro da escola $s$, alunos foram sorteados para turmas pequenas ou regulares. A composição das turmas \textit{entre} escolas não foi randomizada.
\end{itemize}

Essa estrutura hierárquica --- alunos dentro de turmas dentro de escolas --- tem implicações diretas para a estimação e para a inferência, como veremos ao longo desta seção.

\subsection*{Variáveis disponíveis}

O conjunto de dados contém as seguintes informações por aluno:

\begin{itemize}[leftmargin=*]
  \item \textbf{Notas}: pontuação em leitura e matemática para cada ano (jardim de infância ao terceiro ano). São as variáveis de resultado.
  \item \textbf{Tipo de turma designado}: indicador do braço do experimento ao qual o aluno foi alocado (\texttt{small}, \texttt{regular}, \texttt{regular+aide}).
  \item \textbf{Características do aluno}: raça/etnia, gênero, elegibilidade para merenda gratuita (proxy de renda familiar). Essas são covariadas pré-tratamento.
  \item \textbf{Identificadores}: código da escola e da turma.
  \item \textbf{Características da professora}: anos de experiência e titulação, quando disponíveis. Usadas nos testes de falsificação.
\end{itemize}

\subsection*{Rotina de análise em R}

O script de análise carrega os dados do pacote \texttt{AER} do R \citep{Zeileis2008}, que distribui uma versão dos dados do STAR em formato pronto para análise. O script:

\begin{enumerate}[leftmargin=*, label=(\roman*)]
  \item Carrega os dados e seleciona a amostra analítica do jardim de infância, eliminando observações com informação faltante nas variáveis de interesse.
  \item Cria o indicador binário de tratamento $D_i = \mathds{1}[\texttt{stark} = \texttt{\textquotedbl small\textquotedbl}]$, separando turma pequena das duas condições de controle agrupadas.
  \item Padroniza as notas de leitura e matemática (média zero, desvio-padrão um) para facilitar a interpretação dos coeficientes em desvios-padrão.
  \item Define covariadas pré-tratamento: gênero, raça e elegibilidade para merenda gratuita.
  \item Constrói os objetos usados nas estimações e tabelas das subseções seguintes.
\end{enumerate}

\begin{lstlisting}[style=Rstyle, caption={Carregamento e preparação dos dados STAR}]
# -----------------------------
# 1) Configuração inicial
# -----------------------------
rm(list = ls())

library(AER)       # dados STAR
library(dplyr)     # manipulação
library(fixest)    # efeitos fixos (feols)
library(estimatr)  # lm_robust / erros clusterizados
library(sandwich)  # vcovCL
library(lmtest)    # coeftest

# -----------------------------
# 2) Carregar e filtrar dados
# Foco no jardim de infância (ano k)
# -----------------------------
data("STAR", package = "AER")

dados_k <- STAR %>%
  filter(!is.na(stark), !is.na(readk), !is.na(mathk)) %>%
  mutate(
    # Tratamento binário: turma pequena vs. demais
    D       = as.integer(stark == "small"),
    # Notas padronizadas (desvio-padrão = 1)
    read_z  = scale(readk)[, 1],
    math_z  = scale(mathk)[, 1],
    # Covariadas pré-tratamento
    feminino    = as.integer(gender == "female"),
    nao_branco  = as.integer(ethnicity != "cauc"),
    merenda     = as.integer(lunchk == "free"),
    # Identificadores
    escola  = as.factor(schoolk)
  )

cat("Observações na amostra analítica:", nrow(dados_k), "\n")
cat("Escolas:                         ", n_distinct(dados_k$escola), "\n")
cat("Tratados (turma pequena):        ", sum(dados_k$D), "\n")
cat("Controle:                        ", sum(1 - dados_k$D), "\n")
\end{lstlisting}


%% ─────────────────────────────────────────────────────────────────
\section{Identificação: Por que a Aleatorização Resolve o Problema}
\label{sec:star_identificacao}

\subsection*{O estimando e a estratégia de identificação}

Para cada aluno $i$, definimos dois resultados potenciais:
\begin{align}
  Y_i(1) &\equiv \text{nota de } i \text{ se designado para turma pequena,} \\
  Y_i(0) &\equiv \text{nota de } i \text{ se designado para turma regular.}
\end{align}

Seja $s(i)$ a escola do aluno $i$ e $D_i \in \{0,1\}$ o indicador de designação para turma pequena. O estimando de interesse é o \textbf{efeito médio da designação para turma pequena} --- um ITT (\textit{Intent-to-Treat}):
\begin{equation}
  \text{ITT} = \mathbb{E}[Y_i(1) - Y_i(0)].
  \label{eq:itt_star}
\end{equation}

Como a randomização foi realizada \textit{dentro} de cada escola, a hipótese de identificação relevante é a independência condicional na escola:
\begin{equation}
  \bigl(Y_i(1),\, Y_i(0)\bigr) \perp\!\!\!\perp D_i \;\big|\; s(i).
  \label{eq:cia_escola}
\end{equation}

Em linguagem intuitiva: dentro de uma mesma escola, alunos designados para turmas pequenas e alunos designados para turmas regulares são comparáveis em média, porque o sorteio foi conduzido apenas entre os alunos matriculados naquela escola. Portanto, diferenças sistemáticas de desempenho entre esses grupos, após condicionar na escola, podem ser interpretadas como efeito causal da designação para turma menor.

A hipótese \eqref{eq:cia_escola} implica imediatamente que:
\begin{equation}
  \mathbb{E}[\varepsilon_i \mid D_i, s(i)] = 0,
  \label{eq:exog_escola}
\end{equation}
onde $\varepsilon_i = Y_i - \mathbb{E}[Y_i \mid D_i, s(i)]$ é o erro residual após condicionar em escola e tratamento. Não há viés de seleção dentro de escola.

\subsection*{ITT vs.\ LATE}

O ITT mede o efeito da \textit{designação} para turma pequena, não necessariamente o efeito de \textit{frequentar} uma turma pequena. Os dois coincidem se todos os alunos designados para turma pequena efetivamente frequentam uma turma pequena --- o que pode não ser o caso em presença de não-cumprimento (\textit{non-compliance}): alunos podem mudar de turma, abandonar o experimento ou ingressar em turmas diferentes da designada.

Quando há não-cumprimento parcial, e se aceitarmos a designação original como instrumento válido para a turma efetivamente frequentada, o LATE (\textit{Local Average Treatment Effect}) mede o efeito sobre os \textit{cumpridores} --- aqueles que frequentaram a turma pequena justamente porque foram para ela designados. Neste guia, trabalhamos com o ITT como estimando principal por ser mais diretamente identificado pela aleatorização original.

\begin{exemplobox}[ITT e LATE no STAR]
Suponha que 5\% dos alunos designados para turmas pequenas tenham migrado para turmas regulares ao longo do ano. O ITT estima o efeito médio sobre \textit{todos} os designados para turma pequena --- incluindo os que na prática não ficaram em turma pequena. O LATE estimaria o efeito sobre o subgrupo que cumpriu a designação. Como o não-cumprimento no STAR foi relativamente baixo, os dois tendem a ser numericamente próximos, mas conceitualmente distintos.
\end{exemplobox}

\subsection*{SUTVA: uma hipótese que merece atenção}

A hipótese SUTVA (\textit{Stable Unit Treatment Value Assumption}) afirma que o resultado potencial de um aluno depende apenas do seu próprio tipo de turma, e não do tipo de turma dos outros alunos. No STAR, isso é potencialmente questionável: turmas pequenas e regulares coexistem \textit{dentro} da mesma escola.

Alguns canais potenciais de interferência merecem atenção:

\begin{itemize}[leftmargin=*]
  \item \textbf{Realocação de professoras}: se as melhores professoras da escola são sistematicamente alocadas a turmas pequenas, o efeito estimado captura tanto o tamanho da turma quanto a qualidade docente.
  \item \textbf{Efeitos de pares}: alunos em turmas regulares podem ser afetados por terem colegas que, em outro desenho, estariam em turmas pequenas.
  \item \textbf{Recursos compartilhados}: espaço físico, materiais e tempo dos coordenadores pedagógicos são divididos entre turmas dentro da mesma escola.
  \item \textbf{Composição das turmas}: se alunos com maior desempenho prévio forem sistematicamente alocados a certos tipos de turma, a comparação entre turmas dentro da escola pode ser contaminada.
\end{itemize}

Este exemplo serve para reforçar um ponto metodológico central: mesmo em experimentos aleatórios, é necessário explicitar as hipóteses de identificação e avaliar sua plausibilidade no contexto específico do estudo. A aleatorização resolve o problema de seleção entre grupos, mas não elimina automaticamente todas as ameaças à validade interna.

\begin{atencaobox}[SUTVA não é automática]
O fato de um experimento ter sido aleatorizado não garante que o SUTVA seja satisfeito. Sempre que turmas, indivíduos ou unidades do mesmo grupo de randomização interagem --- o que é comum em experimentos educacionais e comunitários ---, a hipótese de ausência de interferência precisa ser defendida com argumentos substantivos, não apenas assumida.
\end{atencaobox}


%% ─────────────────────────────────────────────────────────────────
\section{Estimação: Da Diferença de Médias à Regressão com Efeitos Fixos}
\label{sec:star_estimacao}

\subsection*{Especificação 1 — Diferença de médias simples}

A estimativa mais direta é a diferença de médias entre alunos designados para turmas pequenas e demais alunos:
\begin{equation}
  Y_i = \alpha + \beta\, D_i + u_i.
  \label{eq:spec1}
\end{equation}

O coeficiente $\hat{\beta}$ compara a nota média dos alunos em turmas pequenas com a nota média dos demais, sem qualquer condicionamento adicional. Sob randomização incondicional, esse estimador seria não-viesado para o ITT. No STAR, porém, a randomização foi feita \textit{dentro} de escola --- o que significa que comparar alunos de escolas diferentes pode introduzir diferenças pré-existentes entre escolas, não controladas pelo experimento.

\subsection*{Especificação 2 — Efeitos fixos de escola}

A especificação mais alinhada ao desenho experimental inclui efeitos fixos de escola $\lambda_{s(i)}$:
\begin{equation}
  Y_i = \alpha + \beta\, D_i + \lambda_{s(i)} + u_i.
  \label{eq:spec2}
\end{equation}

Ao incluir uma dummy para cada escola, a comparação passa a ser feita \textit{dentro} de escola --- exatamente como a randomização foi conduzida. O coeficiente $\hat{\beta}$ agora compara alunos designados para turma pequena com alunos designados para turma regular \textit{na mesma escola}, eliminando qualquer diferença entre escolas que pudesse confundir a estimativa.

\subsection*{Especificação 3 — Efeitos fixos de escola e covariadas}

A terceira especificação adiciona covariadas pré-tratamento $X_i$ (gênero, raça, elegibilidade para merenda):
\begin{equation}
  Y_i = \alpha + \beta\, D_i + \lambda_{s(i)} + X_i'\gamma + u_i.
  \label{eq:spec3}
\end{equation}

As covariadas têm aqui um papel específico. Em um RCT bem implementado, $D_i$ é independente de $X_i$ por construção (dentro de escola), de modo que $\gamma \neq 0$ não implica que omitir $X_i$ causaria viés. O benefício dos controles aqui é outro: ao capturar parte da variação nos resultados $Y_i$, as covariadas reduzem a variância residual $\hat{u}_i$, tornando o estimador $\hat{\beta}$ mais preciso sem sacrificar a validade causal. \citet{Lin2013} formaliza esse argumento e mostra que, em RCTs, a especificação preferida inclui covariadas centradas e interações com o tratamento, garantindo eficiência assintótica.

\begin{notabox}[{Controles em RCTs --- precisão sem causalidade}]
Em estudos observacionais, adicionamos controles para tentar satisfazer a CIA e eliminar viés. Em RCTs, adicionamos controles para ganhar precisão. A distinção é fundamental: no primeiro caso, a validade causal depende dos controles incluídos; no segundo, a validade já está garantida pela randomização, e os controles apenas melhoram a eficiência.
\end{notabox}

\begin{lstlisting}[style=Rstyle, caption={Três especificações progressivas no STAR}]
# -----------------------------
# 1) Especificação 1: diferença de médias simples
# -----------------------------
spec1 <- lm_robust(read_z ~ D, data = dados_k, se_type = "HC2")

# -----------------------------
# 2) Especificação 2: efeitos fixos de escola
# -----------------------------
spec2 <- feols(read_z ~ D | escola, data = dados_k,
               vcov = "HC1")

# -----------------------------
# 3) Especificação 3: EF de escola + covariadas pré-tratamento
# -----------------------------
spec3 <- feols(read_z ~ D + feminino + nao_branco + merenda | escola,
               data = dados_k, vcov = "HC1")

# Sumário das três especificações
etable(spec1, spec2, spec3,
       keep    = "D",
       headers = c("(1) Sem EF", "(2) EF escola", "(3) EF + covariadas"),
       title   = "Efeito de turma pequena sobre notas de leitura (STAR, K)")
\end{lstlisting}

% ──────────────────────────────────────────────────────────────────
% TABELA 5.1 — Três especificações: efeito de turma pequena sobre notas
% Executar o script R para obter os coeficientes, erros-padrão e
% estatísticas de diagnóstico (R², N, número de escolas).
% ──────────────────────────────────────────────────────────────────
\begin{table}[H]
\centering
\caption{Efeito de turma pequena sobre nota de leitura: três especificações}
\label{tab:spec_comparacao}
\begin{tabular}{lccc}
\toprule
 & (1) Sem EF & (2) EF Escola & (3) EF + Cov. \\
\midrule
Turma pequena ($D_i$)
  & \textit{[inserir]} & \textit{[inserir]} & \textit{[inserir]} \\
  & \textit{[inserir SE]} & \textit{[inserir SE]} & \textit{[inserir SE]} \\[4pt]
Efeitos fixos de escola & Não & Sim & Sim \\
Covariadas pré-tratamento & Não & Não & Sim \\
$R^2$ & \textit{[inserir]} & \textit{[inserir]} & \textit{[inserir]} \\
Observações & \textit{[inserir]} & \textit{[inserir]} & \textit{[inserir]} \\
\bottomrule
\end{tabular}

\vspace{4pt}
\figmeta{Dados: \citet{Krueger1999}, via pacote \texttt{AER}}{Elaboração dos autores}{Variável dependente: nota de leitura padronizada (média = 0, DP = 1). Covariadas: gênero, raça e elegibilidade para merenda gratuita. Erros-padrão HC robustos entre parênteses. * $p < 0{,}05$; ** $p < 0{,}01$; *** $p < 0{,}001$.}
\end{table}

A comparação entre as três colunas da Tabela~\ref{tab:spec_comparacao} ilustra um resultado esperado: em um RCT bem implementado, as estimativas nas três especificações devem ser numericamente similares. Eventuais diferenças entre a coluna (1) e as colunas (2) e (3) indicam diferenças pré-existentes entre escolas, que a aleatorização dentro de escola não controla. A redução nos erros-padrão entre (1) e (3) reflete o ganho de precisão proveniente dos controles.


%% ─────────────────────────────────────────────────────────────────
\section{Inferência I: Qual Erro-Padrão Usar?}
\label{sec:star_ep}

A escolha do estimador de variância não é um detalhe técnico --- é parte da análise substantiva. No STAR, há uma razão estrutural clara para não usar erros-padrão clássicos: alunos dentro da mesma turma compartilham professora, rotina pedagógica, material didático e dinâmica de sala. Seus resultados não são observações independentes. Ignorar essa correlação intraclasse leva a erros-padrão artificialmente pequenos e a testes com taxa de falsos positivos acima do nível nominal.

\subsection*{Correlação intraclasse}

A dependência dentro de grupos pode ser quantificada pela \textbf{correlação intraclasse} (ICC):
\begin{equation}
  \rho = \frac{\sigma^2_c}{\sigma^2_c + \sigma^2_i},
  \label{eq:icc}
\end{equation}
onde $\sigma^2_c$ é a variância entre turmas e $\sigma^2_i$ é a variância dentro de turmas. O \textbf{efeito de desenho} (\textit{design effect}) indica o quanto o erro-padrão efetivo supera o erro-padrão clássico em uma amostra com $\bar{n}$ observações por turma:
\begin{equation}
  \text{DEFF} = 1 + (\bar{n} - 1)\,\rho.
  \label{eq:deff}
\end{equation}

Com $\rho = 0{,}20$ e turmas de tamanho médio $\bar{n} = 20$, o DEFF seria $1 + 19 \times 0{,}20 = 4{,}80$ --- o que significa que os erros-padrão clássicos subestimariam a incerteza por um fator de $\sqrt{4{,}80} \approx 2{,}19$. Estatísticas-$t$ baseadas nesses erros estariam infladas em mais de 100\%.

\subsection*{Quatro abordagens}

Comparamos quatro estimadores de variância:

\begin{enumerate}[leftmargin=*, label=(\roman*)]
  \item \textbf{Clássico}: assume homoscedasticidade e independência entre todas as observações.
  \item \textbf{HC robusto}: corrige heteroscedasticidade, mas ainda trata observações como independentes.
  \item \textbf{Clusterizado por turma}: agrupa observações dentro da mesma turma, reconhecendo a dependência intraclasse. É a escolha natural dado que o tratamento é atribuído no nível da turma.
  \item \textbf{Clusterizado por escola}: agrupa por escola. Mais conservador, com número menor de clusters (79 escolas vs.\ $\sim$300 turmas).
\end{enumerate}

\begin{lstlisting}[style=Rstyle, caption={Comparação de quatro estimadores de variância no STAR}]
# -----------------------------
# 1) Modelo base: EF de escola, sem outras covariadas
# -----------------------------
mod_base <- lm(read_z ~ D + escola, data = dados_k)

# -----------------------------
# 2) Quatro variantes de erro-padrão
# -----------------------------
# (i) Clássico
se_classico <- sqrt(diag(vcov(mod_base)))["D"]

# (ii) HC2 robusto
se_hc2 <- sqrt(vcovHC(mod_base, type = "HC2"))["D", "D"]

# (iii) Clusterizado por turma
dados_k <- dados_k %>%
  mutate(turma_id = paste(schoolk, stark, sep = "_"))

se_turma <- sqrt(vcovCL(mod_base, cluster = ~turma_id,
                        data = dados_k))["D", "D"]

# (iv) Clusterizado por escola
se_escola <- sqrt(vcovCL(mod_base, cluster = ~escola,
                         data = dados_k))["D", "D"]

# Apresentação
resultados_ep <- data.frame(
  Metodo   = c("Clássico", "HC2", "Cluster turma", "Cluster escola"),
  SE       = round(c(se_classico, se_hc2, se_turma, se_escola), 4),
  t_stat   = round(coef(mod_base)["D"] /
             c(se_classico, se_hc2, se_turma, se_escola), 2)
)
print(resultados_ep)
\end{lstlisting}

% ──────────────────────────────────────────────────────────────────
% TABELA 5.2 — Comparação de erros-padrão
% Executar o script R para obter SE, estatística t e p-valor
% sob as quatro abordagens.
% ──────────────────────────────────────────────────────────────────
\begin{table}[H]
\centering
\caption{Estimativa do efeito de turma pequena sob diferentes estimadores de variância}
\label{tab:erros_padrao}
\begin{tabular}{lcccc}
\toprule
Abordagem & Estimativa & Erro-padrão & Estatística $t$ & Valor-$p$ \\
\midrule
Clássico              & \textit{[inserir]} & \textit{[inserir]} & \textit{[inserir]} & \textit{[inserir]} \\
HC2 robusto           & \textit{[inserir]} & \textit{[inserir]} & \textit{[inserir]} & \textit{[inserir]} \\
Cluster por turma     & \textit{[inserir]} & \textit{[inserir]} & \textit{[inserir]} & \textit{[inserir]} \\
Cluster por escola    & \textit{[inserir]} & \textit{[inserir]} & \textit{[inserir]} & \textit{[inserir]} \\
\bottomrule
\end{tabular}

\vspace{4pt}
\figmeta{Dados: \citet{Krueger1999}}{Elaboração dos autores}{Especificação com efeitos fixos de escola. Variável dependente: nota de leitura padronizada. A estimativa pontual é a mesma nas quatro linhas; apenas o estimador de variância varia.}
\end{table}

O padrão esperado na Tabela~\ref{tab:erros_padrao} é que os erros-padrão clusterizados sejam maiores que os clássicos e robustos. A comparação entre clusterização por turma e por escola ilustra o \textit{trade-off} entre adequação ao processo gerador de dados e estabilidade do estimador: com 79 escolas, o estimador sandwich clusterizado por escola opera próximo ao limite inferior recomendado ($G \approx 30$--50), enquanto as turmas oferecem maior número de clusters.

\begin{notabox}[Regra prática de clustering no STAR]
Como o tratamento foi atribuído no nível da \textit{turma}, a clusterização por turma é a escolha teoricamente correta. Se os dados contiverem identificadores de turma confiáveis, devem ser usados. A clusterização por escola é mais conservadora e pode ser reportada como verificação de robustez.
\end{notabox}


%% ─────────────────────────────────────────────────────────────────
\section{Inferência II: Bootstrap}
\label{sec:star_bootstrap}

A Seção~\ref{sec:clustering} mostrou que, com número razoável de clusters, o estimador sandwich clusterizado tende a ter bom desempenho assintótico. Com cerca de 300 turmas, o STAR está bem acima do limiar mínimo --- o que sugere que os intervalos de confiança assintóticos clusterizados devem ser confiáveis. O \textit{bootstrap} é útil aqui como verificação independente e como método para estatísticas cuja distribuição assintótica é menos direta.

\subsection*{Por que reamostrar turmas, não alunos}

Um erro frequente é aplicar o bootstrap padrão --- reamostrar observações individuais com reposição --- em dados com estrutura de clustering. Ao sortear alunos individualmente, o bootstrap padrão quebra a dependência dentro de turmas: uma amostra bootstrap poderia conter cinco cópias do mesmo aluno e nenhuma cópia de sua turma original, destruindo exatamente a estrutura que torna a clusterização necessária. A solução é o \textbf{\textit{cluster bootstrap}}: reamostrar \textit{turmas inteiras} com reposição, mantendo todos os alunos de cada turma sorteada.

\begin{lstlisting}[style=Rstyle, caption={Cluster bootstrap por turma no STAR}]
# -----------------------------
# 1) Cluster bootstrap: reamostrar turmas inteiras
# -----------------------------
rm(list = ls())
# (assumindo que dados_k e turma_id foram criados anteriormente)

set.seed(42)
B       <- 1999
turmas  <- unique(dados_k$turma_id)
G       <- length(turmas)

boot_coefs <- numeric(B)

for (b in 1:B) {
  # Sortear turmas com reposição
  turmas_b  <- sample(turmas, size = G, replace = TRUE)

  # Reconstruir base com todas as observações das turmas sorteadas
  amostra_b <- bind_rows(
    lapply(turmas_b, function(t) dados_k[dados_k$turma_id == t, ])
  )

  # Estimar modelo com EF de escola
  mod_b <- lm(read_z ~ D + escola, data = amostra_b)
  boot_coefs[b] <- coef(mod_b)["D"]
}

# IC bootstrap percentil (95%)
ic_boot <- quantile(boot_coefs, c(0.025, 0.975))
cat("IC bootstrap 95% (percentil): [",
    round(ic_boot[1], 3), ";", round(ic_boot[2], 3), "]\n")

# Comparar com IC assintótico clusterizado
se_assintotico <- sqrt(vcovCL(lm(read_z ~ D + escola, data = dados_k),
                              cluster = ~turma_id)["D", "D"])
est_pont <- coef(lm(read_z ~ D + escola, data = dados_k))["D"]
ic_assint <- est_pont + c(-1, 1) * 1.96 * se_assintotico

cat("IC assintótico 95% (cluster): [",
    round(ic_assint[1], 3), ";", round(ic_assint[2], 3), "]\n")
\end{lstlisting}

A mensagem didática desta subseção é a seguinte: quando há número razoável de clusters e o desenho é bem comportado, os ICs bootstrap e assintóticos clusterizados devem ser numericamente próximos. Uma divergência grande entre os dois seria um sinal de alerta --- sugerindo que alguma hipótese subjacente (número suficiente de clusters, distribuição simétrica do estimador, ausência de observações influentes) pode não estar sendo satisfeita. No STAR, com cerca de 300 turmas, a convergência entre os dois é esperada.


%% ─────────────────────────────────────────────────────────────────
\section{Inferência III: Quatro Disciplinas, Um Problema --- Múltiplos Testes}
\label{sec:star_multiplos}

O STAR oferece resultados de desempenho em mais de uma disciplina: leitura e matemática, ao menos. Quando testamos o efeito de turmas pequenas em todas as disciplinas disponíveis, deparamo-nos com o problema de múltiplos testes discutido na Seção~\ref{sec:multiplos_testes}: quanto maior o número de testes, maior a probabilidade de encontrar ao menos um resultado estatisticamente significativo por acaso, mesmo que o efeito verdadeiro seja nulo em todos os desfechos.

Suponha que testemos o efeito em quatro resultados (por exemplo, leitura e matemática no jardim de infância e no primeiro ano). Ao nível $\alpha = 0{,}05$ com quatro testes independentes, a probabilidade de ao menos uma rejeição espúria é $1 - 0{,}95^4 \approx 18{,}5\%$ --- quase quatro vezes o nível nominal.

\begin{lstlisting}[style=Rstyle, caption={Múltiplos testes: quatro resultados com correção de Bonferroni e BH-FDR}]
# -----------------------------
# 1) Estimar o efeito nas quatro variáveis de resultado
# Leitura (K), Matemática (K), Leitura (1o ano), Matemática (1o ano)
# -----------------------------
dados_multi <- STAR %>%
  filter(!is.na(stark)) %>%
  mutate(
    D      = as.integer(stark == "small"),
    escola = as.factor(schoolk),
    read_k_z  = scale(readk)[, 1],
    math_k_z  = scale(mathk)[, 1],
    read_1_z  = scale(read1)[, 1],
    math_1_z  = scale(math1)[, 1]
  )

resultados <- c("read_k_z", "math_k_z", "read_1_z", "math_1_z")
nomes      <- c("Leitura (K)", "Matemática (K)",
                "Leitura (1o ano)", "Matemática (1o ano)")

pvalores    <- numeric(length(resultados))
estimativas <- numeric(length(resultados))

for (j in seq_along(resultados)) {
  df_j  <- dados_multi[!is.na(dados_multi[[resultados[j]]]), ]
  mod_j <- lm(as.formula(paste(resultados[j], "~ D + escola")),
              data = df_j)
  se_j  <- sqrt(vcovCL(mod_j, cluster = ~schoolk, data = df_j)["D", "D"])
  estimativas[j] <- coef(mod_j)["D"]
  pvalores[j]    <- 2 * pt(-abs(coef(mod_j)["D"] / se_j),
                            df = df_j$escola %>% n_distinct() - 2)
}

# -----------------------------
# 2) Ajuste por Bonferroni e BH-FDR
# -----------------------------
pv_bonf <- p.adjust(pvalores, method = "bonferroni")
pv_bh   <- p.adjust(pvalores, method = "BH")

tabela_multi <- data.frame(
  Resultado    = nomes,
  Estimativa   = round(estimativas, 3),
  P_bruto      = round(pvalores, 3),
  P_Bonferroni = round(pv_bonf, 3),
  P_BH_FDR     = round(pv_bh, 3)
)
print(tabela_multi)
\end{lstlisting}

% ──────────────────────────────────────────────────────────────────
% TABELA 5.3 — Múltiplos testes: quatro resultados
% Executar o script R para obter estimativas, p-valores brutos
% e p-valores ajustados por Bonferroni e BH-FDR.
% ──────────────────────────────────────────────────────────────────
\begin{table}[H]
\centering
\caption{Efeito de turma pequena: quatro resultados com ajuste para múltiplos testes}
\label{tab:multiplos_testes_star}
\begin{tabular}{lcccc}
\toprule
Resultado & Estimativa & $p$ bruto & $p$ Bonferroni & $p$ BH-FDR \\
\midrule
Leitura (K)          & \textit{[inserir]} & \textit{[inserir]} & \textit{[inserir]} & \textit{[inserir]} \\
Matemática (K)       & \textit{[inserir]} & \textit{[inserir]} & \textit{[inserir]} & \textit{[inserir]} \\
Leitura (1.º ano)    & \textit{[inserir]} & \textit{[inserir]} & \textit{[inserir]} & \textit{[inserir]} \\
Matemática (1.º ano) & \textit{[inserir]} & \textit{[inserir]} & \textit{[inserir]} & \textit{[inserir]} \\
\bottomrule
\end{tabular}

\vspace{4pt}
\figmeta{Dados: \citet{Krueger1999}}{Elaboração dos autores}{Especificação com efeitos fixos de escola. Erros-padrão clusterizados por escola. Bonferroni: nível ajustado $= 0{,}05/4 = 0{,}0125$. BH-FDR: controla proporção esperada de falsas descobertas a 5\%.}
\end{table}

A leitura da Tabela~\ref{tab:multiplos_testes_star} deve responder a três perguntas: (i) quais efeitos são significativos sem qualquer ajuste? (ii) quais sobrevivem ao ajuste de Bonferroni? (iii) o padrão é consistente entre disciplinas e anos, ou concentrado em um único resultado? Uma estimativa que se mantém significativa após Bonferroni sugere evidência mais sólida do que aquela que só resiste ao p-valor bruto. Quando os efeitos são positivos e de magnitude similar em todas as disciplinas, a credibilidade aumenta: seria improvável que o mesmo viés produzisse o mesmo sinal em quatro desfechos distintos.

\begin{notabox}[Bonferroni vs.\ BH-FDR]
Bonferroni controla a probabilidade de \textit{ao menos} um falso positivo no conjunto de testes (FWER), tornando cada teste individual muito conservador. BH-FDR controla a proporção esperada de falsas descobertas entre as rejeições, sendo menos conservador quando $m$ é grande. Para quatro desfechos, a diferença entre os dois é pequena --- Bonferroni com $\alpha^* = 0{,}0125$ vs.\ BH com limiar adaptativo. Com dezenas ou centenas de desfechos, a diferença torna-se mais relevante.
\end{notabox}


%% ─────────────────────────────────────────────────────────────────
\section{Inferência IV: Inferência por Permutação}
\label{sec:star_permutacao}

A inferência por permutação, apresentada na Seção~\ref{sec:design_based}, tem uma propriedade especial no STAR: como o processo de aleatorização está documentado (randomização dentro de escola), a distribuição de referência pode ser construída a partir desse mesmo processo, sem qualquer hipótese distribucional sobre os erros.

\subsection*{Respeitando o desenho na permutação}

Como a aleatorização ocorreu \textit{dentro} de escola, a permutação deve respeitar esse bloco: permutamos os rótulos de tratamento ($D_i$) apenas entre alunos da mesma escola, preservando o número de tratados por escola. Permutar globalmente ignoraria a estrutura por blocos e produziria distribuições de referência inadequadas.

\begin{lstlisting}[style=Rstyle, caption={Inferência por permutação com aleatorização em bloco (escola)}]
# -----------------------------
# 1) Estatística observada
# -----------------------------
mod_obs  <- lm(read_z ~ D + escola, data = dados_k)
t_obs    <- coef(mod_obs)["D"] /
            sqrt(vcovCL(mod_obs, cluster = ~turma_id,
                        data = dados_k)["D", "D"])

cat("Estatística t observada:", round(t_obs, 3), "\n")

# -----------------------------
# 2) Distribuição de permutação
# Permutar D dentro de escola (respeitar o bloco de randomização)
# -----------------------------
set.seed(42)
P <- 4999
t_perm <- numeric(P)

for (p in 1:P) {
  dados_perm <- dados_k %>%
    group_by(escola) %>%
    mutate(D_perm = sample(D)) %>%   # permuta D dentro de escola
    ungroup()

  mod_p   <- lm(read_z ~ D_perm + escola, data = dados_perm)
  se_p    <- sqrt(vcovCL(mod_p, cluster = ~turma_id,
                          data = dados_perm)["D_perm", "D_perm"])
  t_perm[p] <- coef(mod_p)["D_perm"] / se_p
}

# -----------------------------
# 3) P-valor de permutação (bilateral)
# -----------------------------
p_perm    <- mean(abs(t_perm) >= abs(t_obs))
p_assint  <- 2 * pnorm(-abs(t_obs))

cat("P-valor de permutação (bilateral): ", round(p_perm, 4), "\n")
cat("P-valor assintótico (normal):      ", round(p_assint, 4), "\n")
\end{lstlisting}

A comparação entre o p-valor de permutação e o p-valor assintótico clusterizado ilustra um resultado típico de experimentos com amostras grandes: os dois tendem a ser numericamente próximos. A inferência por permutação é particularmente atraente porque usa diretamente o processo de aleatorização descrito no protocolo experimental --- não há necessidade de aproximações assintóticas nem hipóteses sobre a distribuição dos erros. Em amostras menores, com poucos clusters ou distribuições assimétricas, a inferência por permutação pode diferir substancialmente dos métodos assintóticos e, nesse caso, oferece garantias de validade mais robustas.


%% ─────────────────────────────────────────────────────────────────
\section{Testes de Falsificação}
\label{sec:star_falsificacao}

Os testes de falsificação não ``provam'' que o STAR tem identificação válida, mas ajudam a detectar problemas que comprometeriam a interpretação causal. Apresentamos três verificações: balanceamento de covariadas pré-tratamento, atrito diferencial e balanceamento das características das professoras.

\subsection*{Balanceamento de covariadas pré-tratamento}

Em um experimento bem implementado, características pré-tratamento dos alunos não devem prever sistematicamente a designação para turma pequena. Para testar isso, regredimos $D_i$ (indicador de turma pequena) nas covariadas pré-tratamento disponíveis e realizamos um F-teste conjunto de significância:
\begin{equation}
  D_i = \alpha + \lambda_{s(i)} + \delta_1\,\text{feminino}_i + \delta_2\,\text{nao\_branco}_i + \delta_3\,\text{merenda}_i + v_i.
  \label{eq:teste_balanceamento}
\end{equation}

A hipótese nula é $H_0: \delta_1 = \delta_2 = \delta_3 = 0$, ou seja, nenhuma covariada prediz o tratamento dentro de escola. Sob aleatorização bem conduzida, esperamos não rejeitar essa hipótese.

\begin{lstlisting}[style=Rstyle, caption={F-teste de balanceamento de covariadas}]
# -----------------------------
# 1) Regredir tratamento em covariadas pré-tratamento + EF escola
# -----------------------------
mod_bal <- lm(D ~ escola + feminino + nao_branco + merenda,
              data = dados_k)

# F-teste conjunto: apenas para as covariadas (excluindo EF de escola)
library(car)
ftest_bal <- linearHypothesis(mod_bal,
  c("feminino = 0", "nao_branco = 0", "merenda = 0"),
  vcov. = vcovCL(mod_bal, cluster = ~turma_id, data = dados_k))

cat("=== Teste F conjunto de balanceamento ===\n")
print(ftest_bal)

# -----------------------------
# 2) Tabela de médias por grupo (dentro de escola)
# -----------------------------
dados_k %>%
  group_by(D) %>%
  summarise(
    pct_feminino   = mean(feminino, na.rm = TRUE),
    pct_nao_branco = mean(nao_branco, na.rm = TRUE),
    pct_merenda    = mean(merenda, na.rm = TRUE),
    n              = n()
  ) %>%
  print()
\end{lstlisting}

% ──────────────────────────────────────────────────────────────────
% TABELA 5.4 — Teste de balanceamento: covariadas pré-tratamento
% Executar o script R para obter as médias por grupo e a estatística F.
% ──────────────────────────────────────────────────────────────────
\begin{table}[H]
\centering
\caption{Verificação de balanceamento: covariadas pré-tratamento por condição experimental}
\label{tab:balanceamento}
\begin{tabular}{lccc}
\toprule
Covariada & Turma pequena ($D=1$) & Controle ($D=0$) & Dif. padronizada \\
\midrule
Feminino (\%)         & \textit{[inserir]} & \textit{[inserir]} & \textit{[inserir]} \\
Não-branco (\%)       & \textit{[inserir]} & \textit{[inserir]} & \textit{[inserir]} \\
Merenda gratuita (\%) & \textit{[inserir]} & \textit{[inserir]} & \textit{[inserir]} \\
\midrule
\multicolumn{3}{l}{F-teste conjunto ($H_0$: todas as diferenças = 0)}
  & $F = $ \textit{[inserir]} \\
\multicolumn{3}{l}{Valor-$p$}
  & \textit{[inserir]} \\
\midrule
Observações           & \textit{[inserir]} & \textit{[inserir]} & \\
\bottomrule
\end{tabular}

\vspace{4pt}
\figmeta{Dados: \citet{Krueger1999}}{Elaboração dos autores}{Comparação feita dentro de escola (efeitos fixos de escola incluídos na regressão). Dif. padronizada $= (\bar{X}_1 - \bar{X}_0)/\sqrt{(s_1^2 + s_0^2)/2}$.}
\end{table}

\subsection*{Atrito diferencial}

A aleatorização garante comparabilidade no momento inicial do experimento. Ao longo dos quatro anos, no entanto, alunos podem sair do experimento --- por mudança de escola, ausência prolongada ou dados faltantes. Se a \textit{probabilidade de permanecer na amostra} for diferente entre turmas pequenas e regulares, a comparação ao final do período pode estar contaminada: estaríamos comparando grupos sistematicamente diferentes dos originalmente sorteados.

Para verificar atrito diferencial, testamos se $D_i$ prediz a probabilidade de o aluno ter notas disponíveis no ano de interesse:
\begin{equation}
  \mathds{1}[Y_i \text{ observado}] = \alpha + \beta_{\text{atrito}}\, D_i + \lambda_{s(i)} + u_i.
  \label{eq:atrito}
\end{equation}

Um coeficiente $\hat{\beta}_{\text{atrito}}$ significativamente diferente de zero indica que tratados e controles têm taxas de atrito distintas --- o que comprometeria a validade interna mesmo com aleatorização original correta.

\begin{lstlisting}[style=Rstyle, caption={Verificação de atrito diferencial}]
# -----------------------------
# Indicador de permanência na amostra para cada ano
# Testamos se D prediz estar na amostra no 1o ano
# -----------------------------
dados_atrito <- STAR %>%
  filter(!is.na(stark)) %>%
  mutate(
    D              = as.integer(stark == "small"),
    escola         = as.factor(schoolk),
    presente_ano1  = as.integer(!is.na(read1))
  )

mod_atrito <- lm_robust(presente_ano1 ~ D,
                        fixed_effects = ~escola,
                        clusters      = schoolk,
                        data          = dados_atrito,
                        se_type       = "CR2")

cat("=== Teste de atrito diferencial (1o ano) ===\n")
cat("Coeficiente D:  ", round(coef(mod_atrito)["D"], 4), "\n")
se_atrito <- summary(mod_atrito)$coefficients["D", 2]
pv_atrito <- summary(mod_atrito)$coefficients["D", 4]
cat("Erro-padrão:    ", round(se_atrito, 4), "\n")
cat("Valor-p:        ", round(pv_atrito, 3), "\n")
\end{lstlisting}

\subsection*{Placebo de professora: características docentes como teste de implementação}

Um aspecto do STAR que merece atenção é a qualidade da implementação. Se as professoras mais experientes ou mais tituladas foram sistematicamente alocadas a turmas pequenas, o efeito estimado confundiria o tamanho da turma com a qualidade docente. Testamos isso verificando o balanceamento das características das professoras entre condições experimentais:

\begin{equation}
  \text{experiência}_{c} = \alpha + \beta_{\text{prof}}\, D_c + \lambda_{s(c)} + w_c,
  \label{eq:placebo_prof}
\end{equation}

onde $c$ indexa turmas e $D_c$ é o indicador de turma pequena. Se $\hat{\beta}_{\text{prof}}$ for significativamente diferente de zero, há evidência de que turmas pequenas e regulares foram atribuídas a professoras sistematicamente diferentes.

\begin{lstlisting}[style=Rstyle, caption={Balanceamento de características das professoras entre condições}]
# -----------------------------
# Agregar para o nível de turma (cada turma tem uma professora)
# -----------------------------
dados_prof <- STAR %>%
  filter(!is.na(stark), !is.na(experiencek)) %>%
  mutate(D = as.integer(stark == "small")) %>%
  group_by(schoolk, stark) %>%
  summarise(
    experiencia = first(experiencek),
    D           = first(D),
    .groups = "drop"
  ) %>%
  mutate(escola = as.factor(schoolk))

mod_prof <- lm_robust(experiencia ~ D,
                      fixed_effects = ~escola,
                      clusters      = schoolk,
                      data          = dados_prof,
                      se_type       = "CR2")

cat("=== Placebo de professora: experiência docente ===\n")
cat("Coeficiente D:  ", round(coef(mod_prof)["D"], 3), "\n")
se_prof <- summary(mod_prof)$coefficients["D", 2]
pv_prof <- summary(mod_prof)$coefficients["D", 4]
cat("Erro-padrão:    ", round(se_prof, 3), "\n")
cat("Valor-p:        ", round(pv_prof, 3), "\n")
\end{lstlisting}

\begin{atencaobox}[Testes de falsificação: limites e interpretação]
Passar em todos os testes de falsificação é evidência circunstancial favorável à validade do desenho, mas não é prova de causalidade. Uma estratégia pode superar todos os placebos disponíveis e ainda ter identificação inválida por razões não testáveis --- por exemplo, um canal de interferência entre turmas que não deixa rastro nas covariadas observadas. Os testes de falsificação devem ser lidos em conjunto com o argumento teórico sobre o mecanismo de identificação, não como substitutos a ele.
\end{atencaobox}


%% ─────────────────────────────────────────────────────────────────
\subsection*{Fechamento: o que o STAR ensinou neste guia}

O Experimento STAR percorreu, em um único exemplo empírico, toda a cadeia metodológica discutida neste guia. A \textbf{aleatorização dentro de escola} resolveu o problema de identificação: dentro de cada escola, tratados e controles são comparáveis em média, eliminando o viés de seleção. A \textbf{regressão com efeitos fixos de escola} alinhou a estimação ao desenho experimental, melhorando precisão sem comprometer validade. A discussão sobre \textbf{erros-padrão} revelou que a escolha do estimador de variância é uma decisão substantiva: ignorar a correlação intraclasse leva a inferências enganosas. O \textbf{bootstrap} por cluster confirmou os resultados assintóticos quando o número de grupos é suficiente, e mostrou quando os dois tendem a divergir. A análise de \textbf{múltiplos testes} forçou uma resposta honesta sobre quais resultados são robustos: um efeito que sobrevive a Bonferroni em quatro disciplinas é mais convincente do que um único p-valor favorável. A \textbf{inferência por permutação} usou diretamente o processo de aleatorização para construir distribuições de referência exatas, sem aproximações assintóticas. Por fim, os \textbf{testes de falsificação} --- balanceamento de covariadas, atrito diferencial e placebo de professora --- forneceram evidência circunstancial sobre a qualidade da implementação experimental.

O ponto que conecta todos esses elementos é o seguinte: cada ferramenta responde a uma pergunta diferente. Identificação pergunta \textit{o que estamos estimando?} Estimação pergunta \textit{como calculamos isso?} Inferência pergunta \textit{quanta incerteza há?} Bootstrap e permutação perguntam \textit{essa incerteza é bem caracterizada pelas aproximações assintóticas?} Múltiplos testes perguntam \textit{o resultado sobrevive a buscas seletivas?} Falsificação pergunta \textit{as hipóteses por trás da estimação são plausíveis?} Nenhuma dessas perguntas é dispensável, e a resposta a qualquer uma delas pressupõe que as anteriores tenham sido respondidas com rigor.
%%
%%  Dados: AER::STAR (pacote AER do R)
%%  Referência principal: Krueger (1999)
%% ─────────────────────────────────────────────────────────────────

\chapter{Aplicação Empírica: O Experimento STAR}
\label{sec:star}
\label{cap:star}

As seções anteriores construíram o arcabouço conceitual da inferência causal: resultados potenciais, identificação, estimação, erros-padrão, testes de hipótese, bootstrap e falsificação. O objetivo desta seção é usar um experimento clássico como laboratório empírico para ver esses conceitos funcionando juntos. O Experimento STAR --- \textit{Student/Teacher Achievement Ratio} --- oferece um dos melhores exemplos disponíveis de avaliação experimental em políticas educacionais. Ele é suficientemente limpo para ilustrar a lógica da identificação por aleatorização, mas suficientemente rico para revelar as complexidades práticas de inferência, múltiplos testes e falsificação.

O experimento foi conduzido no estado americano do Tennessee entre 1985 e 1989 \citep{Finn1990, Krueger1999}. A pergunta central é direta: turmas menores melhoram o desempenho escolar? Para respondê-la com rigor causal, o estado financiou um experimento em larga escala em que estudantes e professoras foram aleatoriamente designados a três tipos de turma: \textbf{turma pequena} (13--17 alunos), \textbf{turma regular} (22--26 alunos) e \textbf{turma regular com auxiliar} (22--26 alunos com uma professora auxiliar). O acompanhamento ocorreu ao longo de quatro anos, do jardim de infância (pré-escola, \textit{kindergarten}) ao terceiro ano do ensino fundamental (K--3).

O STAR é um caso de referência na literatura de avaliação de impacto por três razões. Primeira, é um dos poucos experimentos aleatórios de grande escala realizados em educação --- campo em que, por razões éticas e logísticas, experimentos são raros. Segunda, sua metodologia está bem documentada, o que permite verificar e replicar cada escolha empírica. Terceira, os dados tornaram-se públicos e foram explorados por dezenas de pesquisadores, criando uma literatura acumulada que serve de referência para discussão metodológica \citep{Chetty2011}. Neste guia, não temos o objetivo de contribuir com nova pesquisa sobre educação: usamos o STAR como plataforma para ilustrar, com um exemplo concreto, os métodos discutidos nas seções anteriores.


%% ─────────────────────────────────────────────────────────────────
\section{Dados e Desenho Experimental}
\label{sec:star_dados}

\subsection*{Estrutura da amostra}

O experimento acompanhou aproximadamente 11.600 alunos distribuídos em cerca de 300 turmas em 79 escolas do Tennessee. As 79 escolas foram selecionadas de forma a representar áreas rurais, suburbanas e urbanas do estado. Dentro de cada escola, alunos e professoras foram aleatoriamente designados a um dos três tipos de turma.

É importante distinguir três unidades de análise que coexistem nos dados:

\begin{itemize}[leftmargin=*]
  \item \textbf{Aluno ($i$)}: unidade de observação. Os dados contêm, para cada aluno, notas em diferentes disciplinas, características demográficas e indicadores de situação socioeconômica.
  \item \textbf{Turma ($c$)}: unidade de tratamento. O tratamento --- tipo de turma --- é atribuído no nível da turma, não do aluno individualmente. Todos os alunos de uma mesma turma recebem exatamente o mesmo tratamento.
  \item \textbf{Escola ($s$)}: unidade de randomização. A aleatorização foi feita \textit{dentro} de cada escola: dentro da escola $s$, alunos foram sorteados para turmas pequenas ou regulares. A composição das turmas \textit{entre} escolas não foi randomizada.
\end{itemize}

Essa estrutura hierárquica --- alunos dentro de turmas dentro de escolas --- tem implicações diretas para a estimação e para a inferência, como veremos ao longo desta seção.

\subsection*{Variáveis disponíveis}

O conjunto de dados contém as seguintes informações por aluno:

\begin{itemize}[leftmargin=*]
  \item \textbf{Notas}: pontuação em leitura e matemática para cada ano (jardim de infância ao terceiro ano). São as variáveis de resultado.
  \item \textbf{Tipo de turma designado}: indicador do braço do experimento ao qual o aluno foi alocado (\texttt{small}, \texttt{regular}, \texttt{regular+aide}).
  \item \textbf{Características do aluno}: raça/etnia, gênero, elegibilidade para merenda gratuita (proxy de renda familiar). Essas são covariadas pré-tratamento.
  \item \textbf{Identificadores}: código da escola e da turma.
  \item \textbf{Características da professora}: anos de experiência e titulação, quando disponíveis. Usadas nos testes de falsificação.
\end{itemize}

\subsection*{Rotina de análise em R}

O script de análise carrega os dados do pacote \texttt{AER} do R \citep{Zeileis2008}, que distribui uma versão dos dados do STAR em formato pronto para análise. O script:

\begin{enumerate}[leftmargin=*, label=(\roman*)]
  \item Carrega os dados e seleciona a amostra analítica do jardim de infância, eliminando observações com informação faltante nas variáveis de interesse.
  \item Cria o indicador binário de tratamento $D_i = \mathds{1}[\texttt{stark} = \texttt{\textquotedbl small\textquotedbl}]$, separando turma pequena das duas condições de controle agrupadas.
  \item Padroniza as notas de leitura e matemática (média zero, desvio-padrão um) para facilitar a interpretação dos coeficientes em desvios-padrão.
  \item Define covariadas pré-tratamento: gênero, raça e elegibilidade para merenda gratuita.
  \item Constrói os objetos usados nas estimações e tabelas das subseções seguintes.
\end{enumerate}

\begin{lstlisting}[style=Rstyle, caption={Carregamento e preparação dos dados STAR}]
# -----------------------------
# 1) Configuração inicial
# -----------------------------
rm(list = ls())

library(AER)       # dados STAR
library(dplyr)     # manipulação
library(fixest)    # efeitos fixos (feols)
library(estimatr)  # lm_robust / erros clusterizados
library(sandwich)  # vcovCL
library(lmtest)    # coeftest

# -----------------------------
# 2) Carregar e filtrar dados
# Foco no jardim de infância (ano k)
# -----------------------------
data("STAR", package = "AER")

dados_k <- STAR %>%
  filter(!is.na(stark), !is.na(readk), !is.na(mathk)) %>%
  mutate(
    # Tratamento binário: turma pequena vs. demais
    D       = as.integer(stark == "small"),
    # Notas padronizadas (desvio-padrão = 1)
    read_z  = scale(readk)[, 1],
    math_z  = scale(mathk)[, 1],
    # Covariadas pré-tratamento
    feminino    = as.integer(gender == "female"),
    nao_branco  = as.integer(ethnicity != "cauc"),
    merenda     = as.integer(lunchk == "free"),
    # Identificadores
    escola  = as.factor(schoolk)
  )

cat("Observações na amostra analítica:", nrow(dados_k), "\n")
cat("Escolas:                         ", n_distinct(dados_k$escola), "\n")
cat("Tratados (turma pequena):        ", sum(dados_k$D), "\n")
cat("Controle:                        ", sum(1 - dados_k$D), "\n")
\end{lstlisting}


%% ─────────────────────────────────────────────────────────────────
\section{Identificação: Por que a Aleatorização Resolve o Problema}
\label{sec:star_identificacao}

\subsection*{O estimando e a estratégia de identificação}

Para cada aluno $i$, definimos dois resultados potenciais:
\begin{align}
  Y_i(1) &\equiv \text{nota de } i \text{ se designado para turma pequena,} \\
  Y_i(0) &\equiv \text{nota de } i \text{ se designado para turma regular.}
\end{align}

Seja $s(i)$ a escola do aluno $i$ e $D_i \in \{0,1\}$ o indicador de designação para turma pequena. O estimando de interesse é o \textbf{efeito médio da designação para turma pequena} --- um ITT (\textit{Intent-to-Treat}):
\begin{equation}
  \text{ITT} = \mathbb{E}[Y_i(1) - Y_i(0)].
  \label{eq:itt_star}
\end{equation}

Como a randomização foi realizada \textit{dentro} de cada escola, a hipótese de identificação relevante é a independência condicional na escola:
\begin{equation}
  \bigl(Y_i(1),\, Y_i(0)\bigr) \perp\!\!\!\perp D_i \;\big|\; s(i).
  \label{eq:cia_escola}
\end{equation}

Em linguagem intuitiva: dentro de uma mesma escola, alunos designados para turmas pequenas e alunos designados para turmas regulares são comparáveis em média, porque o sorteio foi conduzido apenas entre os alunos matriculados naquela escola. Portanto, diferenças sistemáticas de desempenho entre esses grupos, após condicionar na escola, podem ser interpretadas como efeito causal da designação para turma menor.

A hipótese \eqref{eq:cia_escola} implica imediatamente que:
\begin{equation}
  \mathbb{E}[\varepsilon_i \mid D_i, s(i)] = 0,
  \label{eq:exog_escola}
\end{equation}
onde $\varepsilon_i = Y_i - \mathbb{E}[Y_i \mid D_i, s(i)]$ é o erro residual após condicionar em escola e tratamento. Não há viés de seleção dentro de escola.

\subsection*{ITT vs.\ LATE}

O ITT mede o efeito da \textit{designação} para turma pequena, não necessariamente o efeito de \textit{frequentar} uma turma pequena. Os dois coincidem se todos os alunos designados para turma pequena efetivamente frequentam uma turma pequena --- o que pode não ser o caso em presença de não-cumprimento (\textit{non-compliance}): alunos podem mudar de turma, abandonar o experimento ou ingressar em turmas diferentes da designada.

Quando há não-cumprimento parcial, e se aceitarmos a designação original como instrumento válido para a turma efetivamente frequentada, o LATE (\textit{Local Average Treatment Effect}) mede o efeito sobre os \textit{cumpridores} --- aqueles que frequentaram a turma pequena justamente porque foram para ela designados. Neste guia, trabalhamos com o ITT como estimando principal por ser mais diretamente identificado pela aleatorização original.

\begin{exemplobox}[ITT e LATE no STAR]
Suponha que 5\% dos alunos designados para turmas pequenas tenham migrado para turmas regulares ao longo do ano. O ITT estima o efeito médio sobre \textit{todos} os designados para turma pequena --- incluindo os que na prática não ficaram em turma pequena. O LATE estimaria o efeito sobre o subgrupo que cumpriu a designação. Como o não-cumprimento no STAR foi relativamente baixo, os dois tendem a ser numericamente próximos, mas conceitualmente distintos.
\end{exemplobox}

\subsection*{SUTVA: uma hipótese que merece atenção}

A hipótese SUTVA (\textit{Stable Unit Treatment Value Assumption}) afirma que o resultado potencial de um aluno depende apenas do seu próprio tipo de turma, e não do tipo de turma dos outros alunos. No STAR, isso é potencialmente questionável: turmas pequenas e regulares coexistem \textit{dentro} da mesma escola.

Alguns canais potenciais de interferência merecem atenção:

\begin{itemize}[leftmargin=*]
  \item \textbf{Realocação de professoras}: se as melhores professoras da escola são sistematicamente alocadas a turmas pequenas, o efeito estimado captura tanto o tamanho da turma quanto a qualidade docente.
  \item \textbf{Efeitos de pares}: alunos em turmas regulares podem ser afetados por terem colegas que, em outro desenho, estariam em turmas pequenas.
  \item \textbf{Recursos compartilhados}: espaço físico, materiais e tempo dos coordenadores pedagógicos são divididos entre turmas dentro da mesma escola.
  \item \textbf{Composição das turmas}: se alunos com maior desempenho prévio forem sistematicamente alocados a certos tipos de turma, a comparação entre turmas dentro da escola pode ser contaminada.
\end{itemize}

Este exemplo serve para reforçar um ponto metodológico central: mesmo em experimentos aleatórios, é necessário explicitar as hipóteses de identificação e avaliar sua plausibilidade no contexto específico do estudo. A aleatorização resolve o problema de seleção entre grupos, mas não elimina automaticamente todas as ameaças à validade interna.

\begin{atencaobox}[SUTVA não é automática]
O fato de um experimento ter sido aleatorizado não garante que o SUTVA seja satisfeito. Sempre que turmas, indivíduos ou unidades do mesmo grupo de randomização interagem --- o que é comum em experimentos educacionais e comunitários ---, a hipótese de ausência de interferência precisa ser defendida com argumentos substantivos, não apenas assumida.
\end{atencaobox}


%% ─────────────────────────────────────────────────────────────────
\section{Estimação: Da Diferença de Médias à Regressão com Efeitos Fixos}
\label{sec:star_estimacao}

\subsection*{Especificação 1 — Diferença de médias simples}

A estimativa mais direta é a diferença de médias entre alunos designados para turmas pequenas e demais alunos:
\begin{equation}
  Y_i = \alpha + \beta\, D_i + u_i.
  \label{eq:spec1}
\end{equation}

O coeficiente $\hat{\beta}$ compara a nota média dos alunos em turmas pequenas com a nota média dos demais, sem qualquer condicionamento adicional. Sob randomização incondicional, esse estimador seria não-viesado para o ITT. No STAR, porém, a randomização foi feita \textit{dentro} de escola --- o que significa que comparar alunos de escolas diferentes pode introduzir diferenças pré-existentes entre escolas, não controladas pelo experimento.

\subsection*{Especificação 2 — Efeitos fixos de escola}

A especificação mais alinhada ao desenho experimental inclui efeitos fixos de escola $\lambda_{s(i)}$:
\begin{equation}
  Y_i = \alpha + \beta\, D_i + \lambda_{s(i)} + u_i.
  \label{eq:spec2}
\end{equation}

Ao incluir uma dummy para cada escola, a comparação passa a ser feita \textit{dentro} de escola --- exatamente como a randomização foi conduzida. O coeficiente $\hat{\beta}$ agora compara alunos designados para turma pequena com alunos designados para turma regular \textit{na mesma escola}, eliminando qualquer diferença entre escolas que pudesse confundir a estimativa.

\subsection*{Especificação 3 — Efeitos fixos de escola e covariadas}

A terceira especificação adiciona covariadas pré-tratamento $X_i$ (gênero, raça, elegibilidade para merenda):
\begin{equation}
  Y_i = \alpha + \beta\, D_i + \lambda_{s(i)} + X_i'\gamma + u_i.
  \label{eq:spec3}
\end{equation}

As covariadas têm aqui um papel específico. Em um RCT bem implementado, $D_i$ é independente de $X_i$ por construção (dentro de escola), de modo que $\gamma \neq 0$ não implica que omitir $X_i$ causaria viés. O benefício dos controles aqui é outro: ao capturar parte da variação nos resultados $Y_i$, as covariadas reduzem a variância residual $\hat{u}_i$, tornando o estimador $\hat{\beta}$ mais preciso sem sacrificar a validade causal. \citet{Lin2013} formaliza esse argumento e mostra que, em RCTs, a especificação preferida inclui covariadas centradas e interações com o tratamento, garantindo eficiência assintótica.

\begin{notabox}[{Controles em RCTs --- precisão sem causalidade}]
Em estudos observacionais, adicionamos controles para tentar satisfazer a CIA e eliminar viés. Em RCTs, adicionamos controles para ganhar precisão. A distinção é fundamental: no primeiro caso, a validade causal depende dos controles incluídos; no segundo, a validade já está garantida pela randomização, e os controles apenas melhoram a eficiência.
\end{notabox}

\begin{lstlisting}[style=Rstyle, caption={Três especificações progressivas no STAR}]
# -----------------------------
# 1) Especificação 1: diferença de médias simples
# -----------------------------
spec1 <- lm_robust(read_z ~ D, data = dados_k, se_type = "HC2")

# -----------------------------
# 2) Especificação 2: efeitos fixos de escola
# -----------------------------
spec2 <- feols(read_z ~ D | escola, data = dados_k,
               vcov = "HC1")

# -----------------------------
# 3) Especificação 3: EF de escola + covariadas pré-tratamento
# -----------------------------
spec3 <- feols(read_z ~ D + feminino + nao_branco + merenda | escola,
               data = dados_k, vcov = "HC1")

# Sumário das três especificações
etable(spec1, spec2, spec3,
       keep    = "D",
       headers = c("(1) Sem EF", "(2) EF escola", "(3) EF + covariadas"),
       title   = "Efeito de turma pequena sobre notas de leitura (STAR, K)")
\end{lstlisting}

% ──────────────────────────────────────────────────────────────────
% TABELA 5.1 — Três especificações: efeito de turma pequena sobre notas
% Executar o script R para obter os coeficientes, erros-padrão e
% estatísticas de diagnóstico (R², N, número de escolas).
% ──────────────────────────────────────────────────────────────────
\begin{table}[H]
\centering
\caption{Efeito de turma pequena sobre nota de leitura: três especificações}
\label{tab:spec_comparacao}
\begin{tabular}{lccc}
\toprule
 & (1) Sem EF & (2) EF Escola & (3) EF + Cov. \\
\midrule
Turma pequena ($D_i$)
  & \textit{[inserir]} & \textit{[inserir]} & \textit{[inserir]} \\
  & \textit{[inserir SE]} & \textit{[inserir SE]} & \textit{[inserir SE]} \\[4pt]
Efeitos fixos de escola & Não & Sim & Sim \\
Covariadas pré-tratamento & Não & Não & Sim \\
$R^2$ & \textit{[inserir]} & \textit{[inserir]} & \textit{[inserir]} \\
Observações & \textit{[inserir]} & \textit{[inserir]} & \textit{[inserir]} \\
\bottomrule
\end{tabular}

\vspace{4pt}
\figmeta{Dados: \citet{Krueger1999}, via pacote \texttt{AER}}{Elaboração dos autores}{Variável dependente: nota de leitura padronizada (média = 0, DP = 1). Covariadas: gênero, raça e elegibilidade para merenda gratuita. Erros-padrão HC robustos entre parênteses. * $p < 0{,}05$; ** $p < 0{,}01$; *** $p < 0{,}001$.}
\end{table}

A comparação entre as três colunas da Tabela~\ref{tab:spec_comparacao} ilustra um resultado esperado: em um RCT bem implementado, as estimativas nas três especificações devem ser numericamente similares. Eventuais diferenças entre a coluna (1) e as colunas (2) e (3) indicam diferenças pré-existentes entre escolas, que a aleatorização dentro de escola não controla. A redução nos erros-padrão entre (1) e (3) reflete o ganho de precisão proveniente dos controles.


%% ─────────────────────────────────────────────────────────────────
\section{Inferência I: Qual Erro-Padrão Usar?}
\label{sec:star_ep}

A escolha do estimador de variância não é um detalhe técnico --- é parte da análise substantiva. No STAR, há uma razão estrutural clara para não usar erros-padrão clássicos: alunos dentro da mesma turma compartilham professora, rotina pedagógica, material didático e dinâmica de sala. Seus resultados não são observações independentes. Ignorar essa correlação intraclasse leva a erros-padrão artificialmente pequenos e a testes com taxa de falsos positivos acima do nível nominal.

\subsection*{Correlação intraclasse}

A dependência dentro de grupos pode ser quantificada pela \textbf{correlação intraclasse} (ICC):
\begin{equation}
  \rho = \frac{\sigma^2_c}{\sigma^2_c + \sigma^2_i},
  \label{eq:icc}
\end{equation}
onde $\sigma^2_c$ é a variância entre turmas e $\sigma^2_i$ é a variância dentro de turmas. O \textbf{efeito de desenho} (\textit{design effect}) indica o quanto o erro-padrão efetivo supera o erro-padrão clássico em uma amostra com $\bar{n}$ observações por turma:
\begin{equation}
  \text{DEFF} = 1 + (\bar{n} - 1)\,\rho.
  \label{eq:deff}
\end{equation}

Com $\rho = 0{,}20$ e turmas de tamanho médio $\bar{n} = 20$, o DEFF seria $1 + 19 \times 0{,}20 = 4{,}80$ --- o que significa que os erros-padrão clássicos subestimariam a incerteza por um fator de $\sqrt{4{,}80} \approx 2{,}19$. Estatísticas-$t$ baseadas nesses erros estariam infladas em mais de 100\%.

\subsection*{Quatro abordagens}

Comparamos quatro estimadores de variância:

\begin{enumerate}[leftmargin=*, label=(\roman*)]
  \item \textbf{Clássico}: assume homoscedasticidade e independência entre todas as observações.
  \item \textbf{HC robusto}: corrige heteroscedasticidade, mas ainda trata observações como independentes.
  \item \textbf{Clusterizado por turma}: agrupa observações dentro da mesma turma, reconhecendo a dependência intraclasse. É a escolha natural dado que o tratamento é atribuído no nível da turma.
  \item \textbf{Clusterizado por escola}: agrupa por escola. Mais conservador, com número menor de clusters (79 escolas vs.\ $\sim$300 turmas).
\end{enumerate}

\begin{lstlisting}[style=Rstyle, caption={Comparação de quatro estimadores de variância no STAR}]
# -----------------------------
# 1) Modelo base: EF de escola, sem outras covariadas
# -----------------------------
mod_base <- lm(read_z ~ D + escola, data = dados_k)

# -----------------------------
# 2) Quatro variantes de erro-padrão
# -----------------------------
# (i) Clássico
se_classico <- sqrt(diag(vcov(mod_base)))["D"]

# (ii) HC2 robusto
se_hc2 <- sqrt(vcovHC(mod_base, type = "HC2"))["D", "D"]

# (iii) Clusterizado por turma
dados_k <- dados_k %>%
  mutate(turma_id = paste(schoolk, stark, sep = "_"))

se_turma <- sqrt(vcovCL(mod_base, cluster = ~turma_id,
                        data = dados_k))["D", "D"]

# (iv) Clusterizado por escola
se_escola <- sqrt(vcovCL(mod_base, cluster = ~escola,
                         data = dados_k))["D", "D"]

# Apresentação
resultados_ep <- data.frame(
  Metodo   = c("Clássico", "HC2", "Cluster turma", "Cluster escola"),
  SE       = round(c(se_classico, se_hc2, se_turma, se_escola), 4),
  t_stat   = round(coef(mod_base)["D"] /
             c(se_classico, se_hc2, se_turma, se_escola), 2)
)
print(resultados_ep)
\end{lstlisting}

% ──────────────────────────────────────────────────────────────────
% TABELA 5.2 — Comparação de erros-padrão
% Executar o script R para obter SE, estatística t e p-valor
% sob as quatro abordagens.
% ──────────────────────────────────────────────────────────────────
\begin{table}[H]
\centering
\caption{Estimativa do efeito de turma pequena sob diferentes estimadores de variância}
\label{tab:erros_padrao}
\begin{tabular}{lcccc}
\toprule
Abordagem & Estimativa & Erro-padrão & Estatística $t$ & Valor-$p$ \\
\midrule
Clássico              & \textit{[inserir]} & \textit{[inserir]} & \textit{[inserir]} & \textit{[inserir]} \\
HC2 robusto           & \textit{[inserir]} & \textit{[inserir]} & \textit{[inserir]} & \textit{[inserir]} \\
Cluster por turma     & \textit{[inserir]} & \textit{[inserir]} & \textit{[inserir]} & \textit{[inserir]} \\
Cluster por escola    & \textit{[inserir]} & \textit{[inserir]} & \textit{[inserir]} & \textit{[inserir]} \\
\bottomrule
\end{tabular}

\vspace{4pt}
\figmeta{Dados: \citet{Krueger1999}}{Elaboração dos autores}{Especificação com efeitos fixos de escola. Variável dependente: nota de leitura padronizada. A estimativa pontual é a mesma nas quatro linhas; apenas o estimador de variância varia.}
\end{table}

O padrão esperado na Tabela~\ref{tab:erros_padrao} é que os erros-padrão clusterizados sejam maiores que os clássicos e robustos. A comparação entre clusterização por turma e por escola ilustra o \textit{trade-off} entre adequação ao processo gerador de dados e estabilidade do estimador: com 79 escolas, o estimador sandwich clusterizado por escola opera próximo ao limite inferior recomendado ($G \approx 30$--50), enquanto as turmas oferecem maior número de clusters.

\begin{notabox}[Regra prática de clustering no STAR]
Como o tratamento foi atribuído no nível da \textit{turma}, a clusterização por turma é a escolha teoricamente correta. Se os dados contiverem identificadores de turma confiáveis, devem ser usados. A clusterização por escola é mais conservadora e pode ser reportada como verificação de robustez.
\end{notabox}


%% ─────────────────────────────────────────────────────────────────
\section{Inferência II: Bootstrap}
\label{sec:star_bootstrap}

A Seção~\ref{sec:clustering} mostrou que, com número razoável de clusters, o estimador sandwich clusterizado tende a ter bom desempenho assintótico. Com cerca de 300 turmas, o STAR está bem acima do limiar mínimo --- o que sugere que os intervalos de confiança assintóticos clusterizados devem ser confiáveis. O \textit{bootstrap} é útil aqui como verificação independente e como método para estatísticas cuja distribuição assintótica é menos direta.

\subsection*{Por que reamostrar turmas, não alunos}

Um erro frequente é aplicar o bootstrap padrão --- reamostrar observações individuais com reposição --- em dados com estrutura de clustering. Ao sortear alunos individualmente, o bootstrap padrão quebra a dependência dentro de turmas: uma amostra bootstrap poderia conter cinco cópias do mesmo aluno e nenhuma cópia de sua turma original, destruindo exatamente a estrutura que torna a clusterização necessária. A solução é o \textbf{\textit{cluster bootstrap}}: reamostrar \textit{turmas inteiras} com reposição, mantendo todos os alunos de cada turma sorteada.

\begin{lstlisting}[style=Rstyle, caption={Cluster bootstrap por turma no STAR}]
# -----------------------------
# 1) Cluster bootstrap: reamostrar turmas inteiras
# -----------------------------
rm(list = ls())
# (assumindo que dados_k e turma_id foram criados anteriormente)

set.seed(42)
B       <- 1999
turmas  <- unique(dados_k$turma_id)
G       <- length(turmas)

boot_coefs <- numeric(B)

for (b in 1:B) {
  # Sortear turmas com reposição
  turmas_b  <- sample(turmas, size = G, replace = TRUE)

  # Reconstruir base com todas as observações das turmas sorteadas
  amostra_b <- bind_rows(
    lapply(turmas_b, function(t) dados_k[dados_k$turma_id == t, ])
  )

  # Estimar modelo com EF de escola
  mod_b <- lm(read_z ~ D + escola, data = amostra_b)
  boot_coefs[b] <- coef(mod_b)["D"]
}

# IC bootstrap percentil (95%)
ic_boot <- quantile(boot_coefs, c(0.025, 0.975))
cat("IC bootstrap 95% (percentil): [",
    round(ic_boot[1], 3), ";", round(ic_boot[2], 3), "]\n")

# Comparar com IC assintótico clusterizado
se_assintotico <- sqrt(vcovCL(lm(read_z ~ D + escola, data = dados_k),
                              cluster = ~turma_id)["D", "D"])
est_pont <- coef(lm(read_z ~ D + escola, data = dados_k))["D"]
ic_assint <- est_pont + c(-1, 1) * 1.96 * se_assintotico

cat("IC assintótico 95% (cluster): [",
    round(ic_assint[1], 3), ";", round(ic_assint[2], 3), "]\n")
\end{lstlisting}

A mensagem didática desta subseção é a seguinte: quando há número razoável de clusters e o desenho é bem comportado, os ICs bootstrap e assintóticos clusterizados devem ser numericamente próximos. Uma divergência grande entre os dois seria um sinal de alerta --- sugerindo que alguma hipótese subjacente (número suficiente de clusters, distribuição simétrica do estimador, ausência de observações influentes) pode não estar sendo satisfeita. No STAR, com cerca de 300 turmas, a convergência entre os dois é esperada.


%% ─────────────────────────────────────────────────────────────────
\section{Inferência III: Quatro Disciplinas, Um Problema --- Múltiplos Testes}
\label{sec:star_multiplos}

O STAR oferece resultados de desempenho em mais de uma disciplina: leitura e matemática, ao menos. Quando testamos o efeito de turmas pequenas em todas as disciplinas disponíveis, deparamo-nos com o problema de múltiplos testes discutido na Seção~\ref{sec:multiplos_testes}: quanto maior o número de testes, maior a probabilidade de encontrar ao menos um resultado estatisticamente significativo por acaso, mesmo que o efeito verdadeiro seja nulo em todos os desfechos.

Suponha que testemos o efeito em quatro resultados (por exemplo, leitura e matemática no jardim de infância e no primeiro ano). Ao nível $\alpha = 0{,}05$ com quatro testes independentes, a probabilidade de ao menos uma rejeição espúria é $1 - 0{,}95^4 \approx 18{,}5\%$ --- quase quatro vezes o nível nominal.

\begin{lstlisting}[style=Rstyle, caption={Múltiplos testes: quatro resultados com correção de Bonferroni e BH-FDR}]
# -----------------------------
# 1) Estimar o efeito nas quatro variáveis de resultado
# Leitura (K), Matemática (K), Leitura (1o ano), Matemática (1o ano)
# -----------------------------
dados_multi <- STAR %>%
  filter(!is.na(stark)) %>%
  mutate(
    D      = as.integer(stark == "small"),
    escola = as.factor(schoolk),
    read_k_z  = scale(readk)[, 1],
    math_k_z  = scale(mathk)[, 1],
    read_1_z  = scale(read1)[, 1],
    math_1_z  = scale(math1)[, 1]
  )

resultados <- c("read_k_z", "math_k_z", "read_1_z", "math_1_z")
nomes      <- c("Leitura (K)", "Matemática (K)",
                "Leitura (1o ano)", "Matemática (1o ano)")

pvalores    <- numeric(length(resultados))
estimativas <- numeric(length(resultados))

for (j in seq_along(resultados)) {
  df_j  <- dados_multi[!is.na(dados_multi[[resultados[j]]]), ]
  mod_j <- lm(as.formula(paste(resultados[j], "~ D + escola")),
              data = df_j)
  se_j  <- sqrt(vcovCL(mod_j, cluster = ~schoolk, data = df_j)["D", "D"])
  estimativas[j] <- coef(mod_j)["D"]
  pvalores[j]    <- 2 * pt(-abs(coef(mod_j)["D"] / se_j),
                            df = df_j$escola %>% n_distinct() - 2)
}

# -----------------------------
# 2) Ajuste por Bonferroni e BH-FDR
# -----------------------------
pv_bonf <- p.adjust(pvalores, method = "bonferroni")
pv_bh   <- p.adjust(pvalores, method = "BH")

tabela_multi <- data.frame(
  Resultado    = nomes,
  Estimativa   = round(estimativas, 3),
  P_bruto      = round(pvalores, 3),
  P_Bonferroni = round(pv_bonf, 3),
  P_BH_FDR     = round(pv_bh, 3)
)
print(tabela_multi)
\end{lstlisting}

% ──────────────────────────────────────────────────────────────────
% TABELA 5.3 — Múltiplos testes: quatro resultados
% Executar o script R para obter estimativas, p-valores brutos
% e p-valores ajustados por Bonferroni e BH-FDR.
% ──────────────────────────────────────────────────────────────────
\begin{table}[H]
\centering
\caption{Efeito de turma pequena: quatro resultados com ajuste para múltiplos testes}
\label{tab:multiplos_testes_star}
\begin{tabular}{lcccc}
\toprule
Resultado & Estimativa & $p$ bruto & $p$ Bonferroni & $p$ BH-FDR \\
\midrule
Leitura (K)          & \textit{[inserir]} & \textit{[inserir]} & \textit{[inserir]} & \textit{[inserir]} \\
Matemática (K)       & \textit{[inserir]} & \textit{[inserir]} & \textit{[inserir]} & \textit{[inserir]} \\
Leitura (1.º ano)    & \textit{[inserir]} & \textit{[inserir]} & \textit{[inserir]} & \textit{[inserir]} \\
Matemática (1.º ano) & \textit{[inserir]} & \textit{[inserir]} & \textit{[inserir]} & \textit{[inserir]} \\
\bottomrule
\end{tabular}

\vspace{4pt}
\figmeta{Dados: \citet{Krueger1999}}{Elaboração dos autores}{Especificação com efeitos fixos de escola. Erros-padrão clusterizados por escola. Bonferroni: nível ajustado $= 0{,}05/4 = 0{,}0125$. BH-FDR: controla proporção esperada de falsas descobertas a 5\%.}
\end{table}

A leitura da Tabela~\ref{tab:multiplos_testes_star} deve responder a três perguntas: (i) quais efeitos são significativos sem qualquer ajuste? (ii) quais sobrevivem ao ajuste de Bonferroni? (iii) o padrão é consistente entre disciplinas e anos, ou concentrado em um único resultado? Uma estimativa que se mantém significativa após Bonferroni sugere evidência mais sólida do que aquela que só resiste ao p-valor bruto. Quando os efeitos são positivos e de magnitude similar em todas as disciplinas, a credibilidade aumenta: seria improvável que o mesmo viés produzisse o mesmo sinal em quatro desfechos distintos.

\begin{notabox}[Bonferroni vs.\ BH-FDR]
Bonferroni controla a probabilidade de \textit{ao menos} um falso positivo no conjunto de testes (FWER), tornando cada teste individual muito conservador. BH-FDR controla a proporção esperada de falsas descobertas entre as rejeições, sendo menos conservador quando $m$ é grande. Para quatro desfechos, a diferença entre os dois é pequena --- Bonferroni com $\alpha^* = 0{,}0125$ vs.\ BH com limiar adaptativo. Com dezenas ou centenas de desfechos, a diferença torna-se mais relevante.
\end{notabox}


%% ─────────────────────────────────────────────────────────────────
\section{Inferência IV: Inferência por Permutação}
\label{sec:star_permutacao}

A inferência por permutação, apresentada na Seção~\ref{sec:design_based}, tem uma propriedade especial no STAR: como o processo de aleatorização está documentado (randomização dentro de escola), a distribuição de referência pode ser construída a partir desse mesmo processo, sem qualquer hipótese distribucional sobre os erros.

\subsection*{Respeitando o desenho na permutação}

Como a aleatorização ocorreu \textit{dentro} de escola, a permutação deve respeitar esse bloco: permutamos os rótulos de tratamento ($D_i$) apenas entre alunos da mesma escola, preservando o número de tratados por escola. Permutar globalmente ignoraria a estrutura por blocos e produziria distribuições de referência inadequadas.

\begin{lstlisting}[style=Rstyle, caption={Inferência por permutação com aleatorização em bloco (escola)}]
# -----------------------------
# 1) Estatística observada
# -----------------------------
mod_obs  <- lm(read_z ~ D + escola, data = dados_k)
t_obs    <- coef(mod_obs)["D"] /
            sqrt(vcovCL(mod_obs, cluster = ~turma_id,
                        data = dados_k)["D", "D"])

cat("Estatística t observada:", round(t_obs, 3), "\n")

# -----------------------------
# 2) Distribuição de permutação
# Permutar D dentro de escola (respeitar o bloco de randomização)
# -----------------------------
set.seed(42)
P <- 4999
t_perm <- numeric(P)

for (p in 1:P) {
  dados_perm <- dados_k %>%
    group_by(escola) %>%
    mutate(D_perm = sample(D)) %>%   # permuta D dentro de escola
    ungroup()

  mod_p   <- lm(read_z ~ D_perm + escola, data = dados_perm)
  se_p    <- sqrt(vcovCL(mod_p, cluster = ~turma_id,
                          data = dados_perm)["D_perm", "D_perm"])
  t_perm[p] <- coef(mod_p)["D_perm"] / se_p
}

# -----------------------------
# 3) P-valor de permutação (bilateral)
# -----------------------------
p_perm    <- mean(abs(t_perm) >= abs(t_obs))
p_assint  <- 2 * pnorm(-abs(t_obs))

cat("P-valor de permutação (bilateral): ", round(p_perm, 4), "\n")
cat("P-valor assintótico (normal):      ", round(p_assint, 4), "\n")
\end{lstlisting}

A comparação entre o p-valor de permutação e o p-valor assintótico clusterizado ilustra um resultado típico de experimentos com amostras grandes: os dois tendem a ser numericamente próximos. A inferência por permutação é particularmente atraente porque usa diretamente o processo de aleatorização descrito no protocolo experimental --- não há necessidade de aproximações assintóticas nem hipóteses sobre a distribuição dos erros. Em amostras menores, com poucos clusters ou distribuições assimétricas, a inferência por permutação pode diferir substancialmente dos métodos assintóticos e, nesse caso, oferece garantias de validade mais robustas.


%% ─────────────────────────────────────────────────────────────────
\section{Testes de Falsificação}
\label{sec:star_falsificacao}

Os testes de falsificação não ``provam'' que o STAR tem identificação válida, mas ajudam a detectar problemas que comprometeriam a interpretação causal. Apresentamos três verificações: balanceamento de covariadas pré-tratamento, atrito diferencial e balanceamento das características das professoras.

\subsection*{Balanceamento de covariadas pré-tratamento}

Em um experimento bem implementado, características pré-tratamento dos alunos não devem prever sistematicamente a designação para turma pequena. Para testar isso, regredimos $D_i$ (indicador de turma pequena) nas covariadas pré-tratamento disponíveis e realizamos um F-teste conjunto de significância:
\begin{equation}
  D_i = \alpha + \lambda_{s(i)} + \delta_1\,\text{feminino}_i + \delta_2\,\text{nao\_branco}_i + \delta_3\,\text{merenda}_i + v_i.
  \label{eq:teste_balanceamento}
\end{equation}

A hipótese nula é $H_0: \delta_1 = \delta_2 = \delta_3 = 0$, ou seja, nenhuma covariada prediz o tratamento dentro de escola. Sob aleatorização bem conduzida, esperamos não rejeitar essa hipótese.

\begin{lstlisting}[style=Rstyle, caption={F-teste de balanceamento de covariadas}]
# -----------------------------
# 1) Regredir tratamento em covariadas pré-tratamento + EF escola
# -----------------------------
mod_bal <- lm(D ~ escola + feminino + nao_branco + merenda,
              data = dados_k)

# F-teste conjunto: apenas para as covariadas (excluindo EF de escola)
library(car)
ftest_bal <- linearHypothesis(mod_bal,
  c("feminino = 0", "nao_branco = 0", "merenda = 0"),
  vcov. = vcovCL(mod_bal, cluster = ~turma_id, data = dados_k))

cat("=== Teste F conjunto de balanceamento ===\n")
print(ftest_bal)

# -----------------------------
# 2) Tabela de médias por grupo (dentro de escola)
# -----------------------------
dados_k %>%
  group_by(D) %>%
  summarise(
    pct_feminino   = mean(feminino, na.rm = TRUE),
    pct_nao_branco = mean(nao_branco, na.rm = TRUE),
    pct_merenda    = mean(merenda, na.rm = TRUE),
    n              = n()
  ) %>%
  print()
\end{lstlisting}

% ──────────────────────────────────────────────────────────────────
% TABELA 5.4 — Teste de balanceamento: covariadas pré-tratamento
% Executar o script R para obter as médias por grupo e a estatística F.
% ──────────────────────────────────────────────────────────────────
\begin{table}[H]
\centering
\caption{Verificação de balanceamento: covariadas pré-tratamento por condição experimental}
\label{tab:balanceamento}
\begin{tabular}{lccc}
\toprule
Covariada & Turma pequena ($D=1$) & Controle ($D=0$) & Dif. padronizada \\
\midrule
Feminino (\%)         & \textit{[inserir]} & \textit{[inserir]} & \textit{[inserir]} \\
Não-branco (\%)       & \textit{[inserir]} & \textit{[inserir]} & \textit{[inserir]} \\
Merenda gratuita (\%) & \textit{[inserir]} & \textit{[inserir]} & \textit{[inserir]} \\
\midrule
\multicolumn{3}{l}{F-teste conjunto ($H_0$: todas as diferenças = 0)}
  & $F = $ \textit{[inserir]} \\
\multicolumn{3}{l}{Valor-$p$}
  & \textit{[inserir]} \\
\midrule
Observações           & \textit{[inserir]} & \textit{[inserir]} & \\
\bottomrule
\end{tabular}

\vspace{4pt}
\figmeta{Dados: \citet{Krueger1999}}{Elaboração dos autores}{Comparação feita dentro de escola (efeitos fixos de escola incluídos na regressão). Dif. padronizada $= (\bar{X}_1 - \bar{X}_0)/\sqrt{(s_1^2 + s_0^2)/2}$.}
\end{table}

\subsection*{Atrito diferencial}

A aleatorização garante comparabilidade no momento inicial do experimento. Ao longo dos quatro anos, no entanto, alunos podem sair do experimento --- por mudança de escola, ausência prolongada ou dados faltantes. Se a \textit{probabilidade de permanecer na amostra} for diferente entre turmas pequenas e regulares, a comparação ao final do período pode estar contaminada: estaríamos comparando grupos sistematicamente diferentes dos originalmente sorteados.

Para verificar atrito diferencial, testamos se $D_i$ prediz a probabilidade de o aluno ter notas disponíveis no ano de interesse:
\begin{equation}
  \mathds{1}[Y_i \text{ observado}] = \alpha + \beta_{\text{atrito}}\, D_i + \lambda_{s(i)} + u_i.
  \label{eq:atrito}
\end{equation}

Um coeficiente $\hat{\beta}_{\text{atrito}}$ significativamente diferente de zero indica que tratados e controles têm taxas de atrito distintas --- o que comprometeria a validade interna mesmo com aleatorização original correta.

\begin{lstlisting}[style=Rstyle, caption={Verificação de atrito diferencial}]
# -----------------------------
# Indicador de permanência na amostra para cada ano
# Testamos se D prediz estar na amostra no 1o ano
# -----------------------------
dados_atrito <- STAR %>%
  filter(!is.na(stark)) %>%
  mutate(
    D              = as.integer(stark == "small"),
    escola         = as.factor(schoolk),
    presente_ano1  = as.integer(!is.na(read1))
  )

mod_atrito <- lm_robust(presente_ano1 ~ D,
                        fixed_effects = ~escola,
                        clusters      = schoolk,
                        data          = dados_atrito,
                        se_type       = "CR2")

cat("=== Teste de atrito diferencial (1o ano) ===\n")
cat("Coeficiente D:  ", round(coef(mod_atrito)["D"], 4), "\n")
se_atrito <- summary(mod_atrito)$coefficients["D", 2]
pv_atrito <- summary(mod_atrito)$coefficients["D", 4]
cat("Erro-padrão:    ", round(se_atrito, 4), "\n")
cat("Valor-p:        ", round(pv_atrito, 3), "\n")
\end{lstlisting}

\subsection*{Placebo de professora: características docentes como teste de implementação}

Um aspecto do STAR que merece atenção é a qualidade da implementação. Se as professoras mais experientes ou mais tituladas foram sistematicamente alocadas a turmas pequenas, o efeito estimado confundiria o tamanho da turma com a qualidade docente. Testamos isso verificando o balanceamento das características das professoras entre condições experimentais:

\begin{equation}
  \text{experiência}_{c} = \alpha + \beta_{\text{prof}}\, D_c + \lambda_{s(c)} + w_c,
  \label{eq:placebo_prof}
\end{equation}

onde $c$ indexa turmas e $D_c$ é o indicador de turma pequena. Se $\hat{\beta}_{\text{prof}}$ for significativamente diferente de zero, há evidência de que turmas pequenas e regulares foram atribuídas a professoras sistematicamente diferentes.

\begin{lstlisting}[style=Rstyle, caption={Balanceamento de características das professoras entre condições}]
# -----------------------------
# Agregar para o nível de turma (cada turma tem uma professora)
# -----------------------------
dados_prof <- STAR %>%
  filter(!is.na(stark), !is.na(experiencek)) %>%
  mutate(D = as.integer(stark == "small")) %>%
  group_by(schoolk, stark) %>%
  summarise(
    experiencia = first(experiencek),
    D           = first(D),
    .groups = "drop"
  ) %>%
  mutate(escola = as.factor(schoolk))

mod_prof <- lm_robust(experiencia ~ D,
                      fixed_effects = ~escola,
                      clusters      = schoolk,
                      data          = dados_prof,
                      se_type       = "CR2")

cat("=== Placebo de professora: experiência docente ===\n")
cat("Coeficiente D:  ", round(coef(mod_prof)["D"], 3), "\n")
se_prof <- summary(mod_prof)$coefficients["D", 2]
pv_prof <- summary(mod_prof)$coefficients["D", 4]
cat("Erro-padrão:    ", round(se_prof, 3), "\n")
cat("Valor-p:        ", round(pv_prof, 3), "\n")
\end{lstlisting}

\begin{atencaobox}[Testes de falsificação: limites e interpretação]
Passar em todos os testes de falsificação é evidência circunstancial favorável à validade do desenho, mas não é prova de causalidade. Uma estratégia pode superar todos os placebos disponíveis e ainda ter identificação inválida por razões não testáveis --- por exemplo, um canal de interferência entre turmas que não deixa rastro nas covariadas observadas. Os testes de falsificação devem ser lidos em conjunto com o argumento teórico sobre o mecanismo de identificação, não como substitutos a ele.
\end{atencaobox}


%% ─────────────────────────────────────────────────────────────────
\subsection*{Fechamento: o que o STAR ensinou neste guia}

O Experimento STAR percorreu, em um único exemplo empírico, toda a cadeia metodológica discutida neste guia. A \textbf{aleatorização dentro de escola} resolveu o problema de identificação: dentro de cada escola, tratados e controles são comparáveis em média, eliminando o viés de seleção. A \textbf{regressão com efeitos fixos de escola} alinhou a estimação ao desenho experimental, melhorando precisão sem comprometer validade. A discussão sobre \textbf{erros-padrão} revelou que a escolha do estimador de variância é uma decisão substantiva: ignorar a correlação intraclasse leva a inferências enganosas. O \textbf{bootstrap} por cluster confirmou os resultados assintóticos quando o número de grupos é suficiente, e mostrou quando os dois tendem a divergir. A análise de \textbf{múltiplos testes} forçou uma resposta honesta sobre quais resultados são robustos: um efeito que sobrevive a Bonferroni em quatro disciplinas é mais convincente do que um único p-valor favorável. A \textbf{inferência por permutação} usou diretamente o processo de aleatorização para construir distribuições de referência exatas, sem aproximações assintóticas. Por fim, os \textbf{testes de falsificação} --- balanceamento de covariadas, atrito diferencial e placebo de professora --- forneceram evidência circunstancial sobre a qualidade da implementação experimental.

O ponto que conecta todos esses elementos é o seguinte: cada ferramenta responde a uma pergunta diferente. Identificação pergunta \textit{o que estamos estimando?} Estimação pergunta \textit{como calculamos isso?} Inferência pergunta \textit{quanta incerteza há?} Bootstrap e permutação perguntam \textit{essa incerteza é bem caracterizada pelas aproximações assintóticas?} Múltiplos testes perguntam \textit{o resultado sobrevive a buscas seletivas?} Falsificação pergunta \textit{as hipóteses por trás da estimação são plausíveis?} Nenhuma dessas perguntas é dispensável, e a resposta a qualquer uma delas pressupõe que as anteriores tenham sido respondidas com rigor.
%%
%%  Dados: AER::STAR (pacote AER do R)
%%  Referência principal: Krueger (1999)
%% ─────────────────────────────────────────────────────────────────

\chapter{Aplicação Empírica: O Experimento STAR}
\label{sec:star}
\label{cap:star}

As seções anteriores construíram o arcabouço conceitual da inferência causal: resultados potenciais, identificação, estimação, erros-padrão, testes de hipótese, bootstrap e falsificação. O objetivo desta seção é usar um experimento clássico como laboratório empírico para ver esses conceitos funcionando juntos. O Experimento STAR --- \textit{Student/Teacher Achievement Ratio} --- oferece um dos melhores exemplos disponíveis de avaliação experimental em políticas educacionais. Ele é suficientemente limpo para ilustrar a lógica da identificação por aleatorização, mas suficientemente rico para revelar as complexidades práticas de inferência, múltiplos testes e falsificação.

O experimento foi conduzido no estado americano do Tennessee entre 1985 e 1989 \citep{Finn1990, Krueger1999}. A pergunta central é direta: turmas menores melhoram o desempenho escolar? Para respondê-la com rigor causal, o estado financiou um experimento em larga escala em que estudantes e professoras foram aleatoriamente designados a três tipos de turma: \textbf{turma pequena} (13--17 alunos), \textbf{turma regular} (22--26 alunos) e \textbf{turma regular com auxiliar} (22--26 alunos com uma professora auxiliar). O acompanhamento ocorreu ao longo de quatro anos, do jardim de infância (pré-escola, \textit{kindergarten}) ao terceiro ano do ensino fundamental (K--3).

O STAR é um caso de referência na literatura de avaliação de impacto por três razões. Primeira, é um dos poucos experimentos aleatórios de grande escala realizados em educação --- campo em que, por razões éticas e logísticas, experimentos são raros. Segunda, sua metodologia está bem documentada, o que permite verificar e replicar cada escolha empírica. Terceira, os dados tornaram-se públicos e foram explorados por dezenas de pesquisadores, criando uma literatura acumulada que serve de referência para discussão metodológica \citep{Chetty2011}. Neste guia, não temos o objetivo de contribuir com nova pesquisa sobre educação: usamos o STAR como plataforma para ilustrar, com um exemplo concreto, os métodos discutidos nas seções anteriores.


%% ─────────────────────────────────────────────────────────────────
\section{Dados e Desenho Experimental}
\label{sec:star_dados}

\subsection*{Estrutura da amostra}

O experimento acompanhou aproximadamente 11.600 alunos distribuídos em cerca de 300 turmas em 79 escolas do Tennessee. As 79 escolas foram selecionadas de forma a representar áreas rurais, suburbanas e urbanas do estado. Dentro de cada escola, alunos e professoras foram aleatoriamente designados a um dos três tipos de turma.

É importante distinguir três unidades de análise que coexistem nos dados:

\begin{itemize}[leftmargin=*]
  \item \textbf{Aluno ($i$)}: unidade de observação. Os dados contêm, para cada aluno, notas em diferentes disciplinas, características demográficas e indicadores de situação socioeconômica.
  \item \textbf{Turma ($c$)}: unidade de tratamento. O tratamento --- tipo de turma --- é atribuído no nível da turma, não do aluno individualmente. Todos os alunos de uma mesma turma recebem exatamente o mesmo tratamento.
  \item \textbf{Escola ($s$)}: unidade de randomização. A aleatorização foi feita \textit{dentro} de cada escola: dentro da escola $s$, alunos foram sorteados para turmas pequenas ou regulares. A composição das turmas \textit{entre} escolas não foi randomizada.
\end{itemize}

Essa estrutura hierárquica --- alunos dentro de turmas dentro de escolas --- tem implicações diretas para a estimação e para a inferência, como veremos ao longo desta seção.

\subsection*{Variáveis disponíveis}

O conjunto de dados contém as seguintes informações por aluno:

\begin{itemize}[leftmargin=*]
  \item \textbf{Notas}: pontuação em leitura e matemática para cada ano (jardim de infância ao terceiro ano). São as variáveis de resultado.
  \item \textbf{Tipo de turma designado}: indicador do braço do experimento ao qual o aluno foi alocado (\texttt{small}, \texttt{regular}, \texttt{regular+aide}).
  \item \textbf{Características do aluno}: raça/etnia, gênero, elegibilidade para merenda gratuita (proxy de renda familiar). Essas são covariadas pré-tratamento.
  \item \textbf{Identificadores}: código da escola e da turma.
  \item \textbf{Características da professora}: anos de experiência e titulação, quando disponíveis. Usadas nos testes de falsificação.
\end{itemize}

\subsection*{Rotina de análise em R}

O script de análise carrega os dados do pacote \texttt{AER} do R \citep{Zeileis2008}, que distribui uma versão dos dados do STAR em formato pronto para análise. O script:

\begin{enumerate}[leftmargin=*, label=(\roman*)]
  \item Carrega os dados e seleciona a amostra analítica do jardim de infância, eliminando observações com informação faltante nas variáveis de interesse.
  \item Cria o indicador binário de tratamento $D_i = \mathds{1}[\texttt{stark} = \texttt{\textquotedbl small\textquotedbl}]$, separando turma pequena das duas condições de controle agrupadas.
  \item Padroniza as notas de leitura e matemática (média zero, desvio-padrão um) para facilitar a interpretação dos coeficientes em desvios-padrão.
  \item Define covariadas pré-tratamento: gênero, raça e elegibilidade para merenda gratuita.
  \item Constrói os objetos usados nas estimações e tabelas das subseções seguintes.
\end{enumerate}

\begin{lstlisting}[style=Rstyle, caption={Carregamento e preparação dos dados STAR}]
# -----------------------------
# 1) Configuração inicial
# -----------------------------
rm(list = ls())

library(AER)       # dados STAR
library(dplyr)     # manipulação
library(fixest)    # efeitos fixos (feols)
library(estimatr)  # lm_robust / erros clusterizados
library(sandwich)  # vcovCL
library(lmtest)    # coeftest

# -----------------------------
# 2) Carregar e filtrar dados
# Foco no jardim de infância (ano k)
# -----------------------------
data("STAR", package = "AER")

dados_k <- STAR %>%
  filter(!is.na(stark), !is.na(readk), !is.na(mathk)) %>%
  mutate(
    # Tratamento binário: turma pequena vs. demais
    D       = as.integer(stark == "small"),
    # Notas padronizadas (desvio-padrão = 1)
    read_z  = scale(readk)[, 1],
    math_z  = scale(mathk)[, 1],
    # Covariadas pré-tratamento
    feminino    = as.integer(gender == "female"),
    nao_branco  = as.integer(ethnicity != "cauc"),
    merenda     = as.integer(lunchk == "free"),
    # Identificadores
    escola  = as.factor(schoolk)
  )

cat("Observações na amostra analítica:", nrow(dados_k), "\n")
cat("Escolas:                         ", n_distinct(dados_k$escola), "\n")
cat("Tratados (turma pequena):        ", sum(dados_k$D), "\n")
cat("Controle:                        ", sum(1 - dados_k$D), "\n")
\end{lstlisting}


%% ─────────────────────────────────────────────────────────────────
\section{Identificação: Por que a Aleatorização Resolve o Problema}
\label{sec:star_identificacao}

\subsection*{O estimando e a estratégia de identificação}

Para cada aluno $i$, definimos dois resultados potenciais:
\begin{align}
  Y_i(1) &\equiv \text{nota de } i \text{ se designado para turma pequena,} \\
  Y_i(0) &\equiv \text{nota de } i \text{ se designado para turma regular.}
\end{align}

Seja $s(i)$ a escola do aluno $i$ e $D_i \in \{0,1\}$ o indicador de designação para turma pequena. O estimando de interesse é o \textbf{efeito médio da designação para turma pequena} --- um ITT (\textit{Intent-to-Treat}):
\begin{equation}
  \text{ITT} = \mathbb{E}[Y_i(1) - Y_i(0)].
  \label{eq:itt_star}
\end{equation}

Como a randomização foi realizada \textit{dentro} de cada escola, a hipótese de identificação relevante é a independência condicional na escola:
\begin{equation}
  \bigl(Y_i(1),\, Y_i(0)\bigr) \perp\!\!\!\perp D_i \;\big|\; s(i).
  \label{eq:cia_escola}
\end{equation}

Em linguagem intuitiva: dentro de uma mesma escola, alunos designados para turmas pequenas e alunos designados para turmas regulares são comparáveis em média, porque o sorteio foi conduzido apenas entre os alunos matriculados naquela escola. Portanto, diferenças sistemáticas de desempenho entre esses grupos, após condicionar na escola, podem ser interpretadas como efeito causal da designação para turma menor.

A hipótese \eqref{eq:cia_escola} implica imediatamente que:
\begin{equation}
  \mathbb{E}[\varepsilon_i \mid D_i, s(i)] = 0,
  \label{eq:exog_escola}
\end{equation}
onde $\varepsilon_i = Y_i - \mathbb{E}[Y_i \mid D_i, s(i)]$ é o erro residual após condicionar em escola e tratamento. Não há viés de seleção dentro de escola.

\subsection*{ITT vs.\ LATE}

O ITT mede o efeito da \textit{designação} para turma pequena, não necessariamente o efeito de \textit{frequentar} uma turma pequena. Os dois coincidem se todos os alunos designados para turma pequena efetivamente frequentam uma turma pequena --- o que pode não ser o caso em presença de não-cumprimento (\textit{non-compliance}): alunos podem mudar de turma, abandonar o experimento ou ingressar em turmas diferentes da designada.

Quando há não-cumprimento parcial, e se aceitarmos a designação original como instrumento válido para a turma efetivamente frequentada, o LATE (\textit{Local Average Treatment Effect}) mede o efeito sobre os \textit{cumpridores} --- aqueles que frequentaram a turma pequena justamente porque foram para ela designados. Neste guia, trabalhamos com o ITT como estimando principal por ser mais diretamente identificado pela aleatorização original.

\begin{exemplobox}[ITT e LATE no STAR]
Suponha que 5\% dos alunos designados para turmas pequenas tenham migrado para turmas regulares ao longo do ano. O ITT estima o efeito médio sobre \textit{todos} os designados para turma pequena --- incluindo os que na prática não ficaram em turma pequena. O LATE estimaria o efeito sobre o subgrupo que cumpriu a designação. Como o não-cumprimento no STAR foi relativamente baixo, os dois tendem a ser numericamente próximos, mas conceitualmente distintos.
\end{exemplobox}

\subsection*{SUTVA: uma hipótese que merece atenção}

A hipótese SUTVA (\textit{Stable Unit Treatment Value Assumption}) afirma que o resultado potencial de um aluno depende apenas do seu próprio tipo de turma, e não do tipo de turma dos outros alunos. No STAR, isso é potencialmente questionável: turmas pequenas e regulares coexistem \textit{dentro} da mesma escola.

Alguns canais potenciais de interferência merecem atenção:

\begin{itemize}[leftmargin=*]
  \item \textbf{Realocação de professoras}: se as melhores professoras da escola são sistematicamente alocadas a turmas pequenas, o efeito estimado captura tanto o tamanho da turma quanto a qualidade docente.
  \item \textbf{Efeitos de pares}: alunos em turmas regulares podem ser afetados por terem colegas que, em outro desenho, estariam em turmas pequenas.
  \item \textbf{Recursos compartilhados}: espaço físico, materiais e tempo dos coordenadores pedagógicos são divididos entre turmas dentro da mesma escola.
  \item \textbf{Composição das turmas}: se alunos com maior desempenho prévio forem sistematicamente alocados a certos tipos de turma, a comparação entre turmas dentro da escola pode ser contaminada.
\end{itemize}

Este exemplo serve para reforçar um ponto metodológico central: mesmo em experimentos aleatórios, é necessário explicitar as hipóteses de identificação e avaliar sua plausibilidade no contexto específico do estudo. A aleatorização resolve o problema de seleção entre grupos, mas não elimina automaticamente todas as ameaças à validade interna.

\begin{atencaobox}[SUTVA não é automática]
O fato de um experimento ter sido aleatorizado não garante que o SUTVA seja satisfeito. Sempre que turmas, indivíduos ou unidades do mesmo grupo de randomização interagem --- o que é comum em experimentos educacionais e comunitários ---, a hipótese de ausência de interferência precisa ser defendida com argumentos substantivos, não apenas assumida.
\end{atencaobox}


%% ─────────────────────────────────────────────────────────────────
\section{Estimação: Da Diferença de Médias à Regressão com Efeitos Fixos}
\label{sec:star_estimacao}

\subsection*{Especificação 1 — Diferença de médias simples}

A estimativa mais direta é a diferença de médias entre alunos designados para turmas pequenas e demais alunos:
\begin{equation}
  Y_i = \alpha + \beta\, D_i + u_i.
  \label{eq:spec1}
\end{equation}

O coeficiente $\hat{\beta}$ compara a nota média dos alunos em turmas pequenas com a nota média dos demais, sem qualquer condicionamento adicional. Sob randomização incondicional, esse estimador seria não-viesado para o ITT. No STAR, porém, a randomização foi feita \textit{dentro} de escola --- o que significa que comparar alunos de escolas diferentes pode introduzir diferenças pré-existentes entre escolas, não controladas pelo experimento.

\subsection*{Especificação 2 — Efeitos fixos de escola}

A especificação mais alinhada ao desenho experimental inclui efeitos fixos de escola $\lambda_{s(i)}$:
\begin{equation}
  Y_i = \alpha + \beta\, D_i + \lambda_{s(i)} + u_i.
  \label{eq:spec2}
\end{equation}

Ao incluir uma dummy para cada escola, a comparação passa a ser feita \textit{dentro} de escola --- exatamente como a randomização foi conduzida. O coeficiente $\hat{\beta}$ agora compara alunos designados para turma pequena com alunos designados para turma regular \textit{na mesma escola}, eliminando qualquer diferença entre escolas que pudesse confundir a estimativa.

\subsection*{Especificação 3 — Efeitos fixos de escola e covariadas}

A terceira especificação adiciona covariadas pré-tratamento $X_i$ (gênero, raça, elegibilidade para merenda):
\begin{equation}
  Y_i = \alpha + \beta\, D_i + \lambda_{s(i)} + X_i'\gamma + u_i.
  \label{eq:spec3}
\end{equation}

As covariadas têm aqui um papel específico. Em um RCT bem implementado, $D_i$ é independente de $X_i$ por construção (dentro de escola), de modo que $\gamma \neq 0$ não implica que omitir $X_i$ causaria viés. O benefício dos controles aqui é outro: ao capturar parte da variação nos resultados $Y_i$, as covariadas reduzem a variância residual $\hat{u}_i$, tornando o estimador $\hat{\beta}$ mais preciso sem sacrificar a validade causal. \citet{Lin2013} formaliza esse argumento e mostra que, em RCTs, a especificação preferida inclui covariadas centradas e interações com o tratamento, garantindo eficiência assintótica.

\begin{notabox}[{Controles em RCTs --- precisão sem causalidade}]
Em estudos observacionais, adicionamos controles para tentar satisfazer a CIA e eliminar viés. Em RCTs, adicionamos controles para ganhar precisão. A distinção é fundamental: no primeiro caso, a validade causal depende dos controles incluídos; no segundo, a validade já está garantida pela randomização, e os controles apenas melhoram a eficiência.
\end{notabox}

\begin{lstlisting}[style=Rstyle, caption={Três especificações progressivas no STAR}]
# -----------------------------
# 1) Especificação 1: diferença de médias simples
# -----------------------------
spec1 <- lm_robust(read_z ~ D, data = dados_k, se_type = "HC2")

# -----------------------------
# 2) Especificação 2: efeitos fixos de escola
# -----------------------------
spec2 <- feols(read_z ~ D | escola, data = dados_k,
               vcov = "HC1")

# -----------------------------
# 3) Especificação 3: EF de escola + covariadas pré-tratamento
# -----------------------------
spec3 <- feols(read_z ~ D + feminino + nao_branco + merenda | escola,
               data = dados_k, vcov = "HC1")

# Sumário das três especificações
etable(spec1, spec2, spec3,
       keep    = "D",
       headers = c("(1) Sem EF", "(2) EF escola", "(3) EF + covariadas"),
       title   = "Efeito de turma pequena sobre notas de leitura (STAR, K)")
\end{lstlisting}

% ──────────────────────────────────────────────────────────────────
% TABELA 5.1 — Três especificações: efeito de turma pequena sobre notas
% Executar o script R para obter os coeficientes, erros-padrão e
% estatísticas de diagnóstico (R², N, número de escolas).
% ──────────────────────────────────────────────────────────────────
\begin{table}[H]
\centering
\caption{Efeito de turma pequena sobre nota de leitura: três especificações}
\label{tab:spec_comparacao}
\begin{tabular}{lccc}
\toprule
 & (1) Sem EF & (2) EF Escola & (3) EF + Cov. \\
\midrule
Turma pequena ($D_i$)
  & \textit{[inserir]} & \textit{[inserir]} & \textit{[inserir]} \\
  & \textit{[inserir SE]} & \textit{[inserir SE]} & \textit{[inserir SE]} \\[4pt]
Efeitos fixos de escola & Não & Sim & Sim \\
Covariadas pré-tratamento & Não & Não & Sim \\
$R^2$ & \textit{[inserir]} & \textit{[inserir]} & \textit{[inserir]} \\
Observações & \textit{[inserir]} & \textit{[inserir]} & \textit{[inserir]} \\
\bottomrule
\end{tabular}

\vspace{4pt}
\figmeta{Dados: \citet{Krueger1999}, via pacote \texttt{AER}}{Elaboração dos autores}{Variável dependente: nota de leitura padronizada (média = 0, DP = 1). Covariadas: gênero, raça e elegibilidade para merenda gratuita. Erros-padrão HC robustos entre parênteses. * $p < 0{,}05$; ** $p < 0{,}01$; *** $p < 0{,}001$.}
\end{table}

A comparação entre as três colunas da Tabela~\ref{tab:spec_comparacao} ilustra um resultado esperado: em um RCT bem implementado, as estimativas nas três especificações devem ser numericamente similares. Eventuais diferenças entre a coluna (1) e as colunas (2) e (3) indicam diferenças pré-existentes entre escolas, que a aleatorização dentro de escola não controla. A redução nos erros-padrão entre (1) e (3) reflete o ganho de precisão proveniente dos controles.


%% ─────────────────────────────────────────────────────────────────
\section{Inferência I: Qual Erro-Padrão Usar?}
\label{sec:star_ep}

A escolha do estimador de variância não é um detalhe técnico --- é parte da análise substantiva. No STAR, há uma razão estrutural clara para não usar erros-padrão clássicos: alunos dentro da mesma turma compartilham professora, rotina pedagógica, material didático e dinâmica de sala. Seus resultados não são observações independentes. Ignorar essa correlação intraclasse leva a erros-padrão artificialmente pequenos e a testes com taxa de falsos positivos acima do nível nominal.

\subsection*{Correlação intraclasse}

A dependência dentro de grupos pode ser quantificada pela \textbf{correlação intraclasse} (ICC):
\begin{equation}
  \rho = \frac{\sigma^2_c}{\sigma^2_c + \sigma^2_i},
  \label{eq:icc}
\end{equation}
onde $\sigma^2_c$ é a variância entre turmas e $\sigma^2_i$ é a variância dentro de turmas. O \textbf{efeito de desenho} (\textit{design effect}) indica o quanto o erro-padrão efetivo supera o erro-padrão clássico em uma amostra com $\bar{n}$ observações por turma:
\begin{equation}
  \text{DEFF} = 1 + (\bar{n} - 1)\,\rho.
  \label{eq:deff}
\end{equation}

Com $\rho = 0{,}20$ e turmas de tamanho médio $\bar{n} = 20$, o DEFF seria $1 + 19 \times 0{,}20 = 4{,}80$ --- o que significa que os erros-padrão clássicos subestimariam a incerteza por um fator de $\sqrt{4{,}80} \approx 2{,}19$. Estatísticas-$t$ baseadas nesses erros estariam infladas em mais de 100\%.

\subsection*{Quatro abordagens}

Comparamos quatro estimadores de variância:

\begin{enumerate}[leftmargin=*, label=(\roman*)]
  \item \textbf{Clássico}: assume homoscedasticidade e independência entre todas as observações.
  \item \textbf{HC robusto}: corrige heteroscedasticidade, mas ainda trata observações como independentes.
  \item \textbf{Clusterizado por turma}: agrupa observações dentro da mesma turma, reconhecendo a dependência intraclasse. É a escolha natural dado que o tratamento é atribuído no nível da turma.
  \item \textbf{Clusterizado por escola}: agrupa por escola. Mais conservador, com número menor de clusters (79 escolas vs.\ $\sim$300 turmas).
\end{enumerate}

\begin{lstlisting}[style=Rstyle, caption={Comparação de quatro estimadores de variância no STAR}]
# -----------------------------
# 1) Modelo base: EF de escola, sem outras covariadas
# -----------------------------
mod_base <- lm(read_z ~ D + escola, data = dados_k)

# -----------------------------
# 2) Quatro variantes de erro-padrão
# -----------------------------
# (i) Clássico
se_classico <- sqrt(diag(vcov(mod_base)))["D"]

# (ii) HC2 robusto
se_hc2 <- sqrt(vcovHC(mod_base, type = "HC2"))["D", "D"]

# (iii) Clusterizado por turma
dados_k <- dados_k %>%
  mutate(turma_id = paste(schoolk, stark, sep = "_"))

se_turma <- sqrt(vcovCL(mod_base, cluster = ~turma_id,
                        data = dados_k))["D", "D"]

# (iv) Clusterizado por escola
se_escola <- sqrt(vcovCL(mod_base, cluster = ~escola,
                         data = dados_k))["D", "D"]

# Apresentação
resultados_ep <- data.frame(
  Metodo   = c("Clássico", "HC2", "Cluster turma", "Cluster escola"),
  SE       = round(c(se_classico, se_hc2, se_turma, se_escola), 4),
  t_stat   = round(coef(mod_base)["D"] /
             c(se_classico, se_hc2, se_turma, se_escola), 2)
)
print(resultados_ep)
\end{lstlisting}

% ──────────────────────────────────────────────────────────────────
% TABELA 5.2 — Comparação de erros-padrão
% Executar o script R para obter SE, estatística t e p-valor
% sob as quatro abordagens.
% ──────────────────────────────────────────────────────────────────
\begin{table}[H]
\centering
\caption{Estimativa do efeito de turma pequena sob diferentes estimadores de variância}
\label{tab:erros_padrao}
\begin{tabular}{lcccc}
\toprule
Abordagem & Estimativa & Erro-padrão & Estatística $t$ & Valor-$p$ \\
\midrule
Clássico              & \textit{[inserir]} & \textit{[inserir]} & \textit{[inserir]} & \textit{[inserir]} \\
HC2 robusto           & \textit{[inserir]} & \textit{[inserir]} & \textit{[inserir]} & \textit{[inserir]} \\
Cluster por turma     & \textit{[inserir]} & \textit{[inserir]} & \textit{[inserir]} & \textit{[inserir]} \\
Cluster por escola    & \textit{[inserir]} & \textit{[inserir]} & \textit{[inserir]} & \textit{[inserir]} \\
\bottomrule
\end{tabular}

\vspace{4pt}
\figmeta{Dados: \citet{Krueger1999}}{Elaboração dos autores}{Especificação com efeitos fixos de escola. Variável dependente: nota de leitura padronizada. A estimativa pontual é a mesma nas quatro linhas; apenas o estimador de variância varia.}
\end{table}

O padrão esperado na Tabela~\ref{tab:erros_padrao} é que os erros-padrão clusterizados sejam maiores que os clássicos e robustos. A comparação entre clusterização por turma e por escola ilustra o \textit{trade-off} entre adequação ao processo gerador de dados e estabilidade do estimador: com 79 escolas, o estimador sandwich clusterizado por escola opera próximo ao limite inferior recomendado ($G \approx 30$--50), enquanto as turmas oferecem maior número de clusters.

\begin{notabox}[Regra prática de clustering no STAR]
Como o tratamento foi atribuído no nível da \textit{turma}, a clusterização por turma é a escolha teoricamente correta. Se os dados contiverem identificadores de turma confiáveis, devem ser usados. A clusterização por escola é mais conservadora e pode ser reportada como verificação de robustez.
\end{notabox}


%% ─────────────────────────────────────────────────────────────────
\section{Inferência II: Bootstrap}
\label{sec:star_bootstrap}

A Seção~\ref{sec:clustering} mostrou que, com número razoável de clusters, o estimador sandwich clusterizado tende a ter bom desempenho assintótico. Com cerca de 300 turmas, o STAR está bem acima do limiar mínimo --- o que sugere que os intervalos de confiança assintóticos clusterizados devem ser confiáveis. O \textit{bootstrap} é útil aqui como verificação independente e como método para estatísticas cuja distribuição assintótica é menos direta.

\subsection*{Por que reamostrar turmas, não alunos}

Um erro frequente é aplicar o bootstrap padrão --- reamostrar observações individuais com reposição --- em dados com estrutura de clustering. Ao sortear alunos individualmente, o bootstrap padrão quebra a dependência dentro de turmas: uma amostra bootstrap poderia conter cinco cópias do mesmo aluno e nenhuma cópia de sua turma original, destruindo exatamente a estrutura que torna a clusterização necessária. A solução é o \textbf{\textit{cluster bootstrap}}: reamostrar \textit{turmas inteiras} com reposição, mantendo todos os alunos de cada turma sorteada.

\begin{lstlisting}[style=Rstyle, caption={Cluster bootstrap por turma no STAR}]
# -----------------------------
# 1) Cluster bootstrap: reamostrar turmas inteiras
# -----------------------------
rm(list = ls())
# (assumindo que dados_k e turma_id foram criados anteriormente)

set.seed(42)
B       <- 1999
turmas  <- unique(dados_k$turma_id)
G       <- length(turmas)

boot_coefs <- numeric(B)

for (b in 1:B) {
  # Sortear turmas com reposição
  turmas_b  <- sample(turmas, size = G, replace = TRUE)

  # Reconstruir base com todas as observações das turmas sorteadas
  amostra_b <- bind_rows(
    lapply(turmas_b, function(t) dados_k[dados_k$turma_id == t, ])
  )

  # Estimar modelo com EF de escola
  mod_b <- lm(read_z ~ D + escola, data = amostra_b)
  boot_coefs[b] <- coef(mod_b)["D"]
}

# IC bootstrap percentil (95%)
ic_boot <- quantile(boot_coefs, c(0.025, 0.975))
cat("IC bootstrap 95% (percentil): [",
    round(ic_boot[1], 3), ";", round(ic_boot[2], 3), "]\n")

# Comparar com IC assintótico clusterizado
se_assintotico <- sqrt(vcovCL(lm(read_z ~ D + escola, data = dados_k),
                              cluster = ~turma_id)["D", "D"])
est_pont <- coef(lm(read_z ~ D + escola, data = dados_k))["D"]
ic_assint <- est_pont + c(-1, 1) * 1.96 * se_assintotico

cat("IC assintótico 95% (cluster): [",
    round(ic_assint[1], 3), ";", round(ic_assint[2], 3), "]\n")
\end{lstlisting}

A mensagem didática desta subseção é a seguinte: quando há número razoável de clusters e o desenho é bem comportado, os ICs bootstrap e assintóticos clusterizados devem ser numericamente próximos. Uma divergência grande entre os dois seria um sinal de alerta --- sugerindo que alguma hipótese subjacente (número suficiente de clusters, distribuição simétrica do estimador, ausência de observações influentes) pode não estar sendo satisfeita. No STAR, com cerca de 300 turmas, a convergência entre os dois é esperada.


%% ─────────────────────────────────────────────────────────────────
\section{Inferência III: Quatro Disciplinas, Um Problema --- Múltiplos Testes}
\label{sec:star_multiplos}

O STAR oferece resultados de desempenho em mais de uma disciplina: leitura e matemática, ao menos. Quando testamos o efeito de turmas pequenas em todas as disciplinas disponíveis, deparamo-nos com o problema de múltiplos testes discutido na Seção~\ref{sec:multiplos_testes}: quanto maior o número de testes, maior a probabilidade de encontrar ao menos um resultado estatisticamente significativo por acaso, mesmo que o efeito verdadeiro seja nulo em todos os desfechos.

Suponha que testemos o efeito em quatro resultados (por exemplo, leitura e matemática no jardim de infância e no primeiro ano). Ao nível $\alpha = 0{,}05$ com quatro testes independentes, a probabilidade de ao menos uma rejeição espúria é $1 - 0{,}95^4 \approx 18{,}5\%$ --- quase quatro vezes o nível nominal.

\begin{lstlisting}[style=Rstyle, caption={Múltiplos testes: quatro resultados com correção de Bonferroni e BH-FDR}]
# -----------------------------
# 1) Estimar o efeito nas quatro variáveis de resultado
# Leitura (K), Matemática (K), Leitura (1o ano), Matemática (1o ano)
# -----------------------------
dados_multi <- STAR %>%
  filter(!is.na(stark)) %>%
  mutate(
    D      = as.integer(stark == "small"),
    escola = as.factor(schoolk),
    read_k_z  = scale(readk)[, 1],
    math_k_z  = scale(mathk)[, 1],
    read_1_z  = scale(read1)[, 1],
    math_1_z  = scale(math1)[, 1]
  )

resultados <- c("read_k_z", "math_k_z", "read_1_z", "math_1_z")
nomes      <- c("Leitura (K)", "Matemática (K)",
                "Leitura (1o ano)", "Matemática (1o ano)")

pvalores    <- numeric(length(resultados))
estimativas <- numeric(length(resultados))

for (j in seq_along(resultados)) {
  df_j  <- dados_multi[!is.na(dados_multi[[resultados[j]]]), ]
  mod_j <- lm(as.formula(paste(resultados[j], "~ D + escola")),
              data = df_j)
  se_j  <- sqrt(vcovCL(mod_j, cluster = ~schoolk, data = df_j)["D", "D"])
  estimativas[j] <- coef(mod_j)["D"]
  pvalores[j]    <- 2 * pt(-abs(coef(mod_j)["D"] / se_j),
                            df = df_j$escola %>% n_distinct() - 2)
}

# -----------------------------
# 2) Ajuste por Bonferroni e BH-FDR
# -----------------------------
pv_bonf <- p.adjust(pvalores, method = "bonferroni")
pv_bh   <- p.adjust(pvalores, method = "BH")

tabela_multi <- data.frame(
  Resultado    = nomes,
  Estimativa   = round(estimativas, 3),
  P_bruto      = round(pvalores, 3),
  P_Bonferroni = round(pv_bonf, 3),
  P_BH_FDR     = round(pv_bh, 3)
)
print(tabela_multi)
\end{lstlisting}

% ──────────────────────────────────────────────────────────────────
% TABELA 5.3 — Múltiplos testes: quatro resultados
% Executar o script R para obter estimativas, p-valores brutos
% e p-valores ajustados por Bonferroni e BH-FDR.
% ──────────────────────────────────────────────────────────────────
\begin{table}[H]
\centering
\caption{Efeito de turma pequena: quatro resultados com ajuste para múltiplos testes}
\label{tab:multiplos_testes_star}
\begin{tabular}{lcccc}
\toprule
Resultado & Estimativa & $p$ bruto & $p$ Bonferroni & $p$ BH-FDR \\
\midrule
Leitura (K)          & \textit{[inserir]} & \textit{[inserir]} & \textit{[inserir]} & \textit{[inserir]} \\
Matemática (K)       & \textit{[inserir]} & \textit{[inserir]} & \textit{[inserir]} & \textit{[inserir]} \\
Leitura (1.º ano)    & \textit{[inserir]} & \textit{[inserir]} & \textit{[inserir]} & \textit{[inserir]} \\
Matemática (1.º ano) & \textit{[inserir]} & \textit{[inserir]} & \textit{[inserir]} & \textit{[inserir]} \\
\bottomrule
\end{tabular}

\vspace{4pt}
\figmeta{Dados: \citet{Krueger1999}}{Elaboração dos autores}{Especificação com efeitos fixos de escola. Erros-padrão clusterizados por escola. Bonferroni: nível ajustado $= 0{,}05/4 = 0{,}0125$. BH-FDR: controla proporção esperada de falsas descobertas a 5\%.}
\end{table}

A leitura da Tabela~\ref{tab:multiplos_testes_star} deve responder a três perguntas: (i) quais efeitos são significativos sem qualquer ajuste? (ii) quais sobrevivem ao ajuste de Bonferroni? (iii) o padrão é consistente entre disciplinas e anos, ou concentrado em um único resultado? Uma estimativa que se mantém significativa após Bonferroni sugere evidência mais sólida do que aquela que só resiste ao p-valor bruto. Quando os efeitos são positivos e de magnitude similar em todas as disciplinas, a credibilidade aumenta: seria improvável que o mesmo viés produzisse o mesmo sinal em quatro desfechos distintos.

\begin{notabox}[Bonferroni vs.\ BH-FDR]
Bonferroni controla a probabilidade de \textit{ao menos} um falso positivo no conjunto de testes (FWER), tornando cada teste individual muito conservador. BH-FDR controla a proporção esperada de falsas descobertas entre as rejeições, sendo menos conservador quando $m$ é grande. Para quatro desfechos, a diferença entre os dois é pequena --- Bonferroni com $\alpha^* = 0{,}0125$ vs.\ BH com limiar adaptativo. Com dezenas ou centenas de desfechos, a diferença torna-se mais relevante.
\end{notabox}


%% ─────────────────────────────────────────────────────────────────
\section{Inferência IV: Inferência por Permutação}
\label{sec:star_permutacao}

A inferência por permutação, apresentada na Seção~\ref{sec:design_based}, tem uma propriedade especial no STAR: como o processo de aleatorização está documentado (randomização dentro de escola), a distribuição de referência pode ser construída a partir desse mesmo processo, sem qualquer hipótese distribucional sobre os erros.

\subsection*{Respeitando o desenho na permutação}

Como a aleatorização ocorreu \textit{dentro} de escola, a permutação deve respeitar esse bloco: permutamos os rótulos de tratamento ($D_i$) apenas entre alunos da mesma escola, preservando o número de tratados por escola. Permutar globalmente ignoraria a estrutura por blocos e produziria distribuições de referência inadequadas.

\begin{lstlisting}[style=Rstyle, caption={Inferência por permutação com aleatorização em bloco (escola)}]
# -----------------------------
# 1) Estatística observada
# -----------------------------
mod_obs  <- lm(read_z ~ D + escola, data = dados_k)
t_obs    <- coef(mod_obs)["D"] /
            sqrt(vcovCL(mod_obs, cluster = ~turma_id,
                        data = dados_k)["D", "D"])

cat("Estatística t observada:", round(t_obs, 3), "\n")

# -----------------------------
# 2) Distribuição de permutação
# Permutar D dentro de escola (respeitar o bloco de randomização)
# -----------------------------
set.seed(42)
P <- 4999
t_perm <- numeric(P)

for (p in 1:P) {
  dados_perm <- dados_k %>%
    group_by(escola) %>%
    mutate(D_perm = sample(D)) %>%   # permuta D dentro de escola
    ungroup()

  mod_p   <- lm(read_z ~ D_perm + escola, data = dados_perm)
  se_p    <- sqrt(vcovCL(mod_p, cluster = ~turma_id,
                          data = dados_perm)["D_perm", "D_perm"])
  t_perm[p] <- coef(mod_p)["D_perm"] / se_p
}

# -----------------------------
# 3) P-valor de permutação (bilateral)
# -----------------------------
p_perm    <- mean(abs(t_perm) >= abs(t_obs))
p_assint  <- 2 * pnorm(-abs(t_obs))

cat("P-valor de permutação (bilateral): ", round(p_perm, 4), "\n")
cat("P-valor assintótico (normal):      ", round(p_assint, 4), "\n")
\end{lstlisting}

A comparação entre o p-valor de permutação e o p-valor assintótico clusterizado ilustra um resultado típico de experimentos com amostras grandes: os dois tendem a ser numericamente próximos. A inferência por permutação é particularmente atraente porque usa diretamente o processo de aleatorização descrito no protocolo experimental --- não há necessidade de aproximações assintóticas nem hipóteses sobre a distribuição dos erros. Em amostras menores, com poucos clusters ou distribuições assimétricas, a inferência por permutação pode diferir substancialmente dos métodos assintóticos e, nesse caso, oferece garantias de validade mais robustas.


%% ─────────────────────────────────────────────────────────────────
\section{Testes de Falsificação}
\label{sec:star_falsificacao}

Os testes de falsificação não ``provam'' que o STAR tem identificação válida, mas ajudam a detectar problemas que comprometeriam a interpretação causal. Apresentamos três verificações: balanceamento de covariadas pré-tratamento, atrito diferencial e balanceamento das características das professoras.

\subsection*{Balanceamento de covariadas pré-tratamento}

Em um experimento bem implementado, características pré-tratamento dos alunos não devem prever sistematicamente a designação para turma pequena. Para testar isso, regredimos $D_i$ (indicador de turma pequena) nas covariadas pré-tratamento disponíveis e realizamos um F-teste conjunto de significância:
\begin{equation}
  D_i = \alpha + \lambda_{s(i)} + \delta_1\,\text{feminino}_i + \delta_2\,\text{nao\_branco}_i + \delta_3\,\text{merenda}_i + v_i.
  \label{eq:teste_balanceamento}
\end{equation}

A hipótese nula é $H_0: \delta_1 = \delta_2 = \delta_3 = 0$, ou seja, nenhuma covariada prediz o tratamento dentro de escola. Sob aleatorização bem conduzida, esperamos não rejeitar essa hipótese.

\begin{lstlisting}[style=Rstyle, caption={F-teste de balanceamento de covariadas}]
# -----------------------------
# 1) Regredir tratamento em covariadas pré-tratamento + EF escola
# -----------------------------
mod_bal <- lm(D ~ escola + feminino + nao_branco + merenda,
              data = dados_k)

# F-teste conjunto: apenas para as covariadas (excluindo EF de escola)
library(car)
ftest_bal <- linearHypothesis(mod_bal,
  c("feminino = 0", "nao_branco = 0", "merenda = 0"),
  vcov. = vcovCL(mod_bal, cluster = ~turma_id, data = dados_k))

cat("=== Teste F conjunto de balanceamento ===\n")
print(ftest_bal)

# -----------------------------
# 2) Tabela de médias por grupo (dentro de escola)
# -----------------------------
dados_k %>%
  group_by(D) %>%
  summarise(
    pct_feminino   = mean(feminino, na.rm = TRUE),
    pct_nao_branco = mean(nao_branco, na.rm = TRUE),
    pct_merenda    = mean(merenda, na.rm = TRUE),
    n              = n()
  ) %>%
  print()
\end{lstlisting}

% ──────────────────────────────────────────────────────────────────
% TABELA 5.4 — Teste de balanceamento: covariadas pré-tratamento
% Executar o script R para obter as médias por grupo e a estatística F.
% ──────────────────────────────────────────────────────────────────
\begin{table}[H]
\centering
\caption{Verificação de balanceamento: covariadas pré-tratamento por condição experimental}
\label{tab:balanceamento}
\begin{tabular}{lccc}
\toprule
Covariada & Turma pequena ($D=1$) & Controle ($D=0$) & Dif. padronizada \\
\midrule
Feminino (\%)         & \textit{[inserir]} & \textit{[inserir]} & \textit{[inserir]} \\
Não-branco (\%)       & \textit{[inserir]} & \textit{[inserir]} & \textit{[inserir]} \\
Merenda gratuita (\%) & \textit{[inserir]} & \textit{[inserir]} & \textit{[inserir]} \\
\midrule
\multicolumn{3}{l}{F-teste conjunto ($H_0$: todas as diferenças = 0)}
  & $F = $ \textit{[inserir]} \\
\multicolumn{3}{l}{Valor-$p$}
  & \textit{[inserir]} \\
\midrule
Observações           & \textit{[inserir]} & \textit{[inserir]} & \\
\bottomrule
\end{tabular}

\vspace{4pt}
\figmeta{Dados: \citet{Krueger1999}}{Elaboração dos autores}{Comparação feita dentro de escola (efeitos fixos de escola incluídos na regressão). Dif. padronizada $= (\bar{X}_1 - \bar{X}_0)/\sqrt{(s_1^2 + s_0^2)/2}$.}
\end{table}

\subsection*{Atrito diferencial}

A aleatorização garante comparabilidade no momento inicial do experimento. Ao longo dos quatro anos, no entanto, alunos podem sair do experimento --- por mudança de escola, ausência prolongada ou dados faltantes. Se a \textit{probabilidade de permanecer na amostra} for diferente entre turmas pequenas e regulares, a comparação ao final do período pode estar contaminada: estaríamos comparando grupos sistematicamente diferentes dos originalmente sorteados.

Para verificar atrito diferencial, testamos se $D_i$ prediz a probabilidade de o aluno ter notas disponíveis no ano de interesse:
\begin{equation}
  \mathds{1}[Y_i \text{ observado}] = \alpha + \beta_{\text{atrito}}\, D_i + \lambda_{s(i)} + u_i.
  \label{eq:atrito}
\end{equation}

Um coeficiente $\hat{\beta}_{\text{atrito}}$ significativamente diferente de zero indica que tratados e controles têm taxas de atrito distintas --- o que comprometeria a validade interna mesmo com aleatorização original correta.

\begin{lstlisting}[style=Rstyle, caption={Verificação de atrito diferencial}]
# -----------------------------
# Indicador de permanência na amostra para cada ano
# Testamos se D prediz estar na amostra no 1o ano
# -----------------------------
dados_atrito <- STAR %>%
  filter(!is.na(stark)) %>%
  mutate(
    D              = as.integer(stark == "small"),
    escola         = as.factor(schoolk),
    presente_ano1  = as.integer(!is.na(read1))
  )

mod_atrito <- lm_robust(presente_ano1 ~ D,
                        fixed_effects = ~escola,
                        clusters      = schoolk,
                        data          = dados_atrito,
                        se_type       = "CR2")

cat("=== Teste de atrito diferencial (1o ano) ===\n")
cat("Coeficiente D:  ", round(coef(mod_atrito)["D"], 4), "\n")
se_atrito <- summary(mod_atrito)$coefficients["D", 2]
pv_atrito <- summary(mod_atrito)$coefficients["D", 4]
cat("Erro-padrão:    ", round(se_atrito, 4), "\n")
cat("Valor-p:        ", round(pv_atrito, 3), "\n")
\end{lstlisting}

\subsection*{Placebo de professora: características docentes como teste de implementação}

Um aspecto do STAR que merece atenção é a qualidade da implementação. Se as professoras mais experientes ou mais tituladas foram sistematicamente alocadas a turmas pequenas, o efeito estimado confundiria o tamanho da turma com a qualidade docente. Testamos isso verificando o balanceamento das características das professoras entre condições experimentais:

\begin{equation}
  \text{experiência}_{c} = \alpha + \beta_{\text{prof}}\, D_c + \lambda_{s(c)} + w_c,
  \label{eq:placebo_prof}
\end{equation}

onde $c$ indexa turmas e $D_c$ é o indicador de turma pequena. Se $\hat{\beta}_{\text{prof}}$ for significativamente diferente de zero, há evidência de que turmas pequenas e regulares foram atribuídas a professoras sistematicamente diferentes.

\begin{lstlisting}[style=Rstyle, caption={Balanceamento de características das professoras entre condições}]
# -----------------------------
# Agregar para o nível de turma (cada turma tem uma professora)
# -----------------------------
dados_prof <- STAR %>%
  filter(!is.na(stark), !is.na(experiencek)) %>%
  mutate(D = as.integer(stark == "small")) %>%
  group_by(schoolk, stark) %>%
  summarise(
    experiencia = first(experiencek),
    D           = first(D),
    .groups = "drop"
  ) %>%
  mutate(escola = as.factor(schoolk))

mod_prof <- lm_robust(experiencia ~ D,
                      fixed_effects = ~escola,
                      clusters      = schoolk,
                      data          = dados_prof,
                      se_type       = "CR2")

cat("=== Placebo de professora: experiência docente ===\n")
cat("Coeficiente D:  ", round(coef(mod_prof)["D"], 3), "\n")
se_prof <- summary(mod_prof)$coefficients["D", 2]
pv_prof <- summary(mod_prof)$coefficients["D", 4]
cat("Erro-padrão:    ", round(se_prof, 3), "\n")
cat("Valor-p:        ", round(pv_prof, 3), "\n")
\end{lstlisting}

\begin{atencaobox}[Testes de falsificação: limites e interpretação]
Passar em todos os testes de falsificação é evidência circunstancial favorável à validade do desenho, mas não é prova de causalidade. Uma estratégia pode superar todos os placebos disponíveis e ainda ter identificação inválida por razões não testáveis --- por exemplo, um canal de interferência entre turmas que não deixa rastro nas covariadas observadas. Os testes de falsificação devem ser lidos em conjunto com o argumento teórico sobre o mecanismo de identificação, não como substitutos a ele.
\end{atencaobox}


%% ─────────────────────────────────────────────────────────────────
\subsection*{Fechamento: o que o STAR ensinou neste guia}

O Experimento STAR percorreu, em um único exemplo empírico, toda a cadeia metodológica discutida neste guia. A \textbf{aleatorização dentro de escola} resolveu o problema de identificação: dentro de cada escola, tratados e controles são comparáveis em média, eliminando o viés de seleção. A \textbf{regressão com efeitos fixos de escola} alinhou a estimação ao desenho experimental, melhorando precisão sem comprometer validade. A discussão sobre \textbf{erros-padrão} revelou que a escolha do estimador de variância é uma decisão substantiva: ignorar a correlação intraclasse leva a inferências enganosas. O \textbf{bootstrap} por cluster confirmou os resultados assintóticos quando o número de grupos é suficiente, e mostrou quando os dois tendem a divergir. A análise de \textbf{múltiplos testes} forçou uma resposta honesta sobre quais resultados são robustos: um efeito que sobrevive a Bonferroni em quatro disciplinas é mais convincente do que um único p-valor favorável. A \textbf{inferência por permutação} usou diretamente o processo de aleatorização para construir distribuições de referência exatas, sem aproximações assintóticas. Por fim, os \textbf{testes de falsificação} --- balanceamento de covariadas, atrito diferencial e placebo de professora --- forneceram evidência circunstancial sobre a qualidade da implementação experimental.

O ponto que conecta todos esses elementos é o seguinte: cada ferramenta responde a uma pergunta diferente. Identificação pergunta \textit{o que estamos estimando?} Estimação pergunta \textit{como calculamos isso?} Inferência pergunta \textit{quanta incerteza há?} Bootstrap e permutação perguntam \textit{essa incerteza é bem caracterizada pelas aproximações assintóticas?} Múltiplos testes perguntam \textit{o resultado sobrevive a buscas seletivas?} Falsificação pergunta \textit{as hipóteses por trás da estimação são plausíveis?} Nenhuma dessas perguntas é dispensável, e a resposta a qualquer uma delas pressupõe que as anteriores tenham sido respondidas com rigor.
\chapter{Conclusão}
\label{cap:conclusao}

Este guia percorreu a cadeia completa de uma análise causal quantitativa:
da pergunta de política ao número numa tabela. Cada escolha metodológica ao longo desse caminho reflete uma hipótese sobre o mundo, e cada hipótese carrega um custo que precisa ser explicitado e defendido.

O argumento central pode ser enunciado de forma direta: \textbf{causalidade não é uma propriedade dos dados, mas das hipóteses que fazemos sobre como os dados foram gerados}. Dados observacionais não ``falam por si''; eles precisam ser lidos à luz de um argumento de identificação. Esse argumento determina o que é estimado, e a estimação determina com que precisão. Sem o primeiro elo, os outros dois não produzem nada interpretável.


\section{A Cadeia de Análise: Um Balanço}

\subsection*{Identificação}

O Capítulo~\ref{cap:identificacao} mostrou que o problema fundamental da
inferência causal --- a impossibilidade de observar o contrafactual --- não
pode ser eliminado, apenas contornado. O contorno exige hipóteses sobre o
processo de seleção ao tratamento. A exogeneidade incondicional (satisfeita
por randomização) elimina o viés de seleção por construção. A independência
condicional (CIA) exige que todas as variáveis que simultaneamente afetam
tratamento e resultado estejam disponíveis nos dados --- uma hipótese
substantiva, não testável diretamente.

A lição prática: antes de estimar qualquer coisa, o pesquisador precisa
responder à pergunta ``por que comparar esses dois grupos é informativo sobre
o efeito causal?'' A resposta a essa pergunta é o argumento de identificação.
Sem ele, nenhum refinamento econométrico subsequente produz interpretação
causal.

\subsection*{Estimação}

O Capítulo~\ref{cap:estimacao} mostrou que identificação e estimação são
problemas distintos. Mesmo quando um estimando é identificado, o estimador
precisa efetivamente calculá-lo. A forma funcional importa: OLS com controles
lineares pode calcular uma média ponderada adversa dos efeitos locais quando
os pesos implícitos diferem dos pesos de interesse. O IPW oferece robustez à
forma funcional de $\mathbb{E}[Y \mid D, X]$ ao custo de depender do
propensity score. Estimadores duplamente robustos combinam os dois modelos
e são consistentes se pelo menos um estiver correto.

\subsection*{Inferência}

O Capítulo~\ref{cap:inferencia} mostrou que a escolha do estimador de
variância é tão substantiva quanto a escolha do estimador do parâmetro.
Erros-padrão clássicos supõem independência e homoscedasticidade ---
hipóteses raramente plausíveis em dados de políticas públicas. O clustering
é obrigatório quando o tratamento é atribuído em grupo. Com poucos clusters,
o wild cluster bootstrap é a solução preferida.

O p-valor mede uma coisa específica: a probabilidade de observar dados tão
extremos quanto os observados, se a hipótese nula fosse verdadeira. Ele não
mede o tamanho do efeito e não substitui o intervalo de confiança. Reportar
o tamanho do efeito com seu intervalo de confiança é sempre mais informativo
do que reportar apenas o p-valor.


\section{O Experimento STAR como Laboratório}

O Capítulo~\ref{cap:star} demonstrou todos esses princípios em um único
exemplo empírico. O STAR mostra que:

\begin{itemize}[leftmargin=*]
  \item A randomização resolve o problema de identificação dentro dos blocos
    experimentais --- mas não elimina a necessidade de explicitar as hipóteses
    (SUTVA pode ser questionado; a qualidade da implementação importa).
  \item A especificação com efeitos fixos de escola alinha a estimação ao
    design experimental, melhorando precisão sem comprometer validade.
  \item A escolha entre erros clássicos, HC e clusterizados produz diferenças
    substanciais nos erros-padrão --- diferenças com consequências reais para
    as conclusões.
  \item Testar quatro disciplinas sem correção para múltiplos testes infla a
    taxa de falsos positivos. Os ajustes de Bonferroni e BH-FDR revelam quais
    resultados são robustos.
  \item A inferência por permutação usa diretamente o processo de aleatorização
    e produz p-valores exatos em amostras finitas.
  \item Os testes de falsificação fornecem evidência circunstancial sobre a
    qualidade do desenho, mas não provam identificação válida.
\end{itemize}


\section{O Que Este Guia Não Cobre}

Este guia apresentou os fundamentos da inferência causal em contextos
experimentais e observacionais com seleção em observáveis. Os casos em que
a CIA não pode ser mantida exigem estratégias quasi-experimentais que exploram
variação exógena de outras naturezas. Esses métodos são o objeto dos demais
guias da Série Avaliação na Prática:

\begin{itemize}[leftmargin=*]
  \item \textbf{Regressão com Descontinuidade (RDD)}: explora pontos de corte
    em regras de elegibilidade para comparações localmente aleatórias.
  \item \textbf{Diferenças em Diferenças (DiD)}: utiliza variação temporal no
    tratamento para controlar por diferenças fixas não observadas entre grupos.
  \item \textbf{Pareamento (\textit{Matching})}: constrói um grupo de controle
    sintético com covariadas similares às do grupo tratado.
  \item \textbf{Controle Sintético (SCM)}: combina unidades de controle para
    replicar a trajetória pré-tratamento da unidade tratada.
  \item \textbf{Variáveis Instrumentais (IV)}: utiliza uma variável que afeta
    o tratamento mas não afeta diretamente o resultado.
\end{itemize}

\begin{conexaobox}[Série Avaliação na Prática]
Os guias de RDD, DiD (dois volumes), Pareamento, SCM, Variáveis Instrumentais,
Poder Estatístico e Fontes de Dados estão disponíveis na Série Avaliação na
Prática do FGV CLEAR. Cada guia pode ser lido de forma independente, mas os
conceitos de identificação, estimação e inferência desenvolvidos aqui são
pressupostos em todos eles.
\end{conexaobox}


\section{Uma Nota Final}

A econometria causal não é uma caixa de ferramentas que, aplicada
corretamente, produz respostas definitivas. É um conjunto de disciplinas
intelectuais para raciocinar com rigor sobre o que podemos e o que não podemos
aprender a partir de dados. Mesmo o experimento mais bem controlado carrega
hipóteses que precisam ser defendidas, não apenas assumidas.

O hábito que este guia procura cultivar é uma pergunta, repetida a cada estimativa: \textit{o que exatamente estou estimando? Sob que hipóteses esse número tem interpretação causal? Quão plausíveis são essas hipóteses neste contexto?} Um pesquisador que faz essas perguntas sistematicamente produz avaliações mais informativas do que aquele que domina ferramentas sofisticadas sem questionar o que elas calculam.

Quando essas perguntas são levadas a sério --- e as hipóteses são explicitadas e debatidas em vez de varridas para baixo do tapete ---, a inferência causal oferece a melhor base disponível para decisões de política pública informadas por evidência.

% ============================================================
% REFERÊNCIAS
% ============================================================
\backmatter
\bibliography{ref}

\end{document}
